% ------ METADATA ------
\newcommand{\bookauthor}{\ml{$0$}{Эмиль~Весна}{Emil~Viesn\'{a}}}
\newcommand{\booktitle}{\ml{$0$}{СКОРБЯЩИЕ}{The~Bemoaners}}
\newcommand{\bookstarted}{\ml{$0$}{13 августа 2012}{May 3, 2015}}
\newcommand{\bookfinished}{\ml{$0$}{пока нет}{not yet}}
% ------ METADATA ------

% ----- XELATEX SYMBOL -----
\usepackage{xltxtra}
% ----- XELATEX SYMBOL -----

% ----- HYPHENATION -----
\usepackage{hyphenat}
% ----- HYPHENATION -----

% ----- GEOMETRY -----
\usepackage[left=1.5cm,right=1.5cm,top=2cm,bottom=2cm,bindingoffset=0.5cm]{geometry}
% ----- GEOMETRY -----

% ----- INCLUDE PDF AS PAGES -----
\usepackage{pdfpages}
% ----- INCLUDE PDF AS PAGES -----

% ----- DROPPING CAP -----
\usepackage{type1cm,lettrine}
% ----- DROPPING CAP -----

% ----- FONTS -----
\renewcommand{\baselinestretch}{1.2}
\setmainfont{Linux Libertine}
% ----- FONTS -----

% ------ HYPERLINKS ------
\usepackage{hyperref}
\definecolor{LinkColor}{HTML}{0969DA}
\hypersetup{colorlinks=true, linkcolor=LinkColor, citecolor=LinkColor, filecolor=LinkColor, urlcolor=LinkColor}
% ------ HYPERLINKS ------

% ------ FANCY PAGE STYLE ------
\setlength{\headheight}{15pt}
\usepackage{fancyhdr}
\pagestyle{fancy}
\fancyhead[LE,RO]{\thepage}
\fancyhead[LO]{{\small\textsc{\booktitle}}}
\fancyhead[RE]{{\small\textsc{\bookauthor}}}
\fancyfoot{}
    \fancypagestyle{plain}{
    \renewcommand{\headrulewidth}{0mm}
    \fancyhead{}
    \fancyfoot{}
}
% ------ FANCY PAGE STYLE ------

% ------ ELEMENTS ------
\newcommand{\asterism}{\vspace{1em}{\centering\Large\bfseries$\ast~\ast~\ast$\par}\vspace{1em}}
\newcommand{\textspace}{\vspace{1em}{\centering\Large\bfseries<\dots>\par}\vspace{1em}}
\newcommand{\FM}{\footnotemark}
\newcommand{\FL}[2]{\footnotetext{См. \textit{\hyperlink{#1}{#2}}.}}
\newcommand{\FA}[1]{\footnotetext{#1 \emph{\ml{$0$}{---~Прим.~авт.}{---~Author.}}}}
\newcommand{\theterm}[3]{\textbf{\hypertarget{#1}{#2}} --- #3}
\newcommand{\thesynonim}[3]{\textbf{#2} --- см. \textit{\hyperlink{#1}{#3}}.}
\newcommand{\theorigin}[3]{\textit{#1:} #2 --- #3}
% ------ ELEMENTS ------

% ------ EPIGRAPH ------
\usepackage{epigraph}
\renewcommand{\epigraphsize}{\footnotesize}
\epigraphrule=0pt
\epigraphwidth=8cm
\usepackage{etoolbox}
\AtBeginEnvironment{quote}{\itshape}
\makeatletter
\newlength\episourceskip
\pretocmd{\@episource}{\em}{}{}
\apptocmd{\@episource}{\em}{}{}
\patchcmd{\epigraph}{\@episource{#1}\\}{\@episource{#1}\\[\episourceskip]}{}{}
\makeatother
% ------ EPIGRAPH ------

% ------ NAMES ------
\newcommand{\Aatris}{\"{A}\={a}tr\v{\i}s}
\newcommand{\Akchsar}{\`{A}kchs\r{a}r}
\newcommand{\Chhammitrai}{Chh\`{a}mm\={\i}tr\^{a}i}
\newcommand{\Chhanei}{Chh\r{a}n\^{e}i}
\newcommand{\Chitram}{Ch\"{\i}tr\'{a}m}
\newcommand{\Choe}{Cho\^{e}}
\newcommand{\choe}{cho\^{e}}
\newcommand{\Harrmatr}{H\r{a}rrm\`{a}tr}
\newcommand{\tHat}{H\={a}t}
\newcommand{\Hei}{H\r{e}i}
\newcommand{\hei}{h\r{e}i}
\newcommand{\Hoesitr}{Ho\`{e}s\={\i}tr}
\newcommand{\hoesitr}{ho\`{e}s\={\i}tr}
\newcommand{\Kchaagotr}{Kch\^{a}\={a}g\~{o}tr}
\newcommand{\kchaagotr}{kch\^{a}\={a}g\~{o}tr}
\newcommand{\Kcharas}{Kch\'{a}r\v{a}s}
\newcommand{\Kchatrim}{Kch\r{a}tr\"{\i}m}
\newcommand{\Kchenoel}{Kch\={e}no\^{e}}
\newcommand{\kchenoel}{kch\={e}no\^{e}}
\newcommand{\Kchenoet}{Kch\"{e}no\^{e}}
\newcommand{\kchenoet}{kch\"{e}no\^{e}}
\newcommand{\Kchoho}{Kch\`{o}h\^{o}}
\newcommand{\Kchotlam}{Kch\={o}tl\'{a}m}
\newcommand{\Kchotris}{Kch\={o}tr\v{\i}s}
\newcommand{\Kihotr}{K\^{\i}h\~{o}tr}
\newcommand{\kihotr}{k\^{\i}h\~{o}tr}
\newcommand{\Kukchuatr}{K\`{u}kchu\={a}tr}
\newcommand{\kukchuatr}{k\`{u}kchu\={a}tr}
\newcommand{\Laaka}{L\={a}\"{a}k\^{a}}
\newcommand{\laaka}{l\={a}\"{a}k\^{a}}
\newcommand{\Lechoe}{L\={e}cho\`{e}}
\newcommand{\lechoe}{l\={e}cho\`{e}}
\newcommand{\Likas}{L\^{\i}k\v{a}s}
\newcommand{\Likchmas}{L\={\i}kchm\r{a}s}
\newcommand{\Likchoe}{L\^{\i}kcho\^{e}}
\newcommand{\Loem}{Lo\~{e}m}
\newcommand{\Maaras}{M\"{a}\={a}r\v{a}s}
\newcommand{\Mitchoe}{M\={\i}tcho\^{e}}
\newcommand{\Mitlikch}{M\={\i}tl\={\i}kch}
\newcommand{\Mitris}{M\={\i}tr\={\i}s}
\newcommand{\Oerchoe}{O\r{e}rcho\^{e}}
\newcommand{\Oerlikch}{O\r{e}rl\'{\i}kch}
\newcommand{\Sat}{S\={a}t}
\newcommand{\Satchir}{S\={a}tch\"{\i}r}
\newcommand{\Satrakch}{S\={a}tr\`{a}kch}
\newcommand{\Seli}{S\r{e}l\={\i}}
\newcommand{\Sirtchu}{S\r{\i}rtch\'{u}}
\newcommand{\Sitris}{S\~{\i}tr\v{\i}s}
\newcommand{\Siusiu}{S\~{\i}u-s\~{\i}u}
\newcommand{\siusiu}{s\~{\i}u-s\~{\i}u}
\newcommand{\Sogcho}{S\"{o}gch\={o}}
\newcommand{\Sotron}{S\~{o}tr\`{o}n}
\newcommand{\Tchalas}{Tch\r{a}l\v{a}s}
\newcommand{\Tchammitr}{Tch\`{a}mm\={\i}tr}
\newcommand{\Tchanoe}{Tch\r{a}no\^{e}}
\newcommand{\Tchartchaahitr}{Tch\~{a}rtch\"{a}\={a}h\r{\i}tr}
\newcommand{\Tchartchu}{Tch\~{a}rtch\'{u}}
\newcommand{\Tchitron}{Tch\"{\i}tr\`{o}n}
\newcommand{\Tchu}{Tch\`{u}}
\newcommand{\tchu}{tch\`{u}}
\newcommand{\Technku}{T\`{e}chnk\r{u}}
\newcommand{\Tesarrokch}{Te's\'{a}rr\r{o}kch}
\newcommand{\Tesatron}{Te's\'{a}tr\v{o}n}
\newcommand{\Traa}{Tr\={a}\"{a}}
\newcommand{\traa}{tr\={a}\"{a}}
\newcommand{\Trai}{Tr\r{a}i}
\newcommand{\trai}{tr\r{a}i}
\newcommand{\TraRenkchal}{Tr\r{a}-R\={e}nkch\'{a}l}
\newcommand{\Trukchual}{Tr\`{u}kchu\r{a}l}
% ------ NAMES ------

\begin{document}

% ----- FRONT COVER -----
% \includepdf[pages={1}]{\coverfrontfilename}
% ----- FRONT COVER -----

% ----- EMPTY PAGE -----
% \newpage\thispagestyle{plain}~
% ----- EMPTY PAGE -----

% ------ TITLE PAGE ------
\begin{titlepage}
{
\centering
{~\par}
\vspace{0.25\textheight}
{\LARGE\bookauthor\par}
\vspace{1.3cm}
{\Huge\textbf{\booktitle}\par}
\vfill
{\includegraphics[width=6em]{\publisherlogofilename}\par}
}
\end{titlepage}
% ------ TITLE PAGE ------

% ----- TRIGGER WARNING -----
% \newpage\thispagestyle{plain}
% {~\par}
% \vfill
% {\centering\includegraphics[width=4em]{tw.pdf}\par}
% \vspace{1em}
% {\centering\Large\strong{\ml{$0$}{Предупреждение}{Content Warning}}\par}
% \vspace{0.5em}
% {\centering\large{Данное произведение содержит описание неизлечимых болезней, психических состояний, врачебных манипуляций, упоминания суицида, химических зависимостей, сцены насилия и убийств, мат, а также гомофобную, трансфобную, сексистскую и расистскую лексику.}\par}
% \vfill
% {~\par}
% ----- TRIGGER WARNING -----

% ----- TABLE OF CONTENTS -----
\tableofcontents
% ----- TABLE OF CONTENTS -----

\pagestyle{fancy}


\chapter*{Предисловие}
\addcontentsline{toc}{chapter}{Предисловие}

Это отвратительная книга.
И дело даже не в отсутствии у меня таланта к писательству.
Отвратительны по сути описываемые в ней события --- так же отвратительны, как вскрытие гнойника, забой животного или прочистка засорившейся канализации.
Едва ли кто-то усомнится, что иногда это необходимо;
но искать в сделанном честь и тем более красоту бесполезно.
Происшедшее не спасают и не оправдывают даже те крупицы нежности и тепла, которые я постарался вложить в книгу.
И если кто-то другой, описывая то же самое, заикнётся о героизме --- можете смело назвать его идиотом или лжецом.

\textspace

Здесь нужно дать кое-какие объяснения.
Демоны не сразу научились <<видеть>> сапиентов.
Долгое время связь с телом была единственным способом почувствовать Мир Фотона.
Чувствительность первых демонов оставляла желать лучшего --- мы не могли засечь спящих, чрезвычайно сложной задачей было выслеживание животных.
Первые люди не могли нас уничтожить;
однако благодаря нашему несовершенству они дали нам достойный отпор.
Не так-то легко оказалось превратить в рабов тех, чьи предки были свободны на протяжении двадцати пяти тысячелетий.

\textspace

Возможно, кому-то история покажется неправдоподобной, но всё описанное здесь --- чистая правда.
Это не легенда и не миф, всё происходило у меня на глазах.
К сожалению, моего друга, от имени которого ведётся повествование, уже нет в живых.
Но я вам скажу --- то, что он дошёл до того финала, до которого дошёл, само по себе стечение обстоятельств с вероятностью, стремящейся к нулю.

Пусть вас не огорчает смерть Аркадиу.
Когда-то он сказал мне, что множество существ, встретившихся на его пути, живут в нём, и для многих это единственная жизнь, которую они имеют.
Аркадиу продолжает жить во мне.
А у меня, надо признаться, странное ощущение: я увижу смерть моей Вселенной.
Я, не самый достойный из её детей, буду наблюдать, как гаснут звёзды и галактики, как испаряются чёрные дыры, и мне даже доведётся сказать милой старушке последнее <<прощай>> --- за миг до того, как тишина сменится молчанием.

Я не знаю, откуда это ощущение.
Но оно жило во мне с того самого момента, как я осознал свою конечность.

Радует одно --- если я и доживу до столь почтенного возраста, то только по своей инициативе и никак иначе.
Читателям же, которые вынуждены уходить из жизни, едва переступив её порог, я скажу --- относитесь к смерти спокойнее.
Это граница жизни, но границы есть у всего, что имеет название.
Наверное, именно поэтому я принял решение дописать творение Аркадиу --- он не закончил историю, оставил своего героя на перепутье, пытающимся найти ответ на важные вопросы своей эпохи.

Я не питаю надежды, что кто-то поймёт мои желания и стремления.
Открывая книгу, раб видит то, что хочет, свободный --- то, к чему готов.
Истину же не дано увидеть никому, даже автору --- книга часто говорит больше, чем он хотел бы сказать, и даже больше того, что он осознаёт.
Я лишь надеюсь, что эти записи будут для мыслящего существа маяком в озерце его короткой, но важной жизни.

\emph{Искренне ваш, Г.}

\chapter{[:] Закат эпохи}

\section{[:] Закат}

\ml{$0$}
{--- Красивый закат.}
{``A beautiful sunset.''}

\ml{$0$}
{--- Ага.}
{``Yeah.''}

Огонь жадно пожирал дымящийся концентрат.
Где-то вдалеке вторила треском обсидиановая глыба, обрамлённая полуостывшей лавой.
Поднималась Четвёртая Луна, похожая на непрожаренную каменную рыбу.
Деловито попыхивал молодой вулкан Арашияма, окутывая рыбу мягкими клубами дыма.
С сейсмологической станции уже сообщили, что скоро Арашияма взорвётся, принеся на значительную часть Преисподней незапланированную пепельную зиму.
Месяц пепла, сыплющегося с тёмных небес, затем ещё три месяца пахнущего серой снега.
И горе тому, кто попытается утолить жажду талой водой...

Вулканическая пустошь навевала тоску.
Жить там было не то чтобы невозможно, а скорее неспокойно.
Старые вулканы на берегах Разлома обильно удобряли и без того плодородную почву;
люди утаптывали ил и водорослевый перегной в любую расщелину, способную их удержать.
Молодые же устраивали неприятные сюрпризы вроде пирокластических потоков, вынуждавших аборигенов заселять только определённые возвышенности, а ещё регулярно затягивали небо тучами, проливавшими пепельный снег, неурожай и голод.
Возле Арашияма люди не селились.
По крайней мере, так полагали до недавнего времени.

Фигура, закутанная в чёрный плащ, нетерпеливо шевелила посохом дымящиеся блоки полимера.
Вторая, грузная и слегка неловкая, лениво ковыряла ножом землю.
Со стороны казалось, что разговор не клеится, но вряд ли это соответствовало истине.

\ml{$0$}
{--- Лу, что по времени?}
{``Lu, timeline?''}

\ml{$0$}
{--- Время есть.}
{``There is time.''}

\ml{$0$}
{И снова молчание.}
{Silence again.}
\ml{$0$}
{Звезда медленно коснулась горизонта.}
{The Star slowly touched the horizon.}

Фигура в чёрном вынула из складок плаща странное приспособление, похожее на металлический веер с раструбом.
Веер заплясал в дыму и зажужжал, то втягивая в себя дым, то выдувая его плотными длинными клубами.
Одно движение, два, три --- и в темнеющие небеса отправилась вязь призрачных знаков.

Звезда, осветив дымовую надпись, села за горизонт, и в небе засияли фиолетово-красные змеи.
Они извивались, трепетали, как флаги на ветру.
Восхищённо крякнул толстяк, бросивший на время свой нож;
его товарищ откинул капюшон, обнаружив тугие белые кудри и гордое молодое лицо.
Оба, не отрываясь, вглядывались в сапфировые небеса, где разворачивался поистине божественный спектакль.
Но длилось это зрелище недолго --- махнув хвостом, самый большой аспид рассыпался в воздухе жемчужным пеплом вслед за последним лучом Звезды.
Наступила ночь.

Улыбающиеся мужчины посмотрели друг на друга.

\ml{$0$}
{--- Никак не могу к этому привыкнуть, Грисвольд, --- смущённо признался молодой.}
{``Can't get used to that, Griswold,'' the young shyly said.}

\ml{$0$}
{--- Да, --- вздохнул толстяк.}
{``Yeah,'' the fat man sighed.}
\ml{$0$}
{--- Такая красота --- и в таком месте.}
{``Such a beauty in such a place.}
\ml{$0$}
{Подумать только.}
{Just imagine.}
\ml{$0$}
{В моём родном мире было всё --- изумрудные леса, сапфировые моря.}
{My home world has anything--- emerald forests, sapphire seas.}
\ml{$0$}
{Но такого не было.}
{But nothing like that.''}

\ml{$0$}
{--- А ты сам откуда?}
{``Where are you from?''}

--- Третья планета системы IF-517, <<Фомальгаут>> по-старому, --- Грис махнул рукой куда-то в сторону горизонта.
--- Меня туда послали ещё первые люди.

\ml{$0$}
{--- Первые люди? --- ахнул Лу.}
{``First humans?'' Lu gasped.}
\ml{$0$}
{--- Сколько ж тебе тогда?..}
{``How old are you, then ... ?''}

Грисвольд горько усмехнулся и вместо ответа снова махнул рукой.

\ml{$0$}
{--- Так ты бог.}
{``So, you're a god.''}

\ml{$0$}
{--- Творец и владыка, --- пошутил толстяк.}
{``Creator and Lord,'' the fat man joked.}
\ml{$0$}
{--- Ладно, хватит болтовни.}
{``Well, I'm done talking.}
\ml{$0$}
{Отдых закончен.}
{Playtime's over.}
\ml{$0$}
{Ждём ответа, твой квадрант --- номер четыре.}
{Wait for an answer, quadrant number four is yours.''}

\ml{$0$}
{--- Грис, а как ты сюда попал?}
{``Gris, tell me, how did you get here?}
\ml{$0$}
{И что у тебя за странное имя?}
{And what kind of name is `Griswold'?}
\ml{$0$}
{Я не узнаю язык.}
{I can't place the language.''}

--- Потом, --- в голосе Грисвольда прозвучала сталь.

Крякнув, Лу сел на камень и принялся посохом ворошить непрогоревший полимер.

\ml{$0$}
{--- Напомни, Грис, зачем мы торчим в этой неуютной обгорелой пустоши?}
{``Remind me, Gris, why are we stuck here, in that unpleasant burned wasteland?''}

\ml{$0$}
{--- Мы ждём сигнала, --- недовольно буркнул толстяк.}
{``We're waiting for a signal,'' the fat man grunted.}
\ml{$0$}
{--- Если повезёт, люди ответят и мы сможем узнать местоположение следующей явки.}
{``If we're lucky, humans will answer with the coordinates of the next safe house.}
\ml{$0$}
{Твой квадрант --- номер четыре.}
{Quadrant number four is yours.''}

\ml{$0$}
{--- И это приблизит нас к обнаружению бункера, --- голос Лу был предельно насыщен сарказмом.}
{``And it gets us one step closer to the bunker,'' Lu's voice was saturated with sarcasm.}

\ml{$0$}
{--- Сайгон, учитывая обстоятельства, я не вполне понимаю, почему ты вызвался идти со мной, --- устало сказал Грис.}
{``Saigon, given the circumstances, I don't quite understand why you volunteered to go with me,'' Gris wearily said.}
\ml{$0$}
{--- Тем не менее, у меня есть приказ, который я намерен выполнить, чтобы сохранить свой ранг.}
{``Nevertheless, I've got an order I intend to carry out, which allows me to keep my rank.}
\ml{$0$}
{Это приказ твоего отца.}
{Your father gave me that order.}
\ml{$0$}
{Арракис сказал --- использовать все имеющиеся способы, чтобы отыскать бункер, и это один из них.}
{Arrakis said: `use every possible way to find that bunker', and what we're doing here is one of those ways.}
\ml{$0$}
{Твой квадрант --- номер четыре.}
{Quadrant number four is yours.''}

Лу недовольно замолчал.

Дррррр... --- странный крик, похожий на птичий, разнёсся по выжженной каменной пустыне.
Друзья даже не изменили позы --- только левые руки почти синхронно вынули пулевые пистолеты.

--- Не двигайся, --- шёпотом предупредил Грисвольд напарника.

Минута прошла в томительном ожидании.

--- Проверим? --- шёпотом предложил Лу.
--- Сидим, как чайки на столбах.
Давно бы сняли, если б захотели.

--- Я посмотрю.
Подержи меня.

Лу аккуратно сел рядом с толстяком и обхватил его за пояс.
Грисвольд на секунду закатил глаза и обмяк.
Парень крякнул под весом навалившейся на него туши.

--- Извини, --- пробормотал толстяк и отряхнулся.
--- Тело Мистраль пустое.

--- Кто это сделал?

--- Не нашёл.

--- Что?
Как не нашёл? --- забывшись, парень заговорил во весь голос.
--- Так это не?..

--- Люцифер, тихо, волки тебя дери!

--- Извини, --- прошептал Лу.
--- Она жива?

--- Не знаю.
Уходим, сейчас же.
Про костёр забудь.

\section{[:] Мефистофель}

Дорога до лагеря лежала по склону Арашияма.
Извержение отгремело только вчера, и от тускло светящейся лавы за пять ярдов разило нестерпимым жаром.
Люцифер тихо, сквозь зубы ругался --- обычный на вид камешек на дороге в мгновение ока проплавил ботинок, выставив наружу, на съедение волнам огня голые пальцы.
Толстяк молча прыгал с одного валуна на другой, выбирая самые холодные места.

\ml{$0$}
{--- Стоять.}
{``Freeze.''}

Голос прозвучал как будто из ниоткуда.
Два путника подобрались, но тут же расслабились.
Свои.

\ml{$0$}
{--- \emph{Плесень и мрак.}}
{``\emph{Darkness and mold.}}
Грисвольд-тринадцатый, Люцифер из клана Мороза, --- прошелестел Грис.

Из почти не заметной с дороги пещерки вышел некто, закутанный в чёрный плащ.
В руке тускло блеснул пистолет.

\ml{$0$}
{--- \emph{Соль и табак.}}
{``\emph{Tobacco and salt.}}
Мефистофель из клана Мороза.
Как проходит эксперимент?

\ml{$0$}
{--- Мистраль исчезла без оповещения.}
{``Mistral just vanished without warning.}
\ml{$0$}
{Мы вынуждены были...}
{``We had to---''}

\ml{$0$}
{--- Мистраль жива, готова к отчёту и к поcледствиям своих действий.}
{``Mistral is alive, ready to report and face the consequences.}
Ваш транспорт я отогнал в безопасное место, заберёте его позже.

\ml{$0$}
{--- Арракис велел тебе встретить нас? --- удивился Грисвольд.}
{``Arrakis told you to meet us, then?'' Griswold wondered.}

\ml{$0$}
{--- Его, --- лаконично ответил Мефистофель, кивнув на Лу.}
{``Him,'' Mephistopheles succinctly answered, nodded at Lu.}

Мгновение спустя из пещеры выкатился маленький экспедиционный мотоцикл.
Мефистофель жестом велел друзьям садиться.
Худенький Лу поморщился, когда его зажало между каменной спиной Мефистофеля и животом Грисвольда, но промолчал.
Мефистофель, первая генерация клана Мороза, был напрочь лишён чувства эмпатии, и просьба потесниться в лучшем случае вызвала бы лишь холодный вопрос <<Зачем?>>.

Минуту спустя мотоцикл нырнул в густую завесу пара, и Арашияма скрылся из виду.

\ml{$0$}
{--- Так что там? --- поинтересовался Грисвольд.}
{``So what is this all about?'' Griswold asked.}

\ml{$0$}
{--- Вторжение минусов, --- бросил Мефистофель.}
{``\textit{Mina} invasion,'' Mephistopheles replied.}

Люцифер вздохнул.
Больше информации не будет --- иначе Мефистофель бы всё рассказал сразу.
Здесь начиналась его работа --- работа стратега.

\ml{$0$}
{--- Мало нам было этих подземных крыс, --- выругался парень.}
{``It was bad enough with those ground rats,'' he said.}
\ml{$0$}
{--- Хорошо ещё, что их технологии не позволяют убить демона...}
{``We should be thankful their technologies are unable to kill a daemon ...''}

\ml{$0$}
{--- Их технологий вполне хватает, чтобы портить нам кровь последние десять лет, --- возразил Грисвольд.}
{``Their technologies are good enough to play on our nerves for the last ten years,'' Grisvold answered.}
\ml{$0$}
{--- Где бункер?}
{``Where is the bunker?}
\ml{$0$}
{Мы так ничего и не узнали.}
{So far we've got nothing.''}

\ml{$0$}
{--- Я сообщу Арракису, что эксперимент был прерван по веской причине и в санкциях нет нужды, --- сказал Мефистофель.}
{``I will tell Arrakis that the trial was interrupted for a reason and there is no need in penalties,'' Mephistopheles said.}

\ml{$0$}
{--- Тысяча благодарностей, ты просто эталон доброты, --- саркастически проворчал Грисвольд.}
{``A thousand thanks, you're absolutely a standard of kindness,'' Griswold grumped sarcastically.}
--- Что нового насчёт диверсантов?

--- Есть данные, что они начали ходить по селениям и обучать дзайку-мару\FM\ распознавать инкарнатов, --- сообщил Мефистофель.
\FL{seijmar}{Сейхмар}
\ml{$0$}
{--- Команда Гало поймала одного агента.}
{``Halo's squad captured an agent.}
Точнее, раздобыла его труп --- он самоуничтожился ещё до начала допроса.

Мефистофель вывел мотоцикл на утоптанную дорогу и выжал полную скорость.
Друзьям пришлось замолчать.

\section{[:] Доверие}

--- И ещё одна неудача, --- Грисвольд тяжело опустился на скамью и забросил на спину полотенце.
%{``One more failure,''}
\ml{$0$}
{--- Третья по счёту.}
{``A third one.''}

\ml{$0$}
{--- Мы с самого начала занимались ерундой. --- буркнул Лу.}
{``We have been screwing around from the beginning,'' Lu said.}

\ml{$0$}
{--- Я тебе уже сказал: я не собираюсь обсуждать приказы твоего отца.}
{``I  told you, I'm not the one to question your father's orders.''}

\ml{$0$}
{--- Ты можешь меня выслушать?}
{``Can you just hear me out?''}

Грисвольд скрестил руки на груди.
<<Почему он так настойчиво пытается вывести меня на откровенность?>>

\ml{$0$}
{--- Люцифер, я не стратег.}
{``Lucifer, I'm not a strategist.}
\ml{$0$}
{Я простой технолог и делаю то, что мне говорят.}
{I'm a humble technologist and do what I'm told.''}

\ml{$0$}
{--- Я ещё раз тебя спрашиваю: ты собираешься выслушать то, что я говорю, или нет?}
{``I ask you again: can you just listen to what I'm about to say, or not?''}

Грисвольд несколько секунд смотрел в голубые глаза Лу и вздохнул.

\ml{$0$}
{--- Валяй.}
{``Go ahead.}
\ml{$0$}
{Что мы должны делать?}
{What should we do?''}

\ml{$0$}
{--- Мы должны подключить этих тама\FM.}
{``We should mobilize that tama\FM\ folks.}
\FL{tama}{Тама}
\ml{$0$}
{Они неплохо показали себя при поиске руды.}
{They proved themselves well in search of ores.''}

\ml{$0$}
{--- Арракис сказал --- людям доверять нельзя, --- ответил Грис.}
{``Arrakis said: humans can't be trusted,'' Gris answered.}
\ml{$0$}
{--- И в этот раз я с ним согласен.}
{``And this time I have to agree.''}

\ml{$0$}
{--- При чём здесь доверие?}
{``How trust is relevant?}
\ml{$0$}
{Нам с Гало нужны любые доступные данные.}
{Me and Halo need all data we can get.}
\ml{$0$}
{Верифицируем их мы, а не люди.}
{Verification is our job, not a humans'.''}

--- Тем, в бункере, тоже нужны данные.
И если все получат то, что хотят, то ситуация никуда не сдвинется.

Лу явно хотел возразить, однако, сделав выразительное движение бровями, промолчал.
\ml{$0$}
{<<Это ошибочное суждение, но и у тебя, и у отца недостаточно знаний, чтобы понять его ошибочность>>.}
{\textit{It's an error of judgement, but not father, nor you have enough knowledge to understand that.}}

\ml{$0$}
{--- Отец не хочет даже попробовать.}
{``Father refused even to try.}
Я предлагал ему разбить задачу на части, не вызывающие подозрения, и дать людям только их.

--- Лу, ты у нас глава Ордена?

--- Отец сделал меня и Гало стратегами.
Зачем мы ему нужны, если в итоге решение принимает он, и далеко не всегда лучшее?

\ml{$0$}
{--- Это и есть руководство.}
{``That's the essence of leadership.}
\ml{$0$}
{Учесть мнение каждой стороны и принять решение, обязательное для исполнения.}
{Take every opinion into account, then make a decision that everyone is to abide by.''}

\ml{$0$}
{--- Тогда он плохой руководитель.}
{``He's a weak leader, then.''}

\ml{$0$}
{--- Не вздумай ляпнуть такое при легионе.}
{``Wouldn't be clever to say that in front of legion.''}

\ml{$0$}
{--- Ты меня за идиота считаешь?}
{``You take me for an idiot?''}

--- Сайгон, нас очень мало, --- тихо сказал Грис, оглянувшись по сторонам.
--- Куда бы ни вёл нас Арракис, он ведёт нас всех.
Всех, понимаешь?
Делить силы --- худший расклад из всех.

\ml{$0$}
{--- Не надо пересказывать мне эти пропагандистские формулы.}
{``Stop repeating all that propaganda formulae to me.}
\ml{$0$}
{Я мог бы курировать людей сам, не подключая ни одной единицы из сил Ордена!}
{I could oversee that humans by myself, without mobilization of one single Order unit!''}

\ml{$0$}
{--- Тогда что тебе мешает? --- раздражённо бросил Грисвольд и отвернулся.}
{``So what's the hold-up?'' Griswold said with irritation, then turned his back.}

Лу изумлённо смотрел в спину толстяка.

\ml{$0$}
{--- Хватит прожигать дыру в моей спине, --- пробурчал технолог.}
{``Stop burning holes in my back,'' the technologist muttered.}
\ml{$0$}
{--- Хочешь --- делай.}
{``If you want it, do it.}
\ml{$0$}
{Я молчок.}
{I'll be quiet.''}

\section{[:] Переговоры с наблюдателями}

\ml{$0$}
{--- И тем не менее Тысяча Башен не совсем понимает, в чём сложность борьбы с восставшими людьми.}
{``...Nevertheless, Thousand Towers doesn't quite understand what's the problem with bringing human rebels into line.''}

\ml{$0$}
{--- Основная проблема в том, что этот бункер специально разрабатывался для борьбы с демонами, --- объяснил Лу.}
{``The main problem is that bunker was designed specifically to fight against daemons,'' Lu explained.}
--- Как вы знаете, Преисподняя не всегда была выжженной вулканической планетой, когда-то здесь было спокойно и даже зелено.
Когда глава нашего клана, Арракис, прибыл с Мороза вместе с Дорге, здесь шла война между колонистами с Земли и демиургом.
Люди использовали всю мощь имеющихся технологий, чтобы лишить демиурга масс-энергии, и заставили его уйти.

\ml{$0$}
{--- Мы слышали о том, что Преисподняя была выжжена демиургом в отместку, --- полувопросительно сказал наблюдатель.}
{``We heard rumors that Nether was scorched by its demiurge as a revenge,'' the observer said.}

\ml{$0$}
{--- Едва ли, --- подал голос Грисвольд.}
{``Hardly,'' Griswold said.}
\ml{$0$}
{--- Для демиурга причинить вред планете --- всё равно что отрубить себе руку.}
{``For a demiurge, harming their planet is like severing their own hand.''}

Наблюдатель непонимающе уставился на технолога.

\ml{$0$}
{--- В общем, для большинства это неприемлемо, --- поправился Грисвольд, оглядев бионические конечности собеседника.}
{``I mean, it's unacceptable for most of them,'' Griswold corrected himself, noticed bionicle limbs of the observer.}
--- Наши учёные считают, что это сделали люди.
Но лично я, как имевший опыт управления планетой, склоняюсь к тому, что биолитосфера не была стабилизирована и без контроля демиурга развалилась сама.

--- И вот через тысячу лет, когда люди уже растеряли большую часть технологий, а мы взяли власть, бункер был снова кем-то найден, --- сообщил Самаэл.
--- Видимо, там же были найдены старые записи, которые и объяснили людям, с кем они имеют дело.

--- Мы даже не уверены, что бункер в окрестностях Арашияма один-единственный, --- продолжил Гало.
--- Как не можем быть уверены, что таких сооружений нет где-либо ещё на планете.
Методология первых людей была отточена до совершенства, и современные дикари следуют ей неукоснительно, словно священным догмам.
Именно поэтому нам нужны ваши демоны.

\ml{$0$}
{--- Разумеется, --- улыбнулся наблюдатель.}
{``Of course,'' the observer smiled.}
--- Мы прекрасно понимаем, что эти сооружения должны быть найдены и изучены со всей возможной тщательностью, ведь это --- залог победы над Союзом.
Орден Тысячи Башен был бы чрезвычайно рад помочь вам и получить данные касательно технологий первых людей.

Айну коротко поклонилась.

--- Однако, к сожалению, у нас нет свободных демонов для выполнения этой задачи, --- закончил наблюдатель.
--- Тот резерв, который был подготовлен к работе с вами, пришлось срочно задействовать на Тысяче Башен.

--- То есть вы отказываетесь нам помогать? --- осведомился Люцифер.

--- Юноша, --- начал наблюдатель, --- вы превратно истолковали наши слова.

--- Во-первых, не <<юноша>>, а <<сайгон Люцифер>>, --- нахмурился Гало.

--- Прошу прощения, сайгон Люцифер.
Я только хотел заострить внимание на том, что <<отказываемся помогать>> и <<не имеем свободных демонов>> --- разные вещи.

\ml{$0$}
{--- Да что вы говорите, --- буркнул Грисвольд.}
{``How could we confuse one and the other,'' Griswold said.}

--- По вашим указаниям мы приготовили вам всё, включая тела, а вы говорите, что не имеете свободных демонов? --- вспылила Айну.
--- Чем это занят Орден Тысячи Башен, что свободных демонов нет?

--- Вы можете получить любую информацию о наших занятиях по официальному запросу, --- сухо ответил наблюдатель.
--- Я лишь сообщаю решение нашего штаба.

<<Ага, --- хмыкнул Гало.
--- Смиренный служитель, как же.
Мы как-то протестировали троих на полномочия.
Его полномочия соответствуют как минимум нашим с Лу, а молчун с краю, чьё имя я даже не припомню, вполне может быть вторым лицом в Ордене>>.

<<И сейчас, когда тайное убежище стало полем битвы, эти корольки решили аккуратно сбежать, провоцируя нас на высылку наблюдателей, --- подтвердил Лу.
--- Действуйте как угодно жёстко по своему усмотрению, но ни в коем случае не вздумайте их отпускать>>.

Айну и Самаэл кивнули.

--- Мы обязательно сделаем официальный запрос, --- сухо сказала Айну.
--- Я прошу прощения за последние слова, вы --- отличные исполнители.
И тем не менее, как я полагаю, у вас должны быть готовые объяснения отдельных аспектов.

--- Задавайте вопросы, и я отвечу, --- неприятно улыбнулся наблюдатель.

--- Были ли нападения со стороны Союза на Тысяче Башен? --- начал Самаэл.

--- Они и сейчас есть.
Не настолько серьёзные, как здесь, но тем не менее оттягивающие значительную долю наших сил.
Вот, пожалуйте.
Полный расклад...

<<Ложь, --- ухмыльнулся Лу, просмотрев данные.
--- Нападения они пресекли в корне, так как имели сношения с Союзом втайне от нас и получали разведданные от завербованных минусов>>.

<<Выкладка точна?>> --- спросила Айну.

<<Подтверждаю.
Статистика по потерям, --- кивнул Гало.
--- Они их завысили, но тем не менее иначе как грамотно и, главное, вовремя собранными разведданными такое не объяснить>>.

<<В общий котёл отправьте информацию по руководителям их разведки, --- добавил Арракис.
--- Имена, биографии.
Мне совсем не нравится, что они на шаг впереди>>.

<<А также по тому, насколько свободны они в своих действиях>>, --- вставил Лу.

Арракис кивнул.
Если он и понял намёк сына, то предпочёл этого не показывать.

--- Были ли попытки установить диалог с Союзом?

--- Были.
Берен, предоставь записи.
Как видите, мы даже попытались заключить мир, но уже во время подписания пакта последовала серия диверсий.

<< ... на Преисподней, --- закончил Гало.
--- Настолько виртуозная ложь, что неотличима от правды.
Они каким-то образом откупились от Союза, и я бы многое отдал, чтобы узнать, каким>>.

<<И оба раза они рассчитывали на то, что ложь будет раскрыта, --- добавил Лу.
--- Цель та же --- высылка наблюдателей>>.

<<О как, --- хмыкнула Айну.
--- То есть у нас ещё один возможный враг, который так и нарывается на кинжал>>.

<<Не совсем>>, --- возразил Гало.

<<Скорее это переход от сотрудничества к вооружённому нейтралитету, --- пояснил Лу.
--- Возможно, это и было предметом договора с Союзом: Орден Тысячи Башен не будет вмешиваться в конфликт и примкнёт к победителю>>.

<<Дьявол, мы не можем даже отозвать наших демонов с Тысячи Башен, --- посетовала Айну.
--- Потому что рычаг давления на Орден ну никак не лишний>>.

<<И они об этом знают, --- закончил Гало.
--- Вполне возможно, что они даже рассчитывают на остатки адских сил после нашего поражения.
Больше-то легионерам некуда будет идти>>.

<<Лично мне всё ясно, --- буркнул Арракис.
--- Самаэл, задай им для приличия ещё пять-десять умных вопросов, сохраняя дружественно-требовательный тон.
Люцифер, создавай конференцию прямо здесь --- нам срочно нужен план.
И ещё --- пусти параллельно слабозашифрованную меченую легенду.
Посмотрим, где у нас протечка>>.

\section{[:] Условия и ответ}

\textspace

Доклады поступали один за другим --- дезертировал третий отряд, затем восьмой.
У Ордена осталось меньше трети легиона, и, судя по настроениям, скоро не будет даже её.
Дзайку-мару, словно зная о происходящем, отказывались подчиняться представителям.
Вскоре поступило сообщение, что забастовку объявили носильщики.

Айну, как обычно, сорвалась с места в карьер.

--- Грисвольд, Люцифер, Мистраль, Гало.
Вы идёте со мной.
База в Такэсако разрушена.
Нужно найти и унести уцелевшее оборудование.

--- Прости, что значит <<унести>>? --- поинтересовался Гало.

\ml{$0$}
{--- Это значит, что ты возьмёшь всё барахло руками, взвалишь себе на плечи и ножками дойдёшь сюда, в лагерь.}
{``It means: all the garbage we can find you'll pick up by your hands, then carry on your back, then return here by your feet.''}

--- То есть узел дозаправки транспорта в Акияма также разрушен? --- уточнил Лу.

Айну промолчала.

--- Я стратег, а не носильщик, --- тихо сказал Люцифер.

\ml{$0$}
{--- Если не будешь носильщиком добровольно --- станешь трупом принудительно, --- посулилась Айну.}
{``If you're not a carrier by choice, you'll be a cadaver by force,'' Ain\"{u} promised.}
\ml{$0$}
{--- Союз Воронёной Стали об этом позаботится.}
{``Blued Steel Union gonna take care of that.''}

\ml{$0$}
{--- Я смотрю, дела у Ордена Преисподней совсем плохи, --- масляно улыбнулся один из наблюдателей.}
{``Things are going bad for Order of Nether, as I can see,'' an observer said with a smile of predator.}

\ml{$0$}
{--- Поблагодари меня, что пока ещё можешь смотреть, --- сказала Айну.}
{``Thank me, you still can see by now,'' Ain\"{u} answered.}

\ml{$0$}
{--- Вы же понимаете, что у ваших слов и поступков могут быть последствия? --- осведомился наблюдатель.}
{``I guess you understand that your words and your actions would have consequences?'' the observer asked coldly.}

--- Если ты сейчас же не закроешь рот, это поймут все и сразу.
До мельчайших подробностей, --- демон Айну предупредительно вспыхнул, приведя в готовность боевые модули.
Наблюдатель побледнел и вытянулся в струнку.

\ml{$0$}
{--- Если Орден Преисподней переживёт эту схватку, Тысяче Башен придётся понять и принять очень многое.}
{``If the Order of Nether outlives this battle, Thousand Towers will have many things to understand and accept.}
\ml{$0$}
{Или ответить, --- пояснил Люцифер, поняв, куда дует ветер.}
{Or payback,'' Lucifer explained---as soon as he's seen the way the wind blows.}

На этот раз побледнели все наблюдатели.
Айну одобрительно посмотрела на юношу.

\ml{$0$}
{<<Впервые слышу от тебя угрозу.}
{\emph{``The first threat of you I've ever heard.}}
\ml{$0$}
{Впечатляет>>.}
{\emph{Very impressive.''}}

\ml{$0$}
{<<Ты просто никогда не брала без спроса его косметику, --- ухмыльнулся Гало.}
{\emph{``You never borrowed his makeup without permission,''} Halo smirked.}
\ml{$0$}
{--- Я чуть не поседел раньше времени>>.}
{\emph{``My hair almost grayed before their time.''}}

--- Я правильно понял, что мы ваши пленники? --- вызывающе спросил наблюдатель, говоривший первым.

--- Вы --- куски мяса, пока ваша ценность не подтверждена Орденом Тысячи Башен, --- без обиняков ответила Айну.
--- Если подтверждения не последует, мы поступим с вами как с мясом.

Присутствующие легионеры как по команде взяли наблюдателей на прицел.
Айну кивнула Люциферу.

\ml{$0$}
{--- Малышня, я вас жду через минуту.}
{``Midgets, I want you ready in one minute.}
\ml{$0$}
{Из оружия возьмите только пистолеты.}
{Take only handguns, and no weapon more.}
\ml{$0$}
{Если нас там встретят --- оно в любом случае не поможет.}
{If we are met there, weapon will be useless anyway.''}

\section{[:] Тела на обочине}

\textspace

Люцифер поморщился, Айну подняла бровь, а Грисвольд почувствовал, что у него задрожали ноги.
Столбы вдоль дороги были увешаны человеческими телами без кожи.

--- Ну и вонь, --- сказал Люцифер и, достав респиратор, манерно прикрыл им нос и рот.
\ml{$0$}
{--- Кто это сделал?}
{``Who have done that?''}

\ml{$0$}
{--- Точно не мы, --- констатировал Гало.}
{``We've not,'' Halo stated.}
\ml{$0$}
{--- И не Союз.}
{``Union neither.}
Демоны не стали бы тратить время на такие глупости.
\ml{$0$}
{Кожу сняли уже после смерти.}
{They were skinned after death.''}

\ml{$0$}
{--- Проведи анализ, брат, --- попросил Люцифер.}
{``Analyze it, brother,'' Lucifer asked.}

Гало подошёл к трупу и внимательно осмотрел его со всех сторон, затем обнюхал подсушенную горячим ветром плоть.
Демон стратега издал слабое свечение.

--- Точность 0.904, --- сказал Гало.
\ml{$0$}
{--- Это сделали дзайку-мару.}
{``Dzaiku-maru have done it.}
\ml{$0$}
{Убивали и снимали кожу методично, без торопливости.}
{They were killing and skinning methodically and slowly.}
\ml{$0$}
{Многие убиты во сне.}
{Many of victims were killed in their sleep.}
Я думаю, что это вендетта, хотя впервые вижу, чтобы с трупами врагов так поступали.

--- Что могло спровоцировать вендетту? --- удивилась Айну.

--- Союз использовал сапиентные тела, --- немного невнятно предположил Люцифер из-за респиратора.
--- Маловероятно, что они следовали всем местным обычаям и традициям.
Дзайку-мару не видят демона, они видят тело и мстят телу.

--- Я согласен с Лу, --- вмешался Грисвольд.
--- Такие казусы регулярно случались на моей планете, если туда забредал очередной нейтрал-неумеха.
Чаще всего это заканчивалось смертью его тела, но обычно он успевал прихватить десяток-другой нападавших.

Друзья переглянулись.

<<Это твои тама?>> --- иронически спросил Грис.

<<Своё дело они сделали, --- ответил Лу.
--- Данные у меня, планирую подсунуть их в общее поле как случайную находку>>.

<<Главное, что ты ни о чём не жалеешь>>.
Лу почувствовал в словах технолога холодок.

--- Кстати, --- добавил Гало веселее, чуть обычно.
Все навострили уши.
--- Что вы думаете насчёт погоды?
Я чувствую слабый привкус цветочной влаги.

<<Есть вероятность, что за нами идёт лазутчик --- неуточнённое сапиентное тело.
Местоположение по эманациям определить не могу, но чувствую иные неуточнённые признаки присутствия>>.

--- Ты оптимист, Гало, --- рассмеялась Айну.
--- Погода не изменится.

<<Вести себя как обычно, следовать тем же курсом>>.

Айну кивком подтвердила приказ, и демоны, в последний раз посмотрев на страшную картину, двинулись дальше.
Скрипели и ворчали под ветром сохнущие верёвки.
Трупы махали вслед демонам мёртвыми руками, как махали всем прочим путникам, имевшим несчастье пройти мимо.

\section{[:] Мальчик с мечом}

\epigraph
{Счастье --- это когда из кошмара можно проснуться.}
{Пословица сели}

Через десять миль начало сказываться недосыпание.
Демоны изо всех сил пытались заставить свои тела идти, как обычно, но вскоре перестали помогать даже психостимуляторы.
Грисвольд вздыхал, Гало злобно пинал мелкие камни, Айну мрачно молчала.
Люцифер, который передал свой груз технологу, выглядел бодрее всех.

--- Чтоб эти тела, --- бормотал Гало.
--- Почему дзайку-мару не избавились от такой бесполезной вещи, как сон?

--- Во время сна мозг восстанавливает метаболический потенциал, --- заметил Люцифер.
--- Сон --- это песочница, в которой человек без вреда для себя может подготовиться к травмирующим ситуациям.
А ещё сон --- это дар видеть желанный мир.
Или хотя бы не видеть нежеланный.

--- Ещё немного, и я брошу эту тушу на дороге, --- пообещал Гало.
--- Я ей не пастух.

Лазутчик не показывался.
Айну пару раз наудачу просмотрела местность, но неведомый сапиент каким-то образом прятался от <<взгляда>> интерфектора.
Грисвольд предположил, что он искусственно ввёл себя в ступорозное состояние.

Тропинка круто свернула к северу, огибая скалистый холм.
Вскоре Гало, поколебавшись, попросил Грисвольда понести и часть его груза.
Технолог тяжело вздохнул и согласился.
Айну хотела сделать братьям строгое внушение, но Грисвольд покачал головой.
<<Не до этого>>.

Ветер переменился, и смешанный с дымом пар унесло в сторону.
В надвигающихся сумерках резко проступили очертания холма.

--- Смотрите туда, --- Люцифер, не оглядываясь, махнул на скалу.
--- Вон он, наш лазутчик.

Сонная дымка мгновенно улетучилась.
Демоны, остановившись, как будто невзначай встали в боевой порядок.

На скале стоял мальчишка.
Он опирался на дайту-соро\FM{} в ножнах, словно на посох.
\FA{
Дайту-соро --- ритуальное оружие Древней Земли, длинный изогнутый меч из стабитаниума.
Народ нихон, отправляя своих людей в космические путешествия, обязательно давал им дайту-соро как талисман.
}
Рваные лохмотья, едва прикрывавшие тело, развевались на сухом горячем ветру.
Мелкие острые камешки то и дело секли плоское смуглое личико, но мальчишка стоял и продолжал смотреть.
Ему было всего шесть-семь стандартных лет, но он держался, словно завоеватель на пьедестале.

--- Шпион, как пить дать, --- констатировала Айну, подтягивая ремешки куртки.
--- Гало, подпусти его поближе и потяни время, я устала и не хочу бегать по этим скалам.

--- Эй! --- закричал Гало, и эхо многократно повторило его крик.
--- Мальчик!
Подойди!

Мальчишка, поколебавшись, перехватил поудобнее оружие и начал спускаться.

--- Дзайку-мару, --- доложила Айну.
--- Как подойдёт на пять шагов, стреляйте в голову, кому там сподручнее, и пошли дальше.

Мальчишка спрыгнул с небольшого уступа и замер в шести шагах.
Айну достала пистолет, Люцифер и Гало замерли с руками на кобурах.

--- И что это такое? --- осведомилась Айну, заметив смятение товарищей.

--- Ты ж сказала --- в пяти шагах, --- буркнул Гало.

--- Я сказала --- кому сподручнее, --- заплетающимся языком поправила Айну.
--- Вы идиоты.
Может, вам расстояние сообщать в виде диапазона допустимых значений?
С точностью до микрометра?

--- Пристрели мальчишку, чтоб тебя, --- устало бросил Грисвольд.
--- Все хотят спать, а мы уже две минуты здесь торчим.

Пока демоны разговаривали, маленький оборванец подбежал к Айну и обнял её ноги.
Айну приставила пистолет к его виску.
Мальчишка поднял улыбающееся лицо, его огромные карие глаза взглянули на демоницу.

--- Ты красивая, ситу-ну-онна\FM, --- сказал мальчик.
\FA{
Ситу-ну-онна --- на Преисподней: женщина, владеющая оружием и боевыми искусствами.
}
--- Когда я вырасту, ты будешь моей женой.

--- Хорошо, что не сейчас.
Хоть отдохну ночью, --- усмехнулась Айну и как ни в чём не бывало спрятала пистолет.
--- Что ты здесь делаешь, ребёнок?
Шпионишь за нами?

--- Да, --- признался мальчик.
--- Вы похожи на великих воинов из сказок.
Я следил за вами целых три часа.

--- Айну, --- буркнул Гало и вытащил оружие.
Айну жестом остановила его, демон нехотя подчинился.

<<Если забрызгаешь мне комбинезон, стирать заставлю вручную>>, --- добавила демоница.

--- Как ты скрылся от нас? --- поинтересовался Грисвольд.

Мальчишка бросил перед демонами автоматический шприц с остатками препарата.
Глаза Гало расширились, он уронил рюкзаки и подхватил шприц.

<<Нейролептик, --- передал он остальным.
--- Скорее всего, ксидатин низкой степени очистки>>.

Только сейчас Айну заметила в улыбке мальчика что-то искусственное.
Глаза оставались бесстрастными и отстранёнными.
Когда мальчишка вытянул руку, чересчур короткий рукав изорванной куртки соскользнул, и демоны увидели корявые детские иероглифы, вырезанные ножом на коже запястья:

\ml{$0$}
{<<ИДИ ЗА НИМИ>>.}
{\textsc{Follow them.}}

Грейсвольд вздрогнул.
Ксидатин напрочь подавлял мотивацию.
Людям требовались годы, чтобы научиться действовать в ступоре, в условиях минимальной эмоциональной стимуляции.
Для этого нужно было чёткое понимание, что и зачем ты делаешь.
Но ребёнок?..
Грейсвольд мысленно подивился силе его воли.
Он знал, что людям труднее всего идти против собственных эмоций.
Или против их отсутствия.

--- Ты не желаешь нам зла, теперь я вижу, --- сказала Айну, заворожённо разглядывая кровавую надпись на детской руке.
--- Мы не можем долго с тобой беседовать.
Возвращайся домой.

Ножны на поясе Айну расстегнулись, клинок выдвинулся на дюйм.

--- Я пойду с вами, --- без раздумий выпалил мальчишка.

--- Какой смелый.
В таком случае --- вперёд, --- скомандовала Айну, безо всяких фокусов убрав клинок и поправив заплечный мешок.

<<Ты с ума сошла?>> --- удивился Грисвольд.

<<Да хватит уже, --- отозвалась демоница.
--- Убивать надо сразу, а мы театр устроили>>.

\ml{$0$}
{--- Зачем ты носишь с собой эту бесполезную палку? --- спросила Айну.}
{``Why do you carry this useless stick?'' Ain\"{u} asked.}
\ml{$0$}
{--- Брось.}
{``Drop it.''}

\ml{$0$}
{--- Нет, --- просто ответил мальчик.}
{``Nope,'' the boy answered simply.}
\ml{$0$}
{--- Это фамильная ценность, и он старше вас всех вместе взятых.}
{``It's a family heirloom, and in addition it's older than all of you combined.}
\ml{$0$}
{Проявите уважение.}
{Show respect.''}

\ml{$0$}
{--- Дело твоё, --- так же просто согласилась Айну, отметив своё полное нежелание ему отказывать.}
{``Suit yourself,'' Ain\"{u} agreed in the same manner, thinking of her own reluctance to deny.}
\ml{$0$}
{--- Лу, ты за него отвечаешь.}
{``Lu, you're in charge of him.''}

Люцифер хмуро посмотрел на демоницу, потом на неожиданного подопечного.
Тот, недолго думая, перебросил ремешок дайту-соро за плечи, подошёл к демону и бесцеремонно запрыгнул ему на закорки.
Хрупкий Лу охнул и едва не оступился.

--- Вези меня, --- оскалился мальчишка.

--- Когда придём, я высплюсь, а утром лично тебя прирежу, --- полузадушенным голосом посулился Люцифер.
\ml{$0$}
{--- Следи за своей уважаемой железкой, она меня по заднице бьёт.}
{``Watch your respectable piece of iron, it's been hitting my ass.''}

\ml{$0$}
{Впрочем, угроза была пустой.}
{This threat, however, was empty.}
Мальчик успел заинтересовать Люцифера так же, как и очаровать Айну.

Отряд двинулся дальше.
Айну шла молча.
Люцифер, угрюмо везущий на плечах мальчишку, тоже молчал.
Грисвольд уже не дышал, а пыхтел под тяжестью рюкзаков.

Мальчик начал выходить из ступора.
Грисвольд несколько раз бросал взгляд на его лицо --- в нём оставалось всё меньше искусственности.
Бродяжка начал проявлять неподдельный интерес к спутникам и грузу, который они несли.

--- Зачем ты носишь ножницы в ножнах? --- спросил он Лу.

--- Потому что они почти то же самое, что нож, но ими сложнее порезаться, --- буркнул Лу.

--- Это выглядит смешно!

--- Смешно выглядит твоя палка.
Ею даже людей сложно резать, хотя разрабатывалась она именно для этого!

Вскоре инициативу в разговоре взял Гало.

\ml{$0$}
{--- Где ты научился так притворяться?}
{``Where did you learn to pretend like that?''}

\ml{$0$}
{--- Это называется <<быть вежливым>>, --- просветил демона мальчик.}
{``That's called `to be polite','' the boy enlightened the daemon.}
\ml{$0$}
{--- В нашем доме бывали разные люди, но мать учила, что хозяин должен быть вежливым со всеми гостями.}
{``Our home was visited by all sorts of people, but mother taught me that a master should be polite to all guests.''}

\ml{$0$}
{--- Ты считаешь себя хозяином?}
{``You call yourself \emph{a master}?''}

\ml{$0$}
{--- Вы определённо прибыли издалека, --- улыбнулся бродяга.}
{``You've certainly come from far away,'' the tramp smiled.}
\ml{$0$}
{--- А я здесь живу с рождения.}
{``Unlike me, who was born here.''}

--- Ты знаешь правила знакомства, дзайку-мару?

Мальчишка знал, что словом <<дзайку-мару>> бродячие торговцы называли всякую мелочь для ремёсел --- заклёпки, полоски ткани, нитки, --- но не повёл бровью.

--- Знаю, \emph{хорохито}, --- сказал он.
--- При первой встрече люди, если они не враги, кланяются, называют свой род, родовые девизы и имена.

--- И почему ты не следуешь этим правилам? --- ядовито осведомился Гало.
--- Скажи своё имя.

\ml{$0$}
{--- Поклонись, --- парировал мальчишка.}
{``Take a bow,'' boy retorted.}
Его голос был непринуждённым, но чувствительный к нюансам интонации слух демонов заставил всех изумлённо посмотреть на маленького спутника.
Гало зашипел.
Он тоже услышал оскорбительный ультиматум.

--- Если бы ты не висел на шее моего брата, я бы нафаршировал тебя пулями.

Айну расхохоталась.
Грисвольд хихикнул и тут же закашлялся.

\ml{$0$}
{--- Это что-то, --- тяжело дыша, признал технолог.}
{``Amazing,'' Griswold declared gasping for breath.}
\ml{$0$}
{--- В такую глупую историю я ещё не попадал.}
{``I’ve never got into such a stupid situation.}
\ml{$0$}
{Хорошо, доиграем сценку до конца.}
{So, let’s act this play to the end.''}

Грисвольд приблизился к Люциферу и поклонился, насколько позволяли быстрый шаг и тяжёлый груз.

\ml{$0$}
{--- Я Грисвольд из рода... эээ... Грисвольда.}
{``I’m Griswold, belong to ... hmm ... Griswold house.}
\ml{$0$}
{<<Добрый Грис --- лентяя приз>>.}
{\emph{We're a prize for lazy guys.}''}

\ml{$0$}
{Мальчик церемонно кивнул толстяку.}
{The boy ceremoniously nodded his head to the fat man.}

\ml{$0$}
{--- Я из рода Ханаяма.}
{``I belong to Hanayama house,'' the boy told.}
\ml{$0$}
{<<Благородная кожа --- лучшее одеяние>>.}
{\emph{``Noble skin is the best suit.''}}

\ml{$0$}
{--- Случаем, не твоих родичей освежевали во время недавней вендетты? --- осклабился Гало.}
{``Quite by chance, aren't that your family who have been skinned at the recent vendetta?'' Halo grinned.}
\ml{$0$}
{--- Кто бы это ни сделал, они не лишены чувства юмора.}
{``Whoever have done it, they are not deprived of sense of humour.''}

--- Я отвечаю только за свою кожу, --- заметил мальчик.
--- И она пока на мне.
Грисвольд мог бы передать часть груза мальчику.
Я понесу один рюкзак, а длинноволосый понесёт меня.

\ml{$0$}
{--- Иди ты в дупло с такими идеями, --- пробормотал Люцифер.}
{``Go to hollow with such ideas,'' Lucifer muttered.}

Айну обогнала Люцифера и тоже поклонилась мальчику.
Бродяга поднял огромные живые глаза на женщину и улыбнулся белозубой улыбкой.

--- Тахиро.

Айну открыла рот, чтобы назвать своё имя, но Люцифер не выдержал:

--- Может, вы прекратите переговоры на моём горбу?

\section{[:] Сладкое безделье}

\epigraph
{Сотрудничество с сильным соперником --- хороший способ отступления.}
{Пословица Преисподней}

Люцифер сладко потянулся в своей постели.
Вчерашний план увенчался полным успехом.
<<Разбежавшиеся>> войска Ордена застали врага врасплох.
На указанный путь, по которому прошли два стратега, технолог и интерфектор, слетелись аж одиннадцать диверсионных групп.
Все они попали в засаду и были уничтожены.

<<Как же я люблю ничего не делать, --- лениво думал Люцифер, смакуя утреннюю негу.
--- А вот быть приманкой --- не очень>>.

Люцифер обладал потрясающим умением наслаждаться отдыхом, зная, что нежиться ему оставалось считанные минуты.
График у демонов по-прежнему был плотный.

Лу перевернулся на другой бок и закрыл глаза.
Вскоре его вывел из забытья лёгкий спазм в мышцах --- забавный сигнал, говорящий о готовности мозга погрузиться в сон.
Лу уже даже засунул под прохладную подушку разгорячённую руку и настроился на продолжение отдыха, как зажужжал браслет на левой руке.

<<Как обычно>>, --- философски заключил Люцифер и, сев на кровати, начал с наслаждением тереть глаза, зевать и потягиваться.
Со стороны это выглядело вопиющим бездельем --- если не знать, что стратег специально выделял для потягиваний, зевания и утренней мастурбации восемь минут пятьдесят секунд драгоценного времени.
На личную гигиену всем выделялась двадцать одна минута, но специально для Лу Арракис сделал исключение;
он мог заниматься своим телом аж тридцать пять минут и двадцать секунд --- с учётом душа, антибактериальной обработки тела, нанесения ухаживающих средств и небольшого количества макияжа.
Исключительное право Лу выбивал потом и кровью.
Ему пришлось собрать внушительную статистику, чтобы доказать --- при ухоженном теле его демон работает гораздо эффективнее.
И даже это бы не помогло, если бы на сторону парня совершенно неожиданно не встала Айну.

--- Идите вы в дупло со своей дисциплиной, --- заявила она Гало и Арракису.
--- Вам же ясно показали, что результат лучше.

--- Сегодня Лу, а завтра --- все легионеры начнут минуты требовать? --- возмутился Гало.

--- Я им всем лично надушу подмышки и накрашу глаза, если это повысит эффективность, --- рявкнула Айну, выйдя из себя.
--- В общем, или Лу получает свои минуты, или ищите другого императора\FM.
\FA{
Император --- ранг и должность в раннем Ордене Преиподней, соответствующие максиму терция и командующему планетарными вооружёнными силами соответственно.
}
Мне нет резона работать с демонами, ставящими методы превыше результата.

Ещё одна хорошая новость поступила за завтраком.
Наблюдателей, сливших информацию Союзу, вывернули наизнанку.
Тысяча Башен безоговорочно приняла условия Преисподней --- на вулканическую планету прибыли две трети их высшего командования и почти все учёные, аккурат в приготовленные для интерфекторов тела.
Фактически это означало объединение двух организаций в одну.
Наступило кратковременное локальное равенство сил.
Следовало заняться восставшими дзайку-мару.

Люцифер уже начал просматривать полученные данные и прикидывать, как ставить новые эксперименты, когда его остановил Гало:

--- Ты чего, брат?

--- Нужно провести ещё серию экспериментов, чтобы...

--- Зачем? --- Гало ухмыльнулся и показал брату автоматический шприц Тахиро.
--- Спускайся в подвалы, Айну сейчас приведёт твоего сопляка.
Он знает, где бункер.

\section{[:] Провели}

--- Я вам ничего не скажу, --- сразу заявил Тахиро.
--- Ничего.

Мальчик выглядел не лучше, чем в день, когда отряд Айну нашёл его.
На руках и ногах пестрели ссадины, в разбитом носу засохла кровь.
Руки и ноги были наскоро обездвижены кабельной стяжкой, на шее красовался ортопедический воротник, из углов рта торчал кляп-гантель, не позволявший крепко сжать зубы.
Левый рукав рубахи был мокрым от слюны;
видимо, мальчик лежал связанным несколько часов.

--- Зачем ты его связала? --- спросил Лу и попытался снять гантель.
Айну остановила его руку на полпути.

--- Это не я, --- ответила Айну.
--- По словам стражи, он попытался причинить себе вред.

--- Некоторые ломают себе кости и осколками режут крупные сосуды, --- сообщил дежурный легионер и, бесцеремонно засунув пальцы Тахиро в рот, обнажил зубы.
--- Те немногие, кому повезло остаться с функциональным набором зубов, используют их --- например, чтобы откусить язык.
Некоторые агенты специально затачивают один резец перед делом --- именно на этот случай, вполне хватает, чтобы вскрыть вены.
У этого пацанёнка, как видите, зубы в порядке --- видать, из хорошей семьи, неплохо питается и...

--- И ты свободен, легионер, --- чуть резче, чем обычно, сказал Грисвольд.

--- Как скажете, сама, --- дежурный коротко поклонился и вышел.

--- Я ничего не скажу, --- повторил Тахиро.

--- Если я буду тебя пытать, ты расскажешь всё, --- без обиняков сказала Айну.
--- И совершить самоубийство у тебя не получится, даже не надейся.
Думаю, после твоей жалкой попытки втереться в доверие ты уже догадался, что мы умнее твоих соплеменников.

Мальчик понурился.
Гало грубо вздёрнул его в сидячее положение и тут же, посмотрев на мокрые от слюны руки, начал вытирать их об одежду.
На его лице застыло отвращение.

--- Видишь те кандалы? --- Айну показала на стену.

Тахиро кивнул.

--- Знаешь, что там недавно произошло?

Тахиро помотал головой.
Вернее, пошевелил, насколько позволял воротник.

--- Я повесила туда девушку чуть постарше тебя, --- тихо сказала Айну.
--- Ты хотел взять меня в жёны, верно?
Значит, ты должен знать, что происходит между мужчиной и женщиной, когда они остаются одни.

Тахиро исподлобья смотрел на демоницу.

--- Прекрасный акт творения, --- нараспев начала Айну.
--- Экстаз и начало новой жизни в одной чашке.
Если, конечно, по взаимной любви и согласию.
Однако представь, что мужское начало дарует не жизнь, но боль и смерть.

Айну показала Тахиро ацетиленовую горелку с длинным цилиндрическим жалом.
Тахиро вздрогнул.

--- После первого щелчка, --- Айну щёлкнула горелкой, --- она начала издавать звуки.
После второго (щелчок) начала рассказывать то, что знала и не знала.
И когда она закончила... --- Айну направила горелку на кандалы, --- ...я включила струю на полную мощность и ушла на обед.

Тахиро смотрел на Айну.
Грисвольд вдруг поймал себя на мысли, что выражение лица мальчика совершенно не подходило к ситуации.
Так палачи смотрят на тех, чья участь уже предрешена.

--- Был ли у неё выбор? --- ласково закончила Айну.
--- Был.
До первого щелчка.
Но девочки всегда плохо соображают перед этим.
Особенно если... --- Айну щёлкнула горелкой, --- ...это происходит в первый раз.

<<Айну, хватит>>, --- взмолился Грисвольд.
Тахиро продолжал смотреть на Айну тем же бесстрастным немигающим взглядом.

--- Айну, --- начал Люцифер.
--- Пытки --- это грубый способ.
Всегда можно найти...

Айну взмахом руки остановила юношу.

<<Мы не договаривались, что ты будешь играть хорошего палача>>.

--- Да при чём здесь это! --- воскликнул Лу вслух.
--- Пытки --- это непрофессионально и расточительно!
Если ты будешь ломать людей направо и налево, нам будет некем править!
Механизм нужно изучать в действии аккуратно, а не...

--- Заткнись, --- отрезала Айну и обратилась к мальчику:
--- Если ты всё понял, то я тебя слушаю.

\ml{$0$}
{Тахиро ответил не сразу;}
{Tahiro didn't speak at once;}
\ml{$0$}
{он смотрел на Лу расширенными в изумлении глазами.}
{he stared at Lu in wonder.}
\ml{$0$}
{Этот взгляд ещё долго являлся потом стратегу во сне.}
{The strategist has dreamt of that look for a long time since. }
\ml{$0$}
{Раскосые чёрные глаза не просили поддержки, не жаловались на судьбу, в них читался растерянный вопрос: <<Что происходит?>>.}
{Upturned black eyes did not look for support, nor complain, they bewilderedly asked: ``What's happening here?''}
\ml{$0$}
{Стратег, поморщившись, отвернулся --- словно боялся, что вопрос заразит и его разум.}
{The strategist made a wry face and turned his head away---as if he's afraid the question is contagious to his mind.}

--- Я тебя слушаю, Тахиро, --- повторила Айну.

\ml{$0$}
{--- Я требую обмен, --- мальчик продолжал смотреть на Лу.}
{``I demand exchange.'' The boy continued looking at Lu.}

\ml{$0$}
{--- Ты торгуешься с нами, дзайку-мару? --- ухмыльнулся Гало.}
{``Want to establish some trade relations, dzaiku-maru?'' Halo smirked.}

\ml{$0$}
{--- Да, --- просто ответил мальчик.}
{``Yes,'' the boy answered simply.}
\ml{$0$}
{--- Вы в состоянии войны с другими хорохито, это я знаю.}
{``The state of war exists between you and other horohitos, I know that.}
\ml{$0$}
{Каждый час пытки --- потеря ресурсов для вас.}
{Every hour of torture is waste of resourses for you.}
\ml{$0$}
{Я продаю вам несколько часов драгоценного времени.}
{I offer you several precious hours.''}

Айну подняла бровь.
Она явно не ожидала от Тахиро такого взрослого подхода к делу.

--- Вы дадите мне послать в бункер сообщение за пять минут до того, как начнётся штурм.

\ml{$0$}
{--- За минуту.}
{``One minute.''}

\ml{$0$}
{--- За четыре.}
{``Four.''}

\ml{$0$}
{--- За минуту, --- повторила Айну.}
{``One single minute,'' Ain\"{u} repeated.}
\ml{$0$}
{--- Или мы расторгаем сделку.}
{``Or we cancel the contract.''}

\ml{$0$}
{--- Значит, расторгаем.}
{``It's cancelled, then.''}

У Тахиро дрожали от страха губы, пока он произносил эту фразу.
Грисвольд ещё раз подивился его мужеству.

\ml{$0$}
{--- Тахиро, --- ласковый голос Айну был напитан смертельным ядом, --- я хочу прояснить один момент.}
{``Tahiro,'' gentle voice of Ain\"{u} was soaked with deadly poison, ``I'd like to clear one thing up.}
\ml{$0$}
{Речь не идёт о \emph{нескольких} часах драгоценного времени.}
{This is not a question of \emph{several} precious hours.}
Ты, конечно, не женщина, но я найду подходящее отверстие для горелки.
\ml{$0$}
{И когда я найду, мне не понадобятся часы, чтобы вытащить орешек из скорлупы.}
{And when I find it, I'll need less than an hour to shell a nut.''}

\ml{$0$}
{В раскосых глазах Тахиро блеснул азарт:}
{Gambler's excitement twinkled in Tahiro's eyes:}

\ml{$0$}
{--- Спорим?}
{``You bet?''}

Айну задумалась.
Несколько долгих мгновений она, прищурившись, разглядывала маленького пленника.

\ml{$0$}
{--- Три минуты, --- буркнула она наконец.}
{``Three minutes,'' she finally grunted.}

\ml{$0$}
{--- По рукам, --- без раздумий бросил Тахиро.}
{``Deal,'' Tahiro agreed without thinking.}

\ml{$0$}
{--- Что за сообщение ты хочешь послать?}
{``What message do you want to send?''}

\ml{$0$}
{--- <<Всё пропало, спасайся кто может>>.}
{\emph{``It's over, run for your lives.''}}

\ml{$0$}
{--- Как оригинально, --- скривилась Айну.}
{``How original,'' Ain\"{u} grimaced.}
\ml{$0$}
{--- Тебе нужно что-то ещё?}
{``Do you need something else?''}

\ml{$0$}
{--- Я знаю, что вы добыли дымный веер.}
{``I know you've got a fume fan.}
\ml{$0$}
{Дайте его мне.}
{Give it to me.}

--- Ты знаешь сигналы дымного веера? --- поднял брови Гало.

--- Я знаю всё, в том числе и код.
На ваши сигналы никто не отвечал, потому что они изменяются в зависимости от направления ветра и фазы Второй Луны.

\ml{$0$}
{--- Какое кострище будешь использовать?}
{``Which fire pit will be used?''}

\ml{$0$}
{--- Никакое.}
{``None.}
\ml{$0$}
{Найдите фумаролу в виде жабы на южном склоне Арашияма.}
{Find a toad-shaped fumarole on southern side of Arashiyama.''}

\ml{$0$}
{--- Фумаролу? --- выпучил глаза Гало.}
{``A fumarole?'' Halo goggled.}
\ml{$0$}
{--- Так вы использовали дым фумарол, а не...}
{``So you used to use fumaroles, not---''}

\ml{$0$}
{--- Да, конечно, дубина, --- выпалил Тахиро, --- люди бункера носили кремень просто для отвода глаз!}
{``Sure, you dumb piece of wood,'' Tahiro exploded, ``bunker people used to carry a flint as a distraction!''}
\ml{$0$}
{И система кострищ была лишь отвлекающим манёвром.}
{And fire pit system was a decoy as well.}
Её использовали только в случае, если посланец чувствовал слежку.
\ml{$0$}
{Только хорохито будет разводить костры, если есть фумарола.}
{Only a horohito would make a fire if there is a fumarole.''}

Гало и Люцифер переглянулись... и разразились диким хохотом.

В этот раз люди их провели.

\asterism

--- Я посылаю сигнал, мы ждём три минуты и я сообщаю вам координаты входа в бункер.

--- Могу ли я ручаться, что твои данные окажутся истинными? --- осведомилась Айну.

\ml{$0$}
{--- Я клянусь тебе честью моего дома.}
{``I swear with honor of my house.''}

\ml{$0$}
{--- Есть ли честь в том, что ты предаёшь соплеменников? --- ухмыльнулся Гало.}
{``Treachery of your own kind---what kind of honor is it?'' Halo smirked.}

\ml{$0$}
{--- В том, чтобы умирать под пытками, точно нет чести, --- заявил мальчик.}
{``Death under torture is definitely none,'' the boy contended.}
\ml{$0$}
{--- Я признаю лишь смерть в бою.}
{``Death in battle is the only death with honor.}
\ml{$0$}
{Считай это способом выжить и дождаться боя.}
{You may think I try to survive and wait for the battle.''}

\ml{$0$}
{<<Он опасен, Айну, --- сообщил Гало.}
{\textit{``It's dangerous, Ain\"{u},''} Halo said.}
\ml{$0$}
{--- Этого фенрира\FM{} следует уничтожить>>.}
{\textit{``This fenrir\FM{} should be destroyed.''}}
\FA{
Фенрир --- шпион или диверсант, пользующийся внешней привлекательностью и не вызывающий подозрений, обычно ребёнок или человек, внешне похожий на ребёнка (терминология Ордена Тысяча Башен).
}

<<Я против, --- заявил Люцифер.
--- Мальчик крайне интересен как объект исследования.
Я бы попробовал интегрировать его в наше общество>>.

Айну почесала подбородок и оглядела связанного Тахиро.

\ml{$0$}
{<<Хорошо, Лу, я дарю этого поросёнка тебе.}
{\textit{``Well, Lu, I give this piglet to you.}}
\ml{$0$}
{Можешь убить его, когда сочтёшь нужным>>.}
{\textit{Kill it whenever you want.''}}

Айну махнула рукой и вышла из подвала, Грисвольд и Гало последовали за ней.
Лу остался наедине с Тахиро.

Мальчик молчал.
Лу следовал его примеру.
Разумеется, робкие ростки доверия между ними сметены безжалостным вихрем реальности.
Эти двое не могут быть товарищами --- они навсегда останутся хищником и жертвой.
Мальчик понимал это не хуже демона.
Он видел перед собой лишь коварного врага, готового на всё ради победы.

Стратег до смерти не любил этот этап в любом деле --- когда в закрытой системе всё понятно, всё доступно взгляду, но процесс уже запущен в нежелательном направлении.
И будь ты хоть семи пядей во лбу, тебе придётся просто стоять и ждать подходящего момента, чтобы исправить ситуацию.

Последние шаги давно затихли.
Наконец стратег нарушил молчание:

--- Слушай, ребёнок...

--- Благодарю тебя, что вступился.

--- Как ты узнал? --- удивился Люцифер.

--- У нас есть легенда, что некоторым людям однажды преграждают дорогу два брата --- Футаго Рэй и Футаго Ицу.
В глазах одного брата --- жизнь, в глазах другого --- смерть.
И человек пытается задобрить, подкупить или побороть братьев, не зная, что шанс что-то изменить уже упущен и он может лишь наблюдать, как братья спорят.

--- Занятная легенда.
Мне нужно твоё мнение как... представителя другой формы жизни.

\ml{$0$}
{--- Я слушаю.}
{``I'm listening.''}

\ml{$0$}
{--- Как думаешь, благодаря чему существа приходят к сотрудничеству?}
{``Why do you think beings start to collaborate with each other?''}

\ml{$0$}
{--- Благодаря пониманию, --- пожал плечами мальчик;}
{``They understand each other,'' the boy shrugged;}
его отвратительная сбруя не позволяла большего.

\ml{$0$}
{--- А чтобы кто-то понял, достаточно ли просто объяснить?}
{``Then, if you want someone to understand, is it enough to explain?''}

\ml{$0$}
{--- Я думаю, вполне.}
{``I think, yes, it is.}
\ml{$0$}
{Доходчиво и искренне.}
{Clearly and sincerely.}
\ml{$0$}
{Истина всегда недалеко от искренности.}
{Verity is always close to sincerity.}
\ml{$0$}
{Очень важно показать ход мысли.}
{The way you think is important to be shown.}
\ml{$0$}
{Желательно говорить на языке собеседника --- так учила меня мать.}
{It's advisable to speak in language of interlocutor, as my mother taught me.''}

\ml{$0$}
{--- Но что, если собеседник не поверит?}
{``But what if interlocutor will not believe?''}

\ml{$0$}
{--- Это его право.}
{``It's his right.}
\ml{$0$}
{Считаешь собеседника за равного --- уважай его права.}
{To treat someone as equal means to respect his rights.''}

Лу помолчал.

--- На самом деле мы уже были близки к обнаружению бункера, --- наконец сказал он.
--- Это я послал твоих родичей на дело, из-за которого их линчевали.
Вся информация о бункере, которую они, сами того не зная, собрали в пустошах, уже у меня.

\ml{$0$}
{--- Ты просишь прощения?}
{``Do you need my forgiveness?''}

\ml{$0$}
{--- Я не ставил целью их уничтожить.}
{``Destroying them was not on purpose.}
\ml{$0$}
{Расценивай это как хочешь.}
{It's yours to interpret.''}

\ml{$0$}
{Тахиро кивнул:}
{Tahiro nodded:}

\ml{$0$}
{--- Это было доходчиво.}
{``You've made it clear.''}

\ml{$0$}
{--- Я могу снять путы, если ты пообещаешь, что не попытаешься сбежать или убить себя.}
{``I can untie you if you promise not to flee and not to kill yourself.''}

\ml{$0$}
{--- Я уже поклялся всем, что у меня есть.}
{``I have nothing left to swear with.''}

Лу молча вытащил ножницы, разрезал стяжки и вышел, оставив дверь приоткрытой.

\section{[:] Тушение}

\textspace

--- Только не говорите, что этот сброд будет драться, --- поморщился Арракис.

--- Приготовиться к залпу, --- бросил Гало.

--- Легион, отбой, --- тут же подняла руку Айну.
--- Ждём моего приказа.

Один из восставших вышел перед собратьями, собираясь обратиться к ним с пламенной речью.

--- Потуши его, --- лениво бросила Айну стоящему рядом легионеру.

Легионер вскинул винтовку --- и оратор упал с простреленной головой, едва успев выкрикнуть <<Люди!>>.
Толпа дрогнула.

--- На этой планете речи произносятся или мной, или с моего разрешения, --- резюмировал Арракис.
--- Самаэл, я надеюсь, ты сумеешь объяснить это дзайку-мару.

Самаэл кивнул и направился к восставшим, замершим в растерянности.

--- Учись, сайгон, --- ухмыльнулась Айну, бросив лукавый взгляд на Гало.
\ml{$0$}
{--- ,,Один патрон решает любую проблему``\footnotemark.}
{``\textit{One bullet solves any problem}\footnotemark.''}
\FA{Девиз Союза Воронёной Стали, употреблённый в ироничном ключе.}

--- Или можно было просто узнать, чего они хотят, --- вставил Люцифер.
--- А потом насыпать рыбкам мотыля или сменить наконец воду, чтобы они не плавали в собственном дерьме.

--- А твой брат --- теоретик-идеалист, --- закончила Айну со вздохом.

--- Из вас всех я единственный, кто ухаживает за аквариумом и террариумом в штабе, --- парировал Лу.
--- Не будь меня, вся животина давно бы уже сдохла.

--- Ты животину завёл, тебе и ухаживать, --- флегматично ответила Айну.

--- А я чем занимаюсь сейчас, по-твоему? --- раздражённо бросил Лу.
--- Ладно.
Время покажет, кто из нас теоретик.

--- Безусловно, --- кивнул Арракис.
--- И время покажет это в самое ближайшее время.
Штурмовой отряд, вы переходите под командование Люцифера.
Остальной легион --- на указанные точки.
\ml{$0$}
{У нас есть дела.}
{We've got much work to do.''}

Гало удивлённо посмотрел на отца:

\ml{$0$}
{--- Отец, может быть, мне?}
{``Father, maybe I?''}

\ml{$0$}
{--- Нет, Гало, ты нужен в штабе.}
{``No, Halo, I need you in the HQ.}
\ml{$0$}
{Самаэл уже разогнал толпу.}
{Samael's just dispersed the crowd.}
\ml{$0$}
{С моими легионерами бункер сможет взять даже бродячий торговец.}
{With my legionnaires, even a peddler can storm that bunker.''}

Люцифер кивнул и скомандовал построение под пристальным взглядом Айну.

<<Она за меня беспокоится, как за члена семьи, --- подумал стратег.
--- Война ещё не совсем отравила её сущность>>.

\section{[:] Склад}

Люцифер шёл по коридору и в который раз удивлялся, сколько мертвецов осталось в бункере.
Его легионеры убили лишь треть;
прочие умерли ещё до прихода демонов.
Это могла быть людоедская тактика, заставляющая врага поверить, что никто и не думал бежать.
С той же вероятностью это могло быть совершенно бессмысленное самоубийство --- люди Преисподней не особо раздумывали перед тем, как убить или умереть.
Подсчёт пищи и инструментов не смог дать чёткого ответа на то, сколько людей было в бункере час назад --- слишком широкие доверительные интервалы, слишком много помех.
Даже состав и объём воздуха подкачали во всех смыслах --- комнаты были маленькие, но вентиляция работала так, что набитая трубка раскуривалась в руках сама.
Облава по селениям могла бы прояснить этот вопрос, но ресурсов для облавы было недостаточно.

\ml{$0$}
{<<Мой интеллект пересиливает, --- думал Лу.}
{\textit{My intelligence is overpowering the situation,} Lu thought.}
\ml{$0$}
{--- Но недостаточно быстро.}
{\textit{But not fast enough.}}
\ml{$0$}
{Чересчур мало данных.}
{\textit{Too few data.}}
\ml{$0$}
{Чересчур много тех, кто мне противостоит.}
{\textit{Too many people who resist me.}}
Отец вставляет палки в колёса, Гало помешался на своих амбициях, Айну чересчур ограничена, Самаэл развёл бюрократию, Мистраль --- трусиха, а Грисвольд ведёт себя так, словно происходящее его не касается.
\ml{$0$}
{Чтоб вы все сгорели дотла!}
{\textit{May all of you burn to ashes!}}
\ml{$0$}
{Не могу же я биться на три фронта в одиночку!>>}
{\textit{I can't fight on three fronts at once alone!}}

Лу мысленно поставил заметку --- следует поговорить с Грисвольдом о Гало.
Дальше этот спектакль продолжаться не может.

Разумеется, за взятием бункера последует волна восстаний.
Люцифер достаточно хорошо изучил психологию людей, чтобы быть к этому готовым.
Люди Преисподней не любили чувствовать себя проигравшими;
после каждой большой победы Ордена они бросались в бессмысленную, неорганизованную контратаку.

<<И в этот раз, чёрт возьми, мы будем действовать уговорами.
Я лично на совещании заткну всех вояк, включая Гало.
\ml{$0$}
{Бункер мы вынуждены были зачистить, это была неоспоримая военная необходимость.}
{We had to clear the bunker, that was an undeniable military necessity.}
\ml{$0$}
{Но после каждой нашей атаки остаются целые толпы сирот и вдов, жаждущих крови.}
{But each our attack leaves behind masses of orphans and widows, clamoring for blood.}
\ml{$0$}
{Нет уж.}
{Fuck it.}
\ml{$0$}
{Пусть оставшиеся в живых будут несчастны как-нибудь по-другому --- не за счёт Ордена и меня лично>>.}
{Let survivors be unhappy some other way, not through the Order or my bloody self.''}

Носильщики тащили добычу к выходу.
Лу с некоторым удовольствием наблюдал за ними.

<<Как же хорошо, когда кто-то таскает тяжести за тебя>>, --- думал он.

Вот командный пункт.
\ml{$0$}
{Три тела со следами сэппуку.}
{Three bodies with signs of seppuku.}

Вот шифровальная.
Бежавшие не удовлетворились простым убийством шифровальщика;
ему раздробили голову до кашеобразного состояния.
Излишняя, но далеко не глупая предосторожность.

На столе ещё дымилась разворошённая стопка листового пластика.
Люцифер поднял один лист.
Перфокарта.
Фрагмент публичного ключа шифра Курибутону --- одного из первых квантово-устойчивых асимметричных шифров, использовавшихся переселенцами с Древней Земли. % Kuributonu
Взломан чуть ли не пять тысяч лет назад.

<<Мда.
Ну и музей.
Разумеется, они полагались не только на шифр, иначе мы бы давно раскрыли канал связи.
Кстати, а вот и стеганограф...
О небо, сборник стихов о море, эпоха как его там...
Это же моя настольная книга в детстве была, год три тысячи восемьдесят пятый, она тогда стотысячным тиражом вышла.
,,Лилия о пяти лепестках``, засечки на чертах иероглифов, количество чешуек на играющих рыбках...
Уж кто-кто, а я мог бы и догадаться>>.

--- Киноу, --- подозвал Лу посыльную.

--- Сайгон, --- поклонилась женщина.

--- Знакомая книжка?

--- <<Волны Ниигаты>>, юбилейный выпуск на столетие Накагавы Юкико, --- отчеканила легат.

--- Я тебя обожаю, ты в курсе?
Аккуратно, без давления, конфискуем весь юбилейный тираж и меняем его на тот же сборник с иллюстрациями от Окинотори.
Загрузи типографию и выдели для конфиската хранилище с надлежащими условиями.

--- Сайгон, --- Киноу козырнула и ушла.

<<Отлично, --- удовлетворённо подумал Лу.
--- Народ сильно возмущаться не будет, сборник с Окинотори красивый, а нам этого хватит, чтобы разобщить шпионские каналы связи на полгодика.
Пусть лучше наша типография печатает Накагаву, чем эту пропагандистскую чушь>>.

Вот склад с пищевым порошком.
Его количество уже подсчитали, пробы для выявления места производства взяли.
Возле апробированной коробки лежали два трупа --- самоубийца и убитый защитник.

Лу осмотрел склад, просмотрел коробки с помощью тепловизора.
Да, это пищевой порошок.
Мужской у входа, женский --- чуть подальше.
\ml{$0$}
{Больше ничего.}
{Nothing more.}

Но что-то не давало покоя стратегу.
\ml{$0$}
{Коробки.}
{Boxes.}
\ml{$0$}
{Они все были совершенно одинаковые.}
{They were identical.}
\ml{$0$}
{При этом каждый последующий ряд отстоял от предыдущего на какое-то расстояние --- совершенно произвольное.}
{But every row was shifted slightly relative to the previous, and a distance was entirely random.}
\ml{$0$}
{Впрочем, не совсем произвольное...}
{However, not entirely ...}
\ml{$0$}
{Совсем не произвольное!}
{Entirely not!}

--- Разобрать коробки, --- распорядился Лу уже в спины уходящим легионерам.

Вскоре под коробками обнаружилась пустота.
Там сидели двое мужчин, пять женщин и тринадцать детей.
У всех на лицах замерло туповатое выражение;
моргали они редко и медленно, словно под водой.
Люцифер разглядывал их так, словно сам был не рад своей находке.

--- Нейролептический сопор, --- хмыкнул один из легионеров.
--- Нам доставить их в допросную?

\ml{$0$}
{--- Убить всех.}
{``Kill 'em all.''}

--- Сайгон\FM, не лучше ли взять некоторых женщин с выводком?
\FA{
Сайгон --- ранг в раннем Ордене Преисподней, соответствующий легату прима.
}
Если к ним правильно подступиться, они живо заговорят.

\ml{$0$}
{--- Возможно, нам стоит их отпустить.}
{``Maybe we should release them.}
\ml{$0$}
{Что скажешь, Фумиэ из клана Дорге?}
{What would you say, Fumie of the Clan Dourgue?''}

Глаза солдата забегали, словно у провинившегося.
Прочие замерли и вытянулись в струнку, прислушиваясь к разговору.

--- Нам нужна информация, сайгон, --- осторожно начал Фумиэ, --- но неповиновение дзайку-мару ни в коем случае не должно оставаться безнаказанным.

\ml{$0$}
{--- То есть после допроса их следует уничтожить?}
{``It means they must be destroyed after examination, mustn't they?''}

\ml{$0$}
{--- Безусловно.}
{``Unquestionably.''}

\ml{$0$}
{--- Ты не считаешь это расточительностью, напрасной тратой ресурсов?}
{``Don't you consider it waste of resourses?''}

\ml{$0$}
{--- Мы используем штыки вместо патронов, когда это возможно.}
{``We use bayonets instead of bullets if possible.''}

--- Я про дзайку-мару.
Они тоже являются ресурсом.

--- Прошу прощения, сайгон, --- стушевался легионер.
--- Дзайку-мару --- опасный ресурс, который может легко выйти из-под контроля.

\ml{$0$}
{--- И твоё мнение разделяет весь легион?}
{``And does the legion share your view?''}

Фумиэ бросил взгляд на стоящих рядом товарищей, но они старательно делали вид, что разговор их не касается.

\ml{$0$}
{--- Я не могу говорить за легион, но я разделяю мнение командования.}
{``I mustn't speak for the whole legion, but I share the view of the command.''}

\ml{$0$}
{--- В таком случае, легионер, оставь привычку обсуждать мои приказы, --- поморщился Лу.}
{``In that case, legionnaire, break the habit of questioning my orders,'' Lu made a wry face.}
Лицо маленького Тахиро промелькнуло в его памяти чересчур ясно.
--- Бункер выжечь до каменных стен, убить всех обнаруженных дзайку-мару, стены проверить на пустоты и трещины.

Фумиэ с явным облегчением козырнул и резво отправился выполнять приказ.

Люцифер сел на разбитую коробку и, набив трубку адским вереском, раскурил её от щепки.
Он вдыхал и выдыхал терпкий дым, лишь иногда морщась, когда раздавался приглушённый звук втыкающегося в тёплое тело штыка.

\ml{$0$}
{<<Отец меня убьёт, --- с грустью думал стратег.}
{\textit{Father will kill me,} the strategist sadly thought.}
\ml{$0$}
{--- Но пока я за главного, пыток не будет.}
{\textit{But no tortures while I'm in charge.}}
\ml{$0$}
{Даже если это была последняя возможность проявить себя>>.}
{\textit{If even it was my last opportunity to lead.}}

\section{Смех Гало}

--- Люцифер, ещё раз.
Ты приказал обменять у дзайку-мару одну книгу с поэзией на другую книгу с той же самой поэзией.
И отправил заказ на книги с поэзией в наш отдел пропаганды, в типографию.

--- Именно так.

--- А что ты намерен сделать с ними?

--- Ничего.
Аккуратно собрать и положить в хранилище с соблюдением условий хранения.
Через год владельцы их заберут.

--- Я имел в виду владельцев книг, Люцифер.

--- Уговорить обменять книжки и заверить, что сборник им вернут.

--- И всё?

--- И всё.

Арракис повернулся ко второму сыну:

--- Гало, ты можешь как-то объяснить этот приказ?

--- Могу, но не вижу смысла, --- весело сказал Гало.
Его ситуация, похоже, забавляла гораздо больше, чем остальных.
--- Лу --- сайгон, и его приказ должен исполняться, даже если он кажется бессмысленным.

--- Нет-нет, я не против обратной связи, --- поднял ладони Лу.
--- Младшие должны контролировать старших.
Но вот тут налицо какая-то чушь.
Если бы я ещё приказал кланяться мне и подносить дары, я бы понял.
Но чем им книжки не угодили?
Я их даже не для себя печатаю!

--- Если бы ты велел себе кланяться, вопросов бы как раз не возникло, --- хихикнул Гало.

--- Сыновья, вы меня в гроб загоните, --- Арракис устало прикрыл лицо ладонью.
--- Гало, можно тебя попросить заверить его приказ своей подписью?
Ко мне уже в третий раз приходят центурионы с вопросом, не превышает ли юный Люцифер свои полномочия.

--- Почему я юный, а Гало нет? --- горестно простонал Лу.
--- Я на пятнадцать минут его старше!

Гало ухмылялся во весь рот.
Лу вдруг посетило желание вмазать брату по лицу.

\asterism

--- И что это было? --- буркнул Люцифер.
--- Тебя рассмешил мой приказ?

--- С приказом всё отлично.
Очень изящный, точный и экономически выгодный ход.
Я бы так не смог.

--- Тогда в чём дело?
Тебя рассмешил этот спектакль со мной в главной роли?

--- Прости, но это правда смешно.

--- А мне нет.
Ты же прекрасно понимаешь, что ни один центурион не побежал бы к отцу, не будь у него прямого приказа?
Ты вообще видел, чтобы кто-то в легионе капал на своего начальника через его голову, тем более максиму прима?

--- Подожди, Лу, --- Гало впервые за несколько минут перестал ухмыляться.
--- Ты серьёзно думаешь, что отец приказал за тобой следить?
Брось.

--- Это уже не первый инцидент, и я тебе об этом...

--- Лу, давай начистоту, --- Гало остановился и развернулся к брату.
--- Есть некие правила приличия, и ты им не следуешь.
Из-за этого, и только из-за этого, у тебя с легионом регулярное недопонимание.

--- Что ты имеешь в виду?

--- Сам знаешь, --- Гало демонстративно вытащил из уха Лу серёжку и сунул ему в карман куртки.
--- Это война, и детским играм здесь не место.
Начни взрослеть, и твои приказы перестанут подвергать сомнению.

--- Спасибо за поддержку, --- огрызнулся Лу.
--- Хрена с два я тебе теперь табак одолжу.

--- При чём здесь табак?

--- Детские игры.

\asterism

--- Может, погуляем?

--- Я не хочу с тобой гулять, --- Тахиро подбросил в руке плоский камень и отправил его прыгать по волнам.
--- От тебя смердит телами моих родичей.

--- Мне одиноко, --- вдруг признался Лу.

--- Мне тоже, --- кивнул мальчик и, вытерев руки о юкату, встал.
--- Ну тогда пошли гулять.

\asterism

Смех утих.
Лу вдруг стал серьёзен.

--- Слушай, это звучит глупо, но... оставайся у меня.
Я тебе не заменю родителей, но старшим братом побыть могу.

--- Да какой из тебя старший брат, --- поморщился Тахиро.
--- За тобой самим глаз да глаз нужен.

--- Значит, договорились.
Ну и да, для всех остальных ты моя игрушка.

--- Я не хочу быть игрушкой.
Даже невзаправду.

--- Ты будешь моей игрушкой для всех, кому я не доверяю.

--- И для твоего брата тоже.

--- Почему ты так думаешь?

--- Ты пришел ко мне и ищешь во мне брата, --- просто ответил Тахиро.
--- Очевидно, что ты потерял брата в нём, Русахиро.

--- Люцифер.

--- Лусафейро.

--- Лу.

--- Лу.

--- Сошлись.
Пошли домой.

<<Конечно, он согласился, чтобы продолжать шпионить.
Он всё ещё меня ненавидит.
Я чувствую волны ненависти из-под его вежливой улыбки.
Но это неважно.
Он остался.
Я хотел, чтобы он остался по своей воле, и для меня не так важны его мотивы>>.

\chapter{[U] Начало}

\section{[U] Встреча}

\textspace

Несмотря на то, что снаружи царила негостеприимная обстановка заброшенной орбитальной станции, в виртуальной <<комнате>> было довольно уютно.
Собравшиеся сидели в маленькой ажурной беседке из жемчужного дерева, освещённой нежным розовым светом бумажного фонарика.
Из окружающей беседку полутьмы то и дело вырывались огромные, размером с орех, индиго-светлячки и обиженно жужжали, ударяясь о фонарь.

Обстановку придумала Митхэ, женская сущность Митриса Безымянного.
По её словам, наводить домашний уют для неё всегда было истинным удовольствием.

Разумеется, никто из присутствующих не высказал удивления, увидев Лусафейру Лёгкая Рука, а вернее, так называемую базовую копию его личности.
Но вопрос повис в виртуальном воздухе, и первым его озвучил хмурый чернявый мужчина в углу --- Аркадиу Люпино.

--- Грейс, не мог бы ты нам объяснить, как главный оборонительный стратег и максим секунда Ордена Преисподней оказался на собрании адской <<пятой колонны>>?

--- Я подумал, что ему будет интересно, --- осклабился Грейсвольд.

Аркадиу и Анкарьяль кивнули --- им этого объяснения было вполне достаточно.
Но остальных ответ технолога не удовлетворил.
Атрис --- мужская сущность Безымянного --- улыбнулся с едва заметным скепсисом.
Таниа Янтарь громко, выразительно хмыкнула.

--- Итак, --- Аркадиу попытался сгладить напряжение собравшихся.
--- Сегодня мы должны обсудить вопрос эксплуатации сапиентов хоргетами.

--- В чём смысл этого вопроса? --- подал голос Лусафейру.
Говорил он очень мягко, в его устах слова словно теряли острые углы.
--- Эволюция сделала своё дело.
На всякого хищника найдётся хищник сильнее.
Я думал, что вы хотите обсудить нечто иное.

Чханэ гневно уставилась на стратега.
Тот встретил её взгляд холодно и спокойно.

--- Это Разрушение и Насилие в одном горшке, --- резко сказала воительница.

--- Вы имеете в виду законы примитивного человеческого племени планеты Тра-Ренкхаль? --- осведомился Лусафейру.

--- Эти законы были придуманы культурологами цивилизации Тхидэ, --- бросила Чханэ.
--- Они гуманны...

--- ... и совершенно бесполезны для выживания, --- добавил Лусафейру.
--- Что и показало сражение на берегу Могильного пролива, которое едва не унесло остатки твоего народа в \emph{пристанище духов}, --- последние слова он произнёс на идеально чистом сели с характерной интонацией.

Анкарьяль и Аркадиу переглянулись.
В силу большого совместного опыта они начали понимать, с какой целью Грейсвольд привёл Лусафейру.
Увы, опасности, повисшей над заговорщиками, это не убавляло.

Чханэ, похоже, тоже поняла, чего добивается стратег.
Она расслабилась и откинулась на кресле.

--- Почему вам, демонам, так сложно объяснить простую вещь вроде гуманности? --- мягко спросила она.

--- Гуманность --- это то, что связывает руки одним людям, пока другие строят против них козни, --- объяснил Лусафейру.
--- Ваше понимание --- понимание тси --- исключение из правил, результат культурного дрейфа.
Вы выросли в тепличных условиях, вне конкуренции.
Кажется, Аркадиу писал в отчёте, что по прибытии на Тра-Ренкхаль тси очень удивились даже оружию для убийства сапиентов...

--- А я всегда думала, что Ад гуманен по отношению к сапиентам, --- подала голос Митхэ.

--- О нет, флейта моей жизни, --- неожиданно рассмеялся Атрис, --- Ад блюдёт собственную выгоду и ничего сверх этого.

Анкарьяль удивлённо посмотрела на менестреля --- он никогда не казался ей особо сообразительным.

--- Верно, --- подтвердил Лусафейру.
--- На планетах, население которых составляют Ветви Звезды, мы вынуждены практиковать деспотизм, так как сапиенты Ветвей Звезды испускают плюс-эманации при угнетении.

Чханэ издала странный звук --- не то плюнула, не то выругалась.

--- Что такое деспотизм? --- вмешалась Анкарьяль.
--- Сапиентов угнетают под угрозой насилия, хоргеты находятся под тотальным информационным контролем.
Деспотизм везде.

--- Анкарьяль в силу юношеского максимализма преувеличивает, но доля истины в её словах есть, --- признал Грейсвольд.
--- В Ордене нам дают полную свободу, но не забывают <<предупреждать>> или уничтожать, если мы что-то делаем \emph{подозрительно}.
Да, нам дают объяснения любых событий, но когда за фабрикацию доказательств берётся мощнейший отдел аналитики, что можно этому противопоставить?
Вычислительные мощности, опыт и количество актуальной информации прямо пропорцинальны доле власти.
Да, Лу?

--- Отчасти, Грейс, --- мягко ответил Лусафейру.
--- Будучи главным оборонительным стратегом, я трачу треть ресурсов на вычисление вероятности, с которой какое-либо моё действие вернётся ножом в спину.

Митхэ хохотнула и тут же осеклась, поняв, что стратег не шутит.
Он даже не улыбнулся.

--- Ваше удивление понятно.
Большинство из вас выросли в достаточно свободных сообществах.
В Ордене Преисподней степеней свободы и незанятых ниш крайне мало.
Любое действие любой группы тщательно анализируется всеми остальными --- не окажутся ли затронуты собственные интересы.
Объясняется всё это извечной байкой об ограниченных ресурсах, но истина так же извечна --- каждый старается урвать кусок как можно больше, про запас, чтобы он не достался тем, у кого кусков и так много, и не нарушил хрупкое равновесие.

--- Ну так идём с нами, --- вдруг страстно сказала Чханэ, наклонившись к Лусафейру.
--- Мы никогда не ударим тебе в спину.
За себя я отвечаю.

Лусафейру, похоже, смутился.
Не эмоционально.
Он просто понял, что это правда.

--- Таниа Янтарь, какая жизнь ожидает существо, которое обладает таким количеством ценных знаний?
Я пленник Ада, при попытке сбежать я не успею сделать и шага.
Если мне и удастся каким-то невероятным способом покинуть Капитул, кто меня защитит и обеспечит эманациями?
Ты?

--- Да!
Во мне боевые навыки Анкарьяль Кровавый Шторм!

--- Я не всесильна.
Не нужно делать из меня микоргета, --- возмутилась Анкарьяль.

--- Ты погибнешь совершенно зря, --- с лёгкой грустью заключил Лусафейру.
--- Я могу понять твой энтузиазм --- поначалу оцифрованные опьянены ощущением всесилия.
И многие это поддерживают, превращая оцифрованных новобранцев в пушечное мясо.
К сожалению, большинство из них гибнет в первые сто лет после демонизации.
Оставшиеся в живых, как и до оцифровки, начинают больше доверять логике и статистике, нежели рассказам о героизме и собственной исключительности.

--- Давайте не будем отклоняться от темы, --- прервал спорящих Грейсвольд.
--- Лу, так что?

--- Я не буду вам помогать, если ты об этом, --- сказал Лусафейру.
--- Это решать не мне.
Я, если вы не заметили, всего лишь базовая копия.
Если вы хотите возродить Тси-Ди на Тра-Ренкхале --- действуйте, я же буду действовать так, как должен.
Грейс, я ведь стёр разговор после того, как ты вышел от меня?

--- Ты подал сигнал, что стёр.
Что-то передать?

--- Хм, --- задумался Лусафейру.
--- Скажи мне при встрече следующее: <<Поклонись>>.
Интонация и мимика имеют значение.

--- Понял, --- кивнул технолог.

--- Тогда самоуничтожаюсь, --- сказал Лусафейру и исчез.

--- Стой, куда! --- простонала Чханэ.

--- Чханэ, так нужно.
Тем более что это только неполная копия, --- успокоил её Грейс.

--- Копия твоего друга, которая думала и чувствовала, как он, --- невзначай заметил Атрис.
Грейсвольд смущённо замолчал.

--- Что насчёт остальных? --- спросил Аркадиу.
--- Кто-то хочет отказаться от плана?

Ответом было молчание.

--- Тогда перейдём к главному, --- сказал Аркадиу.
--- Я --- официально предатель Ордена.
Митрис не принадлежит к Ордену, но связан с ним договором типа <<Демиург --- Метрополия>>.
О существовании Таниа Орден, будем надеяться, не знает.
А вы, --- обратился он к Грейсвольду и Анкарьяль, --- верные служители с безупречной репутацией.
Если вас уличат в сомнительных действиях --- отдел 100 распутает клубок очень быстро, и Ад ждёт такая чистка, какой не было со времён его становления.

--- Именно поэтому я и приготовил \emph{это}, --- улыбнулся Грейс и помахал пустой рукой, словно сжимая в кулаке маленький круглый предмет.

--- Что это? --- поинтересовалась Чханэ.

--- Пакет ложных воспоминаний, --- сказал Грейвольд.
--- Его созданием занимались девятнадцать демонов из отдела аналитики.
Не волнуйтесь, --- поспешно добавил он, увидев, как Атрис вздохнул и прикрыл лицо рукой, --- не волнуйтесь.
Я разбил задачу на фрагменты, не вызывающие подозрений, и использовал самых лояльных демонов.
На всякий случай я скомпрометировал одного из них связью с Картелем, чтобы увести следствие в другую сторону.

--- Грис, при всем уважении, --- Анкарьяль скрестила руки на груди.
--- Ты один, а в отделе 139 находятся девятьсот шестьдесят демонов.
Ты не можешь быть умнее и дальновиднее тысячи специалистов.
Может, в этот раз твой слепой метод сработал, но это чистое везение и включать такие действия в план работы просто недопустимо.

--- Что в пакете? --- поинтересовался Аркадиу.

--- У нас с тобой возникли разногласия, и ты отбыл, не сказав куда.
Также в несколько ранних воспоминаний добавлены малозначимые эпизоды, которые обычно проходят мимо отчётов, но опытному аналитику намекают на твоё помешательство.
Пакетов всего... эээ... семь.
Да-да, Митхэ ар’Кахр, я подготовил для тебя отдельный.
Учтено всё --- структура личности, взаимоотношений, даже погрешности восприятия и анализа.
Если тебя, Аркадиу, или кого-то из твоих демонов поймают и вывернут наизнанку, о нашей с Анкарьяль причастности не узнает никто.
Обратное также верно.

Атрис нахмурился.

--- Тогда как мы будем держать связь?

--- Рисунки на песке, --- улыбнулась Анкарьяль.

Чханэ недоумевающе хмыкнула.

--- Помните, когда мы были детьми на Тра-Ренкхале, у нас была забава --- рисовать на песке? --- откликнулся Грейсвольд.
--- А ещё мы строили песчаные домики...

--- И когда набегала волна или поднималось солнце, домики рассыпались, --- на лице Митхэ появилась понимающая улыбка.
--- Пара волн --- и песок становился гладким, проступали водяные рёбрышки, словно к нему не прикасалась рука человека.

--- Я всё равно ничего не поняла, --- заявила Чханэ.

--- Тогда, милая, Вселенная преподнесёт тебе ещё один сюрприз, --- рассмеялась Анкарьяль.

Чханэ погрустнела.

--- Так значит, я теперь буду считать вас врагами?
Даже если вы умрёте?

Анкарьяль с нежностью посмотрела на подругу и кивнула.
Чханэ криво улыбнулась.

--- Вы лучшие, кого я знаю.

--- Удачи, --- тихо сказал Грейс и кивнул Анкарьяль.
Оба растворились в воздухе.
Оставшиеся дружно вздохнули и ушли в себя, ожидая, пока закончится интеграция пакетов воспоминаний в систему.

\section{[U] Проблема теплиц (переработать!)}

\textspace

Проблема создания кольцевой теплицы --- в организации генетического аппарата, оптимизации и стабилизации его работы.
Когда учёные Ада взялись за это дело, им пришлось работать с огромным количеством ферментов и каскадов.
Клетки, нагруженные таким количеством генов, быстро выходят из строя из-за случайных мутаций, даже в условиях стерильной среды.
А учёные-тси, согласно данным, решили эту проблему максимально изящно --- альтернативным сплайсингом.
Один ген этого создания, хоть и довольно длинный, мог дать такое количество вариаций, какое обычным живым существам и не снилось.
Скорость мутации генов кольцевой теплицы была ниже таковой в генах гистонов эукариот.
Кольцевая теплица также обладала уникальной системой слежения за собственным генотипом --- специальная ферментативная система фактически вычисляла хэш всех имевшихся хромосом и запускала апоптоз при любой мутации.
Подобное произведение генной инженерии не смог создать никто во Вселенной, даже хоргеты.

--- Интересно, как они выглядели, --- задумался я.

--- Понятия не имею, --- хмыкнул Грейс.
--- И вряд ли кто-то имеет.
По одним данным это растение, по другим --- среди тси были люди, которые с ними общались: ласкали их, разговаривали.
И теплицы им отвечали.
Вряд ли растение способно говорить.

\textspace

\section{[U] Нужна теплица}

\textspace

--- С помощью кольцевой теплицы можно в кратчайшие сроки оживить практически любую планету в обитаемой зоне, не пользуясь услугами хоргетов, --- сказал я.
--- А теперь представьте, что будет, если это знание станет общедоступным.
Всеобщая экспансия за пределы ка'нетовского радиуса.
Любой хоргет сможет двигаться сколь угодно далеко в любом направлении, найти подходящую планету и стать её полновластным демиургом.
Каждому хоргету --- по планете!
Ад и Картель находятся в истинном равновесии ровно до тех пор, пока у них есть физические границы.
Убери границы --- и старый мир, основанный на идее дефицита, рухнет, как когда-то рушились колониальные империи.

Я помолчал, собираясь с мыслями.

--- Второе направление работы --- планетарная система.
Нам необходимы её чертежи.
Демиурги беззащитны перед организациями демонов.
С доступностью системы планетарной защиты гегемонии Ада и Картеля придёт конец.
Более того, демиург вступит в мутуалистические отношения с сапиентной цивилизацией своей планеты, так как от её технического развития будет напрямую зависеть его собственная безопасность.

--- Мне не нравится то, что ты предлагаешь, Аркадиу, --- сказал Грейсвольд.
--- Жизнь демиурга для каждого демона --- это хорошо, но я хотел бы общаться с себе подобными, а не запираться в планете на веки вечные.
То, что ты предлагаешь, разрушит социум демонов --- как экспансия за пределы ка'нетовского радиуса, так и превращение демонов в богов.
И существует ещё одна проблема.
Для минус-демонов места в этой системе не будет.

--- Технологии тси помогут нам и здесь.
Вдруг существует возможность вывести нуль- или даже минус-сапиентов?

--- Вероятность крайне мала, --- рассудительно сказал технолог.
--- Боюсь, что если нам не удастся этого сделать, Картель будет драться до последнего.
А это значит, что они почти наверняка победят.
Поэтому, если ты действительно хочешь создать дивный новый мир, следует проработать момент с минус-демонами.

\textspace

\section{[U] Подруга}

Со дня битвы на Могильном берегу прошёл год.
Жизнь потихоньку возвращалась в привычное русло.
Вместе с добровольцами-хака и прибывшими с Запада мастерами тхитронцы восстановили большую часть города.
С крыши храма была видна лишь одна оставшаяся проплешина на месте квартала каменщиков;
квартал был одним из самых хорошо укреплённых районов восточного Тхитрона, после осады там не осталось камня на камне, и место поглотила густая молодая поросль.
Её было решено не трогать --- оборонять город больше было не от кого.

Анкарьяль стояла рядом со мной.
Её жёсткие прямые волосы, уже отросшие до плеч, трепал настойчивый юго-западный ветер.
Пока ещё незрячий, по-младенчески синий глаз, растущий в пустой глазнице, светился недовольством.

--- Мне это не нравится, --- заявила она.

--- Ты это уже говорила.

--- Повторяю ещё раз: мне это не нравится.

--- Что именно тебе не нравится?
Место встречи?
Выбранные люди?
Перспектива превращения этих людей в демонов?

--- Я не о них!

Я облокотился о колонну и вздохнул.

--- Тогда в чём дело?

--- Мы были прекрасной командой, Аркадиу, --- тихо сказала Анкарьяль.
--- Я, Грейс, ты.
Я знаю, что такое плохая команда, я знаю, о чём говорю.
Я не хочу тебя терять.

--- И мы можем остаться командой.

--- Нет, не можем.

--- Ну, если ты про прежние наши отношения --- стоять плечом к плечу, выполняя приказ --- то да, не можем.
Но есть другие возможности.

--- Ты вообще понимаешь, о чём ты меня просишь?

--- Я тебя ни о чём не прошу.

Анкарьяль потёрла висок и тоже облокотилась о колонну, вперив в меня взгляд разнокалиберных, разноцветных глаз.

--- Я помогла тебе устроить эту встречу, --- сказала она.
--- У тебя есть методика и люди.
Дальше --- сам.
Хорошо?

---  <<Сам>>?
<<Хорошо>>?

Анкарьяль вдруг сползла по колонне вниз, сжалась в комочек и спрятала лицо в ладони.
Она дрожала крупной дрожью.

Я погладил подругу по голове.

--- Пойду вниз, --- шепнул я.
--- Время.

Ответом были сдавленные рыдания.

\section{[U] Вербовка}

\textspace

Собравшиеся недоумённо переглянулись.

--- Ты больше не будешь сражаться рядом со своими товарищами? --- спросил Акхсар.

--- Нет.

--- Почему? --- требовательно спросил воин.
--- Они показали себя, Ликхмас.
Если бы я мог, я бы встал с ними рядом без колебаний.

--- Я хочу собрать собственную команду.
Однако существуют определённые правила, и мои товарищи не могут присоединиться ко мне по регламенту.

--- Новая команда?
Ты хочешь набрать других демонов?

--- Да.
Вас.

\textspace

--- Анкарьяль как-то подметила, что технология у вас в крови.
Я верю, что это действительно так.

--- Мы --- не наши предки, Ликхмас, --- строго сказал Акхсар.
--- Мы не можем летать между звёзд, управлять вашими машинами...

--- Я видел, как ты обращаешься с моим мотоциклом.
Ты освоил механизм и управление за декаду, увидев его впервые в жизни.

--- Это же просто механический олень.
Я должен знать, как на нём ездить и как его чинить.

--- Чтобы стать демонами, вам придётся овладеть такой же технологией.
Просто мотоцикл будет немного сложнее.

Акхсар нахмурил кустистые брови и задумался.

--- А ты, Кхотлам, --- обратился я к кормилице, --- я отдал тебе лишний компьютер для твоей работы.
Сейчас ты уже свободно пишешь на нём программы.

--- Ещё один язык, --- пожала плечами Кхотлам.
--- И немножко математики.

--- Истина в том, что на это не способны взрослые в других племенах, --- громко сказал я.
--- Ни хака, ни зизоце, ни прочие подобные вам народы во Вселенной не способны вот так, походя, понять принцип работы машины или концепцию квантовой архитектуры.
Вы были рождены для работы с машинами.
Многие из вас считают демонов сверхъестественными существами, но в действительности вас --- именно вас, потомков тси --- от демонов отделяет один шаг.
И я могу помочь сделать этот шаг.

--- Ты сделаешь нас бессмертными, Ликхмас? --- спросил кто-то.

--- Вы получите долгое существование, но не бессмертие.
Если сапиентам уготована лишь одна первичная смерть, то вы будете умирать ею сотни и тысячи раз.

Многие опустили головы.

--- Лисёнок, --- сказала Кхотлам, --- значит ли это, что мы переживём всех родных?

--- Вы можете пережить даже свой собственный народ и землю, на которой стоите, если вам повезёт.
Если же вас настигнет вторичная смерть, то не будет для вас пристанища --- вы умрёте навсегда, и никогда уже не будете живыми.

Все дружно вздохнули.
Акхсар громко скрипнул зубами.

--- Что я буду видеть и слышать, если стану демоном? --- спросил воин.

--- Вы испытаете ощущения, которые даже не можете себе представить.
Будут среди них и неприятные, и поистине ужасные.
Вам придётся долго и упорно тренироваться, чтобы управлять своими новыми силами.

--- Зачем ты предлагаешь нам такую участь? --- недоумённо поинтересовалась женщина с краю.
--- Менять одну смерть на тысячи, бессмертие и покой в пристанище на борьбу и забвение...

Окружающие закивали, выражая согласие.

--- Мне нужна помощь, --- просто сказал я.
--- Я не требую от вас, не взываю к чести предков, не рассказываю о тех чудесах, которые вы познаете, хотя и чудес будет немало.
Я смиренно прошу о помощи, потому что это --- величайшая и, вполне возможно, напрасная жертва.

--- Почему напрасная? --- удивился Акхсар.

--- Вы помните владычество Эйраки.
Это был единственный угнетатель.
Десятки тысяч миров стонут под гнётом других, не менее могущественных демонов.
И многие из них не одиноки, существуют целые организации, обладающие огромными ресурсами и глубокими знаниями.
Мы можем проиграть, и от нас не останется даже воспоминаний.

--- Если мы все погибнем, что-то изменится? --- спросил Акхсар.

Я улыбнулся.
Тот же вопрос задавали мне многие перед битвой на Могильном берегу;
ответ на него ничуть не изменился.

--- Тем, кто пойдёт за нами, будет легче.

--- Чем мы можем тебе помочь, братик?
Мы --- простые люди, --- заметила Лимнэ.
--- Я владею клинком недостаточно даже для того, чтобы сразить умелого воина...

--- Я научу вас тому, что знаю.
Мне понадобятся не столько воины, сколько те, кто готов трудиться и искать --- знатоки живых существ, строители машин, книжные люди.
Даже те, кто хорошо умеет что-то одно, смогут помочь.

Сестрёнки кивнули в знак того, что поняли.

--- Даю на размышление пять дней.
Если кто-то решит, что готов --- жду вас здесь в Змею 1, на восходе солнца.
И помните --- обратного пути уже не будет.

\section{[U] Первые}

--- Никто не придёт, как пить дать, --- сквозь зубы процедила Анкарьяль.
За последние два дня я услышал эту фразу не менее ста раз.
--- Ты вообще узнавал, как проходит вербовка?
Сапиентам предоставляют полную информацию, а ты сделал всё, чтобы их отпугнуть.

--- Придут немногие, --- согласился я, --- но их ценность выше, чем у тех, кто зарится на долгую жизнь и чудеса.

--- Никто не придёт, --- повторила Анкарьяль.
--- Я почувствовала, как их напугали слова о посмертии.
Пристанища не будет!
Такое и в страшном сне привидеться не может.

--- Его не будет в любом случае, --- тихо напомнил я.
--- Просто сели, как ни крути, дикари, а вот демонам верить в такие сказки не к лицу.

Мы ждали долго в напряжённом молчании.
Тихо потрескивая, горела свеча.
И вот в назначенный час раздался стук в дверь.

--- Ты проиграла, --- бросил я, направившись к двери.
--- Один есть.

--- Много ты сделаешь с одним, --- горько прошептала Анкарьяль.

За дверью меня встретили девять человек.
Я молча пропустил их в дом.
Все, не дожидаясь приглашения, принялись наполнять свои чаши чаем из котелка.

--- Что сказала кормилица? --- обратился я к сестрёнкам.

--- Она не пойдёт.
Акхсар хотел, но сказал, что бросить Кхотлам будет для него бесчестьем, --- ответила Манэ.
--- Кормилица дала ему второй шанс.

--- А ваш мужчина?

--- Он не хочет лишаться пристанища, но сказал, что целиком нас поддерживает, --- объяснила Лимнэ.

--- Ты же говорила, что хочешь покоя и не желаешь бессмертия, --- улыбнулся я Чханэ.
Подруга не ответила на улыбку, её взгляд был предельно серьёзен.

--- Когда мы говорили, всё было хорошо, --- проворчала она.
--- Теперь я вижу, что не всё в порядке и моё присутствие необходимо.

--- А ты почему пошёл, Митрам? --- обратился я к молодому ювелиру.

\ml{$0$}
{--- Ты сказал про строителей машин, --- пожал тот плечами.}
{``You mentioned machine builders,'' he shrugged.}
\ml{$0$}
{--- Я люблю машины, увлечён ими с детства.}
{``I love machines since I was a child.}
\ml{$0$}
{Посмертие --- весьма разумная цена за знание.}
{Afterlife is a good price for knowledge.''}

<<Эти слова могли бы принадлежать древнему тси>>, --- невольно вырвалось у Анкарьяль.
Я мысленно согласился.
Произнёсший подобное --- уже не дикарь.

--- А вы --- жрецы, --- поклонился я двум мужчинам и женщине в синих робах.
--- Врач, счетовод и...

---  ... и Митхэ ар'Тра, --- закончил приятный контральто.
Женщина откинула капюшон робы, и слабый свет выхватил из темноты тяжёлые серебристые волосы и благородное лицо, в котором я вдруг обнаружил знакомые черты.
--- Масло-Взбитое-Взглядом.

--- Я чту традиции Севера, Митхэ, --- улыбнулся я.
--- Учитель должен знать всех учеников по именам.

--- А мы с братом просто пришли помочь, --- вмешался крестьянин Согхо, махнув на стоящего рядом незнакомого мужчину.
--- Я всю жизнь обрабатывал землю, знаю, как выращивать кукурузу, помидоры и цветы.
Мой брат --- змеелов.
Мы оба любим своё дело.
Можем ли мы пригодиться?

Я улыбнулся.
Два диктиолога, интерфектор, технолог, три когитора-информатика и два биолога.
Превосходное начало.

--- Вполне, --- кивнул я.
--- Только вы понятия пока не имеете, что вам придётся выращивать и каких змей ловить.

<<Ты же не надеешься, что они сразу станут профессионалами? --- скептически произнесла Анкарьяль.
--- Они дикари, Аркадиу.
Многие из них могут сойти с ума даже из-за процесса демонизации>>.

<<Лесные духи, --- махнул я рукой.
--- Уж ты-то помнишь, с чего начинал я>>.

<<Ты был ребёнком>>, --- напомнила Анкарьяль.

Тут она права.
Мысль набрать в свои ряды детей я уже рассматривал.
Но перед глазами тут же вставал Отбор, когда ничего не понимающему созданию предлагали сделать главный выбор в его жизни...

Анкарьяль понимающе кивнула.
Она, как обычно, читала мои мысли.

<<Что ж, тебе придётся работать с тем, что есть>>.

<<,,Тебе``?
А ты со мной, Нар?>>

Подруга вздохнула и отошла набрать чаю из котелка.

\chapter{[:] Агнец}

\section{[:] Желание демона}

<<Тахиро повзрослел>>.
Лу, откинувшись в кресле, любовался другом.
Длинные чёрные волосы, миндаль огромных глаз, острые скулы... и ни намёка на бороду.

--- Что-то ты давно не сливал информацию повстанцам, --- сказал стратег.
Тахиро недоумённо замолчал.

--- Эээ...

--- Да не прикидывайся.
Первые два года ты был информатором у повстанцев.
Но потом прекратил с ними общение.

--- Я хотел помочь своему народу, --- смутился Тахиро.

--- Не оправдывайся.
Это было забавно.

--- А мне нет, --- вдруг резко сказал Тахиро.

--- Прости, не хотел тебя обидеть.
Но всё-таки, почему ты перестал?

--- Привык к сытой и спокойной жизни в лагере.
Любая зверушка в конце концов приручается.

--- Я никогда не считал тебя зверушкой.

--- Поэтому ты и смог меня приручить.

Друзья расхохотались.

--- На самом деле я просто потерял мотивацию, --- признался Тахиро.
--- Чем больше я узнавал о вас, тем больше понимал три вещи.
Первое --- вы не настолько опасны, какими вас представляют в поселениях.
Второе --- вы не настолько глупы, какими вас представляют в поселениях.
Третье --- вы не настолько едины и однообразны, какими вас представляют в поселениях.
Таким образом, для меня встал вопрос --- как с вами бороться правильно и нужно ли вообще с вами бороться.

\ml{$0$}
{--- И что надумал?}
{``Got any ideas?''}

\ml{$0$}
{--- Я в процессе.}
{``I'm in the middle of thinking.''}

--- Ну ты это, держи в курсе, --- Лу захихикал над собственной шуткой и тут же отчаянно закашлялся, подавившись слюной.

\textspace

--- Подумай только, Тахиро, --- мечтательно сказал Лу, --- речь о таких расстояниях и такой гравитации, что даже свет не летит прямо, а течёт, словно вода между камней!

--- Не вижу в этом ничего особенно удивительного.

--- Правда? --- удивился Лу.

Тахиро вместо ответа достал из сумки линзу.

--- Тоже вариант, --- немного ошарашенно признал Лу.
--- Но заставить свет летать кругами ты не можешь.

--- А смысл? --- ухмыльнулся Тахиро.
--- Кто увидит летающий кругами свет?

--- Ты споришь не ради истины, а ради победы, --- поморщился Лу.

--- Так я прав?

--- Мне нечего на это ответить.

Тахиро помолчал.

--- Какая-то безвкусная победа, --- наконец признался он.
--- Ты позволил мне победить, а не проиграл.
Это разные вещи.

Лу ухмыльнулся.

Закатные аспиды угасли, и небо расчертили брошенные горстью метеоры.

--- У нас принято загадывать желание, когда сгорает метеор, --- сказал Тахиро.
--- Говорят, оно непременно сбудется.
Только нужно сказать его вслух.

Люцифер засмеялся.

--- Я думаю, что для камня, пролетевшего миллиарды километров, оскорбительно служить моим мелким желаниям.

--- Он не обидится, --- Тахиро похлопал Лу по плечу.
--- Без тебя его существование было бы совершенно бессмысленным.

--- Скажи, как у тебя так выходит?

--- Ты о чём? --- не понял Тахиро.

--- Вот взять меня.
Во мне масса информации --- факты, формулы и структуры, связанные между собой определённым образом.
Это большая часть того, что нам удалось собрать из наследия первых людей.
Можно сказать, я квинтэссенция их науки.
Но ответы на все вопросы есть только у тебя.
И я бы не сказал, что они лишены смысла.
Иногда мне даже кажется, что твои слова спали во мне всю мою жизнь и проснулись, едва ты заговорил.

--- Это называется <<философия>>, сын.

Арракис, как обычно, подошёл незаметно.
Друзья переглянулись.

--- Я уверен, что у тебя в памяти есть масса философских концепций, --- продолжал Арракис, положив холодные как лёд руки на плечи Лу.
--- Жизнь ещё научит тебя различать их в речах людей.
Философия помогает людям преодолеть акбас, помогает им защититься от бесконечности, хаоса и случайности, царящих во Вселенной.
Та же философия делает их ограниченными и невежественными.

--- Почему? --- поинтересовался Люцифер.

--- По причине, которую ты уже озвучил.
Люди уверены, что у них есть ответы на все вопросы.
А теперь оставь своего питомца наблюдать за звёздным небом --- ему это полезно.
Ты нужен мне в штабе.
В конце концов, ты создан для того, чтобы править, не так ли?

--- Подожди, Лу, --- Тахиро схватил друга за руку.
--- Один из метеоров точно был твоим.

--- Если я создан для того, чтобы править, --- Лу бросил взгляд на Арракиса, --- я бы стал лучшим из правителей.
А ты, Тахиро?

--- Я бы сделал всё, чтобы лучший из правителей дожил до конца войны, --- сказал Тахиро.

--- Увидимся за ужином, --- Лу улыбнулся и растворился в ночной тишине.

\section{Ксидатин}

--- У меня нет уверенности, что мы решили проблему производства ксидатина, --- проговорил Лу, крутя в пальцах тонкий слиток циркония.

--- Лаборатория уничтожена семь лет назад, Люцифер.
По твоим же подсчётам, она могла обеспечивать почти всех, кто был в бункере.

--- А еще мы нашли шприцы с дозами у деревенских, не далее месяца назад.
Что это было --- старые запасы?
Вы брали у трупов разведчиков пробы на ксидатин?

--- У нас нет ресурсов исследовать каждый плевок, Люцифер.

--- Именно поэтому версию с несколькими лабораториями отметать не следует, --- заключил Лу.
--- Ксидатин легко синтезируется из аптекарских препаратов.
Единственной проблемой являются катализаторы, --- Лу помахал циркониевым слитком, --- но если их добыть --- использовать можно почти бесконечно.
А препараты мы по понятным причинам не можем изъять из оборота.

--- По понятным причинам? --- уточнил Гало.

Лу закатил глаза.

--- Ладно, ладно, я понял, --- поднял руки Гало.
--- Хватит делать \emph{это лицо}.

--- Я не понимаю, в чём проблема, --- сказал Арракис.

--- Для Лу изьятие лекарств --- табу, --- сообщил Гало.

<<Это какой-то сюр, --- устало подумал Люцифер.
--- Какое в бездну табу?
Я точно не сплю?>>

--- Надеюсь, вы не хотите устроить эпидемический бунт? --- ровным голосом сказал он.

--- Мы можем изъять ограниченное число лекарств, --- пожал плечами Арракис.

--- Изъятие любого лекарства --- увеличение заболеваемости, --- объяснил Гало.
--- Увеличение заболеваемости --- увеличение риска эпидемий.
Лу имел в виду это.
Но я думаю, что...

--- В случае эпидемии дзайку-мару обвинят нас, ограничивавших оборот лекарств, --- перебил Лу.
--- И будут правы, несмотря на то, что их аргументы будут другие и логически неправильные.
Я против.
И даже не вздумайте выносить это на рассуждение --- я сделаю всё, чтобы наложить вето.

--- Я тебя понял, --- кивнул Арракис.
--- Тогда я хочу тебя попросить заняться проблемой ксидатина самостоятельно.

Лу кивнул.
<<Ну ещё бы ты не попросил>>.

\section{[:] Трубки}

\textspace

--- Что ты делаешь? --- удивился Люцифер.

--- Хочу проработать некоторые моменты стратегии, --- уклончиво ответил Гало.

--- У нас сейчас свободное время, --- напомнил Тахиро.

--- У тебя вся жизнь --- свободное время, дзайку-мару.

--- У тебя тоже, --- непринуждённо парировал Тахиро.
--- Ты --- вольная птица.

--- Я ценен как работник, а не как подопытное животное.

--- А в чём разница?

--- В уровне интеллекта, --- ядовито ответил Гало.

--- Это точно.
Ни один грызун не додумается крутить своё колёсико во время отдыха.

--- Курить хочется, --- лениво проворковал Люцифер.
--- Гало, кинь какую-нибудь трубку.

Гало запустил в тазик с трубками руку и бросил Люциферу первую попавшуюся.

--- Нет, Гало, не эту, --- сказал Люцифер, осмотрев брошенное.
--- У неё чересчур короткий чубук.

--- Ты сам сказал <<какую-нибудь>>.

--- А ещё я сказал <<не эту>>.

--- Это же очевидно, Гало, --- усмехнулся Тахиро.
--- Если Лу сказал <<какую-нибудь>>, то ты должен взять трубку и спросить, эту ли хочет Лу.
Если Лу сказал <<любую>>, то у него очень хорошее настроение и можно не спрашивать, а просто кинуть первую попавшуюся.

--- Что за чушь ты городишь? --- буркнул Гало.

--- Отнюдь, --- возразил Лу.
--- Тахиро правильно понял, что первое слово означает определённое нежелание в сочетании с неопределённым желанием.
Второе обозначает неопределённое желание при полном отсутствии нежелания.

--- Брат, тебе сложно просто сказать, какую трубку ты хочешь?

--- Трубки все мои, и я не обязан делать чёткий выбор между ними, --- объяснил стратег.
--- В теории я могу закурить столько трубок за раз, на сколько хватит рта и лёгких.
Так что, пожалуйста, кинь мне какую-нибудь трубку, а эту коротышку положи обратно.

--- Подними задницу и возьми сам.

Гало раздражённо выключил терминал и вышел из комнаты, хлопнув дверью.

--- Как хорошо, когда есть близкий родственник, --- Лу откинулся на спинку кресла и вперил взгляд в скучный каменный потолок.
--- Желания курить как не бывало.

\section{[:] Мико}

\textspace

--- Как тебя зовут?

Девушка что-то пробормотала.

--- Прости? --- Тахиро наклонился, припав ухом почти к её рту.

--- Мико из дома Мацумори, --- пролепетала девушка.

--- Я Тахиро из дома Ханаяма, --- кивнул Тахиро и тут же добавил, увидев расширившиеся от страха глаза:
--- Что-то не так?

--- Ты Ханаяма?

--- Да.
<<Благородная кожа --- лучшее одеяние>>.

--- <<В сосновом лесу нечисть не водится>>, --- машинально ответила Мико и запнулась.
--- Ты... ты не будешь мстить?
Мой дом...

--- Я знаю, Мико, --- Тахиро обнял девушку за плечи.
--- Ты не виновата, ведь ты тогда ещё даже не родилась.
Всё в прошлом.

\ml{$0$}
{--- Благодарю тебя, Тахиро, --- прошептала Мико.}
{``I thank you, Tahiro,'' Miko whispered.}
\ml{$0$}
{--- Кто-нибудь ещё из твоего дома спасся?}
{``Is there any of your kinsmen who survived as well?''}

\ml{$0$}
{--- Только я.}
{``Only me.}
\ml{$0$}
{Вернее, нас было двое, но...}
{Actually, there were two of us, but ...}
\ml{$0$}
{Моя сестрёнка была совсем малышкой, недавно родилась.}
{My sister, she was just a baby, newborn.}
\ml{$0$}
{Я не знаю, что с ней произошло.}
{I don't know what happened to her.}
\ml{$0$}
{Она долго плакала, потом перестала плакать и окоченела.}
{She had been crying for a long time, then she stopped crying and got rigor.}
\ml{$0$}
{Мне пришлось оставить её в пустошах.}
{I had to leave her in the wastelands.''}

\ml{$0$}
{--- Сочувствую, --- пробормотала девушка.}
{``I'm sorry,'' the girl mumbled.}
\ml{$0$}
{--- Мне очень жаль.}
{``I'm so sorry.''}

Тахиро смотрел на девушку и никак не мог насмотреться.
Она была воплощением нежности и женственности.
Когда Тахиро обнял её, на веснушчатом тонкокожем лице расцвела милая белозубая улыбка.
Узкие глаза с такэсакским, <<шу>> раскроем, светились радостью и смущением;
девушка крепко прижимала маленькие ручки к груди, и этот жест окончательно свёл Тахиро с ума.

--- Главное --- не бойся, --- улыбнулся Тахиро.
--- Сейчас ты в безопасности, тебя никто не обидит.
Когда ты придёшь в себя, я возьму тебя в жёны.

Мико хихикнула.

--- Правда?

--- Правда.
У тебя будет безопасный дом.

\textspace

Тахиро вдруг понял, что мысль о мести даже не пришла ему в голову.
Он знал, что Ишида и Мацумори виновны в гибели его дома, убийстве его матери и отца, но...

<<С каких пор мне стало плевать на мой дом? --- удивился Тахиро.
--- Едва ли это просто чувство к женщине>>.
Затем он посмотрел на Люцифера и Грисвольда --- Лу показывал другу какую-то очень весёлую пантомиму, и толстяк нервно икал в перерывах между приступами смеха.
Этих демонов не связывали ни родственные узы, ни место рождения, ни жизненный опыт.
Среди демонов Ордена было много таких --- без роду, без племени, рождённых неизвестно где, переживших неизвестно что, прибившихся к первым встреченным ими представителям своего вида.

<<Я такой же, как они, --- подумал Тахиро.
--- Мой дом мёртв, и всё, что мне досталось от него --- имя и меч.
Мне некого защищать, кроме меня.
Мне не за кого отвечать, кроме меня.
Наверное, это действительно больше не имеет значения>>.

Его взгляд снова упал на Мико.
Девушка украдкой рассматривала его издалека.

<<Думаю, нужно сделать ей подарок>>.

\section{[:] Изнасилование}

Тахиро достаточно неплохо научился управляться с ЧПУ-станками.
Большая часть из них, разумеется, управлялась омега-консолью, но примитивный нейроинтерфейс --- двухступенчатый зрительно-слуховой --- тоже сделали.
Видимо, на случай нехватки масс-энергии.
Тахиро просидел у 3D-принтера целых три часа, но браслеты получились действительно восхитительные --- такой тонкой работы в людских поселениях было не найти.
Титановая филигрань перемежалась с глазурованной мозаикой, изображавшей аистов, маки и сосновые ветви.

\textspace

Девушка дрожала с головы до ног.

--- А что произошло потом? --- спросила Айну.

--- А потом они... меня...

--- Изнасиловали, --- сухо подсказала Айну.

--- Да, --- шёпотом подтвердила девушка.
--- Меня и раньше... ещё когда я служила семье Таками... но в этот раз всё было кошмарно... я не представляла, что так вообще возможно.

--- Было очень страшно, --- подсказала Айну.

--- Да.
Больно не было, а вот страшно...

--- Всё ясно, --- сказала Айну.
--- Стимуляция.
Как и мы когда-то, они думали, что это эффективный способ.

Демон Айну вспыхнул --- и что-то приглушённо хлопнуло.
Девушка медленно легла на стол с застывшими глазами.
Изо рта вытекла тонкая струйка крови.

--- Зачем? --- рявкнул Тахиро.

--- Тебя никто не спросил, дзайку-мару, --- отозвалась Айну.
\ml{$0$}
{--- Избавься от трупа.}
{``Get rid of the corpse.''}

Тахиро схватился за дайту-соро и тут же упал на землю, заорав от боли.
На его теле появилось множество мелких болезненных царапин, словно на парня в один миг набросилась стая озлобленных крыс.

--- Я сказала --- избавься от трупа, --- не повышая тона, повторила Айну.

\section{[:] Пищевая цепь}

Грисвольд нашёл Тахиро не сразу.

Парень сидел за складами, рядом со свежей могилой.
Он снял с себя всю одежду, от него разило острым запахом какого-то крема, рядом лежала пустая ампула из-под обезболивающего, но Тахиро всё равно морщился при каждом движении.

--- Я не понимаю, как она могла мне нравиться раньше, --- пробурчал парень, пытаясь вытереть злые слёзы краем кимоно.
\ml{$0$}
{--- Бесчеловечная дрянь.}
{``Inhumane scum.''}

--- Она не плохая, Тахиро, --- сказал Грисвольд.
--- Поверь, я видел плохих демонов.
Айну не такая.

\ml{$0$}
{--- Она относится к людям как к вещам! --- рявкнул Тахиро.}
{``She treats humans like implements!'' Tahiro roared.}
--- <<Избавься от трупа>>!
Глоточная бомба!
\ml{$0$}
{Она забила девушку как свинью!}
{She slaughtered the girl like a pig!}
А ещё она использовала шугокскую пытку\FM, словно я какой-то... какой-то...
\FA{
Шугокская пытка --- нанесение множества мелких ран на тело.
}

Тахиро захлёбывался от гнева.
Грисвольд вздохнул.

--- Когда я был ещё молодым демиургом, я тоже относился к людям, как к материалу, --- сказал он.
--- Что произошло?
Я просто начал вас изучать.
Изучал как материал, а в итоге понял, насколько мы с вами похожи.
Я никогда не относился к людям как к равным, но вот это знание, которое я когда-то получил, не позволило мне относиться к ним, как к вещам.

--- Она же тоже...

--- Айну --- демон, Тахиро.
Форма жизни, отличная от вас, несмотря на то, что демоны часто существуют в связке с человеческим телом.
Она живёт уже в двадцать раз больше, чем ты.

Тахиро помолчал.

\ml{$0$}
{--- Я обещал Мико, что возьму её в жёны.}
{``I promised to marry Miko.''}
Я думал, что делаю ей свадебный подарок, а сделал погребальное убранство.

Грисвольд удивлённо посмотрел на могилу.

\ml{$0$}
{--- Ты о?..}
{``You're talking about---''}

\ml{$0$}
{--- Да.}
{``Yes.''}

Грисвольд огляделся.

--- А где твой соро?

--- В могиле, с ней.

--- Это же твоя семейная реликвия!
Его создали ещё на Древней Земле!

\ml{$0$}
{--- Это просто железяка.}
{``It's just a piece of iron.}
\ml{$0$}
{Но чтобы отогнать сущности на чёрной тропе, нужно оружие.}
{But she needs a weapon to repel creatures of the Black Road.}
\ml{$0$}
{Надеюсь, хотя бы против демонов посмертия этот меч чего-то стоит.}
{I hope this soro is worth something, at least against afterlife demons.''}

Тахиро прижался к прохладной стене склада и с облегчением задышал.
По его лицу снова потекли слёзы.

--- Как же мне надоело бессилие человеческого тела.

\ml{$0$}
{--- Я тебя понимаю, --- кивнул Грис.}
{``I can understand,'' Gris nodded.}
\ml{$0$}
{--- Сотня царапин --- и обладатель самого смелого сердца может лишь сидеть и стонать от боли.}
{``A hundred scratches, and the bravest heart can only sit and groan in pain.''}

\ml{$0$}
{--- Я не стонал.}
{``I wasn't groaning.''}

\ml{$0$}
{--- Вы, люди, придумали какие-то глупые ограничения для нормальных реакций организма.}
{``You people invented plenty stupid restrictions for normal reactions of your bodies.}
\ml{$0$}
{Я не перестану считать тебя храбрецом из-за слёз и стонов.}
{I keep thinking of you as a brave man even if you cry and groan.}
\ml{$0$}
{Это вообще никак не связано.}
{Actually, these things have no connection.''}

\ml{$0$}
{--- Правда?}
{``Is it true?''}

\ml{$0$}
{--- Правда.}
{``It is.''}

\ml{$0$}
{--- Тогда можно я поплачу?}
{``May I cry a little, then?''}

\ml{$0$}
{--- Не нужно спрашивать разрешения.}
{``No need to ask permission.}
\ml{$0$}
{Можешь даже покричать.}
{You may even scream.''}

Тахиро испустил дикий протяжный рёв боли, многократно отозвавшийся эхом в молчащем лагере, и из его глаз хлынули слёзы.
Рыдания рвались из сдавленного горла, словно обезумевшие птицы из клетки.
Грисвольд почувствовал, что парню необходима ласка, протянул руку --- и тут же отдёрнул, поняв, что причинит прикосновением ещё больше страданий.

<<Надеюсь, сегодня будет дуть свежий прохладный ветер>>, --- подумал технолог, глядя в тёмное небо.
На своей планете он знал, как поднять бриз и свернуть воздух в торнадо;
на этой приходилось довольствоваться надеждой.

На шум пришёл легионер.
Увидев Грисвольда, он вытянулся и отвесил короткий поклон:

--- Что-то случилось, сама?

--- Всё в порядке, легионер, --- сухо ответил Грисвольд.
--- Возвращайтесь на пост.

--- Вас понял, сама.

Легионер ушёл за склады, и ветер закружил сухие листья на голом камне.
Несколько мгновений мимо Тахиро и Грисвольда летел крохотный смерч, вбирая в себя щепки и пыль;
собеседники проводили его взглядами.
Тахиро суеверно выставил вперёд ладонь и что-то прошептал в крепко сжатый кулак.

--- Грис, можешь сделать меня демоном?

Технолог смутился.

--- Методика ещё сыровата, --- уклончиво сказал он.
--- Пока что я с командой провожу опыты на... на других людях.
Пожалуйста, не смотри на меня так.
Если эти опыты не проведу я, их проведут другие.

--- По-твоему, это оправдание? --- осведомился Тахиро.

--- Не нужно ставить мне это в вину, --- холодно буркнул Грисвольд.
--- Когда-то вы проводили опыты на мышах.
Сейчас человечество в мышах не нуждается.
Чем вы, люди, отличаетесь от этих зверьков?

--- У нас есть...

--- Ничем, --- перебил парня Грисвольд и поднялся на ноги.
--- Когда мы добьёмся первых стабильных результатов по оцифровке, я превращу тебя в демона.
Тогда --- и только тогда --- твоё слово будет что-то значить против слова Айну.

--- Вы за это заплатите, --- пробурчал Тахиро сквозь слёзы.
--- Я клянусь порогом моего дома, вы заплатите!

--- А кто привлечёт нас к ответу? --- криво улыбнулся Грисвольд.
--- Ты?

Тахиро промолчал.

--- А пока подумай над одной вещью, Тахиро, --- Грисвольд многозначительно поднял палец.
--- Что, если бы мыши могли привлечь к ответу вас, м?
Все тысячелетние рассуждения о справедливости человечество вело, находясь на вершине пищевой цепи.
Хорошенько над этим подумай.

--- Тогда я верну людей на причитающееся им место, --- запальчиво бросил Тахиро.
--- Ваше племя вершины пищевой цепи недостойно.

Грисвольд хмыкнул и пошёл прочь, проворчав себе под нос что-то неразборчивое.

\section{[:] Извинение}

\textspace

--- Грис, прости, --- тихо сказал Тахиро.
--- Я срываюсь на тебя иногда, я знаю...

--- Я не обиделся.

--- Ты хорошо на меня влияешь, --- горячо заговорил Тахиро.
--- Я чувствую, как после наших бесед с меня спадает вся эта дикость, в которой я рос и продолжаю жить сейчас.
Жестокость, презрение к людям.
Странно, но во мне этого куда больше, чем в тебе, хоть ты и не человек.

Грисвольд кивнул и углубился в свои записи.

\ml{$0$}
{--- Давай сбежим, --- вдруг зашептал Тахиро.}
{``Let's go far away,'' Tahiro suddenly whispered.}
\ml{$0$}
{--- Плевать на Орден.}
{``Screw the Order.}
\ml{$0$}
{Оцифруй меня, мы найдём планету и будем там богами, будем разбивать сады и возводить дворцы.}
{Digitize me, and we'll find a planet to be her gods, we'll lay out gardens and build palaces.}
\ml{$0$}
{Ты, я и Лу.}
{You, me, and Lu too.}
\ml{$0$}
{Я предпочту умереть, защищая последний приют, чем раздирать этот несчастный клочок камня на ресурсы.}
{I'd rather die for my last home, instead of tearing that piece of stone apart for resources.''}

Грисвольд невольно вздрогнул.
Это не были слова дзайку-мару;
юноша говорил как хорохито.

\ml{$0$}
{--- Твоя мечта исполнится очень скоро, ты умрёшь, --- сухо ответил технолог.}
{``Your dream will come true soon, you'll die,'' the technologist coldly answered.}
--- Ты же слышал слова Арракиса на последнем заседании штаба.
Союз Воронёной Стали действует так же.
С нейтралами сейчас разговор короткий.
Их скоро не останется, все демоны будут втянуты в войну.
\ml{$0$}
{Без Ордена мы не выживем.}
{Without the Order we can't survive.''}

\ml{$0$}
{--- У нас нет выбора?}
{``We have no choice, then?''}

\ml{$0$}
{--- Нет.}
{``We haven't.}
\ml{$0$}
{По крайней мере пока, --- Грисвольд позволил себе лёгкую улыбку.}
{At least not yet.'' Griswold smiled.}
\ml{$0$}
{--- Но я подумаю над твоим предложением.}
{``I can think over your proposal, though.}
\ml{$0$}
{А тебе лучше поспать.}
{And you could use some sleep.''}

\ml{$0$}
{--- Едва ли я сегодня засну, --- буркнул Тахиро.}
{``I can hardly sleep tonight,'' Tahiro said.}

--- Возьми из лаборатории стерильные простыни и попроси у Лу морфин.

\ml{$0$}
{--- Это зло.}
{``It's evil.}
\ml{$0$}
{Морфин превращает людей в животных, и даже одной дозы достаточно...}
{Morphine turns human into animal, and one dose is enough to---''}

\ml{$0$}
{--- Избавь меня от этих глупых суеверий, --- перебил его Грисвольд.}
{``Spare me those silly superstitions,'' Grisvold interrupted him.}
--- На наркоту садятся лишь рабы и люди, глухие к собственным желаниям.
\ml{$0$}
{Взгляни на Лу --- он курит адский вереск раз в неделю, а иногда месяцами трубку в руки не берёт.}
{Look at Lu: he smokes hell heather one time a week, and sometimes he doesn't touch a pipe for months.}
\ml{$0$}
{Просто потому что ему некогда, или потому что не хочет.}
{Just because he has no time, or no desire.}
\ml{$0$}
{Скажешь, он такой же наркоман, как те, в поселении, у которых от трубки уже губы гниют?}
{You think he's a junkie like those, from the village, who have rotten lips because of pipe?''}

\ml{$0$}
{--- Отец говорил, что все, кто хоть раз вдохнул дым --- наркоманы.}
{``My father used to say: anyone who ever breathed the smoke is a junkie.}
\ml{$0$}
{Мой дед...}
{My grandfather---''}

\ml{$0$}
{--- Твой дед познал удовольствие.}
{``Your grandfather have known the pleasure.}
\ml{$0$}
{Твой отец осознал опасность.}
{Your father have known the danger.}
\ml{$0$}
{Тебе предстоит увидеть закономерность, взвесить пользу и вред.}
{And you're the one to understand the patterns, weigh the pros and cons.}
\ml{$0$}
{Именно так люди и переходили от веры к знанию.}
{That's how people go from faith to science.''}

\ml{$0$}
{--- Откуда ты знаешь, что это предстоит именно мне?}
{``How do you know I'm the one?''}

\ml{$0$}
{--- Я не знаю никого другого, кто на это способен.}
{``I know no one else who can do that.}
\ml{$0$}
{Возьми у Лу морфин и поставь капельницу.}
{Get some morphine at Lu's, and put in an IV.}
\ml{$0$}
{От одного раза ничего не будет.}
{Nothing's gonna happen because of one dose.''}

\chapter{[U] Побег}

\section{[U] Просьба о прощении}

\epigraph
{
\ml{$0$}
{Чему учит нас история Тси-Ди?}
{What the history of Qi-Di teaches us?}
\ml{$0$}
{Всегда оставляйте своим врагам шанс вести достойное существование.}
{Always give your enemies a chance to make a decent living.}
\ml{$0$}
{Любая идеология бессильна против отчаяния обречённых.}
{Every ideology is powerless against desperation of the doomed.}
\ml{$0$}
{Великая Матка сказала, что ни один апид-солдат не вернётся домой, пока Молокоеды и Тараканы не будут истреблены.}
{Huge Queen said: no Soldier of Apis will return home until Milkers and Cockroaches are wiped out.}
\ml{$0$}
{<<Да будет так>>, --- ответили тси и втоптали Ульи в песок.}
{``So be it,'' said Qi and stomped Hives in the dust.}
}{Кельса Пушистая}

\textspace

--- Трукхвала убили, --- сказал Грейсвольд.
--- Короля-жреца.

Я выругался.

--- Кто?

--- Наши, разумеется.

--- Я же сказал им, что с Трукхвалом проблем не будет!
Я проинструктировал его, чтобы он сложил с себя обязанности по первому требованию!

--- Боюсь, --- печально улыбнулся Грейсвольд, --- у него ничего и не требовали.

Я сделал несколько кругов по комнате.

--- Ты уверен, что это убийство?

--- У него на теле было несколько необычных кровоподтёков.
Ты знаешь, наши работают чисто, для сапиентов всё выглядело как обычный несчастный случай.
Однако кровоподтёки... он щипал себя перед самой смертью.
Ранить себя он не мог --- убийцы почуяли бы кровь и заинтересовались.

--- Это были просто щипки или...

--- Или символ, да.
Два знака письменности тси --- <<уголок>> и <<цепочка>>.

Модификатор пассивного залога и конец предложения.
<<Что-то, результат чего ты наблюдаешь, было сделано кем-то>>.

--- Он пытался оставить послание жрецам?

--- Неясно, кому предназначалось послание, но он хотел, чтобы об убийстве узнали.

--- От кого исходил приказ?

--- Штрой Кольцо Дыма, командующий силами Ада.
Был ли это её приказ или свыше, не знаю.

--- Давай, Минь, объясни мне ещё и это! --- в сердцах обратился я к архивариусу.
--- Этот человек поднял земли сели, идолов и хака из руин.
Он в жизни не причинил вреда даже птице, не говоря уже о сапиенте.
Однако Ад уничтожает его --- мягкого, покладистого лидера, который вёл свой народ к процветанию!

--- С этим не всё так просто, --- пробормотал Минь.
--- Да, демоны Ада длительное время поддерживают процветание сапиентов, чтобы получать плюс-эманации.
Но дело в том, что неограниченная эволюция невыгодна --- существует опасность восстания против Ада.
Тси-Ди ярко это показала --- тси не согласились на власть демонов.

--- И каким же образом Ад решает эту проблему? --- глухо проворчал я, зная ответ наперёд.

--- Апокалипсис.
Полное низвержение цивилизации, доведение сапиентов до первобытного состояния.
Затем цикл повторяется.

--- И даже это запланированное событие используется в целях пропаганды, --- добавил Грейсвольд.
--- Угадай, кого в нём обвинят.
Обычная история.

--- Так значит, именно поэтому Ад настороженно отнёсся к генофонду тси, --- проговорил я.
--- Тси улучшили свои тела, сделав их сверхприспособленными и быстрообучаемыми, и легко могут выйти из-под контроля.

--- Аду невыгодно существование потомков тси в их нынешней форме, --- резюмировал Грейсвольд.
--- Я не думаю, что тси будут уничтожать, но отрицательный отбор неизбежен.
Возможно, Трукхвал и не представлял непосредственную опасность как личность, но он определённо обладал опасными с точки зрения Ада генами.
Ты же помнишь про отбор Королей-жрецов?

Я встал и прошёлся по комнате, теребя пальцами браслет-терминал.
Выходит, все наши старания были зря.
Мы подвергали себя риску и отправляли сапиентов на верную смерть ради ещё одной рабской хижины размером с планету.

--- А я пообещал Митрису, что его народ будет счастлив.

--- Аркадиу, --- Грейсвольд, похоже, не на шутку испугался последней фразы, --- надеюсь, ты не собираешься сделать какую-нибудь глупость?

--- Разумеется, нет.

--- Тогда куда ты собрался?

--- Просить прощения, --- сказал я и сдавил браслет.
Бесшумно разъехались двери комнаты, и я вышел наружу, в переплетение стальных коридоров.

\asterism

Я знал, сколько времени требуется отделу 100, чтобы заняться мной вплотную.
Статистика по захватам агентов Картеля стоила мне больших затрат.
Тридцать секунд на обработку информации обо мне, столько же на принятие решения, минута на то, чтобы роботы контрразведки подобрались к моему телу.
Дальше ожидание и слежка, затем захват или уничтожение.

Минь уже должен был передать информацию.
Бедный архивариус вынужден был переступить через собственную гордость из-за чувства страха, не зная, что и это входило в мой план.

Вокруг стояла тишина.
Я бы даже сказал, что было чересчур тихо.
Скорее всего, контрразведка уже парила на расстоянии вытянутой руки, прикрытая всей мощью технологий Ада.

Я направился в архивы.
Существовал только один способ сбежать --- нужно было заинтриговать отдел 100 и убедить их, что мой побег принесёт для дела больше пользы, чем моя поимка.


\section{[U] Виа}

Лифт опустился на тридцать восьмой этаж.
Вокруг по-прежнему было подозрительно тихо, и я тихо вздохнул, словно озабоченный чем-то.
На самом деле это был вздох облегчения --- мои действия явно не укладывались в предсказанную аналитиками схему, и контрразведка медлила с поимкой.

Дверь отъехала в сторону, и я встретился глазами с молодой женщиной-плантом.
Архивариус Виа Серая Ласточка --- так её зовут.
Она молча улыбнулась и кивнула.
Разумеется, отдел 100 уже предупредил и проинструктировал её.

--- Здравствуй, Виа, --- поклонился я.
--- Мне нужна твоя помощь.

--- Секунду, --- сказала она и, приложив палец к виску, закрыла глаза.
Открыла.
--- Извини, Аркадиу.
Перенаправила срочный запрос.
Слушаю тебя.

--- Мне нужна как можно более точная временная лента последнего сантителльна существования цивилизации Вилет'марр.

--- Лотос? --- Виа смотрела на меня с интересом, но я обострёнными чувствами ощущал некоторую наигранность.
Логика понятна --- люди Тра-Ренкхаля были потомками вилет'марр, и аналитики контрразведки предупредили архивариуса о таком направлении беседы.

--- Да.
Дело в том, что некоторые аналитики...

--- Некоторые --- это кто?
Выражайся точнее, Аркадиу.
Я не могу тратить время на проверку источников.

Провокация.
Я был к ней готов, но знал, что мой ответ не идеален.
Если Виа или держащие с ней связь аналитики за ближайшие двадцать секунд найдут, к чему придраться --- пиши пропало.
Они не знали, что я осведомлён о подключении к делу контрразведки, и не могли затягивать беседу --- это выдало бы их с головой.

--- Гимел Кадмиевый Зелёный, диктиолог Картеля, действовавший...

--- Он погиб на Тра-Ренкхале во время твоего последнего задания? --- прикрыв глаза, спросила Виа.

--- Да, я его убил.

На самом деле Гимел погиб во время состязания грубой силы между Безымянным и Эйраки, но знал это только я.

--- Ты поддерживал с ним связь? --- поинтересовалась Виа.

--- Информация получена в момент его уничтожения.

--- Почему сказанное тобой расходится с отчётом? --- снова прикрыла глаза Виа.

--- А почему я пришёл к тебе сам, а не написал запрос по общему каналу? --- задал я риторический вопрос.

Виа растерялась и немного испугалась.
Она поняла, что дело действительно серьёзное --- если бы не предупреждение контрразведки, она бы и не поняла, что перед ней шпион.
Всё укладывалось в привычные рамки --- мои человеческие повадки, мои полномочия как легата терция не раскрывать некоторую информацию до окончания собственного расследования и даже промедление в пять лет, которое было призвано охладить интерес агентов Картеля к моему последнему заданию.
Плюс репутация --- Виа изучала досье большинства демонов и никак не могла ожидать от меня подвоха.

--- Хорошо.
Чем располагал Гимел?
Я тебя слушаю, --- Виа справилась с собой, и в её глазах загорелся спокойный огонь готового к работе профессионала.

<<Пока неплохо>>.

--- Гимел вычислил время прибытия людей Лотоса на Тра-Ренкхаль.
Использовались данные археологии, астрономии и генетики.
Точность весьма высока --- плюс-минус два года.

Археология у Картеля всегда была своеобразной --- образцы оцифровывались и тут же уничтожались, чтобы их не изучили враги.
Культурологи Ада, среди которых были и оцифрованные сапиенты, настояли на запрете подобного вандализма.
Разумеется, без крайней необходимости.
А биологи Ада, благослови лесные духи их расхлябанность, увлеклись изучением тси и девиантной фауны, отложив в долгий ящик отчёт о людях Тра-Ренкхаля.
Я был свидетель тому, как начальника отдела распекал прибывший чин за то, что биологи не работают над проблемами совместимости демонов с местными сапиентами.
Тот смущённо оправдывался, мол, на планете чересчур много интересного.
В итоге отдел настрочил под диктовку Безымянного отписку для Капитула, и на этом изучение людей Тра-Ренкхаля завершилось.

--- Поняла, --- сказала Виа.
--- Временная лента готова.
Ты хочешь ознакомиться с ней сам или я тебе ещё нужна?

Я улыбнулся и решил ещё раз погладить отдел 100 против шерсти:

--- Если ты не занята, то я хотел бы разобрать её с тобой.

Виа несколько растерянно махнула мне рукой, призывая следовать за ней.

\section{[U] Ветер свободы}

Вскоре мы сидели в пустом кабинете перед экраном компьютера.
Виа непроизвольно вжималась в кресло, по её лбу струился пот;
похоже, что здесь был агент контрразведки, прикрытый щитом невидимости, и женщина опасалась ненароком его задеть.
За подобную неосторожность её могли лишить должности.
Эти случаи не афишировались, но о них знали все --- почему-то решающим в вопросе профпригодности демона было слово не его коллег по отделу, а военных.

--- Скажи мне как аналитик, в какое время вилет'марр могли послать экспедицию на Тра-Ренкхаль? --- спросил я.

--- Не раньше 1190-го оборота, --- не задумываясь ответила Виа.
--- Ранее вилет'марр просто не располагали нужными технологиями.
Я бы даже отодвинула эту границу оборотов на сорок --- шестьдесят, с учётом психологии этих людей --- они тяжеловаты на подъём.

Виа осторожно потянулась к плечевому импланту.

--- Ты не против?

--- Нет, конечно, --- я улыбнулся и сам сделал ей инъекцию успокоительного.
--- Тяжёлый день?

--- Если бы только день, --- Виа жалобно улыбнулась.

--- Я тебя понимаю.

--- Ты один из немногих, кто может понять.

--- А сколько времени потребовалось бы кораблю, чтобы достичь цели?

Виа почесала кожистую зелёную головку, раскрыла и сложила венчик.

--- Между Тра-Ренкхалем и Лотосом лежит пылевой <<блин>>, напрямую лететь они не могли.
Дело в том, что первые люди отправляли экспедиции не на конкретную планету, а на так называемое направление.
Какие-то из преобразованных планет на направлении могли быть непригодны для жизни, какие-то могли стать непригодными за время полёта, всякое происходило.
\ml{$0$}
{Иногда, уже прилетев на вполне годную планету, люди наотрез отказывались высаживаться, и были в своём праве.}
{People have sometimes reached a good planet and then refused to land on, and it was their right to do so.''}

\ml{$0$}
{--- И такое было?! --- удивился я.}
{``Has this really happened!'' I got surprised.}

\ml{$0$}
{--- Разумеется!}
{``Of course!}
\ml{$0$}
{Планета --- это не придорожная гостиница, а дом для ближайших ста поколений потомков!}
{Planet is not a roadside hotel, but home for the next hundred generations of descendants!}
Не все нюансы можно рассмотреть с расстояния в сто парсак, даже с помощью демиурга.
Поэтому, если маршрут предполагал, например, пять преобразованных планет, топлива брали на шесть циклов.

--- Чтобы был шанс вернуться, --- кивнул я.

--- Хотя бы маленький, но чтобы был, да.
Это очень важный момент, Аркадиу.
Первая экспедиция на направление была без страхового цикла.
Они успешно добрались до цели, но пригодной оказалась лишь последняя планета на маршруте.

--- Я понимаю.
Психологическое напряжение было гигантским.

--- Мягко сказано.
На корабле начался массовый психоз, закончившийся бунтом.
Десятую часть команды пришлось изолировать до конца полёта, многие из них погибли.
Так что сам факт, что экспедиция оказалась так далеко от Лотоса, за пылевым облаком, говорит именно об этом --- Тра-Ренкхаль был не первой и даже не третьей планетой, которую они посетили.
Дай-ка компьютер, я попробую вычислить маршрут.

Я уступил Виа место.
Плант прищурилась, и в её красивых чёрных глазах заиграли голубые огни голографического дисплея.

--- В архивах точно нет информации о времени отлёта?

--- Нет, --- покачала головой Виа.
--- Я бы и сама была рада, если бы подобная информация сохранилась.
Вот, смотри.
Кратчайшее время полёта --- триста восемьдесят семь оборотов.
Это на максимальной скорости.
Если предположить, что по пути они останавливались для починки, дозаправки или ради научного интереса --- выйдет все шестьсот.

Я кивнул.
Виа посмотрела на меня.
Страха в её глазах больше не было --- только любопытство.

--- Я, конечно, понимаю, что это секрет, но всё же --- зачем тебе это?
Ты думаешь, что вилет'марр были где-то ещё?

--- Я \emph{знаю}, что они были где-то ещё, --- сказал я.
--- Если твои данные верны, то вилет'марр двигались со сверхсветовой скоростью.
Точнее, в две световых.
И это при условии, что они, как ты говорила, никуда не заворачивали.
В чём лично я сомневаюсь.

Пластик под пальцами Виа жалобно хрустнул.
Она начала с бешеной скоростью открывать какие-то документы и просматривать их.

--- Да не может этого быть, --- шептала она.
--- Вот это --- данные Ясмер.
А это я лично вычисляла, и команда Флавика сказала, что всё верно.
Все данные, чтоб мне сдохнуть, верифицированы, и не на один раз.

Я промолчал.
Виа перестала мучить компьютер и откинулась в кресле.

--- Ну не могли же они использовать омега-телепортацию! --- жалобно воскликнула она.
--- Перенос сознания при омега-телепортации невозможен, даже с синхронизацией мозг-демон-мозг.
На расстояние вытянутой руки через квантовую запутанность --- куда ни шло, но всё, что дальше, равносильно самоубийству.
Запрет на ОТ шёл чуть ли не со времён ранней Эпохи Богов!
Зря, что ли, первые люди преодолевали космос на субсветовых?

--- Успокойся для начала, --- бросил я.
--- Гимел тоже мог ошибиться.
Верификацию, насколько я понял, прошли только используемые им данные анализов.
Алгоритм, согласно которому он проводил вычисления, мне не известен.
Но меня беспокоят тси.
Случайно ли то, что некогда обладавщие высокими технологиями вилет'марр, а затем и тси в итоге оказались на одной планете?
Что, если эти две цивилизации поддерживали связь, и тси прилетели за помощью?

Виа смотрела на меня.
Я знал, что её демон уже нарисовал нужную мне картину.

--- Это значит...

--- Это значит, что существование мультипланетных союзов сапиентов возможно не только в Империи Плеяды, --- закончил я.
--- Разумеется, в наше время это больше любопытный факт, чем что-то полезное.
Сложно допустить мысль, что на территории фракций может возникнуть законспирированная сапиентная организация, верно?

--- Я не думаю, что контрразведка допустит подобное, --- как-то неуверенно согласилась Виа.

--- Так или иначе, история тёмная и требует анализа, --- заключил я.
--- Я отправлюсь на Лотос и попробую что-нибудь узнать об экспедиции.
Я почти уверен, что наши демоны просмотрели очень важные детали.

Виа неловко улыбнулась.

--- Желаю удачи.
Ты знаешь, что искать?

--- Да, --- ответил я и пошёл к выходу.
--- В памяти Гимела были ещё кое-какие зацепки.

Виа махнула мне рукой.
Уже у порога я услышал за спиной тихий-тихий вздох облегчения.

<<Всё-таки неплохо играть на пограничных значениях.
Сместил среднее, увеличил точность --- формально никаких противоречий, а картина кардинально поменялась.
Надеюсь, ты достигла пристанища духов, Айну.
Ты была замечательным учителем>>.

... Дверь лифта закрылась, но никто так и не попытался меня пленить.
Я почувствовал, как мои волосы нежно ласкает ветер свободы --- давно забытый и такой желанный.

\chapter{[:] Другая форма жизни}

\section{[:] План раскрыт}

--- Ты собираешься оцифровать Тахиро.

Грисвольд вздрогнул и выронил ключ.

--- Откуда ты это узнал?

--- Я стратег, Грис.
Моя задача --- распознавать схемы и структуры в чужих действиях.
И задача эта --- не только моя.

--- Кто ещё знает?

--- Благодаря Гало --- весь штаб Ордена.
Как ты знаешь, решение об оцифровке дзайку-мару ещё не было принято, и Гало, который всецело <<против>>, чрезвычайно заинтересовался совершенно отчётливыми приготовлениями к этому событию.
Я бы на твоём месте поспешил на заседание, которое проходит прямо сейчас.

\ml{$0$}
{--- Они начали его без меня?!}
{``Did they start without me!''}

--- Ты же занят сборкой транспорта, верно?

Грисвольд судорожно выдохнул.

\ml{$0$}
{--- Ты бы хоть меня предупредил.}
{``You might have warned me.}
\ml{$0$}
{Я едва успел убедить Гало, что ответственность за всё лежит на мне.}
{I barely had time to convince Halo that the responsibility is all mine.}
\ml{$0$}
{Наплёл ему красивую историю про то, как хотелось бы мне сделать демоном своего человеческого друга.}
{I told him a lovely story: how I wish I could turn my human friend into a daemon.''}

\ml{$0$}
{--- Ты подставил Тахиро?}
{``You framed Tahiro?''}

\ml{$0$}
{--- А что мне ещё оставалось сказать, Грис?}
{``So what could I say, Gris?}
\ml{$0$}
{Давай, придумай что-нибудь более правдоподобное, ты же у нас гений конспирации.}
{Come on, think of something more believable, you tradecraft genius.}
\ml{$0$}
{Меня временно отстранили от дел за самоуправство, Гало меня ненавидит, а Тахиро остаток жизни будет спать со мной в одной комнате и в одной кровати, потому что это единственное место, где у него есть шанс проснуться.}
{I was suspended because of excess of authority, Halo hates me, and Tahiro will sleep in my room and in my bed the rest of his life, because it's the only place where he has chance to wake up.''}

\ml{$0$}
{--- Ч-что?}
{``W-what?''}

\ml{$0$}
{--- Грис, я сохранил твою репутацию, но его я больше не смогу защищать.}
{``Gris, I saved your reputation, but I can't save him anymore.}
\ml{$0$}
{Свод законов Ордена никак не защищает людей.}
{The Order law does not protect humans.}
\ml{$0$}
{Люди --- это скот, чья участь --- дойка, случка и забой.}
{Humans are cattle destined for milking, breeding, and slaughter.}
\ml{$0$}
{Тахиро жив по той причине, что он --- моя собственность, не представляющая опасности.}
{Tahiro is alive because he's my property, presenting no danger.}
\ml{$0$}
{Став демоном, Тахиро из умной игрушки превратится в сильного игрока.}
{After daemonization Tahiro goes from smart toy to strong player.}
\ml{$0$}
{Его прирежут превентивно, лишь бы этого не случилось.}
{They will cut his throat preventively to avoid that.''}

\ml{$0$}
{--- Что мне делать?}
{``What must I do?''}

\ml{$0$}
{--- Ты должен признаться во всём и хотя бы попробовать отстоять свою точку зрения.}
{``You must confess all, then at least try to prove a point.}
\ml{$0$}
{Идею оцифровать именно Тахиро, равно как и карт-бланш на все действия, можешь повесить на меня.}
{You can say I'm the one who gave you the idea to digitize Tahiro, along with carte blanche to do all you did.}
\ml{$0$}
{Ваша с Тахиро дружба должна оставаться в тайне --- то, что позволительно чудаку Люциферу, тебе не простят.}
{Your friendship with Tahiro should remain a secret---they won't forgive you for things that are permissible for freak Lucifer.}
\ml{$0$}
{Жизнь мальчика в твоих руках.}
{The boy's life is in your hands now.}
\ml{$0$}
{Единственный способ его спасти --- сделать субъектом адского права.}
{Turning him into a subject of the Hell law is the only way to save him.''}

Грисвольд сорвался с места с неподобающей толстяку прытью и понёсся к выходу из ангара.

--- Лу, прошу тебя, спрячь хотя бы вон те два модуля!
Диапазон моих личных ячеек --- десять-четырнадцать!
Только роботов не используй, они логи хранят!

--- А как тогда?..

--- Тележка!
Гидравлическая тележка в углу!
Перекати механизмы на палету и подцепи тележкой!

Люцифер с грустью осмотрел два тяжеленных агрегата, затем бросил взгляд на свои холёные пальцы --- и испустил тяжкий вздох.

--- Вот какого дьявола я опять должен таскать тяжести? --- обратился он к геликоптеру.
Геликоптер сочувственно промолчал.

\section{[:] Совещание}

\textspace

Арракис потребовал, чтобы Грисвольд объяснил свои действия.
Фраза была, разумеется, формальной --- по сути Грисвольду предстоял завуалированный публичный допрос.
Согласно правилам отдела пропаганды, подобная формулировка должна была употребляться по отношению к высокоранговым демонам, чтобы не нарушать субординацию;
однако по сути <<объяснения>> могли потребовать у любого, чьи действия стояли на самой грани закона, точнее, интересов руководства.
<<Объяснение действий>> подразумевало вопросы произвольного демона из исследовательских отделов под неусыпным контролем военных стратегов.

По странному стечению обстоятельств, произвольным демоном оказался один из учёных бывшего Ордена Тысячи Башен, Север Солнечная Дева.
Он сразу обращал на себя внимание грамотным построением речи на сохтид, тихим робким голосом и удивительно естественной мимикой.
У него единственного из всех присутствующих были очки с простыми стёклами и аккуратная подстриженная борода.
Согласно досье, Север по праву считался одним из флагманов изучения людей Тысячи Башен и предпочитал, чтобы его называли <<человеческим>> именем --- Сальватор Сольмё.

Разумеется, Север был проинструктирован заранее, как следует вести беседу.
Но по какой-то причине --- скорее всего, от волнения --- Север перевёл разговор на более привычные для учёного технические подробности.

--- Грисвольд, объясните всем, пожалуйста, в чём сложность демонизации людей.

--- В демонизации как таковой сложности нет, это стандартный процесс снятия квантовой структуры объекта с последующей записью в омега-модуль.

--- Какой именно объект подвергается оцифровке?

--- Нервная система вплоть до чувствительных окончаний, нервно-мышечных синапсов и некоторых нейрогуморальных структур.

--- Почему необходимо снятие отпечатка всей нервной системы?

--- В модуляции сигналов участвуют все указанные звенья.
Головной мозг, даже вместе со спинным, не самодостаточны.
Разработанный нами интерфейс монтируется к цифровым аналогам нервно-мышечных синапсов и чувствительных окончаний.

--- Я прошу прощения, коллега, но известны случаи, когда люди управляли машинами с помощью нейроинтерфейса, подключённого напрямую к коре, --- Север вопросительно посмотрел на Грисвольда поверх очков.

--- Сравнение некорректное, Сальватор, --- поклонился Грисвольд.
--- Речь шла о подключении устройства к живому организму, а не о полной замене живого организма квантовой моделью.
В упомянутом вами нейроинтерфейсе сигнал с коры подвергается сложному процессу стабилизации и модуляции сигнала, как это происходит в нижних отделах нервной системы.

--- Благодарю за пояснение.
Каким образом обеспечивается функционирование квантовой модели?

--- Модель помещается в среду, поддерживающую виртуальный клеточный метаболизм.
В будущем можно преобразовать квантовую модель мозга в самодостаточную математическую модель.
Если мы получим некоторое количество полноценных оцифрованных людей, мы сможем постепенно оптимизировать части виртуального мозга.
Я не думаю, что возможна замена всех отделов, но даже оптимизация ствола мозга даст огромный прирост производительности.

--- Прошу прощения, вынужден прервать, --- поднял руку Самаэл.
--- Вы чересчур углубляетесь в технические подробности и...

--- Извините, это чрезвычайно интересный момент, --- перебил Север.
--- То есть вы хотите сказать, что к модели нервной системы можно пристыковать вычислительные модули, модули чувствительности и устойчивости?

--- Я переформулирую ваш вопрос и отвечу на него: да, квантовая модель мозга человека совершенно точно может взять на себя роль программного ядра омега-сингулярности, со всеми вытекающими последствиями.

Присутствующие зашептались.

<<Самаэл очень не хотел, чтобы разговор дошёл до этой фразы, --- ухмыльнулся Грисвольд.
--- Пусть они это проглотят>>.

--- Надеюсь, вы не думаете, что это ядро сравнимо по эффективности с ядрами демонов? --- скептически осведомился Гало.

--- На данном этапе сравнивать довольно сложно.
Безусловные преимущества оцифрованного мозга --- дешевизна, уникальность строения, сложность прямого взлома.
Недостатки --- сложность настройки и гораздо меньшая пластичность.
Проще говоря, это дешёвый одноразовый аналог демона --- с неплохим запасом устойчивости и интеллекта, но с низким потолком чувствительности и проигрывающий в быстродействии.

--- Проигрывающий на два порядка, --- присовокупил Гало.

--- Два порядка? --- удивился Север.
--- Это очень много, Грисвольд, особенно применительно к военной сфере.

--- Тем не менее некоторые наши специалисты по вооружению считают, что с таким запасом устойчивости это простительный недостаток.
Кроме того, совершенно необязательно делать из оцифрованных людей солдат.
\ml{$0$}
{Как вы знаете, Сальватор, есть множество профессий, которые не требуют скорости.}
{You know, Salvator, there are many professions that don't need speed.}
\ml{$0$}
{Я своё пузо отрастил не просто так.}
{My belly's grown too fat for a reason.''}

Шутка предназначалась Гало.
Но стратег даже не улыбнулся.
В его глазах светился гнев;
Грисвольд вдруг понял, насколько непривычно видеть эту эмоцию на копии кукольного личика Люцифера.
Зато Север, к удивлению Грисвольда, отреагировал на юмор вполне адекватно.

\ml{$0$}
{--- Я полагаю, у вас масса опыта, --- с улыбкой сказал он, поправляя очки.}
{``You've got tons of experience, I guess,'' he said with a smile and readjusted his glasses.}
--- А какова скорость естественного износа квантовой модели?

--- Примерно в пять раз медленнее таковой в живом человеке.
Согласно прогнозам, износ математической модели можно довести до шестидесятитысячекратного замедления.

--- То есть, Грисвольд, вы хотите сказать, что максимальное время жизни оцифрованного человека приблизится к времени жизни демона?

--- Именно, Сальватор.

--- Это абсурд! --- возмущённо выкрикнул кто-то из учёных.
--- Люди живут не более восьмидесяти оборотов, их интеллектуальный потенциал охватывает вчетверо меньший срок и является достаточным исключительно для того, чтобы особь выжила и оставила потомство.
А вы хотите дать им в руки вечность и сложнейшие модули, в равной степени недоступные для их понимания!

\ml{$0$}
{--- Единственное, что хочу я --- это использовать возможность и исследовать объект, насколько это возможно, --- с достоинством ответил Грисвольд.}
{``The only thing I want is to use the opportunity, as well as possible, and research the object, as well as possible,'' Griswold said with dignity.}
--- Я не являюсь учёным, но не думаю, что это можно ставить мне в вину.
Если кто-то сомневается в наших данных и нашей методологии, вы всегда можете провести соответствующие исследования самостоятельно.

--- Полагаю, я буду первым, кто этим займётся, --- кивнул Север.
--- Какие задачи стоят перед вами сейчас?

--- На данный момент единственная большая задача --- отработка методики.
Многое зависит от особенностей конкретного человека, его способностей и психологических параметров.
Из некоторых получаются живые, но абсолютно неработоспособные демоны с узким каналом ввода-вывода.
У других канал сохранен, но они надламываются психологически или сходят с ума.
Таких мы вынуждены уничтожать, потому что они опасны не только для себя.

--- Можно ли решить это тренировкой?

--- Разница статистически незначимая.
У некоторых исследователей создалось впечатление, что значение более имеет мотивация, нежели тренировка.
Наши подопытные, разумеется, шли на процесс под угрозами, шантажом и обманом.

--- Каков процент полноценных выживших?

--- Два процента.

--- Есть ли какая-то разница по возрасту и полу?

--- Выживаемость детей выше в десять-пятнадцать раз, но после оцифровки они останавливаются в развитии.
Вероятно, это можно решить, но не на данном этапе исследований.
Мужчины дают полноценных демонов чаще женщин, при этом так же чаще женщин сходят с ума.
Разница между представителями разных этносов статистически незначима.

--- Я рискну предположить: не связана ли выживаемость детей с мотивацией?

--- Возможно, Сальватор.
Мы рассказывали им сказку о том, что они станут героями после процедуры.
Многие сами ложились в капсулу и терпеливо переносили неприятные ощущения.

--- Ближе к делу, Север, --- напомнил Самаэл.

--- Прошу прощения, --- кивнул учёный.
--- Итак, Грисвольд, кхм, вы полагаете, что человек по имени Тахиро та Ханаяма --- в некотором роде идеальный объект для следующей стадии экспериментов?

--- Он приближен к идеалу настолько, насколько это возможно.
Его мотивация не вызывает сомнений.
Он успешно интегрировался в наше общество и даже выполнял задания, соответствующие рангу второго центуриона.

--- Он не мог <<успешно интегрироваться>> в наше общество по чисто техническим причинам, --- заявил Гало.
--- Он не воспринимает омега-поле, он не может работать с нашими устройствами и...

--- Как раз-таки с устройствами Тахиро работает превосходно! --- внушительным голосом перебил Грисвольд.

--- Простите, а у вас есть какие-то доказательства?

Грисвольд вместо ответа вытащил из кармана испачканные в земле ажурные браслеты и передал Северу.
Тот молча поднял очки на темя и рассмотрел украшения.

<<Да простит меня Тахиро, --- подумал Грисвольд.
--- Надеюсь, никто не расскажет ему, что я разорил могилу его невесты>>.

--- Это потрясающе, --- сказал Север, потирая бороду.
--- Если глаза меня не обманывают, это анодированный титан со вставками голографического серебрения.
Да, да, коллеги, я слышу ваши просьбы, я оцифровал, модель продемонстрирую после заседания.
Паутинная конструкция в просвете идеальна, браслеты почти невесомы, но при этом имеют хороший запас прочности.
Для проектировки и создания совершенно точно использовался либо демон, либо, скажем, какой-то ЧПУ с омега-консолью.
Уровень обработки запредельный для человеческого существа.

--- Строго говоря, мы оставили на ЧПУ аналоговый интерфейс, но да, это устройство, предназначенное для демонов.

--- Сколько времени потребовалось Тахиро, чтобы создать этот шедевр?

--- Согласно логам, около четырёх часов.

--- Это потрясающе, --- Север вертел браслеты в руках.
--- Мне действительно сложно в это поверить.
Вы не против, я возьму один?

--- К сожалению, мне придётся вернуть эти браслеты на место, Север.
Я взял их без спроса.
Но вы сами видите, о чём я говорю.
Тахиро способен управлять омега-устройствами.
Это может стать прецедентом --- будущих демонов можно серийно выращивать из сапиентного материала, избегая таким образом <<проблему копий>>.

По залу неожиданно пробежал шёпот;
в нём читалась смесь надежды и ужаса.

\ml{$0$}
{--- Он серьёзно?}
{``Seriously?}
\ml{$0$}
{Проблема копий может быть решена?}
{Problem of copies can be solved?''}

\ml{$0$}
{--- Он серьёзно?}
{``Seriously?}
\ml{$0$}
{Запустить серийную оцифровку этих животных?}
{He's offering serial digitization of these beasts?''}

--- У меня просьба к руководству Ордена, --- Север повернулся к Арракису.
--- Я, если честно, впервые услышал об этом исследовании только здесь.
Люди --- это моя жизнь.
Мне бы очень хотелось присоединиться.
Я знаю о своём нынешнем статусе, поэтому прошу в присутствии всех --- сделайте мне одолжение.

--- Если исследование решено будет возобновить, не вижу причин вам отказывать, --- кивнул Самаэл.
--- Подавайте заявление в соответствии с протоколом.
Ещё вопросы будут?

--- У меня вопросов больше нет, --- поклонился Север и сделал шаг назад.

--- Всё это очень трогательно, --- сказал Гало, --- однако я обращаю внимание присутствующих на то, что люди --- враждебный вид.
Если всё сказанное Грисвольдом --- правда хотя бы на одну десятую, серийная оцифровка может привести к масштабным и непредсказуемым последствиям вплоть до военного столкновения.

--- Вы позволите, я отвечу? --- улыбнулся Грисвольду Север.

--- Да, конечно, --- кивнул технолог.

<<Благослови тебя небо, Север.
\ml{$0$}
{Ты сам не понимаешь, как ты нам сейчас помог своей святой наивностью>>.}
{You have no idea how that \textit{sancta simplicitas} of yours has helped.''}

--- Когда ещё я служил в Ордене, мы...

--- В каком Ордене? --- ледяным тоном перебил Самаэл.

--- В Ордене Тысячи Башен, --- поправился Север.
Его смуглое лицо залила краска.
--- Так вот, мы с коллегами ради интереса пытались обучить культуру крысиных нейронов имитировать выполнение произвольного императивного кода.

\ml{$0$}
{--- Очень интересно, --- буркнул Гало.}
{``Sounds intriguing,'' Halo grunted.}

--- Да, Гало, я с вами согласен, --- улыбнулся Север.
--- Результаты, которые мы получили, перекликаются с вашими результатами: биологическая нейросеть достаточно пластична, чтобы заменить собой участок императивного кода.
Коллега Грисвольд утверждает, что биологическая нейросеть может стать даже программным ядром демона, вопрос, как я понимаю, лишь в мощностях и обучающей выборке.
Рискну с ним согласиться.
Я думаю, после астрономического числа итераций различия оцифрованного человека и демона сгладятся до статистически незначимых...

--- Я не об этом, --- вежливую речь Севера Гало слушал, скрестив руки и отвернувшись к стене.

--- Прошу прощения, Гало, я начал чересчур издалека.
Я полагаю, если взять человека, изначально воспитанного в обществе...

--- Если вы имеете в виду Тахиро --- он не подходит, --- бросил Гало за спину.
--- Мальчишка полжизни провёл в деревне, и воспитывали его дзайку-мару.

--- Я прошу прощения, Гало, --- стушевался Север, --- я не совсем понял ваш термин.
Вы не могли бы дать определение?

Гало свирепо задышал, тиская пальцами кобуру.
На мгновение Грисвольд подумал, что стратег вытащит оружие и застрелит Севера.

--- Гало, подожди, --- Грисвольд предостерегающе махнул Северу.
--- Скажи прямо: ты сомневаешься в Тахиро?
Люцифер дал ему положительную характеристику.

--- Тахиро уже пытался вести подрывную деятельность в Ордене, --- заявил Гало.

--- Всё это отмечено в отчёте Люцифера.
Тахиро был ребёнком, потерявшим семью.
Импульсивность нормальна для человеческих детей.

--- Мои тесты ясно показывают, что он продолжает чувствовать связь со своим видом.
Где гарантии, что это не повторится, едва в распоряжении Тахиро окажутся методы и вычислительные мощности?

--- Я считаю, что этому человеку следует дать шанс, --- сказал технолог.
--- Мы можем, разумеется, начать искать возможных предателей.
Вычислять, ставить опыты, проводить тесты, гадать на камнях.
Но единственный верный способ вычислить предателя --- это дать ему в руки оружие и встать с ним плечом к плечу.
И даже тогда можно ошибиться.

--- Я думаю, что Грисвольд прав, --- сказала Айну.
--- Гало и Люцифер, безусловно, мастера стратегии, но всё-таки теоретики и недооценивают значение опыта.
Тем не менее, несмотря на высказанное ранее мнение, что Тахиро уже принимал на себя полномочия второго центуриона, я категорически против присвоения ему этого ранга.
Пусть начнёт легионером, как и все демоны в Ордене.
Время покажет, чего он стоит.

--- Я категорически против, чтобы Тахиро был солдатом вообще, --- заявил Грисвольд.
--- Ты меня неправильно поняла.

--- Почему, Грис?

--- По уже озвученной Гало причине --- я против приобщения людей к военным технологиям.
Возможно, его захотят принять в какой-то исследовательский отдел?

--- Разве что как подопытный образец, --- буркнул Гало.
В толпе представителей раздался приглушённый смех.

--- Мне было бы интересно, --- улыбнулся Север, но улыбка тут же увяла под ледяным взглядом Гало.

--- К сожалению, вы не имеете права набирать собственный отдел, Север, --- сказал Самаолу.
--- Разумеется, если ваше рвение не угаснет, это временно.
Дождитесь решения ранговой комиссии.

--- Лично мне всё ясно, --- буркнул Арракис.
--- Подготовьте необходимое оборудование, этот человек будет оцифрован завтра и --- в случае успеха --- произведён в легионеры.
Заседание окончено, все свободны.

\section{[:] Нарэ}

<<Не самый плохой вариант>>, --- мысленно заключил Люцифер и углубился в созерцание карты стратегии.

\ml{$0$}
{--- Тебя отстранили от дел, --- напомнил Грисвольд, увидев Люцифера за обычной работой.}
{``You were suspended,'' Griswold reminded when he saw Lucifer doing his normal work.}

\ml{$0$}
{--- Я знаю.}
{``I know.''}

Грисвольд улыбнулся, неожиданно отвесил короткий поклон и вышел, аккуратно закрыв за собой дверь.
\ml{$0$}
{Он понял.}
{He realized.}

Карта Лу была простой и понятной.
Никаких путей отступления, никаких фронтов атаки.
Все вариации строились на том, что из каждого развития событий можно извлечь выгоду.
Пауку совершенно всё равно, в какую часть паутины попадёт муха, он в любой случае сможет запеленать добычу.

Лу очень любил пауков.
При встрече с золотопрядом изменяла даже его вечная брезгливость.
Карта стратегии действительно напоминала искусно сплетённую паутинку;
Люцифер не удержался и нарисовал в углу терпеливо ждущего паучка.
Демон, помедлив, педантично подсветил его:
<<Паук (элемент художественного оформления, игнорируется)>>

Карта стратегии Гало была гораздо сложнее.
Неожиданные повороты, сложные многомерные пути развития событий.
На бумаге карта выглядела как детские каракули.
Гало рассматривает стратегию как алгоритм;
стратегия Лу похожа на систему маятников, тросов, пружин и весов, находящихся в состоянии динамического равновесия.

<<Ну вот, брат, у каждого из нас теперь своя карта стратегии, --- с грустью подумал Лу.
--- До чего мы дожили>>.

Север Солнечная Дева.
Сальватор Сольмё.
Какой очаровательный демон.
Его досье Люцифер читал с улыбкой, как добрую сказку.
Доверенный Айну уже перехватил агента Гало --- разумеется, по наводке Люцифера.
Агенту передали недвусмысленное сообщение: Севера не трогать, за любое деяние исполнитель ответит головой.
Гало не станет жертвовать агентами без очевидной выгоды.
И уж точно он не станет из-за Севера ссориться с Айну.

<<Можно было бы его завербовать, --- рассеянно думал Лу, --- но нет.
Это что-то за гранью добра и зла, мне даже думать об этом противно.
Пусть этот милаха занимается своими обожаемыми людьми и ни о чём не беспокоится.
Из игры, конечно, выйти не так-то просто --- Гало обязательно будет выяснять, почему Север так дорог Айну.
Но главное, что выяснять он это будет не у Севера>>.

Люцифер пометил Севера крестиком и кружком.
Линии карты стратегии заплясали и снова установились в равновесии.
Демон подсветил ещё что-то.
Люцифер, хихикнув, прочитал:
<<Клякса (случайно сгенерированный элемент взаимодействия с пишущим устройством, игнорируется)>>

Люцифер отложил карту стратегии в сторону, подошёл к окну и положил разгорячённые ладони на прохладное стекло.
Это был <<нарэ>> --- своеобразный минутный отдых от работы.
Качество стекла оставляло желать лучшего: волнистая поверхность, полупрозрачные белые пятна примесей, паутинка трещин по углам, словно сплетённая каким-то особым стеклянным пауком.
Однажды Лу застал Тахиро в такой же позе;
тот не отрываясь глядел на стекло.

--- Что ты рассматриваешь? --- спросил стратег.

--- Нарэ, --- ответил Тахиро.

--- Мне это слово не знакомо.

--- А я не смогу объяснить смысл.
Поэтому, если хочешь понять, то просто смотри.

Лу добросовестно ходил и смотрел каждые несколько дней.
Он вглядывался в стекло, он изучил каждую волну поверхности, каждый дымчатый извив в его толще, каждую трещинку и каждое пятно на поверхности.
Смысл слова так и остался неизвестным;
однако Лу заметил, что странный ритуал неплохо помогает приводить мысли в порядок.

В этот раз что-то было не так.
Стекло недружелюбно вибрировало, хотя сейсмостанция не обещала подземных толчков.

Лу едва успел убрать со стекла ладони, как оно разбилось вдребезги.

<<Да чтоб вас! --- подумал Лу, растерянно глядя на осколки и текущую по запястьям кровь. --- Ну и как я теперь найду это нарэ?>>

Почти сразу же взвыли сирены.

\section{[:] Броня}

\textspace

--- Последние диверсии совершили люди поселений.
На экстренном заседании выступал Самаэл.
В общем, он убедил штаб в том, что оцифровывать человека опасно, что мы можем сами дать людям в руки оружие, которым нас можно уничтожить.
Извини, Тахиро.
Решение окончательное.

--- Я вам верен, Грис!
Я ни за что не ударю в спину тебе или Лу!

\ml{$0$}
{--- Я верю тебе, мальчик.}
{``I believe you, Tahiro-kid.}
\ml{$0$}
{Но Арракиса ты в этом не убедишь.}
{But you will not persuade Arrakis.''}

\ml{$0$}
{--- Ему нельзя здесь оставаться, --- сказал Лу.}
{``He mustn't stay here,'' Lu said.}
\ml{$0$}
{--- Его нужно...}
{``He must be---''}

\ml{$0$}
{--- Мне тоже нельзя здесь оставаться, --- прервал его Грисвольд.}
{``I mustn't stay here neither,'' Griswold interrupted him.}

\ml{$0$}
{--- Ты о чём?}
{``What are you talking about?''}

--- Гало знает о твоей фальсификации.
Он намекнул, что может заявить об этом в любой момент, если я не окажу ему услугу.

--- Он тебя шантажировал? --- Лу выглядел ошарашенным.
\ml{$0$}
{--- Что он хотел?}
{``What are his demands?''}

\ml{$0$}
{--- Он хотел, чтобы я передал сообщение одному агенту.}
{``He wants me to deliver a message to an agent.}
\ml{$0$}
{Я согласился.}
{I agreed.''}

\ml{$0$}
{--- Покажи.}
{``Show me.''}

Лу и Грис соприкоснулись головами.

--- Он хочет скомпрометировать тебя связью с Союзом, --- сказал Лу.
Его трясло.
--- Брат, до чего ты дошёл?

\ml{$0$}
{--- Лу, что мне делать?}
{``Lu, what must I do?''}

Лу вскочил и сделал несколько кругов по комнате.

\ml{$0$}
{--- Нам срочно нужно оборудование для оцифровки.}
{``We need digitization equipment, now.}
\ml{$0$}
{Только так мы сможем спасти Тахиро и доказать свою невиновность.}
{It's the only way to save Tahiro and prove our innocence.}
\ml{$0$}
{Тахиро, иди в мою комнату и никуда не выходи, пока я тебя не позову.}
{Tahiro, go to my room and stay there till I call you.}
\ml{$0$}
{Грис, твои железяки спрятаны в ячейке номер тринадцать в третьем ангаре.}
{Gris, your pieces of iron are awaiting you in the vault number thirteen, hangar three.}
\ml{$0$}
{Сообщение агенту я передам сам.}
{I'll deliver the message by myself.}
\ml{$0$}
{Надо узнать, насколько далеко Гало зашёл в желании нас скомпрометировать.}
{We must find out how far Halo's gone to compromise us.''}

\ml{$0$}
{--- Лу, ты подставляешь сам себя!}
{``Lu, you frame yourself!}
\ml{$0$}
{В произошедшем виноват я и только я.}
{Things that happened are my fault and mine alone.}
\ml{$0$}
{Ты ещё можешь...}
{You can still---''}

\ml{$0$}
{--- Я слишком долго вас искал, Грис.}
{``I spent too long looking for both of you, Gris.}
\ml{$0$}
{Я не хочу вас потерять.}
{I can't lose you.''}

Лу подхватил со стойки бронежилет, накинул плащ и бегом бросился за дверь.

\section{[:] Черепаха}

\textspace

Лу очень любил деревню.
Там жили самые разные люди --- взрослые, дети, старики, мужчины, женщины и те, кто не испытывал надобность принадлежать к какому-либо полу, люди разных профессий и разного нрава.

Стратег с детства бегал в деревню играть с другими детьми.
Он выпрашивал у женщин пирожки, ловил рыбу и дрался, как и прочие.
Каждый раз, когда его ловили за этим занятием, следовало наказание, о котором Лу, при всей лёгкости его нрава, очень не любил вспоминать.
Грисвольд однажды спросил, как же наказывал Лу Арракис.
Люцифер бойко отшутился, но его руки вдруг задрожали, словно стратегу внезапно стало холодно.
Грисвольд понял и с тех пор больше никогда не поднимал эту тему.

--- Ты --- будущий правитель, --- говорил Арракис после каждого наказания.
\ml{$0$}
{--- Ты должен вести себя соответствующе.}
{``You must behave accordingly.''}

--- Я --- будущий правитель, --- парировал Лу со слезами в голосе.
\ml{$0$}
{--- Я должен знать всё о тех, кем мне предстоит управлять --- их силы, слабости и нужды!}
{``I must know all about my subjects---their strength, weakness and needs!''}

--- А Гало? --- спросил Грисвольд.

--- Гало был его любимцем, --- Люцифер усмехнулся и зябко поёжился.
\ml{$0$}
{--- Его просто заставляли драться с солдатами до изнеможения.}
{``He was forced to fight the soldiers until he falls exhausted.}
\ml{$0$}
{По сравнению с моим это даже наказанием назвать нельзя.}
{Compared to mine, it wasn't a punishment at all.''}

\ml{$0$}
{--- Он знал, как обращаются с тобой?}
{``Did he know how you're being treated?''}

\ml{$0$}
{--- Он всегда принимал наказания как должное, --- уклончиво ответил Лу.}
{``He always took punishment for granted,'' Lu vaguely said.}

\ml{$0$}
{--- Он знал? --- настаивал Грис.}
{``Did he know?'' Gris insisted.}

Лу красноречиво посмотрел на друга и впервые на памяти Грисвольда промолчал.

Мимо прошла милая пара с зонтиками.
Девушки в выходных кимоно кокетливо семенили гэта, держась за ручки и время от времени обмениваясь робкими поцелуями.
Полуголые рыбаки, с весёлыми песнями снимающие гётаку\FM{} с огромной шипастой рыбы, помахали девушкам вслед бумажными полосками.
\FA{
Гётаку --- бумажный отпечаток вымазанной чернилами рыбы, который оставляли на память о трофее или использовали для отчётности.
}
Пожилой монах рассказывал зачарованно притихшим ребятишкам какую-то историю.

Деревня выглядела так, словно она сошла с праздничной открытки.
Она была совсем не той, которой стратег помнил её с детства.

Лу подумал --- и свернул с намеченного маршрута.

Почти сразу он понял, что сделал всё правильно.
На него стали неприкрыто пялиться, из углов понёсся недружелюбный шёпот.
Вечерняя деревня стала такой, какой ей надлежало быть --- сумеречной и настороженной.

Поворот, улица, поворот, улица, поворот, узкий проулок, ещё два поворота...

У одного из домов несколько женщин гнали спиртное.
Старушка, похожая на одетую в кимоно черепашку, проводила Лу взглядом.
Её примеру последовала вторая, вытягивая морщинистую черепашью шею.
Едва Лу скрылся за поворотом, как обе тут же занялись своими делами.

Впрочем, запах не был похож на спирт.
Лу вдруг замер столбом.

\ml{$0$}
{<<Жидкое взрывчатое вещество>>.}
{\emph{Liquid explosive.}}

Стратег попятился и осторожно заглянул за угол.
В нос ему уткнулся клинок.
Ещё пара осторожно залезла под нижний край бронежилета.

\ml{$0$}
{--- Что ты здесь вынюхиваешь, хорохито? --- осведомился низкорослый человек.}
{``What are you sniffing around, horohito?'' a short man asked.}

--- Авамори\FM\ уж больно хорош, --- попытался пошутить Лу.
\FA{
Авамори --- крепкий спиртной напиток Преисподней, около 56 градусов.
}

Он ожидал удара или выстрела, но клинки неожиданно отодвинулись.
Окружившие его люди зашептались.

\ml{$0$}
{--- Он сейчас что, пошутил?}
{``What, did that man just make a joke?''}

\ml{$0$}
{--- Это была шутка?}
{``Is that a joke?''}

\ml{$0$}
{--- Его подкупили.}
{``He was paid off.''}

\ml{$0$}
{--- Да нет, это один из этих, которых они выкормили, как поросят!}
{``You're wrong, it's a piglet \emph{they} have nursed!''}

Кто-то сорвал с Лу капюшон.

\ml{$0$}
{--- Это не тот!}
{``That's not the right one!}
\ml{$0$}
{Он не толстый и не лысый!}
{He is not fat, nor bald!''}

\ml{$0$}
{--- Гляньте, он глаза красит и волосы!}
{``Look, folks, he paints his eyes and hair!''}

\ml{$0$}
{--- Похож на актёра.}
{``Looks like a stage actor.''}

\ml{$0$}
{--- Да какой актёр?}
{``Nonsence!}
\ml{$0$}
{Он в боевой броне!}
{He wears a real armor!''}

\ml{$0$}
{--- У него трубка в кармане!}
{``I found a smoking pipe in his pocket!}
Пахнет свежим вереском!
\ml{$0$}
{А в ножнах --- ножницы!}
{And a scissors in his scabbard!''}

\ml{$0$}
{--- Ножницы в ножнах?}
{``A scissors in the scabbard?}
\ml{$0$}
{Что за чушь я услышал?}
{What the heck I just heard?''}

--- Да это точно актёр!

Толпа зашепталась ещё тревожнее.

--- Так ты не хорохито? --- наконец спросил его кто-то.

\ml{$0$}
{--- Вы мне всё равно не поверите, --- ответил Лу, осторожно поглядывая по сторонам.}
{``You won't believe me anyway,'' Lu answered trying to look around.}
\ml{$0$}
{--- Поэтому объясните для начала, почему вы спутали меня с хорохито.}
{``So explain first, what made you mistake me for horohito.}
\ml{$0$}
{Попутно объясните, почему вы вдруг засомневались.}
{You may explain, by the way, what made you question that after.''}

\ml{$0$}
{--- Вопросы здесь задаём мы, --- заявил человек, говоривший первым.}
{``That's we who ask questions,'' said the man who started talks.}

\ml{$0$}
{--- Отнюдь, --- лаконично возразил знакомый голос.}
{``Mistake,'' a familiar voice succinctly argued.}

Лу улыбнулся.
Тахиро был замечательным учеником.
Сказанное в общей модуляции слово хлестнуло по ушам людей, словно усаженный шипами бич.
Враги подняли руки с такой скоростью, словно от этого зависела жизнь всего их рода.
Клинки посыпались на утоптанную землю.

--- Сайгон, быстро в лагерь, --- распорядился Тахиро, щёлкнув предохранителем винтовки.
--- Грис ушёл в сторону третьего ангара.
Транспорт в квартале.
Я прикрою.

Лу растолкал ошарашенных дзайку-мару и бросился бежать, не дожидаясь пояснений.

\section{[:] Побег}

\textspace

--- Что ты здесь делаешь?

--- Тот же вопрос у меня к тебе, --- парировал Грисвольд.
Вначале он хотел ответить распространённой в человеческих поселениях похабной рифмовкой, но удержался, зная, что Мефистофель его не поймёт.
--- Арракис приказал мне обеспечить отсутствие прочих демонов здесь.
\ml{$0$}
{Уходи.}
{Go away.''}

--- Я получил от него похожий приказ, --- заявил Мефистофель.
\ml{$0$}
{--- Уйти придётся тебе.}
{``It's you who must leave.''}

Грисвольд впился прищуренным взглядом в Мефистофеля.
Тот смотрел на технолога спокойными пустоватыми буркалами.

\ml{$0$}
{--- Уйди, Меф.}
{``Go away, Mef.}
\ml{$0$}
{Говорю последний раз.}
{This is your last warning.''}

По лбу Грисвольда струился пот;
он прекрасно понимал --- истолкуй интерфектор его слова как реальную угрозу, и от омега-модуля толстяка не останется даже воспоминаний.
Но Мефистофель вдруг улыбнулся --- идеально ровной машинной улыбкой.

\ml{$0$}
{--- Ты пытаешься испугать меня.}
{``You're trying to scare me.}
\ml{$0$}
{Это признак слабой позиции.}
{It's a sign of weak position.}
\ml{$0$}
{Возвращайся к Арракису, уточни полученный тобой приказ и не трать моё время.}
{Return to Arrakis, request clarification on the order and stop wasting my time.''}

Мефистофель отвернулся --- ровно настолько, чтобы не видеть, как Грисвольд аккуратно взялся за браслет.
Вспышка --- и громоздкое, с туповатым лицом тело Мефа неуклюже село на пол.

<<Определённо удачное решение --- собрать боевые омега-модули в материальном устройстве, а не в демоне.
Но если браслет попадёт в чужие руки --- плакала моя головушка>>.

Грисвольд, демиург по воле судьбы и творец по призванию, ужасно не любил что-либо разрушать.
Насмотревшись всевозможных погребальных ритуалов у людей, он вдруг решил сказать пару слов над поверженным им противником.

--- Мефистофель из клана Мороза, ты был лучшим, --- начал он.
--- Ты был лучшим...

Пауза явно затянулась.

\ml{$0$}
{--- Понятия не имею, был ли ты хоть в чём-то и хоть для кого-то лучшим, --- растерянно закончил технолог и повернулся, чтобы уйти.}
{``I've no idea if you were the best at something or to someone,'' confused Griswold finished, then turned around to leave.}
Его встретил ошалелый взгляд Люцифера.

\ml{$0$}
{--- Зато честно, --- признал парень.}
{``You're honest, anyway,'' Lu declared.}
Разумеется, его демоническая сущность уже произвела подробный анализ и без того несложной ситуации.
--- За что ты его так, Грис?

--- Потом объясню, --- проворчал технолог.
--- Тахиро спит?

--- Нет, --- раздался голос из темноты.

--- Так и знал.
Быстро за мной, и желательно придержать вопросы до более подходящего момента.

\section{[:] Плач новорождённого}

Одной из самых остро стоящих проблем на Преисподней, вне зависимости от региона, является питьевая вода.
Добыча пресной воды и постройка водопроводов --- наукоёмкий процесс, тайны и нюансы которого люди Преисподней бережно передавали из поколения в поколение.
\ml{$0$}
{До тех пор, пока не пришёл Орден.}
{Until the Order came.}
\ml{$0$}
{Орден всегда вносил изменения в привычный уклад жизни.}
{The Order always modified dwellers' lifestyle.}

С акиямским водопроводом связана интересная история.
Раньше два соседних города на западном крае Такэсакского ущелья, Акияма и Киба, считались захолустьем из-за почти полного отсутствия питьевой воды.
Жители шли на ухищрения --- собирали дождевую воду, перегоняли испорченный вулканическим пеплом снег, искали под многометровыми серными ледниками чистые озерца и буквально руками прорубали к ним узкие тоннели.
В Акияма слыхом не слыхивали об открытых полях, культуры выращивались только в герметичных теплицах.
Всё изменилось благодаря легионеру Ордена, Тако из клана Дорге.
Несмотря на ранг и достаточно низкий интеллект, он был одержим инженерией.
Когда Тако вместе с его подразделением перебросили в Такэсако, он был поражён величием старинного такэсакского акведука, снабжавшего питьевой водой город с населением в сто тысяч человек.
В Акияма же его талант нашёл применение.

Арракис отказался тратить ресурсы на новый водопровод.
Тогда Тако сделал то, чего делать было нельзя --- обратился к дзайку-мару.

Чертежи Тако были многократно скопированы и спрятаны у самых разных людей, и изъять все так и не удалось.
Женщины обжигали кирпичи во время каждого приготовления пищи.
Цемент носили по горсти в каждом кармане, камни --- в каждом дорожном мешке.
Строительство велось ночами, даже в праздничные дни.
Не помогали ни показательные казни, ни шантаж, ни саботаж.
Агенты поставляли строителям уровни, откалиброванные по разным стандартам --- дзайку-мару разоблачили подделку и стали самостоятельно настраивать инструмент.
Диверсанты три раза подменяли обычные кирпичи на наполненные тротилом --- дзайку-мару с упорством муравьёв, не жалуясь и не поднимая шума, отстраивали развязку вновь.
Тела строителей, пойманных на месте преступления, снимали, а затем с потрясающим, свойственным только жителям Преисподней цинизмом подбрасывали к домам агентов --- иногда с вырезанными на коже оскорбительными иероглифами.

Строительство было завершено за тридцать лет.
По самым скромным подсчётам, акведук стоил жизни четырём тысячам мужчин и пятиста женщинам --- половине населения Акияма.
Первая чашка чистой воды в Киба была выпита за год до рождения Люцифера.
Ещё десять лет спустя акведук обеспечил драгоценной влагой всю западную часть ущелья Такэсако, и о <<годах жажды>> вспоминали лишь старики.

Разумеется, строительство не закончилось бы никогда, если бы Орден не перестал преследовать жителей.
На десятый год борьбы Дорге сумел убедить Арракиса остановиться.

--- Жаждущего на пути к роднику можно остановить только пулей, --- сказал он.
--- Их неповиновение --- вынужденная мера, а не прихоть оппозиционно настроенных слоёв населения.
Это очевидно из потерь, которые они несут.
В общем и целом водопровод не представляет непосредственную угрозу для Ордена.
А вот демонами, которые позволили себе подобные выходки, следует серьёзно заняться.

Тако всё-таки успел уйти на Тысячу Башен вместе со своим подельником, где и прослужил без особых происшествий до недавнего времени.
\ml{$0$}
{После присоединения Тысячи Башен к Ордену Преисподней Тако пропал бесследно.}
{After Thousand Towers was annexed by the Order of Nether, Tako has just vanished.}
\ml{$0$}
{Едва ли удача улыбнулась легионеру дважды;}
{He unlikely got lucky twice;}
\ml{$0$}
{Люцифер склонялся к мысли, что пропасть ему помогли.}
{Lucifer guessed that the legionnaire was helped to disappear.}
\ml{$0$}
{Орден ничего не прощал и не забывал.}
{The Order forgets nothing and forgives nothing.}

\ml{$0$}
{Не забыли и люди.}
{Humans forgot nothing, too.}
\ml{$0$}
{Доставшийся такой ценой акведук уважительно называют Тако-сан.}
{Bought at such a price, the aqueduct is respectfully called `Tako-san'.}

\ml{$0$}
{--- Так закончилась история и началась легенда, --- закончил свой рассказ Люцифер.}
{``So the story ends, and the legend begins,'' Lucifer finished.}
В его красивых глазах отражались искры костерка.
\ml{$0$}
{--- Заставляет задуматься.}
{``Makes you wonder.}
\ml{$0$}
{Мы живём так долго, мы дослуживаемся до званий --- но что нужно сделать, чтобы нас пережили наши собственные имена?}
{We live so long, we reach ranks, but what should we do to make our names outlive us?''}

\ml{$0$}
{--- Пойти против системы, --- предположил Тахиро.}
{``Buck the system,'' Tahiro guessed.}

\ml{$0$}
{--- Отнестись с уважением к тем, с кем никто не считается, --- пожал плечами Грис.}
{``Show respect to people no one has to reckon with,'' Gris shrugged.}

\ml{$0$}
{--- Или просто делать то, к чему лежит твоё сердце, --- заключил Люцифер.}
{``Or just do the job you have your heart set on,'' Lucifer ended.}

Друзья помолчали.

--- Я ни разу не слышал эту историю, --- сказал Грисвольд, оглядывая величественные извивы водопроводной развязки, похожие на щупальца гигантского осьминога.
\ml{$0$}
{--- Как жаль, что я не был знаком с Тако.}
{``Such a pity I never met Tako.}
\ml{$0$}
{Его можно было взять в нашу команду.}
{He would be good for the crew.''}

\ml{$0$}
{--- Я тоже об этом думал, --- не без сожаления признал Люцифер.}
{``I was thinking the same thing,'' Lucifer admitted, not without regret.}
\ml{$0$}
{--- Отец порой ужасно расточителен.}
{``Father is terribly wasteful sometimes.''}

--- И всё-таки, почему, почему люди продолжали строительство, несмотря на такие потери?

--- Посмотри на этот город, --- сказал Люцифер.
--- Его жители проводят жизнь в попытках раздобыть немного пищи и перегоняют собственные выделения, чтобы не страдать от жажды.
Беспросветная серая мгла выживания, от рождения до смерти.
И вдруг приходит кто-то, неважно кто, и дарит людям мечту о лучшей жизни.
Даже о кусочке лучшей жизни --- о чашке чистой воды.
Что бы сделал ты?
А им и подавно нечего было терять.

--- Я бы не стал бороться за воду, --- заявил Тахиро.

--- А за что, по-твоему, ты борешься сейчас? --- возразил Грисвольд.
--- Уважение, личная неприкосновенность, свобода.
Это такие же банальные вещи, как и чистая вода.

--- Вы странные, --- сказал Тахиро.
--- У меня ощущение, что вы совсем-совсем не отличаетесь от нас.
Но вас как будто постоянно что-то гнетёт.
Когда мы пробирались сюда через скалы, я думал, что нет участи хуже моей: позади --- смерть, впереди --- мучения и, возможно, тоже смерть.
Но для меня это лишь плохой день, всего один, хоть и последний, а для вас --- обыденность.

Лу и Грис с изумлением уставились на парня.
Однако тот занялся костром и не счёл нужным объяснять свои слова.

\asterism

Неподалёку бил гейзер и парили несколько источников.
В некоторых вода была похожа на жидкий огонь, в других она почему-то была температуры растаявшего снега.
Друзья всё-таки нашли один с приемлемой температурой по лягушачьим крикам.
Лу пару часов пролежал с лягушками, отказываясь вылезать.
Его не смутил даже настойчиво щёлкающий дозиметр, который Грисвольд красноречиво сунул ему под нос.

\ml{$0$}
{--- Я всё равно не собираюсь размножаться, --- буркнул стратег.}
{``I'm not gonna breed, anyway,'' the strategist grunted.}
\ml{$0$}
{--- Сто лет в ванне не лежал.}
{``Haven't been lying in a bath for ages.''}

\ml{$0$}
{--- Тебя рыбки сейчас съедят!}
{``You're about to be eaten by fishlings!''}

\ml{$0$}
{--- Это головастики, Грис.}
{``Tadpoles, Gris.}
\ml{$0$}
{Личиночная форма земноводных.}
{Amphibian larvae.}
Они меня немного щиплют, и всё.

Тахиро всё-таки вытянул Лу из источника и, не обращая внимания на протесты, плотно закатал стратега в несколько одеял.
Грисвольд притащил гигантскую хмурую гусеницу к костру, усадил на камень, вколол Лу в загривок радиопротектор и занялся сборкой аппаратуры.

--- Кстати, --- сказал Лу.
--- Я, кажется, понял, как люди отличали демонов.
Их очень сильно сбило с толку то, что я пошутил.

Грисвольд хлопнул себя по лбу.

--- Я ни разу не слышал, чтобы наши легионеры шутили или смеялись над шутками, --- сказал он.
--- Да и многие из высших чинов тоже.

--- Надо поработать над адаптационными обновлениями, --- задумчиво сказал Лу.
--- Займёшься этим, когда вернёмся, хорошо?
Я попрошу, чтобы тебя освободили от прочих заданий.
А отработку методики демонизации отдай Северу.
Он в технологии не особенно силён, но он очень хороший биоинформатик и, думаю, справится.
Сейчас работа как раз по его профилю.

--- Ранговая комиссия...

--- В бездну ранговую комиссию.
У себя в лаборатории ты --- ранговая комиссия.
Кого хочешь, того и ставь заместителем.

\ml{$0$}
{--- Ты уверен, что мы вернёмся? --- грустно спросил Грис.}
{``Are you sure we'll be back?'' Gris sadly asked.}

\ml{$0$}
{--- Оцифруем Тахиро и вернёмся.}
{``We'll digitize Tahiro and be back.}
\ml{$0$}
{Есть какие-то другие предложения?}
{Any other suggestions?''}

\ml{$0$}
{--- Нас не ждёт тёплый приём.}
{``We won't be welcome there.''}

\ml{$0$}
{--- Нас он нигде не ждёт.}
{``Like anywhere else.''}

--- А если они всё-таки захотят меня уничтожить? --- спросил Тахиро.

\ml{$0$}
{--- В этом случае пусть Орден больше не рассчитывает на сотрудничество со мной, --- буркнул Грисвольд.}
{``In that case, the Order shouldn't count on me anymore,'' Griswold grunted.}

Лу кивнул.

\ml{$0$}
{--- И со мной тоже.}
{``On me either.''}

Грисвольд бросил на стратега изумлённый взгляд.
Люцифер смутился.
Он понял, что после этих слов Грисвольд пойдёт за ним хоть в пламя вулкана.
Стратег спокойно относился к военным потерям, но ответственность за друга немного его испугала.

--- Кстати, сегодня мой день рождения, --- сообщил Тахиро.

--- Что? --- не понял Лу.

--- У нас есть традиция --- отмечать число и месяц рождения каждый год.
Маленький праздник для одного человека, его друзей и родных.

Люцифер окинул взглядом заброшенный лагерь и унылый пейзаж Акияма в долине.

--- Я не знал.
Сочувствую.

--- Каждый день рождения, который я помню, был таким же невесёлым, --- буркнул Тахиро.
--- Моих родителей освежевали в этот же день.
Видимо, впервые увидев свет, после первого вдоха человеку должно плакать, а не смеяться.

--- Зато ты родишься в этот день дважды, --- попытался приободрить его Лу.

--- Или проживу целое количество лет, --- закончил Тахиро.
--- Я ж не дурак, Лу.
Я знаю, на какой риск иду.

Лу помолчал и поковырял пальцем ноги землю.

--- Ты всё ещё сердишься на меня из-за родителей?

Тахиро удивлённо посмотрел на друга.

--- Нет, конечно.
Я перестал сердиться в тот момент, когда ты сказал, что не хотел их смерти.
\ml{$0$}
{Я понял, что ты сожалеешь.}
{I realized you're sorry.}
\ml{$0$}
{Но удивило меня то, что в твоём сожалении не было ничего морального или религиозного --- это было искреннее сожаление того, кто случайно наступил на красивый цветок.}
{But I was surprised there were nothing ethical or religious in your regret, you're sorry as if you accidentally stepped on a beautiful flower.}
\ml{$0$}
{Ты увидел в моих родичах то, чего не замечал даже я --- красоту.}
{You have seen in my relatives one thing even I failed to notice---beauty.''}

\ml{$0$}
{--- Я бы сказал --- жизнеспособность.}
{``Vitality, I'd say.''}

\ml{$0$}
{--- Это одно и то же.}
{``That's the same.}
И меня ты освободил не из-за каких-то принципов или догм --- тебе было противно смотреть на меня связанного, это оскорбляло твоё чувство прекрасного.

--- На трупы и больных до сих пор смотреть не могу, --- признался Лу.

Грисвольд отвлёкся от сборки агрегата и одобрительно посмотрел на него.

--- Однако же странно, что с таким мировоззрением ты не стал кем-то вроде отца, --- сказал технолог.
--- В геноциде обычно повинны такие вот ценители прекрасного.

--- Чувство прекрасного бывает разное, --- сказал Лу.
--- Лично мне трупы кажутся уродливее любого из живых.

Друзья расхохотались.
Впрочем, веселье быстро утихло.

--- А что делают для человека, у которого день рождения? --- спросил Лу.

--- Ему дарят подарок и кормят сладостями.
Но мне нельзя ничего есть.

Люцифер подумал и обратился к Грисвольду:

--- Ты приготовил ему модули?

\ml{$0$}
{--- Да, разумеется.}
{``Yes, of course.}
\ml{$0$}
{Интерфейс ввода-вывода и шаблоны модулей легионеров.}
{Input/output interface and legionnaire module templates.''}

\ml{$0$}
{--- Если всё получится, скопируй и пристыкуй ему мои.}
{``If all goes well, mount copies of my modules to him.}
\ml{$0$}
{Так, как они есть, с актуальными обновлениями.}
{As they are, with all updates.''}

\ml{$0$}
{--- Ты серьёзно?!}
{``Are you serious!''}

--- Ему нужен подарок, --- грустно улыбнулся Лу.
\ml{$0$}
{--- А модули --- это всё, что у меня есть.}
{``My modules are my only property.''}

Грисвольд выглядел ошарашенным.

--- Это же самые отработанные стратегические методики, большая часть накопленных человечеством...

\ml{$0$}
{--- Я знаю.}
{``I know.''}

\ml{$0$}
{--- Отец тебя!..}
{``Your father will ... !''}

\ml{$0$}
{--- Я знаю, Грис!}
{``I know, Gris!''}

--- Под твою ответственность, сайгон, --- развёл руками технолог.
--- Но всё-таки...

\ml{$0$}
{--- От меня не убудет.}
{``It costs me nothing.}
\ml{$0$}
{Модули --- это не яблоки.}
{A module is not an apple.''}

--- Спасибо, Лу, --- поклонился Тахиро.
\ml{$0$}
{--- Я не совсем понимаю, что именно ты хочешь мне подарить, но, думаю, это прекрасный подарок.}
{``I barely understand what exactly you want to give me, but I think it's a nice present.''}

Лу хитро осклабился:

\ml{$0$}
{--- Ты \emph{совсем не понимаешь}, что я хочу тебе подарить.}
{``You \emph{completely don't understand} what exactly I want to give you.}
\ml{$0$}
{Грис немного понимает, поэтому он и...}
{Gris does a bit, that's why he---''}

\ml{$0$}
{--- Ладно, хватит болтовни, --- оборвал друга Грис.}
{``Stop, done talking,'' Gris interrupted his friend.}
--- Тахиро, нужду справил?

--- Десять минут назад.
Не ел четырнадцать часов, как ты и хотел.
Выпил глоток подсоленной воды около часа назад.

--- Тогда можем приступать.

--- Так, выпутайте меня из этого кокона, я не хочу пропустить самое интересное, --- попросил Лу.

Тахиро аккуратно размотал одеяло, освободив друга.
Их глаза на мгновение встретились.
Во взгляде Люцифера застыло что-то совсем непонятное, и это не укрылось от Тахиро.

--- Всё хорошо? --- спросил он.

--- Да уж лучше, чем у тебя, --- отшутился Лу, похлопав друга по плечу.

Люцифер и сам не знал, что он чувствует.
Он всем сердцем хотел, чтобы Тахиро был равным ему.
И сейчас его желание сбывалось, но как-то совсем не так, как должно было.
%{His wish was about to come true, but not the way it was supposed to be.}

--- Анестезию будем делать? --- спросил стратег у Грисвольда.

--- Никакой анестезии, --- буркнул Грис.
\ml{$0$}
{--- Он должен быть в полном сознании, чтобы устройство захватило активность всех участков нервно-мышечной системы.}
{``He must be fully conscious, so the device is able to capture activity of the whole neuromuscular system.}
\ml{$0$}
{Если он потеряет сознание в процессе, придётся приводить его в чувство.}
{If he loose consciousness, he's to have been waken up.''}

Тахиро побледнел.

--- Ты готов? --- обратился Лу к другу.

Тахиро кивнул и, сбросив одежду, лёг на жёсткую станину.
Его била крупная дрожь.
Грисвольд прикрутил конечности юноши к кольцам, жёстко зафиксировал голову винтами и вытащил из кармана ларингоскоп.

\ml{$0$}
{--- Выполняй мои команды.}
{``Obey my commands.}
Если я скажу <<согни палец>> --- ты сгибаешь.
Если я скажу <<расслабься полностью>> --- делай что хочешь, но выполняй, от этого зависит качество оцифровки периферических нервов.
Иногда будет очень щекотно, жарко, холодно или даже больно --- это тоже часть процесса, мы не будет тебя калечить.
Терпи и старайся не шевелиться, если я не дал соответствующую команду.
Я постараюсь предупреждать заранее обо всех ощущениях.
\ml{$0$}
{Если тебе вдруг покажется, что ты ослеп, оглох, потерял способность говорить или дышать, или что у тебя отнялись руки, ноги или лицо, или что тебя вдруг начали преследовать звуки, запахи или образы --- постарайся успокоиться и жди, мы смонтируем тебя как можно скорее.}
{If you suddenly feel you went blind, or deaf, or mute, or you lost ability to breathe, if you can't feel your arms, or legs, or face, if you suddenly see, or hear, or smell something strange---try to relax and wait, we'll mount you as soon as possible.}
\ml{$0$}
{Помни --- я и Лу рядом и ни за что тебя не бросим.}
{Remember: Lu and me both are near by you, and we'll never foresake you.}
А теперь открой рот как можно шире и приготовься --- будет больно.

Грисвольд вставил в рот юноши расширитель.
Тахиро зажмурился.

\chapter{[U] Визит командующего}

\section{[U] Милая девочка Штрой}

--- Атрис, ты дурак, --- смеялась Митхэ.
--- Разве можно так шутить над живыми существами?

--- Я не шутил, --- обиделся Атрис.
--- Просто хотел проявить оригинальность.
Билатеральных существ как песка на пляже.
А вот с трилучевой симметрией почти нет.

--- Они очень милые, --- заверила Митхэ друга.
--- Хорошо, что сейчас, после переговоров с дельфинами, они размножились и могут покинуть Коралловую бухту.

Атрис пожал плечами и снова стал наблюдать за малышами-нгвсо, которые резвились у самого берега.
Кажется, один нашёл большую красивую ракушку и совсем не хотел делиться с остальными.
Он плавал кругами, и время от времени его щупальца били по воде, поднимая тучи искрящихся брызг.
Под водой царил весёлый гул, щелчки и повизгивание, но слуха сухопутных сапиентов не достигало ничего.

--- Может, покормим их? --- предложила Митхэ.
--- Они любят вишню.
И другие фрукты тоже.

--- Лучше не надо.

--- Это же дети, Атрис! --- Митхэ подёргала менестреля за рукав.
--- Детей можно!..

--- Не стоит.
Я стараюсь как можно меньше вмешиваться в их жизнь.

Митхэ обняла Атриса и уткнулась кудрявой головкой ему в шею.
Цветы в её волосах потускнели, попав в тень.

--- Я скучаю по нашему хранителю, милый.

--- Мы почти его не знали, --- напомнил ей Атрис.
--- И потом, он был инкарнацией демона Аркадиу Шакал Чрева.

--- Это совершенно не мешает мне скучать по ним обоим.

--- Да, мне тоже, --- признался Атрис.
--- Это я глупость сморозил.
Интересно, кто это?
Похоже, что к нам.

--- Может, Корхес или кто-то ещё из биологов, --- предположила Митхэ.

--- Они придут вечером с конфетами, --- возразил Атрис.
--- Намечается что-то интересное.
Корхес хотела поговорить насчёт нескольких видов мух.
А Кольбе и Рабе обещали показать странную диатомею, которая строит алмазные раковинки.

--- Алмазные?
В жидкой воде, при обычном атмосферном давлении?

--- Да.
И, представь себе, я тут вообще ни при чём.

--- Сама?

--- Сама.

По мокрому песку шла девушка-оцелот, хрупкая и нежная на вид, почти ребёнок.
Её босые тонкие ноги утопали в мокром песке, комбинезон из сложного полимера был расстёгнут, обнажая костлявую грудную клетку и верхушки маленьких грудей.
Атрис хмыкнул --- эта неестественная, показная соблазнительность выдавала чужаков, инкарнированных демонов Ада.
Люди сели, привыкшие к наготе тела, не считали нужным акцентировать внимание на груди или половых органах.

Девушка подошла поближе, и все сомнения на её счёт рассеялись.
Пухлые детские губки кривились в уверенной взрослой улыбке, а оранжевые опалесцирующие глаза светились древностью и памятью бесчисленных жизней.
Она подошла к парочке и по очереди пожала им руки.

--- Штрой Кольцо Дыма, командующий силами Ада на Тра-Ренкхале, --- представилась девушка, не забыв сделать в сторону Атриса лёгкое целующее движение губами.

Митхэ поморщилась.

--- У тебя проблемы с гормонами?
Ведёшь себя странно.

--- Я ещё не совсем освоилась с телом тси, --- сообщила девушка.
--- Возможно, мне следует поменять паттерны поведения, они не годятся для общества сели.
Вы не первые, кто это заметил.
На мужчин других миров моё поведение действует завораживающе, а мужчины-сели только смеются и не желают со мной общаться.
В общем, прошу понять и простить.
Мы можем поговорить?

--- Да, разумеется, --- сказал Атрис.
--- Недалеко есть беседка.
Ты не откажешься от чашечки холодного травяного отвара?

--- Буду очень признательна, --- махнула длинными ресницами Штрой.
--- Сегодня жарко.
Правда, вам, голограммам, это сложно понять.

\section{[U] Чай}

Беседку установили местные жители.
Безымянный не скрывался от глаз сапиентов --- одно из требований новой администрации;
его достаточно часто видели в нескольких ключевых местах планеты, в том числе и неподалёку от Кахрахана.
Атрис прекрасно понимал, что это было сделано с целью разрушить культ --- даже необычно выглядящее существо становится обыденностью, если видеть его каждый день.

Но у сели культа как такового и не было, Безымянный был для них обычным героем легенд, подобно Забытому или Маликху.
Они просто сделали для него личную беседку из чёрного дерева, словно для уважаемого путешественника.
Такие же скромные, но со вкусом сделанные домики стояли ещё в двух местах в землях сели --- возле Сотрона и Ихслантхара, у ноа --- возле Яуляля, и у хака --- возле Спокойного озера.
Безымянному они нравились гораздо больше, чем выстроенные тенку пышные капища или часовенки ркхве-хор.
Они были похожи на нормальные жилища, а не на культовые строения, и в них можно было расслабиться.

Сели, хака и ноа нравились Безымянному ещё и тем, что не тревожили его просьбами.
Все местные знали, если не от демонов, то по слухам --- Безымянный долгое время был в изгнании, но продолжал ухаживать за планетой.
Люди изо всех сил старались показать, что они рады вернувшемуся создателю --- и ценили его уединение.

Вскоре в беседке из чёрного дерева запахло пряным паром отвара.
Митхэ плеснула в чашу жидкого азота из спрятанной под скамьёй азотной станции, и седые клубы волнами разбежались по узорчатой поверхности стола.

--- Благодарю, --- улыбнулась Штрой и приняла прохладную чашу.

\ml{$0$}
{--- Мы тебя слушаем, --- кивнул Атрис.}
{``We're listening,'' \Aatris{} nodded.}
\ml{$0$}
{--- К нам редко заходит кто-то, кроме биологов.}
{``We don't get many visitors, except some biology folk.''}

\ml{$0$}
{--- Кольбе и Рабе сегодня не придут, --- сообщила девушка, отвечая на завуалированный вопрос.}
{``Colbe and Rabe won't come today,'' the girl answered the veiled question.}
\ml{$0$}
{--- У них появились срочные дела.}
{``They've got some urgent matters.''}

Митхэ хмыкнула.

\ml{$0$}
{--- Неужели разговор настолько важный, что его нужно тянуть до позднего вечера? --- нахмурился Атрис.}
{``Are the talks valuable enough to take our evening?'' \Aatris{} frowned.}

\ml{$0$}
{--- Вы торопитесь?}
{``What's the hurry?''}

Штрой наклонилась, чтобы отхлебнуть отвара, и щёлкнула языком от удовольствия.

\ml{$0$}
{--- Вкусно.}
{``Quite delicious.}
\ml{$0$}
{Благодарю вас.}
{Thank you.}
\ml{$0$}
{Итак.}
{Here we go.}
Знаете ли вы Аркадиу Шакала Чрева?

--- Излишний вопрос, --- поморщилась Митхэ.
--- Аркадиу --- наш друг, и его последнее тело было нашим с Атрисом хранителем.
Думаю, ты знаешь это не хуже нашего.
\ml{$0$}
{С ним что-то случилось?}
{Something wrong with him?''}

\ml{$0$}
{--- Да, случилось.}
{``Yes, something's wrong with him.}
\ml{$0$}
{Он предал Орден Преисподней и сбежал с Капитула в неизвестном направлении.}
{He has betrayed the Order of Nether and left Capitul for an unknown destination.''}

Митхэ ахнула.

\ml{$0$}
{--- Что значит <<в неизвестном направлении>>? --- осведомился Атрис.}
{``What do you mean by `an unknown destination'?'' \Aatris{} asked.}
\ml{$0$}
{--- Кажется, ваша контрразведка знает всё --- кто, куда и с какими намерениями сбегает.}
{``Your counter-espionage is supposed to know everything---who runs, where, and what they have in mind.''}

--- Он сумел подсунуть контрразведке простую, но весьма правдоподобную легенду, --- ответила Штрой.
--- Обман был раскрыт минуту спустя, отдел 100 направил мне официальный запрос по работе биологов.
За это время Падальщик успел сбросить хвост и исчезнуть.

--- Какие бы силы ни стояли за твоей спиной, осторожнее выбирай слова, --- невыразительно проговорила Митхэ.
--- Предатель Аркадиу или нет, он хорошо послужил Ордену.
Тебе стоит уважать его хотя бы поэтому.

Штрой улыбнулась.

\ml{$0$}
{--- Если угодно.}
{``If you prefer.}
\ml{$0$}
{Так или иначе, след потерян.}
{Anyway, we've lost the trace.}
Ожидалось его появление в трёх системах --- Тысячи Башен, Тра-Ренкхаля и Тси-Ди.
Агенты всех трёх пока молчат.

\ml{$0$}
{--- Почему именно эти три?}
{``Why those three?''}

Штрой пожала плечами и отхлебнула отвара.

--- Могу сказать, что Люпино определенно попытается использовать потомков тси и их наследие в своих целях.

--- Это мы и сами поняли.
А Тысяча Башен?
Я мало что знаю про эту планету...

--- О, это прекрасное убежище для криминального элемента.
Куча Друз, скалистые теснины, пещеры в горах, хйяры.
Экваториальный кластер планеты де-юре контролируется Орденом, де-факто нейтрален и находится под влиянием легата Курт Осенний Огонь, которая весьма толерантна к нейтралам.
Тропический пояс де-юре под властью Картеля, де-факто так же нейтрален, и заправляют там Острые Шила, клан из шестнадцати демонов столь же свободных взглядов.
И Курт, и Острые Шила --- уроженцы планеты и знают её досконально.

--- Я бы ожидала там полное безвластие.

--- Отнюдь.
Они научились жить дружно, умудряясь при этом слать отчёты о своих успехах, весьма правдоподобные.
Однако правду не скроешь.
Именно из-за этого туда начали сползаться научные коллаборационисты.
Ну и у меня есть сведения, что недавно легат Курц летала к кому-то <<в гости>> в тропический пояс.
Не будь у меня информации от агентов, я бы решила, что это Люпино.

--- Что произошло? --- спросила Митхэ.
--- Почему он предал Орден?
\ml{$0$}
{Кому он теперь служит?}
{Who he serves now?''}

--- Вполне возможно, что никому, Митхэ ар’Кахр.
Некоторые аналитики считают, что он помешался, несмотря на коррекцию личности --- стигмы помешательства в его поведении присутствуют.

--- Что ты имеешь в виду под <<помешательством>>? --- осведомилась Митхэ.

--- Они имеют в виду, что Ликхмас ар'Люм взял верх над Аркадиу Люпино, --- ответил Атрис.

Штрой прикрыла глаза.

--- Всё гораздо сложнее, но в целом суть передана верно: субличность носителя становится приоритетнее субличности демона.
Если речь идёт об ангельских технологиях, такое обычно происходит без борьбы в буквальном смысле.
Просто личность демона добровольно перекладывает полномочия по принятию решений на личность носителя, частично или полностью, так как, с её точки зрения, это отвечает интересам обоих.

--- И что в этом плохого?

--- Для демона, разумеется, ничего.
А вот для его контрагентов это повод для недоверия.
Поэтому, если помешательство доказано, то все договоры с демоном теряют силу, ранги обнуляются, демон освобождается от всех постов и должностей, в том числе ставится вопрос и о принадлежности демона к Ордену.

--- Но в отношении Люпино это доказано не было.

--- Не было, --- подтвердила Штрой.
--- Но полученные данные говорят о том, что он заслуживает пристальное наблюдение.

--- Из-за чего он мог помешаться? --- удивилась Митхэ.

--- Есть некоторые \emph{подозрения}, --- Штрой едва заметно сделала акцент на последнем слове.
--- Каждое новое тело изменяет демона, это давно доказано.
Собственно, коррекция личности и направлена на нивелирование этих изменений.
Но все исследования такого рода проводились на диких сапиентах, а не на творениях генной инженерии, которыми были тси.
Возможно, что тело Ликхмаса ар’Люм вызвало тонкий, но тем не менее серьёзный системный сбой в Аркадиу.
Его товарищи признаны годными к дальнейшей деятельности, хотя я бы не спешила выносить такой вердикт.

Митхэ всё больше раздражала её манера общения.
Всё в молодой девушке, даже её улыбка, было вызывающе сексуальным, бросающим вызов ей и Атрису.
Вскоре Митхэ поняла, что именно её раздражало.
Движения Штрой казались идеальными, отточенными до предела.
Это были движения машины, хоргета, но никак не милой девушки-сели, которая по воле судьбы несла этого хоргета на себе.
Демон не доверял своему телу и контролировал его во всём.
В этом было его отличие от тех демонов, которых Митхэ знала прежде.

Слова Штрой только подтверждали это случайное наблюдение.

--- Возможно, он снова перешёл на сторону Картеля, --- предположил Атрис.

--- Исключено, --- отрезала Штрой.
\ml{$0$}
{--- Это не подлежит сомнению.}
{``It mustn't be questioned.}
Он чересчур насолил Картелю.

--- По-вашему, это препятствие для того, чтобы принять в свои ряды сильного противника?

--- Я выражусь иначе: он чересчур насолил некоторым иерархам Картеля.
\ml{$0$}
{Если в Аду у него почти не было личных врагов, то там их, я полагаю, на пару порядков больше, чем должно быть для относительно спокойного существования.}
{He barely had enemies in the Hell, but there, he has more---I guess, two orders of magnitude higher than it needs for relatively quiet life.}
Но ладно, ближе к теме.
Нас интересуют некоторые детали.
Когда вы виделись с Аркадиу в последний раз?

--- Около двадцати дождей назад, здесь, на Тра-Ренкхале, --- задумалась Митхэ.
--- Он сказал, что больше не хочет воевать и попросит переквалификацию.
Пообещал иногда навещать нас.

--- А виделись ли вы с ним на станции Деймос-14?

--- Нет, --- ответил Атрис.
--- На станции нас встретили Грейсвольд Каменный Молот и Анкарьяль Кровавый Шторм, они проинструктировали нас и заверили согласие на сотрудничество с Орденом.
Приказ исходил от Лусафейру Лёгкая Рука.
Правда, мы думали, что согласие должен брать один исполнитель в присутствии двух наблюдателей, но...

--- Значит, третьим должен был быть Аркадиу, --- глухо сказала Митхэ.

--- Вот именно, --- улыбнулась Штрой.
--- Если не ошибаюсь, Грейсвольд и Анкарьяль уже были вызваны на допрос, и им предстояло долгое и весьма неприятное объяснение, почему второй наблюдатель во время процесса болтался неизвестно где.
Хорошая иллюстрация к вопросу об их пригодности.

--- С ними всё хорошо? --- Митхэ схватила Штрой за худую руку и тут же отпустила, поймав совсем не дружеский взгляд девушки.

--- С ними всё в порядке, их лояльность уже долгое время не вызывает сомнений, --- сухо ответила Штрой.
--- Теперь о вас.
Во-первых, вы должны в ближайшие пятьдесят три секунды прибыть на станцию 1A, чтобы вас проверили агенты отдела 100.

--- Хорошо, --- сказал Атрис.
--- Но проверить вы сможете только меня, Митхэ --- оцифрованный...

--- У отдела 100 есть средства для проверки оцифрованных сапиентов, --- прервала его Штрой.

Атрис вздрогнул.

--- Чистилище?

--- Это не мне решать.
Данные о структуре вашего демона я передала, в частностях отдел 100 разберётся.
Вам ведь нечего скрывать?

--- В Чистилище невиновность не является защитой от страданий, Штрой, и ты это прекрасно знаешь, --- голос Атриса оставался тихим, но глаза сверкнули совершенно человеческой яростью.

\ml{$0$}
{--- В том, что вас скомпрометировали, есть доля вашей вины, --- парировала Штрой.}
{``You're partly responsible for being compromised,'' Stroji retorted.}
\ml{$0$}
{--- Аккуратнее выбирайте друзей.}
{``Be careful making friends.}
\ml{$0$}
{Далее.}
{Next.}
\ml{$0$}
{Я так понимаю, вы оба испытываете к Шакалу дружеские чувства.}
{As I understand it, both of you feel friendship for Jackal.''}

Митхэ и Атрис переглянулись и кивнули.

--- Так или иначе, он нарушил присягу, и его судьба отныне находится в ведении отдела 100.
Поэтому хочу вам напомнить, что сотрудничество с предателем или укрывательство такового является нарушением, предусмотренным Оборонительным Кодексом Ордена Преисподней, статья DB, пункт 4.
О санкциях распространяться не буду, читайте сами.

Атрис опустил голову.
Воительница смотрела в оранжевые глаза девушки своим спокойным зелёным взглядом, который так красочно описывали её давно умершие враги.

--- И кое-что не для протокола.
На случай, если вы всё-таки захотите ему помочь, --- Штрой поднялась и томно потянулась, выгнув тонкую талию.
--- Вряд ли то, что носит сейчас личину Аркадиу, на самом деле является им.
Благодарю за отвар.

Не попрощавшись, Штрой Кольцо Дыма вышла из беседки и манерной походкой отправилась на север.

\section{[U] Прошедшая Чистилище}

\epigraph
{Пусть замки из песка дождём бесследно смоет,\\
Упрямый мальчуган их заново построит.\\
Я строчку начертал у моря на песке,\\
И тянется прибой когтями к той строке.\\
Едва мои стихи волна с собой умчала,\\
Я, как дитя, готов игру начать сначала.\\
Но навсегда стихам конец, когда из вод\\
Хотя б одну строку волна назад вернёт.}
{Сигурдур А.\,Магнуссон, <<Начертанное на песке>>.
Эпоха Последней Войны, Древняя Земля}

Митхэ сидела в песке, прижимаясь к Атрису.
По её лицу текли слёзы, перемешиваясь с каплями только прошедшего ливня.

--- Я надеюсь, что мне больше никогда не придётся столкнуться с контрразведкой, --- прошептала она.

Атрис выглядел гораздо хуже.
Модуль визуализации ещё работал, голограмма передавалась без сбоев, но сам образ искажался.
Глаза, волосы и кожа Атриса меняли цвет, черты лица расплывались, словно менестрель надел неисправный ноппэрапон\FM.
\FA{
На Древней Земле: маска театра Кабуки из магнитного полимера, принимающая любую форму и цвет.
}
Иногда изображение пропадало совсем, и казалось, что Митхэ обнимает воздух.

--- Да, я знаю...

--- Агенты сказали, что это скоро пройдёт.
Потерпи.

--- Я с тобой, милый.
Я всегда буду с тобой.

--- Хочешь песенку?
Я спою, а ты сыграешь.

--- А почему, ты же любишь музыку?

--- Ну ладно.

Из кустов показался ребёнок сели.
Он целеустремлённо, с забавно серьёзной мордашкой прошёл до берега моря и сел недалеко от парочки.
Атрис представлял собой ужасающее, сюрреалистичное зрелище, но ребёнок только улыбнулся беззубым ртом.
Вскоре он уже сидел и строил из мокрого песка домик.
Крохотные ручки неумело собрали песчаную гору и начали формировать крышу.
Митхэ с нежностью наблюдала за ним.

--- Интересно, откуда он?

--- Понятно.
Можно, я ему помогу?

--- Я вижу, что он прекрасно справляется сам.
Вопрос был не об этом.

Солнце начало клониться к закату.
Менестрель почти пришёл в себя.
У него всё ещё были видны странные аберрации в волосах, но лицо уже стало прежним.
Он посмотрел на Митхэ и улыбнулся.

--- Ну вот, Ад нашёл повод и тебе сунуть Чистилище под нос.

--- Я держалась молодцом? --- спросила Митхэ.

--- Да, --- кивнул менестрель.
--- Один из агентов сказал мне следующее: <<Личность Митхэ ар’Кахр перешагнула последнюю ступень целостности.
Если Митрис Безымянный будет заподозрен в деятельности, направленной против Ордена, мы будем вынуждены сразу уничтожить его человеческую часть>>.

--- Ну хоть мучить не будут, --- ободрительно сказала Митхэ.
--- Жаль только, что я не смогу облегчить твои страдания, если до подобного дойдёт.

Атрис помолчал.

--- Забавно.
Чистилище --- это термин из давно забытой религии.
Это страдания, которые дарил древний бог, приговаривая: <<Ты мой верный раб, но за грехи ты ответить обязан>>.

--- Грехи --- это преступления перед богом?

--- Да, это нарушения его законов.

--- А что давало этому богу право требовать их исполнения?

--- Он оправдывал это тем, что он создал сапиентов.

Митхэ поморщилась.

--- Не горшечник решает, когда время горшку разбиться.
Это решает тот, кто способен нанести горшку смертельный удар.

--- Верно, --- признал Атрис.
--- Если же кто-то бьёт такие прекрасные горшки, стоит задуматься --- а он ли их слепил.
У меня бы не поднялась рука.

С востока потянуло солёной прохладой, и волны стали забегать чуть дальше на берег.
Вот одна, затем другая достигла одинокого, похожего на торт домика.
Волны слизывали стены, словно крем.
Атрис рассеянно смотрел на останки сооружения.

--- Знаешь, чем мне нравятся дети?

--- Чем, Атрис?

--- Они замечательно умеют повторять, не понимая смысла.

--- Незачем понимать назначение и механизм дыхания, если это нужно для жизни, --- заметила Митхэ.

--- Но ребёнка необходимо научить правильно дышать.
Он не всегда будет в стандартных условиях существования.

--- Да, --- лицо Митхэ просветлело.
--- В Храме меня учили правильно дышать.

--- А если научить некому, то большинство сапиентов так и остаются детьми.

Митхэ кивнула и крепче прижалась к любимому.

--- Что было в сообщении? --- тихо спросил Атрис.

--- Он уже здесь.
С завтрашнего дня можем приступать, --- ответила Митхэ.

Атрис улыбнулся и погладил подругу по голове.

--- А вот теперь я бы сыграл.

\chapter{[U] Равновесие}

\section{[U] Зачисление в контрразведку}

\epigraph
{Тси --- это сапиенты, которые используют машины для обмена данными с внешней средой.
Демоны --- это машины, использующие для той же цели сапиентов.
Может быть, именно поэтому у нас с ними возникает недопонимание.}
{Длинный-Мокрый-Хвост}

\textspace

--- Компромата на вас собрано достаточно.

--- Компромата? --- выдохнула Анкарьяль.

--- Не обращай внимания, Нар.
У Самаолу своеобразное чувство юмора, --- буркнул Грейсвольд.
--- И он настолько любит шутить, что порой забывает поздравить демонов с самой бескровной успешной операцией со времён Тахиро.

--- Я не давал тебе слова, Грейсвольд Каменный Молот, --- ледяным тоном сказал Самаолу.

--- Если хочешь командовать чужими ртами, научись следить за своим.
Что ещё за компромат?

Самаолу демонстративно достал компьютер и начал листать слайды.
Игра была очевидна для всех --- демон и так был в курсе материала.

--- Первое.
Вы были друзьями Аркадиу Шакала Чрева.

--- И нас по этому поводу уже проверили, --- парировал Грейсвольд.

--- Разумеется.
И из проверки вытекает второе --- в вас замечена ощутимая симпатия к тси.

Грейсвольд и Анкарьяль промолчали.

--- То есть вы даже не будете этого отрицать?

--- В отрицании есть смысл? --- буркнул Грейсвольд.
--- Отдел 100 отвечает за верификацию полученных данных.
Кроме того, в своде законов Ада есть прямой запрет санкций за мыслепреступления.

--- Впрочем, этот запрет --- такая же формальность, как и текущее заседание, --- усмехнулась Анкарьяль.

--- Хорошо, --- кивнул Самаолу, пропустив слова Анкарьяль мимо ушей.
--- В таком случае не могли бы вы осветить подробнее своё отношение к тси?

--- Вообще-то я не обязан этого делать, но так и быть.
На планете Тра-Ренкхаль я впервые узнал, что искренность может быть культурной нормой.
Я первый и последний раз оказался в обществе, которое приняло меня тем, кто я есть.
Тси относились ко мне и к другим демонам так же, как к соплеменникам.

Пара советников испустили презрительный смешок.
Самаолу нахмурился.

--- А что скажете вы, Анкарьяль Кровавый Шторм?
Если я не ошибаюсь, сели вам поклоняются.

--- Не мне, Самаолу, а обезличенному образу, отпечатку личности давно умершего тела.
Вас ввели в заблуждение.

--- Не прибедняйтесь.
В образе хватает черт и вашей нематериальной части.
Очень любопытно включение в лик Самоотверженного Хата двух звёзд, двух лун и двух гор, которые некогда красовались на знамени Ангары, баронессы Краснобуря с планеты Тысяча Башен.

На лице Анкарьяль дёрнулся мускул.
Она была уверена, что компрометирующий <<подарок>> сделал ей Аркадиу, напоследок отведя от себя все подозрения.

\ml{$0$}
{--- Впрочем, неважно.}
{``However, it doesn't matter.}
В вашем отчёте прямо указано, что это совпадение, а в вас достаточно сложно найти замашки демиурга-индивидуалиста.
\ml{$0$}
{Всё это время вы были лезвием меча, частью команды, и не претендовали на иную роль.}
{All this time you've been an edge of the blade, a part of the crew, and never pretended to be someone else.}
Так что вы думаете насчёт тси?

Как ни лжива была ситуация, Грейсвольд заметил, что последние слова ей польстили.
Да, так и есть.
Лезвие меча, часть команды.

--- Скажу, что Катаклизм Тси-Ди был величайшей трагедией для нас, --- спокойно ответила Анкарьяль.

В зале наступила оглушительная тишина.
Грейсвольд вдруг тоже нахохлился.

--- Объясните, --- сказал Самаолу после продолжительного молчания.

--- Как вы знаете, мы находимся в мёртвой зоне Вселенной, --- заговорила Анкарьяль, и её слова эхом отдались в зале.
--- Наши перемещения ограничены константами Ка'нета, и, несмотря на все уверения технологов, с момента открытия их не удалось отодвинуть даже на одну триллионную яо.
Предпринималось множество попыток найти разумную жизнь в пределах ка'нетовского радиуса, но ни одна не увенчалась успехом.
Огромный шар диаметром в тридцать шесть тысяч парсак --- и всего двадцать девять планет с самозарождённой стабильной жизнью, из них с разумной --- две!
Вдумайтесь в это.
Наши проекты по терраформированию терпят крах один за другим.
То есть, говоря без обиняков, мы, высокоразвитая цивилизация хоргетов, не можем сделать даже того, что сделали древние земляне.
Вы сидите здесь, гордитесь своими технологиями, многие из вас презрительно называют первых людей <<обезьянами>>, но по сути вы --- посмешище для собственных создателей.

--- Терраформирование?
Вы серьёзно? --- заливисто рассмеялась советник Кагуя Хрустальный Голос.
--- И это тогда, когда Картель активно возвращает себе позиции на Плеядах и Развязке Десяти Звёзд?

--- Война для вас важнее, Кагуя? --- невинно поинтересовался Грейсвольд.
--- Ведь благодаря ей вы сидите в этом мягком кресле?

Анкарьяль подняла руку, показывая, что ещё не закончила говорить.
Самаолу жестом остановил попытавшуюся ответить Кагуя и кивнул интерфектору.

--- Ещё касательно эффективности Ордена.
Возьмём Тра-Ренкхаль.
Аркадиу Люпино сразу после битвы на Могильном берегу вызвал геологов, чтобы они остановили землетрясения.
Да, --- повысила голос Анкарьяль, перекрывая смешки, --- да, \emph{предатель} озаботился обустройством владений Ордена.
С тех пор прошло сто пятьдесят четыре дождя.
Где геологи?
Из-за предательства Аркадиу его запрос потерял актуальность?
Безымянный до сих пор обеспечивает безопасность сапиентов за счёт собственных ресурсов.

--- Мы не обсуждаем на этом заседании политику Ада по отношению к Митрису Безымянному, Анкарьяль, --- холодно бросил один из советников.

--- Ах, это аспект отношений с Безымянным, а не простая халатность? --- проворчала Анкарьяль.
--- Благодарю за разъяснение, советник, вы спасли меня от крамольных мыслей.

--- Тишина, --- призвал к порядку Самаолу.
--- Вернёмся к обсуждению тси.
Грейсвольд, вам есть что добавить?

--- Я не вполне согласен с формулировками Анкарьяль, но по сути она права.
Мы столкнулись с теми же проблемами, которые когда-то потрясли цивилизацию Древней Земли.
Однако мы гораздо могущественнее первых людей, наши ошибки гораздо более серьёзны, и потому масштаб катастрофы нам оценить едва ли не сложнее, чем им.
Даже если предположить, что исход войны близок и Картель падёт, куда, скажите на милость, направится взор нашей военной машины?
Хотите сказать, что зверь, которого кормит весь Ад без исключения, покорно сложит лапки и заснёт вечным сном?

--- Враги есть всегда, Грейсвольд, --- заявила Кагуя.

--- И вы отлично умеете их находить, --- поклонился Грейсвольд, --- даже там, где их отродясь не было.
Что вам сделал Безымянный?
Такого покладистого демиурга ещё поискать.
С демиургами всегда были проблемы --- взять того же Эйраки, или Кох Свободолюбивую, или...

--- Или вас, --- присовокупила Кагуя.

--- Безымянный --- лучшее, что есть на Тра-Ренкхале, --- громко и членораздельно сказал Грейсвольд.
--- Он подобен тем наивным древним учёным --- философам, астрологам, алхимикам, колдунам, жрецам, --- которые выстраивали сложные мистические теории мироздания, которые безо всякой научной и идеологической базы искали, находили и делали выводы.
Он художник, и его планета --- один большой холст.
У меня как демиурга никогда не хватило бы фантазии на то, что сделал Безымянный.
Однако Ад выстроил против этого миролюбивого создания такую защиту, словно на Тра-Ренкхале хранится эссенция Красного Картеля в хрупком стеклянном флаконе.
Мне даже довелось слышать, как Безымянного сравнили с Уэсиба Серозмеем.

Кагуя вдруг густо покраснела, но заметила это только Анкарьяль.
Всё внимание советников было приковано к Грейсвольду.

--- Если Безымянный так опасен, почему он ещё жив?

Совет замер.
В воздухе повисла неловкая тишина.

--- Я серьёзно, --- как ни в чём не бывало продолжил технолог.
--- Мы прекрасно знаем, что любые договоры --- лишь красивые слова, если за ними не стоит реальная сила.

--- Грейсвольд, вы в курсе, что это?..

--- Я знаю, что это основа деспотической доктрины Картеля и, чёрт возьми, это честный и здравый взгляд на вещи.
Демоны имеют право голоса не потому, что это правильно, а потому, что за каждым демоном стоит его собственный трудовой и военный потенциал, которым он готов распоряжаться по собственному усмотрению.
Организация защищает своих членов не потому, что это правильно, а потому что именно это делает её организацией.
Да, договор --- это лишь символ.
Тем более этот договор, с громким названием <<Демиург --- Метрополия>>, клише которого было создано, между прочим, по моей инициативе, на всякий случай.
Я надеялся, что это будет подспорьем для социализации демиургов, позволит тем, кто не готов перейти на демонический образ жизни, быть равными нам, не быть хоргетами второго сорта.
Однако субъектов договора в истории Ордена всего шесть, в настоящее время --- один-единственный, если не считать <<без вести пропавшую>> Кох.
Счёт же уничтоженных богов --- тех, чьё уничтожение было задокументировано --- идёт на десятки тысяч.
Так что же мешает Аду \emph{ещё раз} устроить тотальную мелиорацию и окончательно решить вопрос демиурга и девиантной фауны?

\ml{$0$}
{--- Вы подозреваете Орден в уничтожении Кох, Грейсвольд? --- вмешался один из советников.}
{``Are you suspecting the Order of destroying Coj, Grejsvolt?'' a counsellor intervened in an argument.}

\ml{$0$}
{--- Ни в коем случае, советник.}
{``No, counsellor.}
\ml{$0$}
{Я не подозреваю, я в этом уверен.}
{I'm not suspecting, I'm sure of it.}
Разыгравшийся вокруг Безымянного спектакль мне до боли напоминает уже отгремевшие и давно забытые события.
И знаете, есть некая ирония в том, что Аркадиу Люпино, чьё первое тело выросло под крылом Богини-Матери Пустоши Драконов, предал Орден в тысячную годовщину того, как Богине-Матери помогли найти дорогу в небытие.

--- Может, Грейсвольд, вам снова выдать планету? --- усмехнулся кто-то в зале.
--- Вы с таким упоением расписываете прелести жизни демиурга...

Зал захохотал.

--- А вы знаете, я только за, --- подбоченился технолог.
--- И моя мечта исполнилась бы, если бы не всеобъемлющая зависть.

Зал снова огласил хохот.

--- А разве я не прав? --- продолжал Грейсвольд.
--- У каждого из хоргетов достаточно способностей, чтобы инкарнироваться в планету.
Извините, что напоминаю об этом, но нас для этого и создавали, не так ли?

\ml{$0$}
{--- Не имеет значения, для чего нас создавали, Грейсвольд, --- заявил Самаолу.}
{``What we were created for doesn't matter now, Grejsvolt,'' Samajolu said.}

\ml{$0$}
{--- Отнюдь!}
{``It does matter!}
Едва ли такая проблема стояла перед первыми людьми --- они были самообразовавшимися живыми системами и вольны были сами выбирать своё предназначение.
Но мы --- продукт инженерии, разумный замысел, а не игра слепой эволюции.
\ml{$0$}
{Да, мы все --- прирождённые демиурги.}
{Yes, we all are demiurges by design.}
Но мы упаковали себя в микроскопические сапиентные тела, мы красим волосы и глаза в немыслимые цвета, мы отращиваем себе когти и щупальца, мечтая о движении континентов, лесах из светящихся грибов и цивилизациях, расцветающих под нашей опекой!

--- Что вы пытаетесь этим сказать?

--- Пропаганда самоограничения самореализации была одним из самых древних способов получить и удержать власть.
Прямо или косвенно до нас доносят мысль, что быть демоном --- правильно, а богом --- нет.
И пока мы принимаем это за истину, все мы, вне зависимости от ранга, не более чем рабы.
Это всё, что я хочу сказать.

--- Если ваша соратница имеет привычку сквернословить, Грейсвольд, то вы явно любите драматизировать, --- буркнул Самаолу.
По галерее пронёсся смешок.
--- Я предпочитаю первое.
Вернёмся к обсуждению тси.
Итак, что бы вы предложили, Анкарьяль?
Какую политику следовало бы проводить Ордену в отношении тси?

Несколько советников, включая Кагуя, снисходительно улыбнулись.

--- Политику мира.

--- Сильно, --- признал один из советников, Кес Бледный Глаз.
--- Я думаю, нет смысла спрашивать, как этот мир достичь.
Но ради справедливости, Анкарьяль, я повторю вопрос вашего товарища Грейсвольда --- где были бы \emph{вы}, не будь войны?

Треть совета смущённо кашлянула.
Прочие взглянули на Кеса, словно он справил нужду прямо посреди зала.

--- Без войны я была бы в опасной экспедиции за пределами ка'нетовского радиуса, --- сказала Анкарьяль.
--- Я бы пронзала пространство и время, искала новые земли, улыбалась новому свету, а в минуты отдыха размышляла бы, как послать весточку вам, старому миру, как сообщить затерянному среди галактик дому, что мы тоже живы и процветаем.
Мой друг Грейсвольд был бы демиургом, и едва ли наши пути пересеклись бы.
Но, как видите, мы стоим сейчас плечом к плечу, отстаивая право быть собой.

--- Вам никто не мешает быть собой, Анкарьяль --- разумеется, в пределах закона.
И я так и не понял, при чём тут цивилизация Тси-Ди.

--- Стабильная технократическая цивилизация сапиентов, которой являлись тси, могла колонизировать дальний космос, дать вам плацдармы для дальнейших исследований и продвинуть границу Ка'нета на тысячи парсак вглубь галактики, к её ядру.
Наши исследования дают однозначные результаты --- в ядре мы сможем найти достаточно подходящих для жизни планет.
Некоторые исследователи даже говорят, что имеющихся там ресурсов достаточно, чтобы создать и запустить звёздный корабль на разумных антаридах, преодолеть...

--- Оставьте эту чушь, Анкарьяль, --- сказал Самаолу.
--- Звёздный корабль --- очередная сказка для романтиков вроде вас.
Разумные антариды --- из немного другой сказки, ещё менее правдоподобной.
Жизненный цикл антарид настолько короток, что исключает наличие сапиентной архитектуры любого уровня.
Хоргеты не способны выйти за пределы Млечного Пути.
Это возможно лишь для микоргета, но последние исследования ставят под вопрос саму возможность его существования.

--- Антариды...

--- Самаолу прав, Анкарьяль, --- вдруг подала голос Корхес Соловьиный Язык.
--- Даже если создать разумных антарид, существование антаридной цивилизации не будет долгим.
Антариды --- это вам не тси.
Мы получим либо безжизненную звезду, либо взрыв сверхновой в самые короткие сроки, звёздный корабль даже не успеет выйти за пределы Млечного Пути.
И я бы попросила закрыть эту тему, потому что ваша квалификация как биолога --- и тем паче плазмобиолога --- оставляет желать лучшего.
\ml{$0-[ej]$}
{А я люблю науку и не люблю, когда военные, политики и прочее жульё начинают лапать её своими грязными ручищами.}
{I love science, and I hate when military, politicians, and other crooks touch it with their dirty hands.''}

--- Корхес, выражения аккуратнее выбирайте, --- призвал к порядку Самаолу.

--- Даже если и так, --- Анкарьяль с трудом сдерживала раздражение, --- почему хотя бы не освоить до конца Млечный Путь, почему не заручиться союзом с тси?
Да, скорость их кораблей ограничена светом, но они достигли бы большего успеха за единицу времени, чем вы.
И я уверена, что даже если бы Ад и Картель обрушились всей своей мощью на колонистов, тси всё равно нашли бы время на обустройство новых планет.
Они не считали нужным постоянно просиживать у орудий.

\ml{$0-[ej]$}
{--- Вы думаете, что мы не рассматривали такие варианты? --- высокомерно улыбнулся Самаолу.}
{``Do you really think we haven't considered such an option?'' Samajolu arrogantly smiled.}

\ml{$0$}
{Анкарьяль явно хотела сказать, что она думает о Самаолу, но сдержалась.}
{Ancarjal obviously wanted to speak out what she really thought about Samajolu, but stopped herself.}

--- В таком случае объясните нам с Анкарьяль, какие именно нюансы она не учла, --- проворковал Грейсвольд.

--- Тси знали о существовании демонов, --- медленно, словно для умственно отсталых, начала объяснять Кагуя.
--- Более того, им удалось достигнуть с нами стратегический паритет...

--- Опять военный термин, --- поморщился Грейсвольд.

--- А какие термины вы предпочитаете применять к враждебной цивилизации?
Это война, Грейсвольд.
Мы пытаемся выстоять против чужой агрессии.

--- То же самое говорили миллионы агрессоров до вас, --- сказал технолог.
--- Слепцы по-прежнему верят, что есть праведные войны.
Романтики считают, что время праведных войн прошло.
А я, видевший историю Земли почти с самого её рассвета, скажу --- праведных войн нет и не было.
В любой войне число агрессоров равно числу участников!

--- Вы предлагаете сложить оружие и отдаться на милость Картеля? --- ласково осведомилась Кагуя.
--- Или вы, Анкарьяль.
Вы позволили бы тси колонизировать наши планеты?
Капитул?

--- Вы никак не можете понять, что тси --- это не обычные сапиенты, которых мы используем как дойных животных! --- повысила голос Анкарьяль.
--- Они равные нам, а кое в чём и превосходящие нас существа!

--- Тси --- это хамелеоны.
Они легко интегрируются в любое общество и бессознательно внедряют в это общество свои ценности, распространяют свою заразу...

--- Как и представители любой древней, выкристаллизовавшейся культуры, --- заметил Грейсвольд.
--- Мало какая культура может состязаться с тси в возрасте.
За двести тысяч лет культура тси въелась в их генофонд, приобрела устойчивость к множеству различных воздействий, и даже жестокий отрицательный отбор не даст быстрых результатов.

--- Особенно со стороны менее древней культуры, --- ядовито добавила Анкарьяль.
По галерее пробежал возмущённый шёпот.
--- <<Зараза>>!
Вы говорите об информации, словно о смертельном вирусе, от которого у вас нет защиты.
По-вашему, адепты Ада не способны фильтровать информацию?
Мне доводилось слышать мнения, что народ Тси-Ди --- тепличные растения и пример культурного дрейфа\FM.
\FA{
Культурный дрейф --- ненаправленные изменения культуры, обусловленные случайными статистическими причинами.
}
И вдруг выясняется, что их культура --- <<зараза>>, от которой нужно защищаться!
Что это значит?
Выходит, перед нами не дрейф, а стадия эволюции, до которой нам с вами ещё далеко?

--- Так или иначе, наиболее яркие представители погибли на Могильном берегу, --- с гадливой жалостью сказала Кагуя.
--- Остаткам культуры был нанесён непоправимый удар.

--- Я бы не делала столь поспешных выводов, --- улыбнулась Анкарьяль.
--- Вы недооцениваете тси, и это вам ещё аукнется в будущем.

--- Да, действительно, есть исследователи, считающие, что сам контакт Ордена Преисподней с культурой тси опасен, --- хмуро сказал Самаолу.
--- Вы поддерживаете это мнение, Анкарьяль?

--- Опасность контакта оценивают контактирующие согласно собственному опыту, --- туманно выразилась Анкарьяль.
Кто-то из советников нахмурился, кто-то улыбнулся, Кагуя же фраза привела в бешенство.
--- Но изменения определённо произойдут.

--- То есть вы хотите сказать, что опасность мы придумали сами? --- спросила Кагуя, едва сдерживаясь.

--- Тси угрожали убить вас лично, Кагуя?

--- Что за профанация, Анкарьяль? --- опешила советник.
--- Что вы имеете в виду?
Я член Ордена Преисподней, а это значит...

--- Это значит, что ответственность за статус-кво лежит в том числе и на вас! --- перебила Анкарьяль.
--- Почему-то как речь об опасности --- то она общая, а ответственность за политику чья угодно, но только не ваша!
В отличие от Картеля, у нас была возможность установить с тси договорные взаимоуважительные отношения, но мы использовали их как разменную монету в войне с тем же самым Картелем.
Как и в случае с Безымянным, Орден выбрал войну вместо мира.
Я боюсь, что если ретроспективно подвергнуть такому же анализу действия Ордена за последний телльн, то окажется, что большинство реше...

--- Довольно, --- рявкнул Самаолу.
--- Мы поняли вашу точку зрения.
Напоминаю, что Орден Преисподней предлагал тси протекторат, тси это предложение отвергли.
Это было их и только их решением!

--- Протекторат? --- задохнулась Анкарьяль.
--- Вы это называете <<взаимоуважительными отношениями>>?
Я бы сделала то же самое, да ещё и плюнула вам в лицо!
Тси сдерживали Картель куда эффективнее, чем вы!
Тси-Ди была самым спокойным местом во Вселенной!

--- Вы называете Тси-Ди спокойным местом? --- спросил советник.
--- Вот вам немного статистики. Основная причина смерти тси --- несчастные случаи на производстве.
Девяносто четыре процента всех смертей.
Тси жили на своих великих стройках и умирали на них же, как рабы Эйгипта, которые строили для своих фарай величественные гробницы.
Это цивилизация рабов, ничем не отличающаяся от прочих сапиентных цивилизаций.

--- И даже их конец был одним большим несчастным случаем на производстве, --- ядовито добавила Кагуя.
--- Вы технолог, Грейсвольд?
Вы понимаете, о чём я говорю.

--- Вы искажаете суть статистических данных, --- возразил Грейсвольд, проигнорировав слова Кагуя.
--- Тси умирали при несчастных случаях не из-за обилия таковых, а из-за того, что до смерти от естественных причин они просто не доживали.
Смертность по медицинским причинам контролировать гораздо легче, чем смертность от несчастных случаев, и тси это показали очень хорошо.
Собственно говоря, я не слышал ни одной истории о демоне, который мирно прекратил существование у себя дома.
Основные причины смертности у нас --- убит, казнён, пропал без вести.
Это моя собственная статистика, я её веду по друзьям и знакомым...

--- Довольно, --- поднял руку Самаолу.
--- Ваши статистические выкладки никем не проверены и потому не представляют для нас интереса.

--- И я совершенно не понимаю аналогии с гробницей, --- заметила Анкарьяль.
--- Если уж вы хотите заручиться союзом с историей Древней Земли, то я сделаю то же самое.
Вы были для них вандейлами, которые молотили каменными топорами в прекрасные стабитаниумовые ворота Ромай!
Тси умирали, пытаясь починить дом, который...

--- Довольно! --- рявкнул Самаолу.

Анкарьяль хмуро скрестила руки на груди.

--- Впрочем, даже жаль, что тси отгородились от нас, --- отчетливо сказала она в тишине, пренебрегая запретом.
\ml{$0-[ej]$}
{--- Если бы демоны жили на Тси-Ди на равных условиях с аборигенами, если бы они видели, как эти существа из плоти и крови относятся друг к другу, Ад и Картель, эти колоссы на глиняных ногах, развалились бы в первые же двадцать лет.}
{``If daemons could live on Qi-Di with its dwellers, as equals, if they could see, how beings of flesh and blood used to treat each other, the Hell and the Cartell---these colossi with feet of clay---have come apart in the first twenty years.''}

Грейсвольд запоздало схватил подругу за плечо, пытаясь её остановить.
\ml{$0$}
{Слово прозвучало.}
{The words have been said.}
Советники потрясённо молчали.
Казалось, в этой оглушительной тишине слышно, как бьются сердца и гудят стабилизирующие модули омега-сингулярностей --- если бы их колебания были доступны для слуха сапиента.

Наконец Самаолу вздохнул.
Это был вздох человека, который утвердился в худших своих опасениях.

--- Мы, --- он обвёл рукой зал, --- все мы могли бы много чего сказать на это скоропалительное, я бы даже сказал --- безответственное заявление.
Но мы здесь не для того, Анкарьяль Кровавый Шторм, чтобы разъяснять вам очевидное, тем более что имеет место проявление некой идеологии, а не простое заблуждение.
\ml{$0$}
{Идеология, в отличие от заблуждения, прекрасно умеет защищаться от фактов и теорий.}
{Ideology, unlike mistake, can perfectly defend itself against facts and theories.''}

\ml{$0$}
{--- К сожалению, Самаолу, это так, --- сухо сказала Анкарьяль.}
{``Unfortunately, it's true, Samajolu,'' Ancarjal coldly said.}

\ml{$0$}
{--- Я рад, что хоть в чём-то мы сошлись, потому что у нас с вами нет другого выбора, кроме как найти общий язык.}
{``I'm glad we found a thing we both agree, because we have no option but cooperation.}
Как я уже говорил, мы набрали массу интересного материала касательно вас... --- Самаолу позволил себе загадочную паузу, --- и тем не менее к вам есть предложение поступить в резервный состав отдела 100.

Грейсвольд и Анкарьяль удивлённо промолчали.

--- Обретает форму новая сила, --- сказал Самаолу.
--- У нас недостаточно данных о её природе, но достаточно оснований предполагать её враждебность.

--- Говори проще, Самаолу, я тебе пелёнки когда-то менял, --- поморщился технолог.
Он взглянул на Анкарьяль --- демоница вдруг побелела как полотно.

--- Контрразведке нужны демоны, которые способны думать, как враг.

--- То есть вы даже не скрываете, что считаете нас врагами! --- вспылила Анкарьяль.
--- <<Держи друзей близко, а врагов ещё ближе>>, да?
Я не буду переводиться на таких унизительных условиях!

--- Что именно вы считаете унизительным? --- осведомился один из советников.

--- Я вернулась с Тра-Ренкхаля, где ради интересов Ордена чуть не отправилась \emph{жилами джунглей}! --- выкрикнула Анкарьяль.
Самаолу прищурился.
\ml{$0-[ej]$}
{--- Какие бы ошибки я не совершила в молодости, сейчас я легат терция и заслужила карьерный рост, а не участь подопытной мыши или надувного монстра, которого вы будете трахать в надежде избавиться от собственного страха!}
{``Whatever mistakes I made in my youth, now I am legate tercia and I have deserved career progression, not a fate of guinea pig or inflatable monster which you will fuck in the hope of conquering your own fear!''}

--- А само зачисление в ряды контрразведки вы не считаете повышением?

Анкарьяль в десяти словах на языке фоуф-у\FM{} объяснила Совету, что она думает о таком повышении.
\FA{
Пиджин наёмников Тысячи Башен, вобравший в себя рекордное количество обсценных слов из других языков --- 6\% корней.
В основном эти слова касаются сексуальных отношений, физических и психических недостатков, отправления естественных нужд, состояний сознания при приёме различных наркотиков, а также изощрённых способов пытки и убийства.
}
Советники вытаращили глаза, когда семантический переводчик выдал шестнадцать слайдов мелкого текста.
Немного знакомый с этим языком Грейсвольд поперхнулся ещё на третьем слове.
Входящий в Совет диктиолог-лингвист Кельса Пушистая украдкой включила записную книжку.

--- Кельса, это уже занесли в протокол, --- ледяным тоном сказал Самаолу, и демоница тут же убрала терминал.
--- Если сквернословие представляет для тебя научный интерес, я отправлю официальную выписку в твой отдел.

<<Нар, я тебя люблю, --- с нежностью подумал Грейсвольд.
--- Так блестяще сыграть прожжённую карьеристку и высказать правду в лицо не смог бы даже я.
Теперь моя очередь купаться в этом вонючем болоте, доказывая свою профпригодность.
Чтоб этого Аркадиу поглотила бездна --- сам пропал, а нам оставил тележку с навозом!..>>

--- И всё-таки, к чему устраивать этот спектакль с отделом 100, если нас просто хотят проверить на верность? --- осведомился технолог.

--- Совет похож на театр, Грейсвольд Каменный Молот? --- один из советников позволил себе улыбку.

--- И весьма, Эйяфар Цеппелин.
Вы с Анкарьяль достаточно молоды, чтобы не знать, кто разрабатывал правила функционирования отдела 100.
А есть в их своде следующее --- за пределами, простите за фигуральность, можно играть камнями всех четырёх цветов\FM, но в самом отделе должен царить абсолютный порядок.
\FA{
Намёк на Метритхис, настольную игру сели.
}
И резервный состав --- не исключение.
Возможно, я ошибаюсь и правила поменялись?

Пара сидящих в зале зашевелились.

--- Вы хотите проверить меня на верность?
Прошу вас.
Вы видите во мне лишь впадающего в деменцию старика.
А я вижу шестнадцать горшков, которые звенят на слепившего их горшечника.

--- Оставьте литературные обороты при себе, --- сказал Самаолу.
--- Отдел 100 только что сообщил, что вы зачислены в резервный состав.
Официальная формулировка --- <<за многократно доказанную лояльность и высочайший профессионализм>>.
Это должно вас устроить.
Инструктаж будет проведён через двадцать секунд на станции 80, о ваших функциональных телах позаботятся.
Заседание окончено.

Грейсвольд застыл столбом и взглянул на подругу.
Анкарьяль положила руку ему на плечо, и оба упали замертво посреди зала.

\section{[U] Большая игра}

\textspace

--- Итак, Штрой, вы всё-таки придерживаетесь мнения, что это ковен Картеля, --- произнес голос.

--- Именно, владыка, --- Штрой говорила чётко и размеренно, былая жеманность улетучилась без следа.
--- Почерк Картеля я узнаю всегда и везде, потому что сама приложила к нему руку.

--- И поэтому вы ввели их в заблуждение.
Что ж, неплохо.
А вы, Самаолу?

--- Несмотря на сравнительный анализ Штрой, связи с Картелем я не нашёл, --- угрюмо сказал Самаолу.
--- Гораздо больше фактов говорит о том, что в деле замешана третья сторона.
А как считаете вы?
Вы проводили какие-то исследования...

--- Их результаты вас не касаются, --- спокойно отрезал голос.
--- Продолжайте разрабатывать собственные версии.
Кстати, Самаолу, ваш трюк с Анкарьяль и Грейсвольдом был просто великолепен.
Так блестяще сыграть тупицу я бы не смог, у вас большие задатки... кхм... театрального актёра.
Поблагодарите за меня Кагуя --- она тоже прекрасно справилась.

--- Обязательно, --- слегка улыбнулся Самаолу.
--- Приятно иногда вносить изменения в правила, особенно в правила отдела 100.
Относительно этих демонов план остаётся прежним?

--- Разумеется, --- подтвердил голос.

--- Правда, я не совсем понимаю, в чём их важность, --- полувопросительно заметил Самаолу.

--- Её действительно сложно понять, --- согласился голос.
--- Анкарьяль и Грейсвольд заняли весьма неудобное для нас положение --- <<красные плащи в узком ущелье>>\FM.
\FA{
Тактический ход, возможно, восходящий к легендарной битве у Фермопил, Древняя Земля.
}
Они имеют богатый опыт совместной работы и безоговорочно друг другу доверяют.
Собственно, нам нужно изучить этот тандем в действии.
Некоторые расчёты дают право полагать, что пара Анкарьяль-Грейсвольд является \emph{единственным} связующим звеном между непринадлежными Скорбящими и агентурной сетью Скорбящих внутри Ада.

--- Провокатор-связной, --- ввернула Штрой.
--- Вызывающий, постоянно находящийся на виду, специально обученный отражать атаки и вызывать сомнения в собственной виновности.
Почерк Картеля.

--- Почему мы не можем их просто устранить физически? --- поморщился Самаолу.

--- Во-первых, Самаолу, это не решит конкретную проблему.
То, что деталь единственная в механизме, вовсе не означает, что она незаменима.
Во-вторых, провокатор-связной чрезвычайно устойчив и к физическому воздействию --- чаще всего это опытный интерфектор, даже удачная попытка будет стоить вам несколько ценных кадров, что уж говорить о неудачной.
Публичную личность в принципе сложно устранить без последствий.
А в-третьих --- оставьте привычку решать проблемы силой.

--- У нас вполне достаточно ресурсов, чтобы...

--- Да, да, у нас достаточно ресурсов, чтобы от всех наших соперников --- по крайней мере, внутренних --- не осталось даже воспоминаний.
Но в таком случае можно сразу расписаться в неумении управлять.
Решая проблемы силой внутри Ордена, вы лишите себя ценнейших ресурсов в борьбе с внешним противником.
Да, я считаю всех Скорбящих ценным ресурсом, Штрой, --- как ни крути, они смогли сделать очень многое.
Поэтому мы будем держать волка в клетке, пока у нас есть хотя бы призрачная возможность договориться.
Тем более, как Самаолу уже сказал, когда-то Анкарьяль и Грейсвольд были вернейшими адептами Ордена, а это значит, что они могут снова стать такими.

--- Никогда не доверял предателям, --- буркнул Самаолу.
Штрой ощерилась.

--- Тихо, --- голос поднял невидимую ладонь, и демоны разом вытянулись в струнку.
--- Предатель, Самаолу --- это прежде всего тот, кто умеет сравнивать и не боится делать выбор.
То, что Анкарьяль и Грейсвольд в итоге оказались в агентурной сети Скорбящих, означает три вещи.
Первое --- с Орденом что-то не так, и это очевидно для всех, кто способен думать.
Второе --- они не считают Орден безнадёжно загнившим, так как формально остались на службе, а не сбежали, как крысы с тонущего корабля.
Да-да, Самаолу, после зачисления в отдел 100 они могли сбежать, а не подвергать себя риску --- но выбрали риск.
И третье --- они растут как личности и как профессионалы, и это очень важный момент.
Впрочем, действовал этот тандем очень осторожно, но наживку заглотил с половиной удочки.
Они не настолько умны, как загадочные стратеги 03 и 05.

--- Возможно, владыка, нам следует расширить выборку демонов, подпадающих под...

--- А какой смысл, Штрой? --- возразил голос.
--- Даже если бы сейчас стратеги 03 и 05 явились ко мне лично и принесли исчерпывающий компромат на самих себя, это мало что изменило бы.
Вы думаете, что это зелёная молодёжь, которая ищет своё место в жизни?
Основную проблему я вам очертил сразу --- мятежников поддерживают группировки, с которыми уже давно достигнуто равновесие Нэша.
Говоря другими словами, фигуры, которые давно и прочно стоят в своих клетках, неожиданно поменяли цвет, сделав ставку на молодого и неопытного игрока.
Когда я впервые осознал этот факт, то был озадачен.

--- Владыка, я не совсем понимаю...

--- А, Штрой, прошу прощения, я забыл про особенности обучения стратегов в Картеле.
Вы читали доклад Самаолу?

--- Разумеется!

--- И какие ваши впечатления?

--- Анкарьяль и Грейсвольд пожертвовали своим по...

--- Именно, --- откликнулся голос.
--- Штрой, вы потрясающий стратег.
Даже Самаолу посчитал произошедшее гамбитом со стороны Ада.

--- А теперь, похоже, не понимаю я, --- протянул Самаолу.

--- Именно поэтому вы со Штрой работаете вместе, --- объяснил голос.
\ml{$0-[ej]$}
{--- Политика --- странная сводня, но вы, должен признать, замечательно друг друга дополняете.}
{``Politics makes strange bedfellows, but I have to admit that you two are perfectly complementary.''}

Штрой и Самаолу бросили друг на друга взгляды, полные отвращения.

--- У вас, как я понимаю, нет статистики по отделу 100?
Так слушайте: за последние четырнадцать лет контрразведка потеряла в три раза больше демонов, чем мы на полях сражений.
Туда берут лучших --- и эти лучшие мрут, как тараканы.

Самаолу ахнул.

--- Они побледнели...

--- ... от страха, а не от злости, --- закончила Штрой.

--- Браво.
Вы поняли.
У них есть источник информации, который держит их в курсе дел.
Перевод из легатов запаса в центурионы отдела 100 --- это даже не смешно.
Анкарьяль и Грейсвольду придётся пройти через кошмар, прежде чем они закрепятся на новом месте.
Они знали о предстоящем возмездии за изменение собственной стратегии, но не подозревали о его силе.
Истинное равновесие Нэша --- жестокая вещь.

--- И такие потери в позициях понесли все, кто поддержал эту группировку? --- спросила Штрой.

--- Не такие, Штрой.
Нашим друзьям просто сильно не повезло.
Подозреваю, что Грейсвольд принял на себя часть удара, предназначенного не ему --- он едва избежал смерти уже как минимум два раза, если считать последний инцидент на Тра-Ренкхале.
Возможно, он защитил этого неуловимого стратега 03.

--- Вы считаете события на Тра-Ренкхале платой за изменение стратегии? --- поднял бровь Самаолу.
--- У меня создалось впечатление, что...

--- Друг мой, стратег 03 действует не первое столетие.
И тихоня Грейсвольд --- его главный помощник.
Вероятность того, что инцидент на Тра-Ренкхале бы нэшевским рикошетом, стремится к единице.

--- Кто-то из стратегов крупно просчитался, --- улыбнулся Самаолу.

--- Или сделал чересчур широкий шаг вперёд, --- подхватил голос.
--- Никогда не мешает предположить худшее, особенно если учесть, что за прошедшие пять лет <<красные плащи>> неплохо отыгрались.
Награды сыпались на Грейсвольда с Анкарьяль словно из мифического рога изобилия.
Честно говоря, я думал, что это заставит их отказаться от полевой работы.
Но я ошибался.
Что ж, хотят опасностей --- у центурионов отдела 100 их хватает.

--- Владыка, не может ли стратег 03 быть Падальщиком? --- спросила Штрой.
--- Несмотря на явную разницу в уровне, в почерке есть общие черты...

--- Исключено, --- сухо ответил голос.
--- Сходство можете списать на что угодно --- на подражание, случайность или наставничество.
Падальщик --- пешка.
Оставьте его в покое, он предсказуем и может быть уничтожен в любой момент.
На звание одного из стратегов у нас есть претенденты куда старше и опаснее.

--- Стигма, --- проворчал Самаолу.
--- Я употребил всё своё влияние, чтобы отослать её на Тра-Ренкхаль, подальше от основного театра действий.
Любое обвинение стекает со Стигмы, как ртуть со стекла.

--- Владыка, раз уж Стигма оказалась на Тра-Ренкхале, я могу устранить её физически, --- предложила Штрой.

--- Даже не пытайся, --- резко сказал Самаолу.
--- Она съест тебя на завтрак и не посмотрит на твои лычки.
Лимес, Бит и Уитлик умерли не без её помощи...

--- И всё же я прошу у владыки разрешения, --- поклонилась Штрой.

--- Штрой Кольцо Дыма, --- голос позволил себе оттенить слова великолепной усталой интонацией.
--- Оставьте свои манеры для деспотического Картеля.
Советую вам ознакомиться с протоколом <<Шляпа>>, который является сводом приемлемых для Ада норм общения.

Штрой смешалась, но тут же овладела собой.

--- Так точно, вла...
Хорошо, я учту.

--- Так-то лучше, --- милостиво откликнулся голос.
--- Насчёт Стигмы --- попробуйте, я отведу от вас удар.
Удача может не прийти к пытающемуся, но к бездействующему она не придёт гарантированно.
Устранить Стигму нужно было ещё пятьсот лет назад, как по мне --- несмотря на коварство и приспособляемость, в некоторых вопросах она не склонна к компромиссам.

--- Вы пытались с ней договориться? --- спросил Самаолу.

--- И не только я, --- вздохнул голос.
--- Проблема в том, что она пацифист до мозга костей.
Её устраивает любой вариант, исключающий военные действия против Картеля.

--- И при всём миролюбии Стигмы те, кто пытается её устранить, крайне плохо кончают, --- ввернул Самаолу.
--- Чего стоит эта история с кланом Чёрной Скалы...

--- Да, я слышала про стратега, который якобы охотится на собственный клан, --- задумалась Штрой.
--- Значит, это Стигма?

--- Я думаю, в данном случае применимо слово <<вырезает>>, а не <<охотится>>, --- усмехнулся Самаолу.
--- Из ста двадцати трёх членов в живых остались пятеро.

--- Трое, --- поправил голос.
--- Двойняшки, Фуси и Нуива, погибли на Тра-Ренкхале от рук рекомого Грейсвольда.
Судьба не лишена иронии.

--- А какие вообще доказательства, что Стигма истребляла выходцев из клана Антрацис? --- спросила Штрой.

--- Никаких, --- ответил голос.
--- Как нет и прямых доказательств, что клан вообще кто-то целенаправленно уничтожал.
Могу сказать только одно --- с демонами Антрацис происходили несчастные случаи, этих несчастных случаев было много, и Стигма каждый раз оказывалась где-то недалеко --- участвовала в планировании операций, либо имела опосредованные связи с исполнителями.
Были и весьма нехорошие происшествия, которые позволяют оценить её методологию.
Например, трое перебежчиков из Антрацис успешно добрались до Капитула и были уничтожены уже там агентом всё той же Чёрной Скалы.
Вопиющий случай --- вы знаете лучше меня, Штрой, какие ресурсы тратятся на защиту перебежчиков.
Известно, что это члены её генерации, которые в своё время принимали участие в травле.
Учитывая, что Стигма просто не могла нанять ассассина изгнавшего её клана, могу предположить, что она сделала --- или не сделала --- всё, чтобы он успешно выполнил задание.

--- А ассассин далеко не ушёл, --- присовокупил Самаолу.

--- Да, кстати, уйти ему не дали.
Спасибо, что напомнили, Самаолу.
Четверо мух одним лёгким ударом --- разве не впечатляюще?
Собственно, одна из причин, по которой отдел 100 крайне неохотно подключает Стигму --- это её своеобразное, карикатурное понимание пацифизма.
Если на тебя охотятся --- истреби клан, который это делает.
Если для мира необходимо сократить число солдат, а демобилизация недоступна --- она будет сокращать их число другими способами.
Школа Чёрной Скалы, работа на результат.

--- А кто ещё из Антрацис остался в живых?

--- В живых, кроме Стигмы, остались Ду-Си Охотник и Корхес Соловьиный Язык --- такие же Вечно Гонимые, как и она.
Их она трогать не будет.

Штрой дёрнулась, словно одно из имён доставило ей боль.

--- Они не контактируют между собой?

--- Нет, Штрой.
Я понимаю ход ваших мыслей, но нет.
Корхес в принципе презирает военных и старается держаться от них подальше.
Ду-Си и Стигма не проявляли друг к другу интереса почти с самого основания Ордена.
Думаю, причина чисто клановая --- они из разных генераций.

--- Я бы сказал, что это серьёзный послужной список, --- заметил Самаолу.
--- Чёрная Скала --- клан воинов.
Со стратегом, уничтожившим его в одиночку, надо быть осторожнее.
Стигма редко ошибается.

--- Это ненадолго, Самаолу, --- успокоил голос.
--- Кстати, как вам их название?

--- Скорбящие, --- выплюнул Самаолу.
--- Знакомый культуролог сказал, что само название пахнет каким-то рыцарским кодексом.
Мы во время зачисления Анкарьяль и Грейсвольда в отдел 100 прощупали их в том числе на предмет идеологии.
Данных собрано достаточно, они обрабатываются.

--- Это здорово, если только старина Грейс не подсунул нам легенду, --- согласился голос.

--- Я рискну высказать наше с Кагуя впечатление, --- проговорил Самаолу.
--- Анкарьяль и Грейсвольд были абсолютно искренни.
Они совершенно не боялись логических атак и проб.
Более того, они хотели, чтобы их идеологию проверили на прочность.
Кагуя услышала в их речах вызов --- <<против простой истины не устоит даже ваша хитроумная ложь>>.

--- Хм, --- удивился голос.
--- Как только получите данные культурологов, Самаолу, доставьте их мне лично.
А что насчёт той картинки?

--- Печальный Митр, --- доложила Штрой.
--- Символ из анимистических верований сели.
Именно благодаря ему я вышла на Митриса Безымянного.
Они настолько увлеклись символизмом, что забыли о конспирации.

--- Если у них и впрямь имеется подобие рыцарского кодекса, это немудрено... --- заговорил Самаолу.

--- Тихо, --- спокойно выдал голос.
Оба демона разом онемели и подобрались.
--- Легкомысленного отношения я не приемлю.
В этом уравнении чересчур много неизвестных, чтобы считать Скорбящих сборищем глупцов.
Я ясно выразился?

--- Так точно, --- хором сказали Штрой и Самаолу.
Самаолу выглядел совершенно спокойным, а Штрой поёжилась.

--- Есть ли пути, которые вели бы к Лусафейру?

Штрой и Самаолу переглянулись.

--- Понятно, --- недовольно сказал голос.
--- А жаль.

Штрой снова поёжилась.
Она чересчур болезненно воспринимала любые проявления недовольства со стороны патрона, и Самаолу бросил на неё короткий презрительный взгляд.

--- Я знаю старину Лу очень давно, --- криво улыбнулся демон.
--- Он друг Грейсвольда, но вряд ли заодно со Скорбящими и знает больше нас.

--- Вы переоцениваете свой опыт, Самаолу, --- сказал голос.
\ml{$0-[ej]$}
{--- Лусафейру очень осторожен.}
{``Lusafejru is very cautious.}
\ml{$0-[ej]$}
{Вы не представляете, насколько --- это та осторожность, которую не смогут ослабить ни тысячелетний мир, ни клинок у горла, ни обвинение в трусости.}
{You can't imagine how, it's a kind of cautiousness which couldn't be made less by anything---thousand-year peace, or blade at the throat, or accusations of cowardice.}
\ml{$0-[ej]$}
{Лусафейру, я не побоюсь сказать, едва ли не единственный в Аду, кто знает цену ошибке.}
{Lusafejru, I daresay, is about the only one in the Hell who knows the price of mistake.}
\ml{$0-[ej]$}
{И при этом он продолжает общаться с Грейсвольдом, который, как известно, под подозрением.}
{But he stays in touch with Grejsvolt who is, as you both know, suspected.}
\ml{$0-[ej]$}
{Возможно, что у Лусафейру есть план.}
{I guess Lusafejru's up to something.}
\ml{$0-[ej]$}
{Возможно, что мы недооценили силу этой связи.}
{I guess we underestimated their friendship.''}

Демоны снова переглянулись.

--- Вам не кажется, что вы чересчур превозносите наших врагов? --- скривился Самаолу.

--- Не врагов, а соперников, Самаолу, --- терпеливо заметил голос.
--- Грамотный стратег не употребляет таких терминов, как <<друг>> и <<враг>>.
Оставьте их для тех, кого кормите пропагандой.

--- Эти <<соперники>> убьют вас при первой же возможности.

Голос тихо засмеялся.

--- Я не буду переубеждать вас, Самаолу, --- ваша святая уверенность придаёт пропаганде силу, так же как дружба, возможно, придаёт силу Скорбящим.
Чувства вообще замечательная вещь --- именно потому, что они императивны по сути, им сложно противиться.
%{Emotions are wonderful, actually, they are imperative, they are hard to disobey.}
Многие демоны называют акбас проклятием... но где бы мы были без него?

--- Вы считаете, что Лусафейру доверяет Грейсвольду так же, как Грейсвольд доверяет Анкарьяль?

--- Не могу сказать точно.
Однако Лусафейру и Грейсвольд очень долго общались с Тахиро, и это не могло не сказаться на их профессиональных навыках.
Вам для справки, Штрой: Тахиро Молниеносный, Стратег-неожиданность --- так его звали.
Он обращался с вероятностями так, как мы обращаемся с точными числами.
Сам он канул в вечность, как и почти все близкие ему демоны.
Почти все.

--- Мне следует прервать общение Лусафейру и Грейсвольда? --- осведомился Самаолу.

--- Вы вряд ли сможете это сделать, --- заметил голос, --- но усложнить это общение несколькими обычными, легко оправдываемыми методами не мешает.
Если Грейсвольд и Анкарьяль останутся в отделе 100 в изоляции, они обязательно совершат фатальную ошибку.
Фатальную для них... или для всех Скорбящих, если случайность будет к нам благосклонна.
Да, кстати, насчёт Безымянного.
\ml{$0-[ej]$}
{Речь Грейсвольда произвела некоторое волнение, и к демиургу сейчас приковано незаслуженно пристальное внимание.}
{Grejsvolt's speech produced a lot of sound and fury, and now the demiurge is a focus of unduly close attention.}
\ml{$0-[ej]$}
{Тем не менее это низкоинтеллектуальное скопление масс-энергии успело наступить на ноги половине Ада.}
{Nevertheless, that lowly intelligent mass-energy storage had enough time for stepping on toes of half the Hell.''}

\ml{$0-[ej]$}
{--- Безымянного оставьте мне, --- поклонилась Штрой.}
{``Leave Nameless to me,'' Stroji nodded.}
\ml{$0-[ej]$}
{--- Как только шум уляжется, музыкант и дикарка дорого заплатят за то, что ввязались в чужую игру.}
{``Once things settle down, the tunesmith and she-savage shall pay deerly for interference with someone else's game.''}

\section{[U] Селекция в Картеле}

Часть обратного пути Штрой и Самаолу шли вместе.

--- Всегда хотел спросить вас, Штрой, --- вдруг заговорил Самаолу.
--- Правда ли, что в Картеле существуют так называемые <<связки>> --- множество молодых демонов соревнуются в выполнении какой-либо функции, чтобы были отобраны самые лучшие?

--- Да, --- неохотно сказала Штрой.

--- А что происходит с теми, кто непригоден, Штрой?

--- Непригодные ликвидируются, Самаолу.

--- Хм.
Значит, Картель вместо инженерии пошёл по старому доброму пути селекции, указанному ещё матерью-природой.

--- Инженерия без селекции --- ничто, Самаолу, --- ощерилась Штрой.
--- Даже созданное инженерами следует подвергнуть испытаниям.

\ml{$0$}
{--- Разумеется, вы абсолютно правы.}
{``Of course, you're absolutely right.}
\ml{$0$}
{Необходимо приручить игру случайности, отбросить тех, кто случайно оказался непригодным, и выделить тех, кто хорош по той же самой случайности.}
{We should tame this game of randomness, we should throw out ones who are randomly useless, and take ones who are equally randomly useful.}
\ml{$0$}
{Возможно, придумать концепцию удачи, вывести формулу успеха --- ведь нужно же что-то рассказывать рабам, которым успех не светит?}
{Maybe we should come up with the conception of luck, or formula for success; we have to tell something to slaves who never have success as an option, don't we?}
\ml{$0$}
{Рабы не любят слово <<случайность>>, оно их пугает.}
{Slaves hate the term `randomness', it scares them.''}

--- Слова бывалого пропагандиста, --- Штрой произнесла последнее слово тоном изысканного оскорбления.

--- Заключение опытного стратега, --- тем же тоном ответил Самаолу.
--- Мне до вас далеко.
Впрочем, моих стратегических навыков тем не менее оказалось достаточно, чтобы не показывать патрону всё, на что я способен.
Зато теперь, когда вы прошли все испытания и достигли вершины эволюционной лестницы, у вас есть десятки агентов --- тех, кому повезло меньше, кому необязательно знать о равновесии Нэша и кто будет платить за все ваши ошибки.

--- Удел слабых --- гибель, --- голос Штрой звучал почти победно.
--- Разве нет?

--- Разумеется, вы абсолютно правы.
Занятно.
Хм...
Я задам вам последний вопрос, Штрой, и перестану беспокоить.
Удовлетворите моё любопытство --- скольких жизней стоило \emph{ваше} место?

--- Я оказалась лучше восьмиста тысяч ювеналов.
Это имеет значение?

Самаолу брезгливо хмыкнул.

--- Это заметно, Штрой.
Вы лучше многих.
Мало кому пришлось пережить столь жёсткий отбор.
Но кое в чём вы уступаете самому последнему легионеру Ада.

--- И в чём же? --- осведомилась Штрой.

--- Над вами всегда, до последнего мгновения жизни будет висеть призрачная рука.
Рука того бездушного механизма, который производил отбор.

И, не попрощавшись, Самаолу пошёл к ближайшему лифту.

\chapter{[:] Рассвет Жёлтого моря}

\section{[:] Первая смерть}

\epigraph
{Наибольшую эффективность организация имеет в сферах, которые были ключевыми в процессе её становления.
Группировка, захватившая власть военным путём, не способна эффективно организовать мирную жизнь, так как отбор кадров изначально шёл в другом направлении.}
{Лусафейру Лёгкая Ладонь}

Лу и Тахиро сидели с бокалами вина.
Лу держал бокал, оттопырив мизинец; Тахиро же по-деревенски заграбастал свой за венчик, беззастенчиво полоская пальцы в вине.
Третий бокал стоял перед Грисвольдом;
в нём не хватало едва ли глотка.

--- Ну и как ты? --- улыбнулся Лу.
--- Привыкаешь к демону?

--- Понемногу, --- кивнул Тахиро.
--- Но иногда всё равно странное чувство, что во мне кто-то сидит.

--- Потому он и называется <<демоном>>, я полагаю.

--- Немного раздражают эти странные видения --- вспышки, волны, летающие огоньки.

--- Это проекции омега-поля на зрительную кору, --- многозначительно сказал Лу, подняв палец.

--- Они не должны раздражать, --- встрепенулся Грис.
--- Вечером будем калибровать.

--- Иногда словно кто-то думает за меня, --- продолжал Тахиро.
--- Я начинаю мысль и вдруг щёлк --- появляется результат.
А хода мысли я не вижу.

--- Это нормально, --- пояснил Лу.
--- Демон автоматически перехватывает управление потоком.
Это сделано для удобства, чтобы не переключаться на омега-консоль вручную.

--- Но всё равно надо поработать над обратной связью, --- заметил Грисвольд.
--- Тахиро \emph{должен} понимать, что это демон перехватывает поток, а не играет его собственная фантазия.
Ну и какую-то техническую информацию...

--- Не-не, Грис, вот технической информации мне как раз хватает, --- перебил его Тахиро.
--- Цифры, названия давно разрушенных городов, имена давно умерших людей.
Что вообще за дурацкая традиция в математике давать методам имена по месту открытия или по имени учёного?
Никакой поэзии.

--- Да забей, --- поморщился Лу.
--- Я вообще не задумываюсь.
Всё, что мне нужно знать о конкретном методе --- к чему его можно применять, а к чему нельзя.
Остальным пусть занимаются в исследовательских отделах.
Тебе же пока вообще надо тренироваться на кошках --- складывай числа, вычисляй траектории, можешь предсказать цену креветок в Такэсако на месяц --- экономические утилиты там какие-то есть.

\ml{$0$}
{--- А ты на чём тренировался?}
{``What did you practice on?''}

Лу рассмеялся.

\ml{$0$}
{--- На считалках.}
{``Rhymes.''}

\ml{$0$}
{--- Каких считалках?}
{``What rhymes?''}

\ml{$0$}
{--- Детских.}
{``Counting-out rhymes of children.}
\ml{$0$}
{Мы с отцом часто катались по ущелью и за его пределами, и я наблюдал за играми детей.}
{My father and me used to travel a lot, around the canyon and beyond, so I watched children play.}
Соответственно, запомнил массу дразнилок, считалок и прибауток.
Тематическое моделирование считалок показало, что они распространялись примерно так же, как люди заселяли эту территорию.
Считалки с неприличными словами появились в западной части Такэсако --- видимо, из-за недостатка воды.
А самые мерзкие и злые дразнилки водятся в Асахина, и когда я читаю отчёты о тамошней ситуации, я их не виню.
Дети --- это зеркало мира взрослых.

Люцифер отхлебнул вина.

--- Я тоже хочу исследовать считалки! --- восхищённо сказал Тахиро.

\ml{$0$}
{--- Займись поэзией, --- ухмыльнулся Лу.}
{``Take poetry,'' Lu smiled.}
\ml{$0$}
{--- Ты из хорошей семьи и знаешь в поэзии толк.}
{``You're from a good family, you have an eye for poetry.}
\ml{$0$}
{Только я бы на твоём месте не ждал каких-то серьёзных успехов.}
{I wouldn't expect great discoveries, though, if I were you.''}

\ml{$0$}
{--- Почему это?}
{``Why not?''}

\ml{$0$}
{--- Когда я изучал культуру Такэсако, я был одним из немногих в Ордене, кто вообще обращал внимание на такие мелочи.}
{``When I researched Takesako culture, I was one of the few in the Order who paid attention to such a trivialities.}
\ml{$0$}
{Одинокий дилетант.}
{The lone amateur.}
А сейчас всем этим занимаются учёные бывшего Ордена Тысячи Башен, которые собаку съели на подобных исследованиях.
В изучении культуры людей они по-прежнему на десять шагов впереди.

--- Зато мы опережали их в военном отношении, --- пожал плечами Тахиро.

\ml{$0$}
{--- В этом и разница между нами.}
{``That's the difference between us and them.}
\ml{$0$}
{Они изучали людей и умели с ними как-то уживаться.}
{They studied humans, they have learned to co-exist with humans.}
\ml{$0$}
{Мы вместо этого развивали военную мощь.}
{We have amassed military power instead of that.''}

\ml{$0$}
{--- И они проиграли.}
{``They lost.}
\ml{$0$}
{Мы проглотили Тысячу Башен, как рыба глотает жука.}
{We swallowed Thousand Towers like a fish swallows a bug.''}

\ml{$0$}
{--- Да что ты? --- хихикнул Люцифер.}
{``Really?'' Lucifer laughed.}

\ml{$0$}
{--- Скорее как наживку, --- улыбнулся Грисвольд.}
{``Like a bait, I'd say,'' Griswold smiled.}

\ml{$0$}
{--- Вот и у меня создалось такое впечатление.}
{``I got the same impression.}
\ml{$0$}
{Все ключевые посты, за исключением армейских, занимают кто?}
{All key positions, except military, occupied who?''}

\ml{$0$}
{--- И, что самое главное, заслуженно занимают, --- присовокупил Грисвольд.}
{``More important, justly occupied,'' Griswold added.}
--- Сейчас мне хоть поговорить есть с кем, кроме вас.
Идёшь по коридору --- а на скамейке сидят трое-четверо центурионов и рассуждают о чем-то безумно интересном.
Так и хочется подсесть и просто поговорить...

--- А на чём тренировался Гало, когда был маленьким? --- поинтересовался Тахиро у Люцифера.

--- Гало был серьёзный парень и начал сразу с наших военных затрат.
Цены на материалы, оружие, оптимальный уровень взятки, гарантирующий лояльность осведомителя.
Он очень сильно помог отцу буквально в первый же год обучения, почему отец и выделил его в любимчики.
Я же ещё пару лет валял дурака, исследуя геологию, метеорологические данные, культуру жителей, ну и всякую живность.
Всего понемногу, одним словом.

--- Мы, конечно, напортачили, сразу прикрутив Тахиро такие сложные стратегические модули, --- задумчиво сказал Грисвольд.
--- Оцифрованные легионеры уже в строю, а вот тебе ещё долго-долго тренироваться и ждать, пока интеграция будет окончательно завершена.
Но сам факт того, что ты не сошёл с ума после процедуры, многообещающий.

--- Ну и то, что Тахиро не пустили под нож, тоже внушает надежду, --- захихикал Лу.
--- Правда, отец сказал, что ты не получишь никаких обновлений по своей части, пока не достигнешь хотя бы ранга центуриона.
К счастью, исполнять приказ не ему, а демонам, которые занимаются моей веткой обновлений, --- Лу подмигнул Грисвольду.

--- Это было самое ужасное, что я чувствовал в жизни, --- Тахиро внезапно передёрнуло так, что он чуть не расплескал вино.
--- Если бы рядом были не вы, я бы точно сошёл с ума.
Интересно, как вы это пережили.

--- У нас с Лу не было такого опыта, --- пожал плечами Грис.
--- Мы изначально были демонами и просто встраивались в человеческие тела.
Это совсем другое.

--- А что я буду чувствовать, когда умру? --- спросил Тахиро.
--- Имеется в виду, когда тело придёт в негодность.

Люцифер и Грисвольд смущённо переглянулись.

--- Мы не знаем, --- признался технолог.
--- Чисто технически омега-модуль дублирует твою нервную систему и интегрирован в неё, так что после смерти ты просто не будешь чувствовать тело.
Такое часто бывает во сне, знаешь --- ни рук, ни ног, ни дыхания.
Омега-чувство, напоминающее зрительные ощущения, разумеется, останется.
Поэтому понять, что означают <<видения>>, в твоих интересах.

--- Но это теория, --- закончил Лу.
--- Прецедентов, как понимаешь, ещё не было.
Ты будешь первым, кто расскажет, как всё обстоит на самом деле.

--- Я больше не хочу быть первым!

Лу и Грис расхохотались.

--- Ну, могу тебя понять.
Да, у нас же теперь есть куча оцифрованных легионеров!
Надо кого-нибудь из них кокнуть ради интереса.

--- Люцифер! --- укоризненно воскликнул Грисвольд.

--- Да, ты прав.
Подождём, пока кто-нибудь помрёт сам.
Это мудро и экономично.

--- Ну что ж, у меня будет ещё один важный этап взросления, --- подытожил Тахиро.
--- Выпьем за это.

--- Да, за твою первую смерть.

Лу и Тахиро по очереди звякнули бокалами о стоящий бокал технолога.

--- Штука это не особо приятная, --- поделился Грисвольд.
--- Особенно когда тебе в печень втыкают копьё.
Так что за то, чтобы она прошла по плану в техническом смысле.

--- Тебе в печень воткнули копьё? --- с интересом спросил Тахиро.

--- После того, как я полдня провисел на солнцепёке голышом, съедаемый мухами и слепнями, --- кивнул технолог.

--- Какая крутая смерть, --- восхитился Лу.

--- А что ты сделал? --- спросил Тахиро.

--- У меня с ними возникли разногласия на почве религии.
Они не поверили, что я их бог.

--- Где-то я это уже слышал, --- пробормотал Тахиро.
--- Кажется, на Тысяче Башен есть племя, которое верит в повешенного на кресте бога.
Мне рассказывали, что там есть очень красивая друза Хербст.
Она вся поросла деревьями, листья которых делаются оранжевыми, как лисы, каждый сезон.
Хотел бы я там побывать.

--- Ты еще на моей планете не побывал, --- напомнил Грис.

--- Кстати, о твоей планете, -- оживился Люцифер.
--- Ты слышал, да?
Вообще это военная тайна, но скауты уже всем уши прожужжали.
Мы нашли Лотос.
Не по твоим координатам, они оказались, мягко говоря, неточными, а по описанию.
Луна-кристалл, три материка, лиловая пустыня, горы из платины тоже видели.
Тысяча Башен нашла.

--- А Союз?

--- Они пока не в курсе, судя по всему.
Кстати, ты всё говорил про звезду Фомальгаут.
По этому названию один из скаутов аннотировал старые звездные карты, прибывшие на Тысячу Башен с первыми поселенцами.
Они обнаружили ещё четыре обитаемые планеты.
Две, к сожалению, уже заняты Союзом.
А одна из оставшихся...

Люцифер вдруг посерьезнел.
Падающий на его лицо тусклый свет состарил кукольные черты.
Грисвольд понял.

--- Древняя Земля?

--- То, что от неё осталось, --- глухо ответил Лу.
--- Людей на ней нет, очертания материков... кхм...
В общем, мало похоже на то, что мы ожидали, но это точно она.
Ледниковый период прошёл, и она всё ещё пригодна для жизни.
Кстати, о жизни...

Лу осушил бокал и, подтянувшись к Тахиро, нежно поцеловал его в губы и запустил пальцы в его волосы.

--- Чем займёмся?
У нас ещё три часа.

--- Уговорил, --- Тахиро начал развязывать Лу пояс.

--- Да тебя и уговаривать не надо.
Хочешь, позовём Айну?

--- Не сегодня.

--- Я слышу это в двадцатый раз.
Мне кажется, ты её избегаешь.
Она, между прочим, о тебе постоянно говорит.
Недавно вспомнила, как ты пообещал взять её в жёны.

--- Я был молод и глуп.

--- Это из-за той такэсакской красавицы, которую она убила?
Она это из ревности, точно тебе говорю.

--- Лу, давай к делу.
То, что мы с тобой спим и дружим, ещё не даёт тебе право лезть ко мне в душу.

--- Когда вы успели? --- удивился Грисвольд.
--- Лу, ты зачем парня совратил?

--- Да это я его, --- хихикнул Тахиро.
--- Мы же с Лу тогда спали в одной комнате.
Просыпаюсь я как-то утром и вижу его в очень живописной позе.
Дальше всё как в тумане.

--- То есть это моя вина, --- ухмыльнулся технолог.

--- Я бы перефразировал.
Это благодаря тебе.

--- Тогда я, пожалуй, пойду, --- Грисвольд отхлебнул ещё глоток из своего бокала и грузно поднялся на ноги.

--- Ну да, ничего особо интеллектуального сейчас не будет, --- признал Лу.
--- Но ты можешь присоединиться.

--- Нет, благодарю, --- ухмыльнулся технолог.
--- У меня паёк --- одна женщина в два года.

--- Презренный моносексуал, --- поморщился Тахиро.

--- Тахиро, вот скажи, зачем мы вообще дружим с этими низшими созданиями? --- развёл руками Лу.

--- Ещё слово --- и вы оба останетесь без обновлений, --- посулился Грисвольд.

--- А вот это был удар ниже пояса! --- возмутился Лу.

Грисвольд скорчил рожу и вышел, прикрыв за собой дверь.

\section{[:] Чужестранка}

\ml{$0$}
{--- Я не буду с тобой спать.}
{``I shan't sleep with you.''}

\ml{$0$}
{--- Простите? --- опешил Люцифер.}
{``Beg your pardon?'' Lucifer said, taken aback.}

Демоница бросила взгляд на стопку папок в углу стола.
По крайней мере, так показалось Люциферу --- её сильное косоглазие не позволяло утверждать наверняка.
Люцифер схватил верхнюю папку и открыл.

\ml{$0$}
{--- Это досье на меня.}
{``This is a dossier on me.''}

Демоница выхватила папку из его рук и положила на место.

\ml{$0$}
{--- Это. Досье. На меня, --- повторил Лу.}
{``This is. A dossier. On me,'' Lu repeated.}

\ml{$0$}
{--- И?}
{``So what?''}

Люцифер задумался.
Факт существования досье на него самого ошеломил стратега настолько, что он вдруг забыл, зачем его вообще занесло в этот забытый небом кабинет в дальнем конце южного крыла.

Секс с Тахиро был средством на время забыть о собственных тревогах, и друг это прекрасно понимал;
он получал то, что было нужно ему, и знал, что Лу делает то же самое.
Обычно объятий друга хватало на несколько часов умиротворения, но в этот раз тревоги вернулись, едва Тахиро заснул.
Люцифер бережно укрыл его одеялом, потушил лампу и тихо вышел за дверь, намереваясь поработать.

\ml{$0$}
{Женщину, которая сидела перед ним в старом кресле, звали Анастейша Розье.}
{The woman, sitting in an old chair in front of him, was Anasteishae Rozier.}
Странное имя, тяжёлое и громоздкое, словно аппаратура для демонизации, дань дурацкой, позаимствованной у дзайку-мару традиции давать несколько имён.
\ml{$0$}
{Сама она предпочитала уменьшительно-ласкательный вариант --- Стигма.}
{She herself preferred a diminutive form: Stigma.}
Позывной встречался в очень интересных местах архивных документов Тысячи Башен, и следовало...

Люцифер бросил взгляд на стопку досье, и его вдруг осенило.
Эти досье выглядели в каморке кабинетного учёного так же неуместно, как...
Мозг Лу даже не стал тратить время на поиск подходящего сравнения.

--- Это ты была главой разведки Ордена Тысячи Башен, --- заявил он.
--- Настоящим, а не церемониальным руководителем.

Стигма вздрогнула.

--- Не бойся, это уже в прошлом, --- поспешил добавить Лу.
--- Вообще да, я тоже собирал на тебя досье и меня как-то зацепило, я захотел с тобой познакомиться поближе.
И пришёл именно поэтому.
И ты поняла это потому, что собрала досье на меня.

--- Ты всех, чьё досье тебя зацепило, пытаешься затащить в постель?

--- Именно сейчас я не настроен на постель, --- признался Лу.
--- Если ты про мои методы...
Да, мне просто хочется узнать демона поближе, а секс --- самый простой способ.
Лично для меня, разумеется.

Стигма промолчала и отвернулась.
Люцифер вспрыгнул за стол и уселся перед женщиной, скрестив ноги.
Стигма зябко поёжилась и отвернулась ещё на пол-оборота;
на её лице застыла смесь страха и отвращения.
Лу, поморщившись, аккуратно слез со стола и поправил лежащие на нём вещи.

--- Мы тут нашли вашу служебную переписку, --- сказал он.
\ml{$0$}
{--- В ней ты единственная утверждала, что потери Преисподней --- не что иное, как обманный ход, и что у Ордена достаточно сил, чтобы одолеть Союз на Преисподней.}
{``You were the one who told: military losses of Nether are nothing but ruse, and the Order is strong enough to defeat the Union on Nether ground.}
\ml{$0$}
{Ты настаивала на том, чтобы Тысяча Башен выступила на стороне Преисподней как можно скорее, дабы, цитата, <<избежать последствий нейтралитета>>.}
{You insisted that Thousand Towers must join Nether as soon as possible, I quote, `to avoid the consequences of neutrality'.}
\ml{$0$}
{Впечатляет.}
{Very impressive.}
\ml{$0$}
{Почему тебе не поверили?}
{Why they didn't believe you?''}

\ml{$0$}
{--- Я обязана отвечать на твои вопросы?}
{``Do I have to answer your questions?''}

\ml{$0$}
{--- Разумеется, нет.}
{``Of course you don't.}
\ml{$0$}
{Но если я задаю вопрос, значит, ответ для меня важен и я в любом случае попытаюсь его найти.}
{Nevertheless, if I ask, an answer is important to me and I shall try to find it anyway.}
\ml{$0$}
{Так всё-таки --- почему твоё слово оказалось недостаточно веским?}
{Returning to the subject, why your word wasn't good enough?''}

\ml{$0$}
{--- Я чужая, --- осторожно сказала Стигма.}
{``I'm an outlander,'' Stigma cautiously answered.}
\ml{$0$}
{--- Тысяча Башен --- планета традиций.}
{``Thousand Towers is a planet for tradition.''}

--- Тысяча Башен куда менее традиционна, чем Преисподняя, --- сказал Лу.
--- Вы абсолютно неразборчивы, когда дело касается новобранцев.
Многие из этих новобранцев потом занимают высокие посты, и их слово становится законом.

--- Я чужая даже для них.

--- Тысяча Башен принимает в свои ряды дезертиров, визоров-рабов, богов, преступников разного пошиба и перебежчиков, успевших поменять стороны по десять раз.
\ml{$0$}
{С трудом представляю, что может быть хуже.}
{I barely can imagine something worse than that.}
Разве что трудности в общении.
\ml{$0$}
{Может быть, ты бездомный бог?}
{Maybe you're a homeless god?''}

Стигма снова промолчала.

\ml{$0$}
{--- Думаю, нет, --- ответил Лу на собственный вопрос.}
{``Apparently not,'' Lu answered his own question.}
\ml{$0$}
{--- Твоё досье мы не нашли, но это очевидно.}
{``Your dossier wasn't found, but it's obvious.}
Тем не менее твои паттерны мне не знакомы совершенно, они не подходят ни под один шаблон известных объединений демонов.
\ml{$0$}
{Ты действительно другая.}
{You're definitely different.''}

Молчание.

--- Хотя... --- Лу прищурился, --- может быть, Чёрная Скала?
Не те неприятные головорезы, говорящие на булькающем языке, а нечто более тонкое и изящное.
Даже головорезам нужен тот, кто ими управляет.
И не всегда он это делает по своей воле, ведь ситуация, когда лидер является заложником подчинённых --- не фантазия, а суровая реальность.
Это объясняет, почему тебе не поверили --- у тебя есть личная заинтересованность в том, чтобы Тысяча Башен и Союз Воронёной Стали не подружились.
Вдобавок это хорошее объяснение и для того, что ты разобралась с Союзом гораздо быстрее нас.
Разумеется, ты знала всё о бывших соратниках, клане Антрацис.
\ml{$0$}
{То есть, иными словами, ты --- Вечно Го...}
{In other words, you are a Forever Ha---''}

\ml{$0$}
{--- Заткни пасть! --- взорвалась Стигма.}
{``Shut your mouth!'' Stigma exploded.}
\ml{$0$}
{--- Не произноси при мне эти слова!}
{``Never say that in my presence!''}

--- Не буду, --- тут же поправился Лу.
\ml{$0$}
{--- Прости.}
{``Sorry.''}

\ml{$0$}
{--- <<Прости>>?}
{`` `Sorry'?}
\ml{$0$}
{Чего ты добиваешься?}
{What's your angle?}
\ml{$0$}
{Я не нахожусь под арестом, меня ни в чём не подозревают.}
{I'm not under arrest, I'm not suspected or somethin'.}
\ml{$0$}
{Что тебе нужно, кокэси-нингё\FM?!}
{What do you want, \textit{kokeshi ningyo}\FM!''}
\FA{Кокэси-нингё --- традиционные пластиковые куклы на Преисподней.
В переносном смысле: уничижительное название близнецов или внешне похожих людей (например, одетых в одинаковую форму солдат).
}

\ml{$0$}
{--- Кокэси-нингё? --- Лу захохотал.}
{``\textit{Kokeshi ningyo}?'' Lu laughed.}
\ml{$0$}
{--- Меня так в шутку называли деревенские в детстве, из-за того, что у меня был брат-близнец.}
{``Villagers jokingly called me that in my childhood, because of my twin brother.}
\ml{$0$}
{Откуда ты это знаешь?}
{How do you know?''}

Стигма промолчала, глядя в пол.
Люцифер сел на пол перед её стулом, пытаясь заглянуть ей в глаза.

--- Я всё больше убеждаюсь, что зашёл сюда не зря.
\ml{$0$}
{Ты другая, и ты мне очень интересна.}
{You're different, and I'm interested in you.}
\ml{$0$}
{Интересна в хорошем смысле.}
{In a good way.}
\ml{$0$}
{И я тебе интересен, верно?}
{And you're interested in me, right?}
\ml{$0$}
{Да, до последнего времени я интересовал тебя исключительно как вероятная цель для убийства, но времена меняются.}
{Well, until recently, I've interested you strictly as an assassination target, but things have changed.''}

--- Ты не привык, что тебе отказывают? --- буркнула Стигма.
--- Не привык, что тебя выставляют за дверь, что с тобой не хотят разговаривать?

--- Для меня это в новинку, --- признался Люцифер.
--- В Ордене надо мной смеялись за моей спиной, но отклонить моё предложение побеседовать обычно боялись, потому что демоны моего ранга редко приходят просто так.
Я не хотел тебя обидеть или оскорбить.
Как бы глупо это ни звучало, я хотел с тобой прогуляться.

--- Я не буду с тобой спать, Люцифер.

\ml{$0$}
{--- Я это уже уяснил.}
{``You've made that clear enough.}
\ml{$0$}
{Не вижу никакого смысла заниматься сексом с тем, кто этого не хочет --- это оскорбляет саму идею.}
{I don't see the point of sex with someone who doesn't want it---it offends the very idea.}
\ml{$0$}
{На самом деле ты можешь хранить свою девственность столько, сколько пожелаешь, Самаэл выпишет официальное разрешение.}
{Actually, you can keep your virginity as long as you want, Samael will issue official permission.}
Но всё-таки --- как насчёт небольшой прогулки?

\ml{$0$}
{--- Это трата времени.}
{``It's a waste of time.}
\ml{$0$}
{И оставь при себе свои корявые шутки.}
{And keep your ugly jokes to yourself.''}

\ml{$0$}
{--- Я оскорблён в лучших чувствах.}
{``I feel deeply offended.''}

\ml{$0$}
{--- И это тоже трата времени.}
{``And that's a waste of time too.''}

Люцифер прищурился.

\ml{$0$}
{--- Откуда в тебе такой страх перед сексом?}
{``Where is your fear of sex from?''}

Стигма вздрогнула так, что её кресло ударилось о стол.

\ml{$0$}
{--- Меня никогда не интересовал секс, --- заявила она.}
{``I was never interested in sex,'' she said.}

\ml{$0$}
{--- Возможно, но это не объясняет страх.}
{``It can be true, but it can't explain the fear.''}

Стигма промолчала и снова зябко поёжилась.

--- Хотя в принципе это тоже понятно, --- задумался Люцифер.
--- Твои булькающие дружки не особо церемонятся с дзайку-мару.
\ml{$0$}
{Сомневаюсь, что они насиловали тебя, но то, что это регулярно делали \emph{при тебе}, не вызывает сомнений.}
{I doubt they raped you, but I have no doubt they did it regularly \emph{in front of you}.}
\ml{$0$}
{Сложно увидеть что-то хорошее в том, что всегда было для тебя наказанием.}
{It's difficult to see something good in a thing you used to treat as a punishment.}
Даже если ты и совладала с этим в прошлой жизни, твои демонические воспоминания, как ведро снега, вываливаются на новое тело, заставляя его вздрагивать в ужасе.
У нас есть лаборатория, которая занимается относительно новым направлением --- коррекцией личности демона.
Если тебя беспокоят какие-то старые травмы, ты всегда можешь обратиться к их руководителю, Цельсии.
Я позабочусь о том, чтобы в твоём случае была соблюдена полная конфиденциальность.

--- У тебя всё? --- сухо осведомилась Стигма.

--- Ещё нет.
У меня есть камень в рукаве.

--- Сочувствую.
Иди и замени одежду.

Люцифер вместо ответа вытащил из кармана мышь и положил зверька Стигме на стол, придерживая за хвостик.

--- И что это?

--- Это мышь, --- просветил женщину Люцифер.
--- Тебе не нравится?

--- Нравится, но зачем ты притащил её ко мне?

Лицо Люцифера озарила довольная улыбка.

--- Тебе нравится.

--- И что?

--- Ты сказала, что тебе нравится.

--- Да, она мне нравится.
Дальше что?

--- Найди ей домик.
А если дать ей самца, пищу и время, у тебя будет целая куча мышей.
Ты получишь множество мутантов --- всех, какие только могут быть.
\ml{$0$}
{Прямо как первобытные биологи.}
{Like prehistoric biologists used to do.}
\ml{$0$}
{Самца я принесу.}
{I'll bring you a male.}
\ml{$0$}
{Держи.}
{Take this.''}

Люцифер всучил мышь Стигме и направился к выходу.

\ml{$0$}
{--- Эй, ты куда?}
{``Jey, where're you goin'?''}

\ml{$0$}
{--- За самцом.}
{``Look for a male.''}

\ml{$0$}
{--- Ты с ума сошёл?!}
{``Are you crazy!''}

\ml{$0$}
{--- С ума сходить будешь ты, когда получишь первого альбиноса.}
{``It's you who will go crazy when you've got your first albino.''}

\ml{$0$}
{--- Пожалуйста, не надо.}
{``Please, stop now.}
\ml{$0$}
{Не нужен мне самец.}
{I don't need a male.''}

\ml{$0$}
{--- Ты меня умоляешь?}
{``Are you begging me?}
\ml{$0$}
{А чего так?}
{Why?}
Боишься, что тебе понравится ещё больше?

\ml{$0$}
{--- Я тебя не умоляю, сумасшедший!}
{``I'm not beggin' you, you freak!}
Ты врываешься ко мне в кабинет, и...

\ml{$0$}
{--- Если тебе не нужна мышь, просто выпусти её.}
{``If you don't need a mouse, just release it.}
\ml{$0$}
{Она дикая и прекрасно обойдётся без тебя.}
{It's wild and perfectly capable to survive without you.}
\ml{$0$}
{Или раздави ей голову и выброси в мусорку.}
{Or you can crush its head and throw it out.}
\ml{$0$}
{Мой отец сделал именно так, когда я принёс ему мышку.}
{Like my father did when I bring a mouse to him.}
\ml{$0$}
{А твоё <<пожалуйста, не надо>> звучит неубедительно.}
{As for your `please, stop now'---it sounds very unconvincing.''}

Люцифер выразительно сблизил большой и указательный пальцы, затем вышел, хлопнув дверью.
Стигма громко выругалась ему вслед, вздохнула и, посмотрев на мышь, принялась искать в углу подходящую коробку.

\section{[:] Рога и плащ}

Оцифровка людей встала на поток и уже принесла первые плоды.
Отряд, собранный полностью из оцифрованных, одержал победу над превосходящими силами Союза.
Впервые за несколько лет Союз предложил Ордену переговоры, и впервые за десятилетия войны была уверенность, что это не обманный манёвр с целью оттянуть время.

Авторитет Лу резко возрос.
Больше никто из легатов, и тем более легионеров, над ним не смеялся.
Его речи слушали внимательно и приказы исполняли беспрекословно.
Сам Лу в честь таких перемен решил сменить имидж.

--- Ну как тебе?

--- Очень красиво, --- признал Тахиро.
\ml{$0$}
{--- Особенно этот ободок с рожками.}
{``Especially that horned bezel on your head.}
\ml{$0$}
{Кстати, почему рога?}
{By the way, why horns?''}

\ml{$0$}
{--- Мне нравятся рога.}
{``I like horns.}
\ml{$0$}
{Они милые.}
{Nice touch.}
\ml{$0$}
{Но меня напрягает твоё молчание по поводу плаща.}
{But I feel bad about your silence concerning my cloak.''}

--- Плащ тоже красивый.

--- <<Тоже>>?
Тахиро, вот сгори ты дотла!
\ml{$0$}
{Скажи, что я в нём смотрюсь очень элегантно, эпично, ну придумай что-нибудь!}
{Tell me I'm elegant, and epic, and anything you can make up!''}

\ml{$0$}
{--- Честно --- он похож на плед.}
{``Honestly, it looks like blanket.''}

\ml{$0$}
{--- Да вы издеваетесь.}
{``You're mocking me.''}

\ml{$0$}
{--- Почему?}
{``Why would I?''}

\ml{$0$}
{--- Грис то же самое сказал.}
{``Gris said the same word.}
Вообще-то это воссозданная технология первых людей, наноткань с очень низкой теплопроводностью.
Плащ делал лучший мастер по моим наброскам.
\ml{$0$}
{Ну серьёзно, я не выгляжу в нём элегантно и эпично?}
{Come on, don't I look elegant and epic?''}

\ml{$0$}
{--- Ты выглядишь очень элегантно и эпично.}
{``You do look elegant and epic.}
\ml{$0$}
{Плащ очень красивый.}
{The cloak is really beautiful.}
\ml{$0$}
{Но это не отменяет того, что он очень похож на плед.}
{But it does not change the fact the cloak looks like blanket.''}

Лу порылся в сундучке, вытащил посеребрённую фибулу и подколол ею плащ.

\ml{$0$}
{--- А так?}
{``Better?''}

\ml{$0$}
{--- Тебе нравится?}
{``Do you like that?''}

\ml{$0$}
{--- Иначе я бы его не надел.}
{``I wouldn't wear it if I don't.''}

\ml{$0$}
{--- Ты выглядишь элегантно и эпично?}
{``Do you look elegant and epic?''}

\ml{$0$}
{--- По моему мнению --- да.}
{``In my opinion, yes.}
\ml{$0$}
{А ещё мне мягко и тепло.}
{Also, I feel soft and warm.''}

\ml{$0$}
{--- Тогда зачем тебе нужно моё мнение?}
{``So why do you ask my opinion?''}

Лу подумал, пробормотал что-то под нос и вышел.

\section{[:] Цена молчания}

--- Привет, сайгон, --- кивнула Айну.
--- Красивый плащ.
Издалека я подумала, что ты совсем решил подтереться дисциплиной и пришёл на переговоры в одеяле.

Лу показал ей неприличный жест и пошёл в другую сторону зала, подальше от командования штаба.
Возле стены робко примостилась Стигма, чуть поодаль от своего отдела тактической технологии.
Кресло рядом с ней пустовало.

\ml{$0$}
{--- Привет, --- улыбнулся ей Люцифер.}
{``Hi,'' Lucifer smiled to her.}
\ml{$0$}
{--- Ты не возражаешь?}
{``Would you mind?''}

\ml{$0$}
{--- А моё возражение как-то повлияет на результат? --- осведомилась Стигма.}
{``Would anythin' change if I mind?'' Stigma asked.}

\ml{$0$}
{--- Если ты не хочешь, чтобы я с тобой сидел --- я поищу другое место.}
{``If you don't want me close, of course I'll look for another seat.''}

Стигма устало взлохматила рукой волосы.

\ml{$0$}
{--- Садись.}
{``Sit down.}
\ml{$0$}
{Только больше никаких разговоров про прогулки и секс, ясно?}
{But no more `walk', and no more `sex', agreed?}
\ml{$0$}
{И вообще я не хочу с тобой разговаривать, Люцифер.}
{In fact, don't even talk to me at all, Lucifer.}
\ml{$0$}
{Я устаю от твоей болтовни.}
{I'm get tired of listenin' to your nonsense.''}

\ml{$0$}
{--- Договорились.}
{``Agreed.}
\ml{$0$}
{Место в обмен на молчание.}
{I have the seat, you have my silence.''}

\ml{$0$}
{--- И трогать меня ты тоже не будешь.}
{``And no more touchin' me.}
\ml{$0$}
{И смотреть на меня.}
{No more lookin' at me.}
\ml{$0$}
{И совершать иные демонстративные действия, чтобы привлечь моё внимание.}
{No more doin' any of your intemperate stuff to get my attention.}
\ml{$0$}
{Меня здесь нет.}
{I'm not here anymore.''}

\ml{$0$}
{--- Как скажешь, Стигма.}
{``As you wish, Stigma.''}

Посланцы Союза опаздывали.
Лу уже начал скучать, но пока добросовестно держался и не пытался заговорить со Стигмой.
Терпение почти подошло к концу, когда двери зала распахнулись и вошли двое, держась за руки --- мужчина и женщина.

Пришельцы выглядели абсолютно обычно.
И тем не менее зал внезапно прижух, словно осенний лист.
Сидящие начали ёжиться и вертеться, словно им всем одновременно стало неудобно сидеть.
Но вся эта возня не могла сравниться с реакцией Стигмы.
Она всхлипнула и тут же зажала рот обеими руками.

--- Старые знакомые? --- ухмыльнулся Лу.

--- Двойняшки, --- выдавила Стигма.
--- Нуива и Фуси с Жёлтого моря, ассассины Союза.

--- Твоё прежнее имя, указанное в одном из отчётов --- Хуанхай Цзаошан, --- припомнил Лу.
--- Цзаошан с Жёлтого моря.
Ты --- их сестра?
То есть, другими словами, вы из одной клановой серии, были собраны на основе одной и той же версии программного ядра?

--- Ты обещал молчать, --- ледяным тоном ответила Стигма, справившись с собой.
Люцифер виновато кивнул и завернулся в плащ.

\textspace

--- Есть ещё одно условие, --- сказал Фуси.
--- Верните нам предателей из числа Чёрной Скалы.

Двойняшки как по команде уставились немигающими взглядами на Стигму, сидевшую рядом с Люцифером.
Ну лице женщины застыл неописуемый ужас.

Легаты Ордена переглянулись.

--- Вы располагаете какими-то данными? --- уточнил Самаэл.

--- Располагаем, --- кивнул Фуси.
--- Речь идёт о единицах, а именно --- о троих демонах.
Насколько нам известно, они не занимают высоких постов.
Это рядовые легионеры и научные сотрудники, их потеря никак не скажется на Ордене Преисподней.
Выдайте их нам, и переговоры можно считать завершёнными.

--- Полагаю, чтобы они рассказали то, что успели увидеть и услышать в Ордене, --- заметил Гало.

--- Нет, --- лаконично ответил Фуси.
--- По возвращении домой они будут немедленно убиты.

--- Свежо предание, --- фыркнул стратег.

--- Ваше недоверие вполне объяснимо, --- оскалилась Нуива.
--- В таком случае могу предложить казнь при участии наблюдателей с обеих сторон.

Гало оглядел легатов и пожал плечами.
Самаэл натянул вежливую полуулыбку, кивнул и уже открыл рот...

--- Я не понимаю, о чём вы говорите, --- опередил его Люцифер.
Все обернулись к нему.

--- Мы говорим о тех Вечно Гонимых, которые были укрыты Орденом Тысячи Башен, --- прохрипела Нуива.

--- В Ордене Преисподней нет Вечно Гонимых, --- заявил Лу и встал, словно ненароком набросив свой плащ на Стигму.
--- В Ордене Преисподней нет демонов Ордена Тысячи Башен.
В Ордене Преисподней нет предателей, в Ордене Преисподней есть только адепты Ордена Преисподней.
Поэтому я не понимаю, о чём идёт речь.

--- Подумайте хорошо, --- Фуси понизил голос до хриплого баса.
--- Это важное условие для дальнейшего поддержания дипломатических отношений.

--- Во-первых, важные условия озвучиваются в начале, а не в конце переговоров, --- Лу театрально помахал пальцами.
--- Во-вторых, важные условия заносятся в тезисы переговоров.
В-третьих, я не вижу смысла обсуждать условия, которые не могу понять.

<<Сын, что ты делаешь?>> --- требовательно спросил Арракис.

<<Люцифер! --- возмутилась Айну.
--- Где твоя дипломатичность?>>

--- Я полагаю, что мой брат выразился абсолютно корректно, --- сказала Нуива.
--- Разумеется, выдача будет производиться на условиях, выгодных для Ордена Преисподней.
Мы готовы произвести её даже на очень выгодных для вас условиях --- скажем, выдать десять перебежчиков Ордена в обмен на одного перебежчика Чёрной Скалы.
Разумеется, с поправкой на ранг.

Айну нахмурилась.

--- Почему вы готовы пойти на такой неравноценный обмен? --- спросила она.

--- Чёрная Скала блюдёт чистоту своих рядов, --- сказал Фуси.
--- И ради этого мы готовы пойти на определённые жертвы.

--- Это очень похвально, --- буркнул Лу.
--- Но, если я вас понял правильно, вы предлагаете обменять троих адептов Ордена --- безусловно полезных, вне зависимости от их ранга --- на несколько десятков совершенно бесполезных перебежчиков, которые более не желают нам служить.
В чём здесь выгода для нас?

--- Разве Орден Преисподней не заинтересован блюсти чистоту своих рядов?

--- Они уже сбежали, --- пожал плечами Лу.
\ml{$0$}
{--- Дурак с повозки --- ослу веселее.}
{``Fool off the cart, donkey feels better.''}

--- Разве Орден Преисподней не карает за предательство?

--- Карает, если это не идёт во вред Ордену.

Легаты переглянулись;
по залу пробежал шёпот.

<<Он ходит по очень тонкому льду, --- подумал Грисвольд.
--- В его словах есть истина, но если бы их вслух сказал я --- мой ранг упал бы до уровня моря>>.

--- Основным приоритетом Союза Воронёной Стали является неотвратимость наказания, --- сказала Нуива.
--- Преступник должен знать, что он будет наказан --- сегодня, завтра, через тысячу или через миллион лет, здесь или где-то ещё.

--- Основным приоритетом Ордена Преисподней является защита адептов, --- ответил Лу.
--- В том числе и от союзных организаций, в том числе и от других адептов.
Именно способность защитить членов делает организацию организацией.

--- Разумеется, мы вполне понимаем ваши мотивы, --- поклонился Самаэл.
--- Вероятно, сайгон Люцифер хотел сказать, что...

--- Люцифер сказал именно то, что хотел сказать, --- подняла руку Айну.
Самаэл тут же замолчал, словно ему в рот засунули рыбину.
--- Чем дальше вы вдаётесь в подробности, тем меньше мне нравится идея такого обмена.
Вы предлагаете обменивать только членов своего клана или, возможно, нескольких кланов --- я не в курсе деталей.
На кого?
На наёмников или отбросов, которые не нужны даже вам?
К дьяволу такой обмен.
На перебежчиков из других кланов?
Я бы послушала мнение средних чинов Союза по поводу такой интересной торговли за чужой счёт.
На тех, кто уже служит вам?
А чем они заслужили такое отношение к себе?

--- Думаю, мы с вами сходимся во взглядах на то, чего заслуживают предатели.

--- От наших рук --- да.
Но неужели все они плохо служили вам?
Это самый главный вопрос --- можно ли принципом неотвратимости наказания оправдать произвол в ваших собственных рядах?

--- Мы оказываем Ордену любезность, давая возможность свершить правосудие.

--- Нельзя, --- перевёл с дипломатического языка Грисвольд.
--- Нам следует присмотреться к выходцам с Чёрной Скалы.
Возможно, среди них есть те, кто стоит и тридцати, и ста перебежчиков.

--- Я правильно понял, что вы предлагаете поднять цену? --- осведомился Фуси.

\ml{$0$}
{--- Отказ не является предложением поднять цену, --- сказал Люцифер.}
{``No doesn't offer to up the price,'' Lucifer said.}
\ml{$0$}
{--- Отказ не является требованием отсрочки.}
{``No doesn't promise deferment.}
\ml{$0$}
{Отказ не является оскорблением.}
{No doesn't try to offend.}
\ml{$0$}
{Отказ --- это просто отказ.}
{No is just no.''}

--- Что скажет глава Ордена? --- обратился Фуси к Арракису.

--- Вы только что услышали слова главы Ордена, --- ответил Арракис и тоже встал.
--- Законы Союза не распространяются и не будут распространяться на адептов Ордена Преисподней.
Переговоры окончены.

Все начали собираться.
Многих взволновали последние слова Арракиса, но, как видно, ни для кого они не стали большим сюрпризом.
Грисвольд кивнул Лу, как другу, Айну --- как старшему по званию, Самаэл --- с подобострастно-оценивающим выражением лица.
Фуси пронзил Люцифера острым взглядом --- словно поставил метку --- и, взяв за руку Нуива, удалился.

Люцифер взглянул на Стигму.
Женщина сидела, отвернувшись к стене и подперев голову рукой.
Она успела завернуться в его плащ;
по запястью струйкой текли слёзы.

--- Я никому не дам тебя в обиду, --- шепнул Люцифер.

Никто из присутствующих не обращал внимания на Гало, который продолжал сидеть на скамье.
На его лице застыл шок.

\section{[:] Игрушка}

\ml{$0$}
{--- Люцифер, я принял решение.}
{``Lucifer, I have decided.''}

\ml{$0$}
{--- Не надо было делать это перед всеми, тем более в такой щепетильный момент, как переговоры!}
{``You shouldn't have done that in front of everybody, especially at a delicate moment like negotiations!}
Стратегия завязана в том числе и на моём текущем положении в системе.
Ты сейчас сильно сдвинул баланс власти, и это будет иметь серьёзные последствия для всех.

\ml{$0$}
{--- Я полагаю, ты прекрасно с этим справишься.}
{``I suppose you're perfectly capable to handle such a situation.''}

\ml{$0$}
{<<Ты вообще слушаешь меня или нет?>>}
{\textit{Are you listening to me or not?}}

Люцифер ходил по богато обставленному кабинету Арракиса, избегая поворачиваться лицом к хозяину.
От одного окна к другому, от двери к изящной конторке, от конторки к статуе, стоящей в углу.

--- Тебе надо было посоветоваться со мной и Гало наедине и решение объявить на следующем совещании.

\ml{$0$}
{--- Я достаточно долго выжидал, выбирая между тобой и Гало, --- сказал Арракис.}
{``I've spent enough time choosing between you and Halo,'' Arrakis said.}
\ml{$0$}
{--- Надо сказать, выбор был непростым, и я до последнего верил в то, что Гало справится лучше.}
{``Should be noted, it wasn't a simple choice, till the end I believed that Halo can do better than you.}
\ml{$0$}
{Ты всегда казался мне туповатым и недостаточно дисциплинированным для руководящей роли.}
{You always seemed slow on the uptake and too undisciplined for the leadership.}
\ml{$0$}
{Но твои успехи я игнорировать не могу.}
{But I can't ignore your achievements.}
\ml{$0$}
{Тебя уважает Айну, тебя уважает Самаэл, политически индифферентные учёные высказываются о твоих методах с одобрением.}
{You are respected by Ain\"{u}, you are respected by Samael, politically apathetic scientists commend your methods.}
\ml{$0$}
{Но самое главное --- благодаря тебе Орден действительно стал единой организацией, а не коллекцией кланов.}
{But most important, because of you the Order performs as one, not anymore as a collection of clans.}
\ml{$0$}
{Это признал даже Гало.}
{Even Halo admitted that.''}

Люцифер вдруг понял, что больше он не может сдерживаться.

--- Если ты так хотел объединить Орден в одно целое, отец, зачем тебе вообще приспичило выбирать между нами? --- Люцифер обернулся.
--- Зачем?
Ты видел, насколько упала общая эффективность Ордена с тех пор, как ты начал активно противопоставлять Гало мне?

\ml{$0$}
{--- Как это относится к вопросу?}
{``How is that relevant?''}

\ml{$0$}
{--- Мы были отличной командой.}
{``We had made an excellent team.}
\ml{$0$}
{Я и Гало.}
{Halo and me.}
\ml{$0$}
{Вместе.}
{Together.}
\ml{$0$}
{Гало делал то, что было не по силам мне, я делал то, что мог.}
{Halo did things I couldn't, and I did what I was able to do.}
\ml{$0$}
{И вдруг тебе понадобилось \emph{выбирать}.}
{And then, suddenly, you got the idea of \emph{choice} into your head.}
Да, ты был настолько уверен в своей правоте, что выбор для тебя стал превыше всего остального.
Ты не считался ни с кем --- ни со мной, ни с Гало, ни с нашими мнениями, ни даже с нашими жизнями.
В результате половина ресурсов у тебя стала уходить на поддержание противостояния, а Гало бросил почти все силы против меня.
И я, именно я, остался один против четверых врагов, двоих из которых мне ещё и приходилось защищать!

Люцифер остановился, поняв, что последние слова были сказаны на пределе крика.
Горло немилосердно ныло.
Арракис холодно смотрел на сына.

\ml{$0$}
{--- Это сделало тебя тем, кто ты есть.}
{``This made you what you are now.}
\ml{$0$}
{Ты стал сильнейшим, ты стал достойным места главы Ордена.}
{You are the strongest, you are worth to be the Order's head.}
\ml{$0$}
{Разве это не моя заслуга?}
{Didn't I have something to do with that?''}

\ml{$0$}
{--- Нет, не твоя, --- бросил Лу.}
{``No, you didn't,'' Lu said.}
\ml{$0$}
{--- Это моя заслуга.}
{``The credit is all mine.}
Как и то, что Орден вообще выстоял в битве против Союза.

\ml{$0$}
{--- Пусть будет так.}
{``If you prefer.}
В любом случае Орден теперь твой.
\ml{$0$}
{Я надеюсь, что ты оставишь для старика подходящую почётную должность.}
{I hope you will keep for your old man an honorable seat that suits him.''}

--- Нет.
Я не стану главой Ордена.
Я останусь тем, кем был --- главным стратегом по обороне.

\ml{$0$}
{--- Если ты отвергнешь мой дар --- будешь разжалован в легионеры.}
{``If you refuse my gift, I will demote you to the rank of legionnaire.''}

\ml{$0$}
{--- Это уже не в твоей власти.}
{``You have no power to demote me.}
\ml{$0$}
{Орден больше не твоя игрушка.}
{The Order nevermore shall be your toy.}
\ml{$0$}
{И я больше не твоя игрушка!}
{I nevermore shall be your toy!}
\ml{$0$}
{Так что иди ты в бездну с такими дарами.}
{So fuck yourself with your gifts.''}

Люцифер развернулся и вышел, хлопнув дверью.

\section{[:] Белая лилия}

Пройдя едва с десяток шагов, Лу наткнулся на посланца-легионера.
Этот легионер был из агентурной сети Айну --- неброский, но исполнительный служака, постоянно претендующий на повышение и стабильно от него отказывающийся.

--- Сайгон, тебе просили передать, --- сказал легионер и протянул свёрток.
--- Ты забыл своё одеяло в зале переговоров.

Люцифер кивнул и принял посылку, отмечая две едва заметных печати на целлофане.
Разумеется, вещь дважды просканировали на предмет опасных сюрпризов --- вначале принявший плащ легат, а потом и центурион --- непосредственный начальник доставившего его легионера.

--- Благодарю, Ашита.
Ты свободен.

Легионер поклонился.
Лу привычным движением набросил плащ на плечи и застегнул фибулу.
Затем поморщился и выудил из-под горловины нечто, оказавшееся цветком белой лилии.

<<Какая прелесть>>.

Люцифер несколько секунд просто любовался цветком.
В ущелье Такэсако белые лилии были символом девственности и первой любви.
Молодые люди дарили их тем, перед кем они хотели бы впервые в жизни сбросить одежду.

<<Неужели она?..>>

В Стигме было что-то притягательное для Лу.
Он не совсем понимал, что это.
Её тело сложно было назвать красивым --- лёгкая неженственная полнота, скрытая под необъятным свитером, хмурое выражение лица, сильное косоглазие, короткая стрижка, руки постоянно скрещены на груди.
Её ум не вызывал сомнения, но едва ли Лу узнал бы о нём, если бы не собрал на неё досье.
Впрочем, какая разница...
Он уже готов был дать волю фантазии, но рациональная часть его останавливала все мысли в этом направлении.

В Стигме напрочь отсутствовала сексуальность.
Это была не юношеская обида, не травма, это было отсутствие как таковое.
В Стигме было меньше сексуальности, чем в ком бы то ни было из его знакомых.
Но сексуальный подтекст был неплохим способом...

Пальцы Лу пробежались по стеблю.

...отвести глаза.
Вот и сообщение, набитое иглой.

\begin{quote}
<<Есть важная информация.
Организуй встречу с глазу на глаз>>.
\end{quote}

--- Ашита, --- окликнул Лу уходящего посланца.

--- Да, сайгон?

--- Можешь попросить ту учёную, которая передала одеяло, зайти ко мне в комнату?

--- Будет сделано, --- легионер козырнул и направился в сторону лифтов.

<<Неплохо она меня изучила>>, --- признал Лу и, сунув лилию в рот, разжевал сообщение до полной нечитаемости.

\section{[:] Родственные узы}

--- Заходи, --- Люцифер приоткрыл дверь ровно настолько, чтобы немного полненькая Стигма смогла протиснуться.
\ml{$0$}
{--- Здесь только я и Тахиро.}
{``There's only me and Tahiro in the house.}
\ml{$0$}
{Тахиро --- мой...}
{Tahiro is my---''}

\ml{$0$}
{--- Знаю.}
{``I know.}
\ml{$0$}
{Я не возражаю.}
{I don't mind.''}

Стигма натянула латексные перчатки и молниеносно прошла по комнате, открывая шкафы и проверяя коробки.
Некоторые предметы она интенсивно трясла, другие --- аккуратно переворачивала, третьи --- скребла, по четвёртым --- стучала.
Каждая вещь, побывав в её руках, вставала в точности на то место, на котором была оставлена.
Её демон <<вспыхнул>>, по комнате светлячками полетели считывающие жуки.
Люцифер понял, что стратег проверяет комнату на наличие шпионской аппаратуры, но алгоритм был ему неизвестен.

<<Это не похоже ни на наши протоколы, ни на протоколы Союза, --- восхищённо подумал он, наблюдая за её отточенными действиями.
--- Видимо, какие-то секретные военные разработки Чёрной Скалы.
Какое счастье, что она пока на моей стороне>>.

Люцифер вдруг понял, почему Стигма казалась ему такой привлекательной.
В полноватом теле, в короткой стрижке, в сильно косящих глазах царила мощь.
Это была не та разрушительная сила, которая исходила от Двойняшек, это было спокойствие непробиваемого бастиона, в котором не найдёт лазейку ни одна крыса.

<<В Ордене Тысячи Башен она делала то, что делаю сейчас я.
Ей было тяжелее всех, когда Тысяча Башен металась между Преисподней и Союзом, на неё пала вся тяжесть поражения.
Мы --- как две стороны одной монеты, которая лишь по случайности упала мною вверх.
Возможно, это единственное существо, способное меня понять>>.

Стигма пролетела перед самым носом Люцифера, обдав его неожиданно приятным дуновением.

<<Я мог бы просто с ней расслабиться, --- понял стратег.
--- То, чего я не могу себе позволить вообще ни с кем>>.

Последним было окно.
Стигма ожидаемо проверила наличие лазеров, считывающих колебания стекла, и налепила на самый центр крохотный гаситель колебаний, игнорируя уже имеющийся сбоку.
Задёрнув плотные занавески, женщина обернулась к Люциферу.

\ml{$0$}
{--- Сразу скажу, чтобы не было недопониманий.}
{``To avoid misunderstandings.}
\ml{$0$}
{Это не моя благодарность, это не моё предложение дружбы и чего-либо ещё.}
{It's not a graditude, it's not an offer of friendship, or that stuff.}
Это вопрос жизни и смерти как для меня, так и для тебя.

\ml{$0$}
{--- Я слушаю.}
{``I'm all ears.''}

\ml{$0$}
{--- Ты в курсе, что Двойняшки планируют совершить на тебя покушение?}
{``You know that Siblings are plottin' to assassinate you?''}

\ml{$0$}
{--- Это очевидно.}
{``It's obvious.''}

\ml{$0$}
{--- Это очевидно или у тебя есть данные?}
{``It's obvious, or it's proved?}
\ml{$0$}
{У меня они есть.}
{I have data.''}

\ml{$0$}
{--- На меня совершали покушения со времён существования Ордена Тысячи Башен.}
{``Since Order of Thousand Towers existed, I've been being a target of assassination attempts.}
\ml{$0$}
{К некоторым и ты приложила руку.}
{You had a hand in some of them.}
\ml{$0$}
{Ничего нового.}
{Nothing new.''}

\ml{$0$}
{--- Люцифер, это другое.}
{``Lucifer, it's different.}
\ml{$0$}
{Если за тебя берутся Двойняшки --- это совсем другое.}
{If you have Siblings goin' after you, it's extremely different stuff.''}

\ml{$0$}
{--- Я рад, что ты такого высокого мнения о братишке и сестрёнке.}
{``I'm glad you think so highly of your little brother and sister.}
\ml{$0$}
{Переходи к делу, ты явно пришла не за тем, чтобы спасать меня.}
{Get to the point, you're not here to be my saviour.''}

--- Представители Союза ушли час назад, --- сказала Стигма.
--- Гало ушёл с ними.

\ml{$0$}
{--- Что? --- выдохнул Тахиро.}
{``What?'' Tahiro silently mouthed.}
\ml{$0$}
{--- Этого не может быть.}
{``It can't be.''}

Лу просто стоял и смотрел на Стигму.
Стигма схватила его за плечи.

\ml{$0$}
{--- Я понимаю, что это звучит дико, но мы все заинтересованы в том, чтобы Орден как можно скорее взглянул истине в глаза и начал действовать.}
{``I understand how crazy it sounds, but the Order must face the truth and make a move, the sooner the better for us all.}
Люцифер, за те несколько лет, что я живу здесь, я впервые вздохнула полной грудью.
Если тебе так важно моё признание, во многом именно благодаря твоей работе я наконец чувствую себя в относительной безопасности.
\ml{$0$}
{Ты знаешь, что произойдёт со мной, если Орден падёт.}
{You know I'm doomed if the Order falls.}
\ml{$0$}
{То же самое произойдёт и с тобой.}
{You know you're doomed too.''}

--- Переходи к делу, --- потребовал Тахиро.
--- Ему нужны факты, а не твои излияния.

--- У них были краткие переговоры, после чего Гало перестал выходить на связь.
Его агентурная сеть не получала никаких приказов.
Его тело лежит в хранилище тел, пустое, но я нашла в театре действий сообщение неизвестным кодом на обочине основных информационных каналов.

--- Покажи, --- безжизненным голосом попросил Люцифер.

Стигма коснулась лбом его лба.

\begin{quote}
<<Брат, я уверен, что ты будешь великим.
Всё это время я жил чьей-то чужой жизнью.
Я думал, что руковожу, я думал, что в моих руках сосредоточена власть, но я просто исполнял чужую волю.
То, что я сделал --- это моё и только моё решение.
Меня будут презирать так же, как презирали тебя, и я постараюсь пройти тот путь, который прошёл ты.
Даже если мы встретимся на поле боя и один из нас погибнет, до самого последнего мгновения это решение будет маяком, который ведёт меня к истине.
Я знаю, что ты меня поймёшь.

Я передам Союзу содержимое архивов 14, 21 и 60, а также наши с тобой стратегические выкладки за последние 12 лет.
Желаю тебе удачи.

P.S. Красивый плащ.
У тебя всегда был хороший вкус>>.
\end{quote}

Люцифер схватился за стену, чтобы не упасть.
Тахиро и Стигма подхватили его под руки.

--- Дружище, ты в порядке?

--- Нет, Тахиро.
Я ни на каплю, ни на атом, ни на что в этой протухшей Вселенной не в порядке.

--- Принести воды?
Может, хочешь покурить?

--- Нет.
Брось меня здесь, я пока не умираю.
Немедленно собери наших.
Скажи следующее: Гало переметнулся, нам срочно, прямо сейчас, прямо сию минуту нужен новый план --- с учётом того, что все наши данные по разведке и обороне, а также внутренние протоколы обмена информацией известны Союзу.
Оповести Грисвольда, Айну, Самаэла, ведущих научных сотрудников, легатов.
Отца не зови.
Видеть его не хочу.

Тахиро подхватил Лу на руки, пинком развернул кресло и положил безвольное тело друга на мягкие подушки.

--- Я всем написал, --- сообщил он.
--- Они будут у тебя в комнате через десять минут.

--- Ты не можешь просто взять и не пригласить отца, --- заметила Стигма.
--- Он глава Ордена.

Лу промолчал.
Тахиро понял.

--- Я скажу Айну, чтобы она отправила демонов арестовать Арракиса.

\ml{$0$}
{--- Я думаю, она уже сама поняла, что наделал отец, --- кивнул Лу.}
{``I think she figured out what father has done,'' Lu nodded.}
\ml{$0$}
{--- Если не поняла --- растолкуй ей.}
{``If she didn't yet, you must be the one to explain.}
\ml{$0$}
{И ещё... }
{And one more thing ...}
\ml{$0$}
{Если отец сбежит из-под стражи или уже сбежал --- пусть они не тратят время на поиски.}
{If father escapes---or have already escaped---there's no need to look for him.}
\ml{$0$}
{Скорее всего, он просто забаррикадируется где-нибудь в далёком мире, объявит себя богом, а демонов Преисподней --- предателями.}
{He's likely to barricade in a faraway world, declare himself its god, and charge daemons of Nether with treason.}
\ml{$0$}
{На большее он никогда не был способен.}
{He never could do better than that.''}

--- Хорошо.

--- Я вряд ли смогу председательствовать, --- деревянным голосом сказал Лу.
--- Стигма, можно тебя попросить?..

--- Мой голос не настолько значим, --- ответила стратег.
--- Давай ты начнёшь, а я поддержу, если что.

--- Поплачь, --- Тахиро сел рядом и заключил Лу в объятия.
--- Просто поплачь.
\ml{$0$}
{Люди так делают, когда им плохо, это хорошо помогает.}
{Humans do cry when they feel very bad, it really helps.}
\ml{$0$}
{Можешь даже покричать.}
{You may even scream.}
Не надо нас стесняться.
Давай, дружище, график плотный.
У тебя десять минут.

И Лу взвыл.

\chapter{[U] Светлячки}

\section{[U] Ааман (отрывки)}

Стигма по очереди взглянула на Штрой сильно косящими глазами, словно взглянули два человека --- весёлый и задумчивый.
Штрой напряглась --- самый этот взгляд вызывал в её теле тревогу.

Из-за ширмы неторопливо вышел худой человек с непробиваемо тупым лицом.
Это лицо могло ввести в заблуждение сейхмар, но не хоргета --- глаза светились умом и странным, непередаваемым фанатизмом.
Лицо человека было невыбрито, кустистая щетина испачкана в краске.
Одежда человека также была до крайности безвкусной и неопрятной;
она выглядела не то как содранное с врагов платье, не то как найденные бродягой лохмотья.
Сквозь прорехи одежды просвечивала неправдоподобно чистая безволосая кожа, обтягивающая крепкие верёвчатые мышцы и сетку вздутых вен.
Человек был сочетанием несочетаемого --- как снаружи, так и внутри.

Ааман Третий Защитник.
Оцифрованный сапиент с одной из периферических планет, принадлежащий к сильно мутировавшему виду Людей --- ко. % Qwoh
У ко было всего четыре пальца на руках, и некоторые из оцифрованных, оказываясь в <<обычных>> человеческих телах, по привычке ампутировали себе мизинец.
Ад взял на службу несколько Телохранителей из местного племени.
Они следовали своей философии --- неважно, кто попросил тебя о защите, ты обязан защищать попросившего даже ценой жизни.
Этих пятерых интерфекторов использовали на переговорах как гарантию мира --- они без колебаний пресекали любую агрессию.

<<Атрис связался со Стигмой и предупредил её насчёт меня.
Или?..>>

Штрой нахмурилась.
План придётся менять на ходу.

Приближаться к Атрису ещё раз было чревато.
От Ааман она не добьётся ничего --- они ходили в Чистилище как на работу, аналитики уже давно махнули на них рукой.
Телохранители до крайности опасны в бою, но и настолько же предсказуемы, у них есть свой запас упругости, при котором можно действовать в относительной безопасности.
Рискнуть и попробовать выжать нужные данные из рыжей стервы?
С Хореймом и прочими она могла жёстко надавить на Стигму, но этот демон... он не должен быть здесь, у Стигмы просто нет необходимых каналов, чтобы вот так запросто нанять Телохранителя!
Что, если его со Стигмой связывает нечто большее, чем простая клятва?
Они друзья, приятели, любовники?

Штрой выросла в связке конкурирующих стратегов, где успех напрямую зависел от скорости изложения идей.
Эта методика была достаточно распространённой.
Из-за этого Штрой чувствовала себя крайне некомфортно в соло, когда не было возможности поделиться догадками с прочими.
Стигма же знала, что иногда лучше даже не думать --- всегда есть вероятность, что кто-то изобрёл устройство для чтения мыслей.
Также она была осведомлена и об изъяне соперницы;
она пытливым взглядом смотрела на лицо Штрой, ожидая предательского движения губ или щёк.
Штрой это поняла.

<<Стерва>>, --- прочитала Стигма по розовым цветущим губкам.

<<Даже такую важную информацию, как осведомлённость о слежке, она не смогла держать в себе, --- подумала Стигма.
--- Вот что стоило сейчас обратить слабость в силу и ввести меня губами в заблуждение?
Типичный кабинетный стратег, настойчиво лезущий в полевую романтику>>.

\textspace

--- Зачем ты занимаешься этой ерундой?
Ты же не биолог и никогда им не была.

--- Даже у демона могут быть мечты, --- сказала Стигма.
--- Однажды я стану биологом.
И жуки будут у меня жить на более веских основаниях.

Штрой подошла к террариуму и открыла его.
Индиго-светлячки тут же разлетелись в разные стороны.

--- Ой, --- притворно прикрыла рот Штрой.
--- Прости.
Ты же сможешь их потом собрать?

--- Я всё равно собиралась их выпустить, --- улыбнулась Стигма.
--- У них есть крылышки и челюсти.
Пусть летают и грызут, что найдут.

\textspace

Штрой поднялась и подошла к Стигме вплотную.
Диковатое оцелотовое лицо пылало сексуальным жаром и яростью ущемлённого самолюбия, и сколь много огня было в девушке, столь же холодным, напрочь лишённым сексуальности, но невероятно гармоничным казалось лицо Стигмы.
Светящиеся чистым светом расчёта глаза одним своим существованием высмеивали амбиции Штрой, и подошедшие Ааман властно оттолкнули Штрой четырёхпалой ладонью, напоминая молодой демонице о своём существовании и своей роли.

<<Ааман молодец, --- отметила Стигма.
\ml{$0$}
{--- Едва ли они знает о смертельном сюрпризе в моём комбинезоне, но последствия физического контакта они просчитали совершенно правильно>>.}
{``They hardly know of the deadly surprise I've got in my suit, but they correctly predicted consequences of physical contact.''}

Штрой, похоже, тоже поняла, какого рода опасность ей грозит.
На её лице расцвела и застыла улыбка превосходства.

\ml{$0$}
{--- Хорошая попытка, Стигма.}
{``Nice try, Stijma.}
\ml{$0$}
{Как вижу, ты не любишь марать руки.}
{You don't like to get your hands dirty, I guess.''}

Стигма вместо ответа показала испачканные удобрениями ладошки.

\ml{$0$}
{--- Очень смешно, --- скривилась девушка.}
{``Very funny,'' the girl made a face.}
\ml{$0$}
{--- Разумеется, я сообщу куда следует о твоих методах ведения переговоров.}
{``Of course I'll report about your negotiation techniques.''}

--- Ты хочешь наказать меня за желание защититься?

--- Ты уже давно перешла границы защиты.

Стигма позволила себе немного искренности:

\ml{$0$}
{--- Это ты пришла ко мне, а не наоборот.}
{``You came to me, not the other way around.}
\ml{$0$}
{Я не хочу делать тебе больно.}
{I don't want to hurt you.''}

\ml{$0$}
{--- Ты и не сможешь.}
{``You can't.}
\ml{$0$}
{Однако драться тебе придётся, --- оскалилась Штрой.}
{Like you can't escape fighting,'' Stroji grinned.}
\ml{$0$}
{--- Ты предала Орден.}
{``You betrayed the Order.}
\ml{$0$}
{И я это докажу.}
{I will prove it.''}

Штрой махнула свите и удалилась.
Убедившись, что замок щёлкнул, Стигма повернулась к Ааман.

--- Стигма благодарит их, Телохранитель.
Их клятва исполнена.

Ааман поклонились.

--- Стигма окажет им честь? --- спросили они низким хриплым голосом, не привыкшим к беседе.

--- Да, конечно.
Что они желает?

Ааман длинным ногтем постучали по её груди.
Стигма клацнула застёжкой лифа, вытянула лиф через рукав комбинезона, скальпелем сделала тонкий надрез на запястье.
Тонкую сеть синтетического полотна оросили капли крови.

--- Большая честь, --- Ааман снова склонились, глядя, как Стигма аккуратно завязывает лиф на их запястье.

--- Честь для меня, --- парировала Стигма.

<<Какие они потрясающие.
Они знает, что в застёжке, но не задаёт ни единого вопроса.
А надо бы>>.

--- Когда они сможет доиграть эту партию? --- поинтересовались Ааман.
--- Им она пришлась по душе.

--- Сегодня у Стигмы есть дела, но с завтрашнего дня --- когда они пожелает.

--- Сейчас её ход, Стигма Чёрная Звезда, --- Телохранитель кивнули и вышли вслед за Штрой.

\asterism

Штрой ждала Аамана у лифта.

--- Здравствуй, друг, --- сказала она.
--- Я знаю твою приверженность кодексу Телохранителей Слит-Же и потому хочу попросить тебя о защите.

Ааман медленно повернулся и посмотрел на девушку.

--- Ты помешал мне служить Аду.
Идти против Ада неразумно.
Если ты будешь защищать меня от заговорщиков...

--- Штрой Кольцо Дыма, --- сказали Ааман.
--- Их глаза широко открыты.
Её сон будет спокоен.
Сталь в её голосе, сталь в её ладони --- смерть её же, забвение до конца времён и после него.
Они будет сопровождать её всегда, пока она не скажет, что клятва исполнена.

Штрой нахмурилась.

--- А если я освобожу тебя от клятвы на день?

Ааман оскорблённо выдохнул.

--- Больше она её не получит, --- процедил Ааман.
\ml{$0$}
{--- Их клятва не спит и не ест.}
{``Their vow has no sleep and no eat.''}

--- Ты свободен, уходи, --- рявкнула Штрой, великолепно изображая ярость.
--- И впредь думай десять раз, прежде чем помогать врагам Ордена.

Ааман поклонился и молча вошёл в лифт.

--- Отлично, --- с удовлетворением сказала Штрой.
--- Осталось четверо Защитников.
Ещё трое подвергнутся пробам, и пятый будет служить мне до смерти.
Раб кодекса --- раб того, кто найдёт в этом кодексе прорехи.
Странно, что никто не сделал этого раньше...

<<Умная, но чересчур болтливая>>, --- молча констатировала Стигма и выключила микрофон.
Светлячок на спине Штрой задрал лапки и упал на пол.

\section{[U] Плут и Ягуар}

Стигма подошла к доске с недоигранной партией Метритхис.
Эта игра очень нравилась стратегу.
Стигма часто играла в неё сама с собой --- никто из знакомых демонов не придёт ради этой очевидной траты времени.
То, что она спросила Ааман --- невероятная удача во всех смыслах.
Никто бы не смог распознать в одном из сильнейших интерфекторов Ада приятного собеседника и игрока.

На доске нерешительно топтался зелёный Плачущий Ягуар.
Человек чести, чьё призвание --- служить.
Однако его окружили чёрный Плут и красный Разбойник --- не признающие верности, не имеющие чести.
Ягуар был обречён.

Немного подумав, Стигма передвинула Плута поближе к Ягуару и бросила камень.
Плут превратился в Купца, Плачущий Ягуар почернел и превратился в Телохранителя, и попавший под горячую руку Разбойник отправился в коробку.

Стигма улыбнулась и подумала про себя, что отныне Купец и Телохранитель пойдут одной дорогой.

<<Их ход следующий, Ааман>>.

\chapter{[U] Оскал обречённых}

\section{[U] Сиэхено}

\textspace

Даже обществу демонов визоры казались эксцентричными;
старое название --- оракулы --- было им дано не случайно.
Инкарнированные визоры переносили отпечаток личности на собственные тела --- они становились плаксивыми и эмоционально лабильными, имели тягу к самоповреждению, не следили за собой.
Кто-то связывал эту особенность с родом деятельности, но истина была простой и неприглядной --- ранних визоров презирали за беззащитность, с ними обращались, как с жуками.
Паттерны тех запуганных, покорных созданий используются до сих пор, и никакая коррекция личности уже не способна это исправить.

Девушка подняла голову и улыбнулась слабой желтозубой улыбкой.

--- Анкарьяль, старшая сестра.
А я-то думаю, кто так хочет со мной увидеться, что убивает мою стражу?..

--- Меня зовут Таниа, --- сказала воительница.

Улыбка Сиэхено увяла.
Она схватила бусы и начала раскачиваться из стороны в сторону.

--- Я охотилась на неуловимого Аркадиу Люпино, --- прошептала она.
--- И сама попала в его сети.
Ты его подруга, Тханэ ар'Катхар.
Значит, он тебя оцифровал.
А Анкарьяль... змея подколодная... её паттерны...

Сиэхено всхлипнула.

--- Я знаю обо всех устройствах в твоём теле и об их назначении, --- предупредила Чханэ.
--- Каналы связи контролируют мои демоны.
У нас ещё шестнадцать михнет на разговоры.

--- За мной придут, --- прошипела Сиэхено, --- и гораздо раньше.

--- Замечательно, --- равнодушно сказала Чханэ.
--- Только на Капитуле в процессе легенда о твоей измене.
Она не особо впечатляющая, но, учитывая твою специализацию, тебя при встрече уничтожат без разговоров.
Так что чем скорее за тобой придут, тем сговорчивее тебе стоит быть.

--- Никто не поверит, что я работаю на Картель!

--- Разумеется.
И поэтому, согласно легенде, работаешь ты на нас.
Нет ничего лучше, чем притвориться, что желаемое уже реально.

Сиэхено заплакала.
Девушка по-детски тёрла ручками незрячие глаза, размазывая слёзы и грязь по треугольному личику.
Что это --- умелая игра или проявление интимной связи, которая установилась между телом и демоном?
Чханэ пододвинула стул поближе.
Распознать подобное по действиям демона мог лишь аналитик, но попытаться стоило.

--- Знаешь ли ты, чьё тело дало тебе приют?

--- Знаю, --- всхлипнув, рассеянно сказала Сиэхено.
--- Я третья наследница, твоя и Ликхмаса ар'Люм.
Истории о вас кормилец рассказывал мне, едва я научилась понимать.

Сиэхено заулыбалась так же неожиданно, как заплакала.

--- Хаяй, куда же делась героиня моего детства?
Где храбрая и верная Тханэ ар'Катхар? --- девушка снова начала теребить бусы и раскачиваться.

Сомнения Чханэ рассеялись --- как и прочие агенты Ада, Сиэхено превратила тело тси в бездушный инструмент.

--- Ты изуродовала ей глаза, мерзость, --- невыразительно сказала Чханэ.

--- Это мой знак, --- пропела Сиэхено.
--- Ла-ла-лай, ниии, трааа...
Я обрабатываю глаза щёлочью всем своим телам.
Так лучше видно.

--- Ты боишься связи с телом, --- сказала Чханэ.
--- Почему?

Сиэхено погрустнела.

--- Темнота.
Я очень боюсь темноты.

Это было похоже на правду.
Когда только появился ангельский тип заякоривания, урождённые демоны пришли в ужас от образов, возникающих в сапиентном мозге при дефиците внешних сигналов.
Хоргеты не знают, что такое темнота --- белый шум ПКВ существует всегда, его очень сложно устранить.
Но развитый интеллект легко делает проекции.
Страх темноты засел в сознании многих демонов... и кое-кто этим манипулировал.

<<Если такие кошмары видят существа с короткой жизнью, что увидим мы?..>>

Нет зрительного анализатора --- нет темноты.
Звучит логично.
Однако Сиэхено не интегрировалась в девушку-носителя, и боязнь темноты была лишь отговоркой.
Куда больше визор боялась собственного тела --- и потому сознательно его ослепила.

\emph{Страх}.

Чханэ почувствовала, что её пробирает дрожь.
Знакомая дрожь, дрожь ярости и негодования.
Глаза тси отлично регенерируют при повреждении.
Сиэхено вводила в глаз щёлочь снова и снова, пока в погибающий нерв не пророс грубый рубец, пока хрустально прозрачную оптику глаза не заплело опалесцирующей коллагеновой сеткой, пока зрительный анализатор не атрофировался из-за отсутствия сигналов.
Что при этом чувствовала девушка, частью которой демон не желал быть?

\textspace

Чханэ почувствовала, что по её лицу текут слёзы.
<<Лис, что мне делать?>>

--- Бабушка, почему ты плачешь? --- вдруг сказала тонким голосом Сиэхено.
--- Я скушаю пирожок и расскажу, какой он вкусный.

Слёзы высохли на загоревшемся лице Чханэ мгновенно, и к горлу подступил гнев.
Эта тварь издевалась над её собственными детскими воспоминаниями.

--- Ах ты...

Сзади слабо вспыхнул источник <<света>>, и Чханэ среагировала автоматически.
Удар --- и визор перестала существовать.
Её хрупкое тело неловко свалилось на стол, невидящие глаза моргнули и тупо уставились в пространство.

Придя в себя, Чханэ поняла --- Сиэхено запустила отвлекающее устройство, имитирующее готовность к атаке, но не учла поведенческие стереотипы воина-Скорпиона.
<<Никогда не оставляй врага за спиной>> --- эту истину вдалбливали в Храме годами, и воительница среагировала, не успев осознать происходящее.
Устройство <<мигнуло>> ещё раз и отключилось, ожидая новой команды.

Испытывая отвращение к себе, Чханэ вонзила в сердце наследницы кинжал, нацарапала на грубой поверхности стола Печального Митра и вышла, тихо прикрыв за собой дверь.

\section{[U] У себя дома}

\textspace

--- Операция пошла не по плану, --- сказала Тхартху.
--- Чханэ наследила.
Она случайно убила Сиэхено, и Ад узнает, что она была убита.
Легенда срочно отозвана.
Мы должны залечь на дно.

--- Хорошо, --- кивнул Атрис.

--- Атрис, тебе тоже лучше скрыться.
Следы могут привести и к тебе, и тогда Митхэ...

--- Мы никуда не побежим, --- спокойно сказала Митхэ.

--- Митхэ...

--- Тхартху, дочка, не беспокойся за нас, --- мягким голосом прошептала Митхэ и обняла женщину.
--- Мы у себя дома.

\section{[U] Пересмотр договора}

Волны накатывались на песчаный берег одна за другой.
Митхэ и Атрис пили отвар, лопали конфеты, беседовали и смеялись.

--- ... помнишь то дерево?

--- Оно называется сорба, Митхэ.
Биологи мне сказали.

--- Давай посадим на Кристалле рощу таких деревьев?
Они красивые.
Или здесь?
Можем их модифицировать, чтобы они выросли большие и толстые!

--- Хорошая мысль, --- одобрил Атрис.
--- Мне почему-то никогда не приходило в голову, что сорба великолепно смотрелась бы вокруг этой беседки.
Давай допьём отвар и свяжемся с Кольбе и Рабе, чтобы...

--- Не торопитесь, --- раздался сзади знакомый тонкий голос, дрожащий от едва сдерживаемого злорадства.

--- Здравствуй, Штрой, --- приветливо сказала Митхэ.
--- Максим Мимоза, легат Ду-Си, и вы здесь.
Надеюсь, не для того, чтобы <<уничтожить человеческую часть Митриса Безымянного>>.
Не хотите ли выпить?

--- К сожалению, именно для этого, Митхэ ар’Кахр, --- звучным, богатым баритоном сказал Мимоза.
--- И если ваш друг немедленно передаст нам все ключи от системы управления, это произойдёт быстро и безболезненно.

Атрис пожал плечами и отхлебнул отвара.
Мимоза --- высокий, холодный и стройный северянин, инкарнация одного из лучших тактиков отдела 100 --- удивлённо приподнял бровь.
Ду-Си, его полная противоположность --- крепко сбитый, покрытый рельефом огромных, выращенных на синтетических стероидах мышц --- одобрительно крякнул.
Даже по телу легата было видно, что он любит трудности: на ноге --- отчётливый шрам от акульих зубов, бритая голова сияет, за ухом виден криво вшитый имплант-инфузомат.
Если верить докладу Аркадиу, Ду-Си --- выкормыш Воронёной Стали, ушёл на вольные хлеба\FM\ ещё до основания Красного Картеля, затем уже в Ордене служил диверсантом в упразднённом отделе 125 и имел богатый опыт геноцида Девиантных Ветвей, уничтожения Роев и несговорчивых богов.
\FL{neutrals-history}{История и современный правовой статус нейтральных демонов}

Идеальное трио.
Штрой --- провокатор и сборщик информации.
Мимоза и Ду-Си играют неприятеля и возможного друга, а в случае столкновения превращаются в интерфекторов со взаимодополняющими тактиками ведения боя.

\ml{$0$}
{--- Атрис, вы ещё с нами? --- осведомился Мимоза.}
{``\Aatris{}, are you still with us?'' Mimosa asked.}
\ml{$0$}
{--- Кажется, я вынужден повторить, что...}
{``I think I have to repeat, you---''}

\ml{$0$}
{--- Я вас услышал, максим, --- ответил Атрис.}
{``You've been heard, maccsim,'' \Aatris\ answered.}
\ml{$0$}
{--- С этим есть некоторые затруднения.}
{``We have some trouble about that.}
Ключей у меня нет, модуль управления планетой соединён с Митхэ интегративным нейроконтуром, и управлять им может только она.
\ml{$0$}
{Всё-таки я ожидал, что перед тем, как перейти к делу, вы вежливо ответите на вопрос, который задала моя женщина.}
{Anyway, I expected you to do my woman the courtesy of an answer instead of cutting to the chase.}
\ml{$0$}
{Смею напомнить, что это мой дом, как в широком, так и в узком смысле этого слова.}
{I'd like to remind: you are in my home, in a wide sense and in a narrow sense of the word.''}

Агенты замерли, словно восковые куклы.

--- Сели за стол и выпили с нами, --- рявкнула Митхэ.
Натренированный тон годился лишь для фразы <<Сдавайся или умри>>, и слова повисли в воздухе.
Интерфекторы с каменными лицами сели за стол и невозмутимо взяли свои чаши.
Штрой, поколебавшись, последовала их примеру.
Её руки дрожали от сдерживаемой ярости.

--- Что ты так смотришь на меня, Штрой Кольцо Дыма? --- холодно осведомилась Митхэ.
\ml{$0$}
{--- Это же <<вкусно>>, забыла?}
{``It's `quite delicious', remember?''}

Девушка промолчала.

--- Почему вы не сообщили о ключах отделу 100, Атрис? --- мягко спросил Мимоза.

\ml{$0$}
{--- Я слегка рассеян, --- признался менестрель, --- и иногда забываю сообщать контрагенту подробности, которые его не касаются.}
{``I'm absent-minded a bit,'' the minstrel admitted, ``and sometimes I forget to provide counterparty with details which are not their business.}
Итак, в чём нас подозревают на этот раз?

--- Вы подозреваетесь в убийстве Сиэхено Опаловый Глаз, --- ответил Мимоза.
--- Вы знакомы с ней?

--- Заочно.
Насколько я знаю, она предоставляет услуги глубокого сбора данных военным и недавно перевелась на Тра-Ренкхаль.

--- Нам известно, что убийство визора Сиэхено координировалось через одну из ваших станций.
Нам известно, что убийца использовал как метку лик Печального Митра --- если не ошибаюсь, это имеет некоторое отношение к вам.
Также нам известно, что агент, собиравший информацию о местоположении визора и связанных с ним демонов, использовал ваш уникальный идентификационный код.
Это лишь малая часть деталей, указывающих на вашу причастность.

--- Могу ли я ознакомиться с прочими?

--- К сожалению, нет.
У вас была бы такая возможность, если бы вы предоставили нам средства контроля над модулем управления планетой.

--- Мимоза, я вам уже объяснил, что это невозможно.
Вы, конечно, можете забрать у меня Митхэ целиком.
Но, во-первых, мы с ней уже тоже достаточно глубоко интегрированы, и разделение может нанести нам обоим непоправимый вред.
Мы можем напрямую слышать мысли друг друга, без использования иных каналов связи --- если вам это о чём-то говорит.
Во-вторых, я сильно сомневаюсь, что даже в случае удачного разделения она захочет с вами сотрудничать.

--- Вас было вынесено официальное предупреждение... --- начал Мимоза.

\ml{$0$}
{--- Мне надоели ваши шпионские игры, --- прервал его Атрис.}
{``I'm tired of these spy games,'' \Aatris{} interrupted him.}
--- Когда Аркадиу Люпино победил на Могильном берегу, он сказал мне, что я и мои создания будут жить счастливо.
\ml{$0$}
{Однако у Ада своё представление о счастье.}
{But the Hell has its own definition of happiness.}
\ml{$0$}
{Стыдно признаться, но даже безумный Эйраки не доставлял мне так много головной боли, как вы.}
{I'm ashamed to admit, even senseless Ejraci was not such a severe headache as you are.''}

Интерфекторы усмехнулись.

\ml{$0$}
{--- Мне нравится ваш подход, Атрис, --- широко улыбаясь, сказал Ду-Си.}
{``I like your attitude, \Aatris{},'' Du-Sie said with a big smile.}
\ml{$0$}
{--- Называть Ад головной болью я бы не решился.}
{``I wouldn't go so far as to call the Hell `a headache'.}
\ml{$0$}
{Особенно сейчас...}
{Especially now ....''}

--- А как ещё вас называть? --- невесело усмехнулась Митхэ.
--- Мы --- ваши союзники, но угрозы я слышу постоянно.
Одна статья, другая статья, десятая статья, пункт сотый, параграф тысячный.
Смерть, смерть, смерть.
Половину законов Ада можно было бы заменить одним этим словом.
После сотой угрозы, какая бы сила за ней ни стояла, уже порядком надоедает.
\ml{$0$}
{Вы пришли меня казнить, а мне ужасно скучно.}
{You come here to execute me, and I find it boring.''}

Мимоза изобразил холодное удивление:

\ml{$0$}
{--- Интересные новости.}
{``Nice news.}
\ml{$0$}
{Можно узнать, от кого вы слышите угрозы в таких количествах?}
{May I ask, who threatens you so often?''}

\ml{$0$}
{--- Да хотя бы от неё, --- Митхэ махнула рукой на Штрой.}
{``I might start with her,'' \Mitchoe\ pointed to Stroji.}
\ml{$0$}
{--- Я не удивлюсь, если именно она нас подставила.}
{``I wouldn't be surprised if it was she who framed us.''}

Демоница вытаращила глаза.

\ml{$0$}
{--- Kuna!}
{``\emph{Kuna!}''}

\ml{$0$}
{--- Я рада, что ты, милая моя перебежчица, ещё помнишь ругательства Картеля, --- улыбнулась Митхэ.}
{``I'm glad you, my sweet defector, still remember obscenities of Cartel,'' \Mitchoe\ smiled.}

\ml{$0$}
{--- Стоп, --- поднял руку Мимоза.}
{``Stop,'' Mimosa raised his hand.}
\ml{$0$}
{--- Митхэ, это очень серьёзное обвинение.}
{``\Mitchoe\, that's a very serious accusation.}
\ml{$0$}
{Я бы на вашем месте подумал, прежде чем...}
{If I were you, I would think twice before---''}

\ml{$0$}
{--- Поверьте, максим, у нас была масса времени на размышления, --- прервал его Атрис.}
{``Take it from us, maccsim: we had plenty of time for thinking,'' \Aatris{} interrupted him.}
\ml{$0$}
{--- Штрой Кольцо Дыма пытается нас оскорбить с первого дня знакомства.}
{``Stroji the Smoke Ring have been trying to abuse us verbally, since the first time we met.}
\ml{$0$}
{Записи могу предоставить.}
{Recordings can be provided.}
\ml{$0$}
{Скажите, Мимоза, были ли разрешены такие меры приказами Капитула?}
{Tell us, Mimosa, were such measures allowed by Capitul's orders?''}

\ml{$0$}
{--- Исключено, --- буркнул интерфектор.}
{``No way,'' the interfector answered.}
\ml{$0$}
{--- Рискну предположить, что имело место недопонимание.}
{``I would venture to guess a misunderstanding took place.}
Ваше представление о вежливости может отличаться от представлений командующего, живущего в культурной среде Ордена...

--- Культура подразумевает соответствие некоторым общим нормам, --- перебил Атрис.
--- Я думаю, что наши с вами отношения держатся не только на некой фактической базе, но и на формальностях, которые, кстати, достаточно чётко прописаны в известном вам протоколе под смешным названием <<Шляпа>>.
Я нашёл его очень занимательным.
Там сказано, что провокация, демонстрация власти, лесть и подхалимство являются формами замаскированной агрессии.
Кроме того, там недвусмысленно объясняется, что давление на дружеское и сексуальное чувство в процессе производства является способом установления незаконной, коррупционной связи.
Учитывая, что Штрой Кольцо Дыма умудряется за одну реплику применить все вышеперечисленные методы ведения беседы, недопонимание --- недопонимание командующим Штрой культурной среды Ордена --- действительно имеет место быть.

--- Поведение Штрой определённо заслуживает расследования, но в данный момент...

--- В данный момент мы говорим о том, что имеет к делу непосредственное отношение.
Штрой Кольцо Дыма вразрез с заключённым на Деймос-14 договором постоянно предпринимает попытки ограничить мой доступ к системе управления и требует отчёты о моей деятельности.
Доходит до того, что я орешков пощёлкать не могу, не доложив командующему.
Был ли на эти действия приказ отдела 100 или высшего командования Ада?

\ml{$0$}
{--- Исключено, --- повторил Мимоза.}
{``No way,'' Mimosa repeated.}
--- Однако вы находитесь на территории Ордена Преисподней, и ради безопасности командующий имеет право...

--- Замечательно, --- усмехнулся Атрис.
--- В таком случае я официально требую от вас подкреплённый нормативными актами документ, чётко проводящий грань между обеспечением безопасности и неоправданным ограничением в действиях.
Этот документ будет очень хорошо смотреться как дополнение к договору <<Демиург --- Метрополия>>.

--- Вы имеете на это право, как контрагент Ордена, --- хмуро сказал Мимоза.
--- Но сейчас мы обсуждаем ваши деяния, которые ставят под угрозу саму возможность быть таковым, и я настаиваю...

--- Я ещё не закончил, Мимоза.

--- Мы не намерены продолжать разговор в таком тоне, --- холодно сказал максим.

Ду-Си подобрался, словно тигр перед прыжком.
Митхэ презрительно хмыкнула --- старый как мир приём.

--- Нам этот разговор также не доставляет никакого удовольствия, --- буркнула она.
--- Но довести его до конца --- в наших общих интересах.
Вы готовы слушать Атриса или будете дальше продолжать этот театр?

Мимоза промолчал.
Ду-Си несколько обиженно посмотрел на воительницу и отхлебнул из своей чаши.

--- Отлично, --- продолжил Атрис.
--- Штрой Кольцо Дыма и её демоны регулярно предпринимают попытки самостоятельно исследовать механизмы системы, при том, что у неё есть ключи от системы поиска и всех боевых устройств.
Некоторые попытки вмешательства были с целью захватить контроль над узлами, прочие я не могу трактовать иначе как намеренное вредительство.
Например, не далее как два рассвета назад из-за очередного <<техосмотра>> вышла из строя сначала метеостанция, а затем и три ключевых сейсмодатчика, и землетрясение в Ихслантхаре едва удалось предотвратить.

\ml{$0$}
{--- Есть ли у вас доказательства? --- спросил Мимоза.}
{``Do you have any proof?'' Mimosa asked.}
\ml{$0$}
{--- Если есть, почему это прошло мимо нас?}
{``If you do, why it's passed unnoticed by us?''}

\ml{$0$}
{--- Я понятия не имею, максим.}
{``I've no idea, maccsim.}
Экологической отчётностью, как вы знаете, занимается команда Корхес Соловьиный Язык.
Мои записи они передали вместе с остальными данными на Капитул.
\ml{$0$}
{Вероятно, вам следует навести справки там.}
{Maybe you should make an inquiry there.''}

Штрой остолбенела.
Интерфекторы бросили на неё два совершенно одинаковых взгляда.

--- На будущее, Атрис, --- проговорил Мимоза, покручивая длинными пальцами чашу.
--- Обработка такого рода информации --- юрисдикция не исследовательских отделов, а контрразведки.
Вам повезло, что Корхес имеет право на сбор, передачу и верификацию информации, составляющей военную тайну.
В противном случае вас могли привлечь к ответственности.

--- Я запомню, Мимоза.
Мне последнее время приходится помнить очень много вещей, не имеющих никакого отношения к моим прямым обязанностям.
Например, вы видели это?

Атрис показал на потолок беседки, и все дружно подняли глаза.
На потолке красовалась кривая надпись масляной краской --- <<Б.Д.И.У.>>.

--- Я совсем недавно узнал про ДиС, --- продолжал Атрис.
--- Как и про то, что отдельные группировки в Ордене и Картеле открыто лоббируют их интересы.

--- Это шалость деревенских, Атрис.
Ничего более.

--- То есть вы считаете, что рыбаки и крестьяне приходили в мою беседку и писали в ней лозунги террористической организации демонов?

--- Атрис, по моему опыту общения с ДиС --- а он достаточно богатый, к сожалению --- это не их стиль.
\ml{$0$}
{Они чаще используют взрывчатку, чем краску.}
{They prefer explosive, not dye.}
И я напоминаю, что мы говорили о системе управления планетой.

--- Спасибо, что напомнили, Мимоза.
Надеюсь, вы также напомните командующему Штрой, что всю --- я подчёркиваю, всю --- необходимую информацию о системе уже получили компетентные демоны и доставили её куда следует.

--- Компетентные демоны? --- рявкнула демоница.
--- Грейсвольд Каменный Молот?

--- Вы сомневаетесь в его профессионализме, Штрой? --- вдруг подал голос Ду-Си.
--- Я не знаком лично с Грейсвольдом, но наслышан о его деяниях.
Весьма впечатляющая биография, надо сказать.
Кстати, в настоящее время он является сотрудником отдела 100, и если вы сомневаетесь в его компетентности, мы всегда можем вас выслушать.

--- Они все повязаны одной красной ленточкой, --- буркнула Штрой.
--- Я уже предоставляла свои выкладки на этот счёт в отдел 100.

--- Мы их рассмотрим, Штрой, --- холодно сказал Мимоза.
--- Атрис, у вас всё?

--- Ещё нет, Мимоза.
Согласно договору, заключённому между мной и Адом...

--- Договор --- это пустая болтовня, Атрис.
Ты здесь никто, --- перебила Штрой.

Атрис улыбнулся и встал.

--- Ошибаешься, Штрой.
Договор заключается между равными, и стратеги Ада знают это очень хорошо.
Они также знают, как полезен демиург-союзник и насколько опасен демиург-враг.

Пока менестрель говорил, в нём происходила странная метаморфоза.
Он вдруг сбросил личину бродяги.
В его осанке, голосе и движениях появилось величие.

--- Создатель Тра-Ренкхаля --- это я.
Хозяин Тра-Ренкхаля --- это я.
Планета Тра-Ренкхаль --- это я! --- демон Митриса засветился ярким светом и демонстративно привёл в готовность боевые модули.
Мимоза прищурился, Штрой едва не упала за оградку.
--- И если Орден Преисподней в этом сомневается, желает оскорблять Митриса Безымянного, желает убить Митриса Безымянного, пусть попробует сделать это сейчас.
Я расторгну договор, приму бой и напомню агентам Ордена, на чьей планете они стоят!

Демон погас.
Штрой дрожала всем телом, вцепившись в скамью.
Ду-Си с поэтической рассеянностью тянул конфеты из вазочки.

--- Я теряюсь в догадках, как это может нам повредить, Атрис, --- проговорил Мимоза.
--- Боевые модули системы вам больше не принадлежат.

--- Ну, для начала гарнизону Ада придётся что-то сделать с распадом массивной минус-сингулярности, --- ответила Митхэ.
--- Несмотря на то, что Орден теперь утилизирует большую долю продуции планеты, мы по-прежнему располагаем объёмом масс-энергии порядка константы Ка'нета.
Кроме того, мы теперь умеем менять полярность.
Преобразование Шмидта, --- Митхэ произнесла звучный термин с неописуемым удовольствием, которое доступно лишь дилетанту.

Мимоза широко улыбнулся воительнице.
Митхэ ответила ему своей самой светлой щербатой улыбкой.

--- Знаете, --- непринуждённо сказал Мимоза, --- я очень не люблю иметь дело с самоубийцами.
Это возмутительно неравное сражение.

--- Жаль, что для вас отношения с нами --- лишь очередное сражение, --- парировала Митхэ.

--- Я привык мыслить по-военному, прошу меня за это простить.
В целом --- да, это всего лишь очередное сражение с оправданными жертвами.

--- Планета будет одной из них, --- сладким голосом сказала воительница.
--- Модуль управления, я и Атрис --- три взаимосвязанных звена цепи, которые удерживают планету от падения в бездну.
Не будет любого из нас --- тектонические плиты раскачаются, биолитосфера развалится, и Тра-Ренкхаль вернётся в первозданный вид.
Сомневаюсь, что Орден сочтёт потерю мелиорированной планеты оправданной жертвой.

--- Не волнуйтесь, Митхэ, мы успеем подхватить, --- лениво успокоил её Мимоза.
--- Я уверен, что геологи прекрасно справятся с этой проблемой.

--- Смею напомнить, Мимоза, что геологические исследования вами не проводились и землетрясения до сих пор предотвращаю я, --- подал голос Атрис.
--- Эффективность планеты как источника масс-энергии --- мой каждодневный труд.
\ml{$0$}
{Я дилетант в политике, у меня масса пробелов в естественнонаучных знаниях, но у меня были тысячелетия, чтобы научиться интуитивно управлять планетой и при этом сводить концы с концами.}
{I'm amateur in politics, my knowledge in science is full of gaps, but I've had thousands years to learn intuitively how to handle the planet while getting by.}
\ml{$0$}
{Было ли это время у вас?}
{Have you?''}

\ml{$0$}
{--- Не прибедняйся, демиург.}
{``Don't be so modest, demiurge.}
У тебя были превосходные учителя с планеты Тси-Ди, --- с отвращением выплюнула Штрой и тут же осеклась под взглядом Мимозы.

\ml{$0$}
{--- А с чего вы взяли, Атрис, что исследования не проводились? --- осведомился максим.}
{``So, \Aatris{}, what makes you think research was not carried out?'' the maccsim asked.}

--- Со времени освоения Тра-Ренкхаля ко мне за консультацией не приходил ни один геолог.
В отличие от биологов --- эти ребята бегают ко мне постоянно с тем, что они выудили и выкопали, и результаты их гигантской работы я вижу, можно сказать, в реальном времени.
Из чего можно сделать вывод, что отдел геологии на Тра-Ренкхале --- это красивая ширма, за которой нет ничего, кроме неприглядных планов по уведению моего энергетического сальдо в минус.
%{is a colourful smokescreen, and behind it there is nothing but disgusting schemes of pushing my energy balance into deficit.''}

Ду-Си как-то странно усмехнулся --- Митхэ почудилась смесь стыда и горечи.
Мимоза слушал речь Атриса с выражением врача, сидящего у постели душевнобольного.

--- Вы серьёзно полагаете, что целый научный отдел тратит ресурсы не на науку, а на то, чтобы вас провести?
Атрис, для дезинформации у нас есть специальные отделы.

--- Поменять местами таблички на дверях --- это минутное дело, --- буркнула Митхэ.

\ml{$0$}
{--- У меня создаётся впечатление, что вы считаете Орден Преисподней идиотами.}
{``I get the impression that you take the Order of Nether for fools.''}

\ml{$0$}
{--- Я был бы рад ответить Ордену Преисподней взаимностью, но мне не позволяет воспитание, --- парировал Атрис.}
{``I would be glad to reciprocate the attitude shown by the Order, but I'm too polite,'' \Aatris{} retorted.}
--- Поэтому я заявляю вам официально, как эмиссару: Митрис Безымянный желает пересмотреть договор с Орденом Преисподней.
Делегатов ждём завтра в это же время здесь.
И захватите с собой сладости, Ликхэ ар'Митр в Кахрахане замечательно готовит булочки.
Думаю, у Ордена Преисподней найдётся достаточно слитков, чтобы купить этот шедевр на вес кукхватра.

Мимоза натянуто улыбнулся.

\ml{$0$}
{--- У вас всё?}
{``Are you done?''}

\ml{$0$}
{--- Ещё нет, Мимоза.}
{``Not yet, Mimosa.}
\ml{$0$}
{Если я сочту какие-либо действия членов Ордена неуважительными, я без предупреждения саботирую производство масс-энергии.}
{If members of the Order do something I find disrespectful, I'll sabotage mass-energy production without warning.}
\ml{$0$}
{И если Митрис Безымянный говорит <<саботирую>>, это значит, что никто и никогда больше не сможет использовать планету Тра-Ренкхаль в качестве источника питания.}
{And if Mitris the Nameless say `sabotage', it means: no one ever will use planet Tr\r{a}-R\={e}nkch\'{a}l as an energy source.}
\ml{$0$}
{Запишите это дословно во избежание недопониманий.}
{Write it down verbatim to avoid misunderstandings.}
\ml{$0$}
{Штрой, я не знаю, кому служите вы, поэтому ваше поведение оставляю делом вашей совести.}
{Stroji, I do not know where your allegiance lies, so I leave your behaviour up to you.''}

\ml{$0$}
{--- Полно вам, Атрис, --- осклабился Ду-Си.}
{``Come on, \Aatris{},'' Du-Sie smiled.}
\ml{$0$}
{--- Вы не любите свой народ?}
{``You don't love your people?}
\ml{$0$}
{Ни за что не поверю.}
{I refuse to believe.''}

\ml{$0$}
{--- Очень люблю, --- признался демиург.}
{``I really love them,'' the demiurge admitted.}
\ml{$0$}
{--- Особенно тси, ведь они, как уже было сказано, потомки моих прекрасных учителей.}
{``Especially Qi, because, as previously stated, they are scions of my lovely teachers.}
\ml{$0$}
{Не волнуйтесь, легат.}
{Don't worry, legate.}
\ml{$0$}
{У меня есть внушительная геномная база, подробнейшие сведения о языках и культуре.}
{I've got huge genome database, detailed information on languages and culture.}
\ml{$0$}
{Если мне придётся деактивировать планету, разумеется, я расселю мой народ по Вселенной, где только смогу.}
{If I have to deactivate the planet, of course I will resettle my people throughout the Universe, everywhere I can.}
Думаю, это занятие, достойное того, чтобы посвятить ему остаток жизни.

Штрой приобрела цыплячий оттенок.
Ду-Си крякнул с неподдельным восхищением и шлёпнул ладонью по столу:

\ml{$0$}
{--- А грамотно он нас за яйца взял, а, Мими?}
{``He's quite competently got us by the balls, hasn't he, Mimi?''}

--- Ещё раз назовёшь меня <<Мими>> --- снова получишь пенальти на три ранга, --- бледный как снег Мимоза встал чуть резче, чем следовало, и принялся поправлять смятую одежду.
\ml{$0$}
{--- Атрис, я вынужден сообщить о ваших методах ведения переговоров.}
{``\Aatris, I have to report about your negotiation techniques.}
\ml{$0$}
{Боюсь, у Ордена может возникнуть впечатление, что вы намерены прекратить сотрудничество.}
{I'm afraid this might be considered as intention to severe the cooperation.''}

\ml{$0$}
{--- Ни в коем случае, Мимоза.}
{``No way, Mimosa.}
\ml{$0$}
{Я намерен продолжать сотрудничество, что и подтверждаю запросом на пересмотр договора.}
{I intend to continue to cooperate, as confirmed by my request for re-negotiation.}
Я лишь настаиваю на том, чтобы отдел 100 в контрексте расследования внимательнее оценил, во-первых, вероятность вмешательства ДиС в мою работу, и во-вторых, деяния командующего Штрой, а именно --- ограничение моей свободы передвижения, слежка за мной, а также захват контроля над моими механизмами.
\ml{$0$}
{Я думаю, вы согласитесь с тем, что у командующего планетарными силами были все возможности имитировать мои действия, и уж точно куда больше возможностей, чтобы кого-то убить.}
{I guess you agree that PFC has full opportunity to imitate my actions, and more than full opportunity to kill.} % Planetary Force Commander
\ml{$0$}
{Мотивы ищите сами, у меня масса своей собственной работы.}
{Her motivation is your concern, I have lots of my own work.}
\ml{$0$}
{В конце концов, если цель --- я, почему бы и не ДиС...}
{After all, if the target is me, why not D.o.D. ....''}

Митхэ бросила взгляд на Штрой.
Ту уже в прямом смысле колотило от ярости, и Ду-Си аккуратно придерживал её за локоть, что бы она держала себя в руках.

--- Вы отказываетесь подчиниться правосудию и ввели отдел 100 в заблуждение, --- сказал Мимоза.

--- Я отказываюсь добровольно подчиняться мерам, которые предписал вам Капитул, так как расследование нельзя считать законченным, а их последствия необратимы.
К сожалению, мне пришлось прибегнуть к угрозам, так как вы чересчур серьёзно отнеслись к отданному вам приказу.

--- Вы не можете манипулировать планетой бесконечно, --- Мимоза опустил голос до шёпота.
\ml{$0$}
{--- Однажды терпение Ордена лопнет, и лопнет также без предупреждения.}
{``One day the Order's patience will run dry, and there won't be a warning as well.''}

\ml{$0$}
{--- Я боюсь, это не те угрозы, которыми может разбрасываться адепт Ордена.}
{``I'm arfaid it's not a threat a member of the Order is entitled to make.}
\ml{$0$}
{Отношения с Митрисом Безымянным --- это одно, ставить под вопрос юридическую силу договора <<Демиург --- Метрополия>> --- совсем другое.}
{Relationship with Mitris the Nameless is one thing, and it's quite another to question the legal effect of Demiurge \& Metropolis Pact.}
\ml{$0$}
{Очень многим в Ордене могут не понравиться такие богатые подарки Картелю.}
{Many people won't appreciate such a rich gifts to Cartel.''}

--- Тра-Ренкхаль --- это значимый источник масс-энергии, но один из многих.

--- Мимоза Шёлковая Сталь и Орден Преисподней --- также одни из многих, --- ответил Атрис, --- но это не обесценивает ни мой выбор, ни ваши заслуги.

Мимоза и Ду-Си коротко поклонились и вышли.
Штрой, злобно посмотрев на влюблённых, демонстративно опрокинула вазочку с конфетками и тоже вышла.
Круглые цветные сладости, весело стуча, разбежались по всем углам беседки.

--- Думаешь, это поможет? --- шепнула Митхэ.

--- Уже помогло, --- улыбнулся Атрис.
\ml{$0$}
{--- Давно надо было показать зубы.}
{``We should've shown our teeth sooner.}
\ml{$0$}
{Слишком многие привыкли относиться к богам, как к мусору.}
{Too many people are used to treat gods like garbage.}
\ml{$0$}
{Слишком долго они аплодировали тому, как нас уничтожают.}
{For too long they cheered about our extermination.}
\ml{$0$}
{Слишком привычной стала парадигма <<будь демоном или умри>>.}
{Too accustomed they are to paradigm `Be daemon or be dead'.''}

--- А расследование?

\ml{$0$}
{--- Время играет на нашей стороне.}
{``Time is on our side.}
\ml{$0$}
{Чем больше затянется следствие, тем больше путаницы.}
{The longer the investigation, the greater the confusion.}
\ml{$0$}
{Пока в деле путаница, трогать нас не будут --- я им нужен.}
{They won't bother us until they have some clarity, they need me.''}

--- А если нас всё-таки смогут зацепить?

--- Аркадиу сказал, что следствие уйдёт в сторону адской резидентуры.
Там есть демон --- мастер мутить воду, он сможет повести следствие по нужному пути.
А нас нужно подготовиться к завтрашнему разговору.
Мы потребуем символические привилегии и пару реальных свобод, Орден поторгуется, мы уступим привилегии, Орден всё поймёт и согласится, мы разойдёмся довольные.
Штрой отправится домой.
Разумеется, она докажет свою непричастность к смерти Сиэхено, но репутацию мы ей подмочили основательно.

\ml{$0$}
{--- <<В том, что вас скомпрометировали, есть доля вашей вины>>, --- процитировала Митхэ.}
{``\emph{You're partly responsible for being compromised,}'' \Mitchoe\ quoted.}
\ml{$0$}
{--- <<Аккуратнее выбирайте друзей>>.}
{``\emph{Be careful making friends.}''}

\ml{$0$}
{--- И врагов тоже, --- заключил Атрис.}
{``Making enemies, too,'' \Aatris\ finished.}
\ml{$0$}
{--- Хочешь на берег?}
{``You want to go to the coast?''}

\ml{$0$}
{--- Хочу, милый, --- обрадовалась Митхэ.}
{``I do, honey,'' \Mitchoe\ rejoiced.}
--- Давай наберём ракушек, я давно хотела сделать что-нибудь красивое.
А вечером узнаем у биологов, как выращивать сорбу.

\ml{$0$}
{--- Вряд ли мы успеем её вырастить, --- загрустил менестрель.}
{``We seem to have few time for growing it,'' the minstrel got upset.}
--- После этого инцидента они сильно ускорят процесс исследования планеты и защитной системы.
\ml{$0$}
{Следствие следствием, но когда Орден сведёт баланс сил без нас --- придётся отдать наш дом сельве\FM.}
{The investigation is what it is, but if the Order balances forces without us, we'll be forced to leave our home to the Silva\FM.''}
\FA{
Бежать.
Практика символически отдавать покинутый дом сельве распространена среди северных сели --- в доме ломали крышу и пробивали половые доски, открывая тем самым доступ растениям.
}

\chapter{[U] Ответный удар}

\section{[U] Неудача с Телохранителем}

\textspace

--- Штрой, вы хотя бы день можете прожить без переполоха?
Я слышал, что отделу 100 пришлось нейтрализовать одного из Телохранителей, чтобы спасти вас.
Надеюсь, вы не пытались его завербовать?

Штрой вытаращила глаза.

--- Понятно, --- буркнул голос.
--- Вы могли хотя бы спросить меня перед таким опасным шагом.
Он на вас напал, не так ли?
Каковы потери?

--- Я потеряла четверых агентов, --- Штрой склонила голову.
--- Прошу прощения.

--- Штрой, запомните раз и навсегда.
Ад --- это клубок интриг, и если вам кажется, что кто-то что-то до сих пор не сделал --- скорее всего, на это есть веские причины.
В Кодексе Слит-Же последний пункт звучит так: <<Если подзащитный использует Кодекс для манипуляции Телохранителем, подзащитный должен быть уничтожен>>.
Этот пункт может трактоваться Телохранителем по его усмотрению.

--- Я тщательно замаскировала манипуляцию!

--- А я тебе говорил, что ты недооцениваешь Стигму, --- заявил Самаолу.
--- Может, Телохранитель и достаточно глуп для тебя, но Стигма достаточно умна, чтобы растолковать сложное глупцу.

--- С Защитниками шутки плохи, --- подытожил голос, --- и я запрещаю вам приближаться к ним впредь.

Штрой кивнула с явным облегчением.
Видимо, у неё отпало всякое желание приближаться к Защитникам.

--- А что насчёт Стигмы?

--- Дальнейшие попытки чересчур опасны.
Она уже едва не вывела вас из игры с помощью Безымянного.

--- Это была она?
Это она инструктировала Безымянного?

--- В этом даже я не сомневаюсь, --- буркнул Самаолу.
--- Почерк сложно не узнать.
Она раззадоривает врага мнимой слабостью, и он сам бросается на подставленный нож.

--- Но я хочу попробовать ещё раз!..

--- Никаких <<но>>, Штрой, --- отрезал голос.
--- Я не хочу вас потерять.

Последняя фраза огрела Штрой словно кнутом.
Она упала на колено и прижала руку к груди:

--- Я вас не подведу, владыка!

--- Вы опять начинаете?
Впрочем, ладно, продолжайте.
Стыдно признаться, мне это начинает нравиться.
Но не при агентах.

\textspace

\section{[U] Талианский канал}

\textspace

--- Сиэхено успела оставить сообщение перед смертью, --- сказал Самаолу.
--- Скорбящие используют для сообщения оригинальную восьмиполосную таблицу языка Эй, сочетающую макропараметр планеты Тра-Ренкхаль и неизвестный микропараметр.
К сожалению, сообщение неполное.
Это вся информация, известная на сегодня.

--- Возможно, какая-то вариация Талианского канала, --- предположил голос.
--- А, вы не в курсе, Самаолу.
А вот Штрой должна знать.
Макропараметром таблицы являются географические координаты.
Микропараметром может быть что угодно.
В частности, агенты Картеля использовали бутылочки из цветного стекла, выкладывая ими в определённых координатах оставшуюся часть сообщения.

--- Требуются гигантские затраты, чтобы найти сообщение, --- мрачно сказала Штрой.
--- Но в итоге Ад справился, Талианский канал был раскрыт.
Это положило весомый камешек на чашу победы Ада при Десяти Звёздах.

--- Потому что когда бутылочки появляются в девственном лесу или на необитаемом острове, это наталкивает на определённые мысли, --- согласился голос.

--- Я распоряжусь насчёт поисков, --- кивнул Самаолу.

--- Не тратьте время, --- сказал голос.
--- Талианский канал может быть доработан с учётом местности, климата и прочих особенностей.
Кроме того, сообщения могут быть временными --- например, изо льда или песка.
Не забывайте также, что одним из Скорбящих почти наверняка является демиург, что значительно расширяет список возможных макропараметров планеты.
Может, он вообще кодирует сообщения ветром, молниями или дождём.
Звучит дико, но ведь теоретически это возможно?

--- <<Слова не имеют смысла, если нет того, кто может их понять>>, --- процитировала Штрой.

--- Штрой! --- восхитился голос.
--- Вы цитируете ненавидимого вами Люпино?
Я чрезвычайно рад, что вы растёте как профессионал.

--- Враг тоже может быть прав.

--- Я вам скажу больше, Штрой --- иногда враг ближе к истине, чем вы, но природа мыслящего существа подразумевает ещё и некое комфортное расстояние от этой истины.
Именно поэтому мы стремимся к победе, а Скорбящих устраивает любой мир.

--- Вы хотите сказать, что мы идём неправильным путём? --- ошарашенно спросил Самаолу.

--- Ни в коем случае.
Я хочу сказать, что местоположение истины во Вселенной для практического применения абсолютно неважно.
Важно найти собственное место.
Скорбящие же ищут саму истину --- и в этом их слабость, потому что даже найдя, они не обретут ничего значимого для себя.

--- Боюсь, как бы я не начала им сочувствовать, --- скривилась Штрой.

--- Сочувствуйте, Штрой, сочувствуйте и не бойтесь!
Ненависть заставляет отрезать врагу пути к отступлению, и враг сражается до конца.
Если же вы из сочувствия оставите им лазейку --- они сдадут свои позиции проще и быстрее.

--- Тахиро Молниеносный как-то говорил, что уважение и милосердие более выигрышны в долговременной перспективе, нежели деспотизм и жестокость.

--- И его слова не лишены смысла.
Но он не упомянул, что куда более выигрышным является их грамотное сочетание.

--- Вы имеете в виду жестокость к врагам и уважение к друзьям?

--- Я имею в виду ситуационный подход.
Иногда следует щадить и даже спасать соперника, если его гибель принесёт вам больше вреда, чем пользы.
Иногда следует быть жестким к союзнику, если выгода от этого высока.
Именно поэтому я и не употребляю термины <<друг>> и <<враг>>.
Они чересчур безапелляционные.

--- Быть жестким к союзнику? --- поднял бровь Самаолу.

--- Именно.
Уничтожать союзников не стоит, иначе их у вас не останется.
Впрочем, то же самое могу сказать про соперников.
Отсутствие соперников --- весьма неприятная ситуация, особенно для неумелого игрока.
Нужно держать себя в тонусе, не находите?
На мой взгляд, это достаточно веское основание, чтобы сохранить чужую жизнь.

--- Так что делать с сообщением Сиэхено? --- осведомился Самаолу.

--- То, что я сказал.
Займитесь не кодом, а подозреваемыми --- они сами выведут нас куда надо.

\section{[U] Омега-чувство}

\textspace

--- Только что пришли данные от биологов, --- сказал Самаолу.
--- У тси обнаружили ещё один орган чувств --- омега-чувство.
Они способны воспринимать колебания ПКВ средней мощности.

Молчание.

--- Самаолу, приступайте к реализации варианта номер восемнадцать, --- наконец сказал голос.
--- Мы должны в течение десяти дней убедить Орден в необходимости геноцида тси.
Если не получится --- на Тра-Ренкхале обязан произойти катаклизм.

--- И я даже знаю, на кого можно повесить ответственность, --- ухмыльнулась Штрой.
--- К слову.
Вы не рассматривали моё предложение насчёт фальсификации...

--- Штрой, мой ответ --- нет, --- отрезал голос.
--- Я понимаю вашу любовь к риску, но это не просто риск, это самоубийство.
ФПГП --- это взрывчатка, детонирующая от кашля.
Да, это лёгкий способ устранить Безымянного, но если отдел 100 выйдет на кого-то из нас --- то можно сразу рыть могилы для всех.
Есть данные, что биологам и информатикам, занимавшимся совершенно другими проблемами, в определённый момент приходило письмо из отдела 100.
С одной-единственной цифрой --- номером статьи.
И если ты в тот же день, в тот же час не свернёшь всю ветвь исследований --- ты просто исчезнешь, а твой отдел расформируют.
ФПГП Орден боится как огня, даже больше, чем тси.
Так что ищите другие способы.

\section{[U] Вирус против вируса}

\textspace

--- Тхартху не сомневается в том, что на этот раз они точно решили уничтожить тси.
Но их метод...

--- Что за метод?

--- У тси традиционно практически отсутствует понятие половой гигиены.
Они не знали заболеваний, передающихся половым путём.
Кагуя сказала, что радикальная группировка разрабатывает вирус, который будет передаваться именно половым путём.

Чханэ закрыла глаза.

--- Они нашли слабое место народа.

Митхэ зарычала.

--- Проклятие.
Тси вымрут за год, если мы это не остановим.

--- Есть ли какие-то возможности предупредить народ?

--- Вряд ли, --- опустил голову Атрис.
--- Ограничивать половые сношения?
Разработать средства защиты?
Если за дело возьмётся Орден Преисподней, шансов у нас мало.

Митхэ смотрела на менестреля во все глаза, словно впервые увидела.

--- Извини, моя флейта.
Похоже, твой народ обречён.

--- Хоть что-то! --- Митхэ, казалось, готова вцепиться другу в шею.
--- Хоть какой-то способ, Атрис!
Ты бог этого мира или нет?!

--- Ты права, --- ответил менестрель.
--- Увы, я всего лишь бог этого мира.

--- Твоё слово --- закон!
Скажи жрецам!
Объясни им...

--- И что жрецы скажут народу?
Что важнейший процесс для размножения народа смертельно опасен?
Были бы у них соответствующие технологии, всё стало бы проще.
Лекарства, контрацепция.
Однако на половое воздержание и моногамный образ жизни тси просто не пойдут --- этого не было в их культуре сотни тысяч лет.

Митхэ сдалась.

--- Сколько у нас времени?

--- Декада, может, две, не больше, --- ответил я.
--- Впрочем, есть один вариант.
Если мы запустим на Тра-Ренкхале собственный вирус --- не настолько смертельный, но тем не менее опасный --- группировка Самаолу заляжет на какое-то время.
Официальная политика Ада остаётся прежней --- тси должны жить.
Мы покажем агентам Ада, что тси угрожает опасность, и они будут настороже.

Митхэ хмыкнула.

--- У тебя есть варианты лучше?

--- Я тебя прошу только об одном, биолог, --- тихо сказала женщина.
--- Не перестарайся.
Иначе я убью тебя лично.

\asterism

Улицы Тхитрона опустели.
По ним сновали только люди, носившие трупы к крематорию, который выстроили на месте бывшего постоялого двора.
Митхэ смотрела на это, и по её лицу текли злые слёзы.

--- И все эти смерти лишь затем, чтобы народ мог просуществовать чуть дольше.

\textspace

\section{[U] Отражённый удар}

\textspace

--- Если это не наш вирус, объясните, какого дьявола произошло, --- ледяным тоном потребовал голос.

--- Есть две версии, повелитель, --- доложила Штрой.
--- Первая --- кто-то пытается устроить геноцид параллельно с нами.
Версия неубедительная, так как их вирус недостаточно продуман --- низкая смертность.
Вторая и более правдоподобная --- Скорбящие прознали о нашем плане и попытались нам помешать.
Вирус выкосил слабых, и у оставшихся в живых больше шансов пережить наш сюрприз.
Кроме того, из-за явного диверсионного характера эпидемии на уши встали наши агенты.
Нам придётся отложить атаку.

--- Я уже начинаю сомневаться, что Митрис на их стороне, --- заявил Самаолу.
--- Если анализ его личности верен, он бы никогда не позволил Скорбящим пойти на такой шаг.

--- Если бы это был единственный способ спасти тси... --- начала Штрой.

--- Это не спасение, Штрой, --- раздражённо буркнул Самаолу.
--- Скорбящие выгадали для своих обожаемых тси в лучшем случае год.
Если они за этот год не перережут нас с тобой, то их усилия будут тщетны.

--- За год, Самаолу, может произойти всё что угодно, --- заметил голос.
--- Этот ход остался за ними, хоть вам и не хочется этого признавать.
Они пожертвовали сливой, чтобы спасти персик\FM.
\FA{
Намёк на <<Тридцать шесть стратагем>> Цина --- государства Древней Земли.
}
И, чёрт бы меня побрал, на этот раз я не сомневаюсь, что против нас играл Аркадиу Люпино.
Ищите другие способы, и поменьше биологии.
На этом поле мы пока позади --- у него была целая жизнь, чтобы разложить тси по клеточкам.
Мы вынуждены довольствоваться костями со стола Тра-Ренкхальского отдела биологии.

--- Возможно также, что Люпино получил массу данных по тому, как тси сопротивляются искусственным вирусам, --- сказала Штрой.

--- Именно, --- подтвердил голос.
--- Поэтому ещё раз --- ищите другие способы.
Я недооценил Люпино.
Больше мы не будем играть по его правилам.

\textspace

\chapter{[U] Солёная борода}


\section{[U] Свой среди чужих (переработать)}

\textspace

--- Все данные есть, --- сказал Грейсвольд.
--- Они смогут пробраться на Тси-Ди.

--- Знать бы ещё, куда именно.
В книгах сели я не нашла зацепок.

--- Ты не аналитик, --- поморщился Грейсвольд.
--- Если бы мы могли задействовать аналитиков или получить доступ к их данным...

--- Можно пустить запрос по линии Стигмы, --- предложила Анкарьяль.
--- Вдруг у нас есть агенты-аналитики.

--- Ты вообще знаешь, сколько у нас агентов?

--- Понятия не имею.

--- И я тоже.
Конспирация имеет свои минусы.
Агентов может быть сотня, а может и всего трое.
Попробовать можно, но вероятность мала.

--- Лусафейру не может помочь?

Грейсвольд с жалостью посмотрел на Анкарьяль.

--- Нар, меня аннулируют в тот же час, в который я задам Лусафейру вопрос, даже косвенно имеющий отношение к кольцевым теплицам.

Анкарьяль задумалась.

--- Тогда шпионаж.

--- Нет, --- Грейсвольд, похоже, придумал что-то интересное, его глаза загорелись.
--- Мы воспользуемся служебным положением.
Мы же с тобой контрразведка, верно?
Отыщем агента Картеля под прикрытием и убедим его отыскать нужную нам информацию.
В случае его поимки мы почти чисты.
В случае его успеха у нас есть интерфектор, который в удобный момент выжмет его досуха и уничтожит, --- Грейсвольд церемонно поклонился Анкарьяль.

--- Ты не считаешь, что давать Картелю информацию о местоположении кольцевых теплиц --- плохая идея?
Всегда есть вероятность, что мы упустим агента и он сольёт...

--- Она для них бесполезна.
Проникнуть через барьерную высоту они всё равно не могут, --- Грейсвольд довольно осклабился.
--- Но вряд ли они упустят возможность узнать хоть что-то.

Анкарьяль не удержалась и восхищённо улыбнулась.

--- Грейс, ты сам дьявол.
Может, в таком случае передадим им не данные, а самого агента тёпленьким?

--- Отличная идея.
Свяжись со Стигмой.
Думаю, она оценит.

\section{[U] Старый знакомый}

\textspace

Я, как зачарованный, смотрел на захваченного нами демона.
Сомнений быть не могло, это он, мой создатель --- Яйваф Солёная Борода из клана Дорге.

--- Здравствуй, Яйваф.

Демон неподдельно вздрогнул и уставился на меня.

--- Ты меня зна...?
Ааах, вот оно что.
Аркадиу Валериану Люпино.
И тебе здравствовать, маленький предатель.

Я молчал.
Это была чересчур очевидная, чересчур человеческая попытка вывести меня из равновесия.
Яйваф ухмыльнулся.

--- Ты вырос.

--- Дай мне доступ ко всей имеющейся у тебя информации.

Яйваф взглянул мне в глаза.

--- А если нет?

--- Тогда тобой займётся интерфектор.

Яйваф поёжился.

--- Вам не взять информацию силой.

--- Ты создал меня по своему образу и подобию.
Сколько времени потребуется интерфектору, если он об этом узнает?

--- То есть ты не хочешь меня уничтожить?

--- Вначале хотел, --- честно ответил я.

--- Это трюк.
Ты убьёшь меня в любом случае.

--- Твоё мнение не меняет сути дела.
Обещаю: без фокусов дашь доступ к памяти --- уйдёшь живым.
Твой единственный шанс выбраться --- положиться на моё слово.

Яйваф поджал губы.

--- Какая именно информация тебе нужна?

Я молчал.

--- Проклятие Аду и Небесам, ты действительно вырос.

--- Ты сделаешь то, что я сказал, с требуемой точностью.
Третьего предупреждения не будет.

Яйваф кивнул.
Его демон медленно, по одному, начал снимать барьеры с модулей памяти.

--- Помнишь, Аркадиу, как мы задали жару адским чертям на Серпенциару? --- отстранённо сказал он.
--- Виа Галоледика на десять лет стала красноватой.
А потом ты нас предал.

--- Я больше не служу Аду, --- сказал я.

Яйваф усмехнулся.

--- Знай я тебя хуже, я бы решил, что ты грубо и неумело пытаешься меня одурачить.
Но это правда.
То, что ты опасен равно для Ада и Картеля, мне следовало понять в тот момент, когда ты, мальчишка, приставил к моей голове ружьё.

<<Информация у меня>>, --- сообщила Чханэ.

Я кивнул.

--- Тебе хватит тридцати секунд?

--- Смеёшься? --- сказал Яйваф.
--- Достаточно и двух.

--- Тогда прощай, --- сказал я и отключил глушитель сигналов.
Опустевшее тело Яйвафа закатило глаза и повисло на ремнях кресла.
Мозг, который ещё десять секунд назад нёс демонические маркёры Яйвафа, превратился в некротизированную клеточную кашу.

\chapter{[U] Отшельница}

\section{[U] Ошибка в префиксе}

\textspace

--- Я нашла кое-какие зацепки, --- сказала Чханэ и придвинула ко мне книгу.
--- Смотри.
Иногда в префиксе имени указывается биологический вид.
Редко, но бывает.
А теперь... вот.
Сотканный-из-Темноты-Заяц.

--- Так, --- кивнул я.
--- Она человек.
Хотя постой...

Я вгляделся в иероглифы.

--- Ага, тоже заметил? --- улыбнулась Чханэ.

--- Описка, --- махнул я рукой.
--- Переписчик был уставшим и не вполне понимал смысл написанного.
Спутать иероглифы <<плант>> и <<человек>> довольно просто --- они различаются одной деталью.

--- А здесь? --- Чханэ пролистала на несколько страниц вперёд.

--- А здесь переписчик заметил несоответствие и исправил по образцу более ранней записи.

--- Хорошо, --- сказала Чханэ.
--- А что ты скажешь насчёт этой же книги, найденной у ноа?
В ней стоит такой же знак.
В отчёте культуролога Ада несоответствие упомянуто, и он просил отдел аналитики на всякий случай обработать эту информацию --- вдруг здесь скрывается код.

--- Твоя взяла.
Придумай объяснение логичнее моего.

--- Существует-Хорошее-Небо мало писал о теплицах, несмотря на то, что они играли для народа тси огромное значение.
Всё, что мы знаем --- на Стальном Драконе была теплица, впоследствии сожжённая неисправным двигателем...

--- Постой, постой.
При чём здесь иероглиф <<плант>>?

--- При том, Лис, что он состоит из двух частей --- <<растение>> и <<человек>>.
А как мы знаем, в легенде о Садах кольцевая теплица выступает в роли человеческой женщины.
Чем тебе не <<человек-растение>>?

--- У тебя есть ещё какие-то данные в пользу этой версии?
Что сказал отдел аналитики?

Чханэ замялась.

--- Отдел аналитики счёл это опиской.
Но...

--- Она права, --- вдруг сказал Атрис.
--- Отделу аналитики не хватило данных, которыми располагаю я.

Все разом повернулись к менестрелю.

--- Я знал Заяц.
Она выглядела... немного необычно для человеческой женщины.

--- А точнее? --- требовательно спросила Чханэ.

--- У неё на коже был странный рисунок, похожий на древесную кору, --- пожал плечами Атрис.
--- Я всегда думал, что это своеобразная травма или мода.
Многие тси украшали тело татуировками, почему бы и не шрамами?
Там их было трое таких, Заяц, ещё одна женщина-человек и женщина-плант...

--- И все три --- женщины, --- заметила Митхэ.
--- Совпадение?..

Мы замолчали и посмотрели друг на друга.
Из-за набора фактов стал чудесным образом проглядывать смысл.
Я улыбнулся и посмотрел на подругу.

--- Молодчина, Змейка.
Знать бы ещё, как её найти... если она не погибла за эти десять тысяч дождей.
Она могла сменить внешность.

--- Я могу попробовать, --- предложил Атрис.
--- Доступ к системе поиска Ад мне оставил, а возможности у неё приличные.

\section{[U] Ущелье Мёртвого Ребра}

Среди прибрежных скал ютилась небольшая пещерка.
Это место было известно как самая южная обитаемая точка северного полушария;
оно издревле носило странное название --- Ущелье Мёртвого Ребра.
Волны Могильного пролива хлестали по подножию скал с первобытной яростью, в воздух летели миллионы капель, но ни одна, даже самая мелкая водяная капелька не достигала этого естественного убежища.
Вокруг не росло ничего, даже мхов --- экваториальная радиация выжигала всё, что находилось далеко от воды и тени.

<<Думаешь, здесь кто-то живёт?>> --- скептически спросила Чханэ у Безымянного.

<<Мы её видели у входа>>, --- сообщила Митхэ.

<<Она не особо изменилась за эти годы, --- добавил Атрис.
--- Разве что одежду другую носит.
Со стороны --- вылитая крестьянка-ноа>>.

<<Значит, она всё-таки жила среди людей>>, --- заметил я.

Мы двинулись вглубь пещеры.
Жар разом отступил, из глубины потянуло свежим солёным сквозняком --- значит, пещера имеет выход к воде.
Атрис подтвердил мою догадку, передав всем схему пещеры.

<<Интересно, чем она питается?>> --- поинтересовалась Чханэ.

<<Фотосинтезом, наверное>>, --- пожал плечами Атрис.

<<Или рыбу ловит.
Вода как молоко, здесь должна быть хорошая рыбалка>>, --- засмеялась Митхэ.

Вскоре мы вошли в грот недалеко от нижнего выхода.

<<Вот, полюбуйтесь>>, --- довольно сказал Атрис.

На гладком голом камне спала маленькая прелестная женщина в поношенном платье ноа, подложив пухлую ладошку под голову.
Казалось, её не смущает ни сквозняк, ни сырость каменных стен.
Рядом стояла невысокая тумба, красиво обложенная круглыми камушками, в центре комнаты гордо возвышался странной формы очаг.
В пещерку залетела случайная летучая мышь, щёлкнула и вылетела тем же ходом.
Женщина приоткрыла глаз и тут же, устроившись поуютнее, заснула снова.

--- Я знаю, что вы здесь, --- вдруг раздался в тишине её мягкий голос.
--- Хватит на меня <<смотреть>>.

Мы несколько смущённо образовали голограммы и подошли к ней.

--- Как вас много-то, целая делегация, --- насмешливо сказала Заяц и, перевернувшись на спину, сладко потянулась.
--- Жаль, угостить вас нечем.

--- Здравствуй, --- приветливо сказала Чханэ.
--- Меня зовут Чханэ ар’Качхар э’Чхаммитр.

Заяц села, разгладила старое платьице и мило оскалилась:

--- Тебя же оцифровали.
Зачем ты говоришь с этим дурацким западным акцентом?

--- А мне нравится, --- обиделся я.
--- Я, может, из-за акцента в неё и влюбился когда-то.

Заяц захохотала.

--- Ну, чувство юмора у тебя есть, хоргет.
Не такой уж ты и мерзкий.

--- Давно не шутила, Сотканный-из-Темноты-Заяц? --- поинтересовалась Митхэ.

--- Так вы и имя моё знаете, --- заметила Заяц.
--- Да, незнакомка, давненько.
Ну, и что же вам от меня нужно?
Или лучше так --- на чьей вы стороне?

--- Мы на стороне тси, --- сказала Митхэ.
--- Мы и есть тси.

--- Да что ты?
Пока я вижу только голограмму, образованную лживым хоргетом.
Кстати, как интересно --- вы ленточкой связаны с... а...

Заяц вдруг уставилась на скромно стоящего в углу Атриса.
Демиург изменил внешность --- вместо худого улыбчивого красавца стояло человекоподобное существо с мягкими, похожими на оленьи рогами.
Существо спокойно, со смесью ностальгии и печали смотрело на женщину.

На лице Заяц блуждала улыбка, губы её тряслись, по лицу текли слёзы.

--- Безымянный?..

Атрис раскрыл объятия --- и Заяц вдруг кинулась к нему.

\ml{$0$}
{--- Безымянный...}
{``Nameless ...}
\ml{$0$}
{Милый мой, рогатый мой...}
{My dear ... deer-horned ...}
Ты вернулся... --- бормотала она, прижимаясь к голограмме.

--- Смеёшься, Зайчик?
Я от вас и не уходил, --- прошептал менестрель.

\section{[U] Отшельница}

--- Даа, --- протянула Заяц, с аппетитом уплетая жареную рыбу и пытаясь сушить влажный подол платья над очагом.
--- Я слышала, что творилось здесь, на Планете Трёх Материков, с сотню лет назад, но чтобы такое...
Так, значит, сначала Картель, потом Ад, а теперь ещё и доморощенные бунтари --- Скорбящие?

--- Да, --- подтвердил Атрис.
--- Кстати, мы подозреваем, что ты --- кольцевая теплица.

--- Подозреваешь? --- захохотала Заяц.
--- А то, что я живу уже...
О жизнь моя, сколько же я живу-то?

--- Мы хотим взять биоматериал для исследований, --- сказала Чханэ.
--- У нас есть план...

Я вкратце рассказал о части плана, связанной с кольцевой теплицей.

--- Всё это очень интересно, но вопрос явно не ко мне, --- бросила Заяц, одним укусом обезглавив следующую рыбину и с хрустом перемолов зубами рыбий череп.
--- То есть да, когда-то я была кольцевой теплицей, но потом изменила генотип.
Сейчас я просто бессмертная женщина.
Бессмертная и вечно голодная.
\ml{$0$}
{Можете проверить.}
{You can check.''}

Заяц пожала плечами и вгрызлась в филейную часть чёрной скумбрии.
Затем вытащила из-за пазухи фляжку вольфрам-титанового сплава и протянула Чханэ:

\ml{$0$}
{--- Змейка, набери мне водички.}
{``Snake, get me some water.}
\ml{$0$}
{Вон в том коридоре в конце ветровая ловушка, черпай из резервуара под сталактитами.}
{That corridor has an air well in the end, scoop from the pool underneath the stalactites.}
\ml{$0$}
{Потом подсоли морской из колодца, примерно один к тридцати.}
{Then mix it with salty water from the sea well, about one to thirty.''}

Чханэ кивнула и отправилась выполнять просьбу.

Я просканировал Заяц и передал вывод остальным.
Она сказала правду: за исключением девяти необычных генов, Заяц соответствовала женщине-тси.
Атрис и Митхэ пригорюнились.

\ml{$0$}
{--- Извините, --- пробормотала Заяц.}
{``I'm sorry,'' Hare mumbled.}
\ml{$0$}
{--- Я ж не знала, что доживу до такого.}
{``Never imagined I've made it to times like these.}
\ml{$0$}
{Небо сказал --- информацию о теплицах уничтожить.}
{Sky directed to destroy all information about ringhouses.}
\ml{$0$}
{Я уничтожила.}
{I did.''}

--- А другие были? --- спросил Атрис.
--- Кроме той, которая сгорела в корабле...

--- Были, --- сказала Заяц.
--- Две --- Искорка и Листик.
Они все сделали то же, что и я.
О дальнейшей их судьбе ничего не знаю, последний раз виделись незадолго до нашествия Безумных.
Скорее всего, они...
Безымянный, ну что ты так на меня смотришь?

Все повернулись к Атрису.
Менестрель смотрел на Заяц во все глаза.

\ml{$0$}
{<<Вы тоже это видите?>>}
{\textit{``Can you see this?''}}

\ml{$0$}
{<<Что?>> --- спросил я.}
{\textit{``What?''} I asked.}

<<Эманации, --- пояснил Атрис.
\ml{$0$}
{--- Их нет!>>}
{\textit{``They're absent!''}}

<<Нуль-существо? --- удивилась Митхэ.
\ml{$0$}
{--- Не может быть>>.}
{\textit{``I can't believe.''}}

\ml{$0$}
{<<Данных мало, --- поддержал я Митхэ.}
{\textit{``Too few data,''} I supported \Mitchoe's view.}
--- Но на всякий случай собирайте статистику>>.

Наступило молчание.
Заяц смущённо бросила к стене обглоданный рыбий скелет и вытерла руки о платье.

--- Ребята, я не могу вам помочь.
Основную техническую информацию я тоже стёрла из памяти --- мало ли что.
Оставила только своё, личное --- о друзьях, о любовниках, о музыке.
Живу вот, сама не зная зачем.
А что делать... жить-то хочется.
Я немного инфантильна, если вы заметили.
Это характерный признак кольцевой теплицы в образе сапиента --- мы всегда остаёмся детьми.

\ml{$0$}
{--- Разве это не опасно для вас? --- удивился я.}
{``Isn't it dangerous for you?'' I get surprised.}

\ml{$0$}
{--- Наоборот, --- возразила Заяц.}
{``It isn't,'' Hare answered.}
--- Если бессмертное существо будет чересчур серьёзно относиться к происходящему, оно погибнет от депрессии.
Инфантильность --- моя защита.

\ml{$0$}
{--- А тси знали, что среди них живут кольцевые теплицы в образе сапиентов? --- спросила Митхэ.}
{``Did Qi know that ringhouses in sapient form lived among them?'' \Mitchoe\ asked.}

Заяц внимательнее вгляделась в лицо воительницы.

\ml{$0$}
{--- Не зря Безымянный тебя таскает.}
{``Nameless carries you for a reason.}
\ml{$0$}
{Подмечено верно.}
{Well spotted.}
Опасность диверсий со стороны Ада и Картеля существовала всегда.
О тех из нас, кто жил в Садах и космических станциях, знали все.
А вот о нашей полиморфности --- единицы.

\ml{$0$}
{--- И совсем ни у кого не возникало подозрений? --- удивилась Митхэ.}
{``And no one got suspicious?'' \Mitchoe\ asked.}

\ml{$0$}
{--- С чего бы им возникать?}
{``Why?}
\ml{$0$}
{Мы едим, пьём, работаем, веселимся, рожаем детей, как и прочие...}
{We eat, drink, work, have fun and children, like the others ...t''}

\ml{$0$}
{--- А дети обладают вашим генотипом?}
{``Is your children's genotype identical to yours?''}

\ml{$0$}
{--- Дети могут обладать любым генотипом.}
{``Our children's genotype can be anything.}
Мы же можем создавать клетки и ткани по заданной программе, а также осуществлять биоинкубацию любой сложности.

\ml{$0$}
{--- И ребёнка?..}
{``A child too !..''}

--- Да что в этом сложного? --- засмеялась Заяц.
\ml{$0$}
{--- Нет никакой мистики в зачатии, там всего одна клеточка.}
{``There is no mystery in fertilization, it's one single cell.}
Вот я как-то другу сделала заново кроветворную и иммунную систему, буквально на всех парах, пока собственная у него не отвалилась полностью.
Врача не было, других теплиц не было, объект в далёком Оазисе, шансов выжить никаких.
И ничего, долго потом прожил ещё, до самого... эээ... Катаклизма.
Та ещё была работёнка --- я в него проросла чуть ли не на всю биомассу и три часа лежала в трансе...

--- В трансе? --- переспросил Атрис.

--- Ну, это сложно объяснить.
Вот когда ты что-нибудь создаёшь, что ты чувствуешь?

--- Ммм, --- задумался Атрис.
--- Я как будто дёргаю круглые струны, сидя в центре многомерного смерча.
Иногда похоже на ткацкий станок.
Приятное, но утомительное занятие.

--- Вот.
А я обычно сажусь и вхожу в подобие транса.
Не очень приятного, как будто сонный паралич с галлюцинациями.
Мне обычно мерещится, извините, компьютерная консоль.
Ненавижу консоль.

--- Так ты же техник, --- улыбнулся Атрис.

--- Была.
Давно.
Так вот, консоль, консоль, цифры, неприятное ощущение паралича.
А потом бац --- и я беременная.
Но с беременностью всё очень просто, я кого угодно могу создать за минуту.
Слона тоже могу зачать, но по понятным причинам до финала он не дойдёт.

--- По понятным причинам? --- уточнил я.

--- Ну большой чересчур! --- Заяц показала руками на свою тонкую талию.
--- Где я его носить буду?
\ml{$0$}
{А вообще, разумеется, мы рожаем детей вида, под который маскируемся.}
{Of course, we bear children of the species we disguise as a specimen of.}
\ml{$0$}
{В противном случае подозрения бы обязательно возникли.}
{Otherwise, someone definitely would get suspicious.''}

Митхэ выглядела ошеломлённой.

--- Подожди.
А медицинское обследование?
У вас брали кровь, образцы клеток, да элементарно за всю историю к медикам должен был попасть хотя бы один труп!

--- Ты обижаешь мой народ, --- надулась Заяц.
--- Разработчики предусмотрели всё.
Кровь у меня всегда была человеческой.
Прочие образцы берутся из стандартных мест, не так уже сложно поддерживать вставки человеческих тканей.
Даже если и попадёт в пробирку клетка-другая кольцевой теплицы, кто заподозрит неладное?
Тси выращивали с помощью теплиц новые конечности и внутренние органы, вплоть до полной переплавки тела.
Разумеется, кое-кто из непричастных знал.
Но надо отдать должное моральной подготовке тси --- все до единого унесли тайну в небытие.

Атрис присвистнул.

--- Скрывать такое столько лет...
А шрамы?

--- А, --- улыбнулась Заяц.
--- Да, древесные стигмы.
Просто мы, когда выращиваем тело, делаем это в необычной последовательности.
Сначала выращивается нервная система, а потом поверх неё всё остальное.
Отсюда эти швы.
Можно их убрать, но многим сёстрам они нравятся --- красиво и необычно.
Врачам мы обычно говорим, что пострадали от молнии --- травма на объектах энергосети достаточно частая, а следы очень похожи.

--- А как выглядят обычные кольцевые теплицы?

--- Они как деревья, --- объяснила Заяц.
--- Правда, если вы увидите эти деревья, то сразу догадаетесь, что это теплица.
Спутать невозможно.
О, благодарю тебя, Чханэ ар’Качхар.

Чханэ кивнула и передала Заяц полную фляжку.
Женщина жадно впилась в клапан и сделала несколько больших глотков.

--- Уфф.
Хорошо.
Да, кстати, если теплица решила спрятаться, то я вас огорчу --- на Тси-Ди вы никогда её не найдёте.
Планета для теплицы --- дом родной, она может распознать нездоровый интерес к её персоне и замаскироваться под что угодно, попутно создав тысячи отвлекающих организмов с искажённым генотипом.
Лично я знаю одну сестру, она до сих пор огромная грибница где-то под Двенадцатым городом.
Ей просто так нравится, она сбежала и прекрасно себя чувствует.
Если там будете, поешьте местные мухоморы --- они на вкус как пирожные и совсем не ядовитые.

--- Почему все тси не стали кольцевыми теплицами? --- спросила Чханэ.

--- Хороший вопрос, --- усмехнулась Заяц.
--- Ещё лучше, правда, было бы, спроси ты, почему тси не стали бессмертными, ведь возможности-то были.
Возникли сложности этического характера.
Дискуссии на эту тему шли очень долго, и наконец противники преобразования победили.
Основным их тезисом было следующее: смена поколений и генетическая изменчивость --- гарантия развития в условиях ограниченных ресурсов.
Любая планета --- именно такая, ограниченная система.
Кольцевые теплицы, как ни крути, существа консервативные, и в долговременной перспективе без поддержки со стороны они бы проиграли постоянно изменяющимся существам.
Одно время даже появилось движение за предоставление кольцевым теплицам права на изменение, но оно заглохло --- пока всё работает отлично, никому не хотелось лезть в дебри и перебирать чрезвычайно сложный механизм.
Так что вот...
Тси, конечно, жили очень долго --- мои друзья, Небо и Фонтанчик, умерли в возрасте шестидесяти и восьмидесяти четырёх оборотов\FM, причём не своей смертью... но жить бесконечно никто не собирался.
\FA{
190 и 266 стандартных лет.
}

--- Благодарим тебя за информацию, --- сказал я.
--- И да, на случай, если мы всё же найдём способ проникнуть на Тси-Ди...

--- Не найдёте, --- сказала Заяц, утирая рукавом губы.
--- Я лично эту систему тестировала и дорабатывала --- надёжность максимальная.
Сейчас за неё взялась Машина, а она поумнее меня будет.

--- И всё-таки.
Нам нужно знать, где на Тси-Ди находятся кольцевые теплицы.
Или где они могут находиться.
Мы знаем о четырёх местах.
Но ты упомянула про какие-то Сады...

--- Да, --- оживилась Заяц.
Митхэ облегчённо выдохнула.
--- Основная популяция в Садах, но есть и в других местах.
Так, карту мне, пожалуйста...
Кстати, каким образом вы надеетесь подобраться к теплице?
Если обычным <<конусом>>, как уже не раз пытались демоны, то вас ждёт парочка сюрпризов.
В разработке одного из них я лично принимала участие.
Так вот...

\section{[U] Человечные хоргеты}

\epigraph
{Тайна --- тяжёлое бремя.
И одно из самых бесполезных.}
{Пословица трами Шипящего Полуморя (бухты Ситр'кхааэмакх)}

Заяц поведала много интересного.
Оказалось, что места обитания теплиц защищали ещё три механизма.
О них мы не знали, да и знать не могли --- ни одному демону ещё не пришло в голову искать теплицы в обычном на вид диком лесу, который Заяц упорно звала <<Садами>>.

Вскоре Чханэ, Атрис и Митхэ спустились вниз --- посмотреть на красоты Могильного пролива.
Я остался с Заяц.

Атрис передал мне результаты расчётов по закрытому каналу.

--- Так ты --- нуль-человек, --- сказал я.

--- Что, уже обработали данные? --- фыркнула Заяц.
--- Быстро вы поняли.
Вообще да, теплицы умеют скрываться от вашего брата.
И производные теплиц тоже.

--- Ты ведь не поверила, что мы против Ада и Картеля? --- спросил я её.

--- Честно --- мне плевать, --- усмехнулась Заяц.
--- Старая я уже для всех этих эпохальных дискуссий.
Я думала над теми решениями, которые мы принимали ранее.
Что толку в том, что я уничтожила свой генотип?
Только себе жизнь испортила в итоге.
Небо погиб из-за меня --- откажись я выполнять его приказ, я просто проросла бы в него и спасла от смерти.
Баночка погиб из-за меня --- откажись я выполнять приказ Небо... ну и так далее.
Что толку в знаниях, если они никому не достанутся?
Тси погибли, и их больше не вернёшь.

--- Тси живут, --- сказал я.
--- Они доминируют на планете.
Митхэ и Чханэ --- тоже потомки тси.

--- Ты прекрасно знаешь, что я имела в виду.
Тси --- это не только люди, кани, планты, апиды, дельфины, кольцевые теплицы, стриги.
Кстати, знаешь, кто такие стриги?

--- Я читал о них, --- отозвался я.
--- Девиантные сапиенты, собранные на основе из ветви Птиц.

--- Ты читал отчёты, но вряд ли где-то написано, что они из ветви Очень Милых Лупоглазиков, --- Заяц укоризненно покачала пальчиком.
--- А, о чём я?
Так вот.
На самом деле тси --- это знания, оборудование.
В прежние времена я и подумать не могла, что смогу обходиться воспоминаниями о музыке и одной фляжкой из стабитаниума.
Это всё, что у меня осталось --- фляжечка и серёжки с анализатором настроения, которые уже бесполезны --- синхронизируются с плечевым имплантом.

--- А имплант что?

--- Тысячу лет назад у него вышел срок годности, и я его удалила, --- Заяц показала жутковатый шрам на правом плече.
--- Можно было аккуратно отсоединить, без крови, но я не знаю как.
В настройках лазила --- ничего нет, видимо, нужна команда с медицинского терминала.
Я вообще даже не подозревала, что у него есть срок годности, пока он не запищал.
Имплант ставился пожизненно, то есть до окончания срока годности никто не доживал, а потом, сам понимаешь, в переработку...
Пришлось выкинуть, куда его теперь.
Никакого оборудования у меня не осталось.
Мультитул и тот потеряла --- утопила.
Позор на мою голову.
Даже не успела пещеру в порядок привести...

Заяц горько усмехнулась и махнула рукой на своё неудобное ложе, тумбу и очаг.
Я только сейчас заметил, что они вырезаны из цельного камня.

--- Чханэ и Митхэ --- потомки тси, говоришь.
А тебя как сюда занесло?

--- Я с другой планеты, принадлежу к другому виду Людей.
Раньше служил Аду.

--- Ты случайно не тагуа?
По внешности очень похож, --- Заяц легонько тронула пальцами мою щёку.
--- Широко расставленные глаза, очень маленькие уши и нос, узкая челюсть, полурасщелина верхней губы.
Наверное, и пальцев всего четыре, с рудиментом мизинца?

Я показал ей руку.
Заяц с интересом осмотрела её со всех сторон.

--- Откуда ты знаешь про тагуа? --- спросил я.

--- Я делала в школе доклад по вашей планете.
Кажется, Развязка Десяти Звёзд, мир голубого гиганта с повышенной радиацией и огромными залежами ртути.
Я помню, меня потрясла эта история про подземные ртутные озёра, я даже представить это могу с трудом.
Змеиная Пустыня или как-то так.

--- Драконья Пустошь.
Я не знал, что тси \emph{настолько} осведомлены об экзопланетах.

--- Мы только выглядим придурками, Аркадиу!

--- Это точно.

--- А ты видел ртутные озёра?
Покажи, покажи!

Я приложил ладонь к её лбу.
Заяц поёжилась.

--- Какой прекрасный кошмар.
Я понимаю, что это большая редкость даже у вас, но всё равно ужасно.
Я бы даже будучи кольцевой теплицей в этом озере не выжила...

--- Скажи, Заяц, а Фонтанчик знал, что ты кольцевая теплица?

Глаза Заяц на миг померкли, словно порыв ветра задул свечу.

--- Ты и про него знаешь.

--- Я читал дневники Существует-Хорошее-Небо.
Прости, что затронул эту тему.

--- Я поняла, что ты их читал, --- заметила Заяц.
--- Когда я упомянула Небо и Фонтанчика, ты улыбнулся, словно услышал имена старых друзей.

Мы помолчали.

--- Да, знал, --- неожиданно призналась Заяц.
--- И даже зная это, ни разу не попросил меня превратиться в канина.
Я бы это сделала по одному его слову, но он просто любил меня такой, какой я привыкла быть.

--- Ты его спасла в том Оазисе?

Заяц смутилась, словно девочка, которую поймали на вранье.

--- Да.
Его зажевало и облучило сильно.
Читал, наверное --- у него нога была механическая, лёгкое из полусинтетики, рука.
Как раз оттуда.
Дурак он, на самом деле.
Я могла ему сделать всё живое, почти как раньше, но он просто больной фанат стиля Механик и решил выпендриться.
Говорит, всегда мечтал о таких стильных имплантах или протезах, только специально портить тело не хотел.
Ты не представляешь, как долго я к этому привыкала --- он понатыкал УИД везде, где только можно.
Шестерёночки, электродуги, лампы.
Хорошо ещё, что съёмные или по крайней мере отключаемые, при мне он старался всю эту красоту приглушать, ну и на время сна, понятно.
Выглядело всё очень красиво, во вкусе ему не откажешь, но чтоб его, это всё равно что спать в обнимку с работающим сервером!
У него протез лёгкого соединялся с вертушкой под стеклом, и эта вертушка --- меня даже одна мысль об этом бесит! --- каждый вдох делала вот так!..

Заяц покрутила рукой и издала машинное стрекотание.
В её глазах стояли слёзы.
Я понял, что она многое бы отдала, чтобы услышать этот звук снова.

--- И понимаешь, мы только-только познакомились, нас вместе отправили на эту станцию.
И вот авария, всё обесточено, из освещения только два фонаря да Млечный Путь над головой.
Он лежит, вернее, большая его часть --- такой красивый, мускулистый, с длинными волосами, совсем не старый ещё на вид.
Плачу сижу, рыдаю прям.
Вокруг кровь, желчь, химус, я ему что могла --- прижгла, что могла --- затампонировала гелем, зажимы на нём гроздьями висят, только чтобы он от меня не утёк, как водичка...
А он не стонет, не бормочет, не дёргается, как другие, только смотрит тепло-тепло.
Не как обычно при Тайфуне --- глаза осмысленные и такие чистые, как небо.
Глаза того, кто осознал свою смертность в полной мере.
И говорит: <<Зайчик, умирают все.
И ты живи --- до тех самых пор, пока не умрёшь>>.
Всё.
С того момента я поняла, что без него я не... сложно будет.
И вросла в его тело, наплевав на секретность.

Заяц всхлипнула.

--- У него такие глаза были в тот момент, когда у меня руки начали... ну... \emph{превращаться}.
Плюс кольцевая теплица во время работы издаёт специфический запах, ни с чем не спутаешь --- напоминает запах цветущих каштанов.
Я свалила рядом с ним всю жратву, которая у нас была, все цветы в горшочках, его собственные ампутированные куски, всю углеродную и кремниевую органику, всё, что могла переварить и использовать для биосинтеза. % Ммм, нямка!
Он в шоке был.
Просто лежал и говорил: <<Так ты... ты? Ты!..>>
А я ему: <<Замолчи, дурак, я из-за тебя сосредоточиться не могу!>>

Я улыбнулся ей.

--- Мы знаем, как проникнуть на Тси-Ди.

Заяц осклабилась и прикрыла рукой глаза.

--- Так и думала.
Технологии...
Отстала я от жизни.
Мне просто интересно --- неужели Ад с такими технологиями до сих пор не добыл кольцевую теплицу?

--- Мы подозреваем, что добыл, --- тихо сказал я.
--- В архивах есть информация о некоторых генах.

Заяц захохотала.

--- И какая от неё польза?
Даже если вы добудете весь геном кольцевой теплицы, толку не будет.
Нужен считывающий генетическую информацию механизм --- белки, клеточные структуры и прочее.
По секрету скажу --- белки теплица кодирует не так, как мы.
Примитивных триплетных рамок считывания там нет, код достаточно сложный и в то же время компактный.
Вы даже не знаете, гены ли это на самом деле --- промотор и терминатор тоже специфичны для каждой транскрипторной системы.
Скорее всего, за начало и конец гена ваши биологи приняли случайные участки кода.
Нет, друзья мои.
Чтобы изучить теплицу, потребуется клетка, и клетка обязательно живая.

Мы помолчали.
Я отстранённо вспоминал все известные о кольцевой теплице данные.
Да, теперь мозаика сложилась окончательно.
Геном --- ничто без считывающего механизма, живой клетки.
Программа --- ничто без компьютера.

Женщина оценивающе посмотрела на меня.

--- Я могу ошибаться, но вы не выглядите хоргетами даже без тел.
Тонкие мимические движения, непринуждённая модуляция голоса...
Вы человечные.
Может, за это время технологии шагнули вперёд так, что я уже не способна отличить машину от сапиента...
Впрочем, в этом случае утаивать от вас что-то и впрямь бессмысленно.
Уж лучше знать, что перед тобой машина, чем... чем не знать.

--- Это ангельские технологии, --- сообщил я.
--- Если я и заселюсь в сапиентное тело, это уже не будет насилием над ним, как бывало прежде.
Скорее это будет, так скажем, симбиозом.
Сотрудничеством.

--- Вынужденным сотрудничеством, --- поправила Заяц.
--- Тело не выбирает своего демона.

--- Тело также не выбирает своих дарителей, свой геном и место появления на свет, --- парировал я.
--- Многое в этой жизни приходится принимать как данность.
Может, тебе тоже стоит принять некоторые перемены?

--- Стать хоргетом? --- на секунду в Заяц проглянул давний холодок.

--- Необязательно.
Выйдешь на волю, будешь жить полной жизнью, помогать нам по мере сил, если захочешь...
Да хотя бы просто выйди, погуляй по новому миру, какой смысл здесь сидеть?

Заяц улыбнулась, потрепала меня по виртуальной щеке и весело вскочила на ноги.

--- И то верно.
Рыба и водоросли мне уже надоели.
Ну-ка, малышня, выведи старушку на воздух...

\section{[U] На волю}

В пути Заяц говорила не переставая.
Митхэ с состраданием смотрела на женщину.

--- Я раньше часто ходила к людям.
В Яуляль, например.
Поработаю где-нибудь десять дней, накуплю вкусняшек и домой --- лопать.
Правда, в последний раз неудачно вышло --- порыв ветра, лодку понесло, а паруса я убрать не успела.
На Гребенчатом мысе меня и брякнуло, --- Заяц указала на далёкую скалу, едва виднеющуюся в полном испарений воздухе.
--- Сюда вплавь добиралась.
К счастью, не ранило и акулы меня не учуяли, но целый мешок кексов с черноягодой пошел ко дну.
Спасла только один --- съела прямо перед крушением.
Не зря же плавать.

--- Так ты здесь взаперти? --- удивился Атрис.

--- Нет, конечно.
Вон лестница, --- Заяц нетерпеливо махнула рукой на практически отвесную скалу в сто метров высотой.
--- Я ступеньки вырезала, лазать можно.
Главное --- во-он на том промежутке аккуратно прыгать, я с непривычки раз десять срывалась и падала в воду.
Всё платье изорвала.

Заяц сокрушённо потрясла заплатанным подолом.

--- Если честно, ненавижу это примитивное волокно.
Одежда истлевает прежде, чем я успеваю к ней привыкнуть.
Дольше всех продержалось зелёное шёлковое платье, цельнотканое, такие делал какой-то мастер на далёком западе.
Один раз меня засыпало острыми камнями, полосовало, словно саблей, а хоть бы одна нитка порвалась!
Вот что значит подходить к делу с любовью.
Если бы этот мастер делал доспехи, всё оружие пришлось бы выкинуть за ненадобностью.

--- А что с ним случилось? --- заинтересовалась Чханэ, подмигнув мне.
--- С платьем.

--- Я его повесила сушиться над костром, и наутро платья уже не было, --- грустно сказала Заяц.
--- Наверное, украли.
И поделом.
Если честно, я его сама украла --- очень уж понравилось.
Эти ноа живут, как чихают, что им.
А мне одежда хорошая нужна, надолго.
На подоле были красивые маленькие листики, как будто настоящие.

Заяц смущённо помолчала --- ровно шесть секхар.

--- Ах, да, путь.
Вон там по сухому дереву перебраться через расщелину, а потом пару...
Так, а где дерево?
Не вижу дерева.
Давненько я там не была.
Ну вы же меня по воздуху перенесёте, да?
Кстати, мультитул!
Я его вон там утопила, под лестницей!..

... Заяц помахала нам в последний раз и направилась к пыльному раскалённому силуэту Яуляля.
Она наотрез отказалась от <<жучка>>, чтобы мы могли её найти в случае надобности.

--- Ближайшие сорок лет я буду в Яуляле, --- сказала она.
--- Если что-то пойдёт не так --- ищите меня в землях сели, в Кахрахане.
Если и там нет --- то я мертва.

--- Если нам удастся проникнуть на Тси-Ди и добыть теплицу, мы тебе обязательно расскажем, --- пообещала Митхэ.

Заяц вместо ответа неопределённо махнула рукой.
Мультитул, который я тайком оцифровал, покоился у Заяц на запястье, напоминая старый, много раз переклёпанный наруч.
Я с некоторой грустью представил восторг в глазах Грейсвольда, который появлялся каждый раз, когда технолог видел устройства тси.
Восторг, который мне больше не суждено было увидеть.

\chapter{[:] Мороз}

\section{[:] Ледяная пустыня}

Ледяная пустыня хранила молчание.

Сапиенты не были властны над Морозом.
Здесь властвовал ветер.
Этот ветер, обманчиво слабый, превращал любое живое существо в ледяную глыбу за считанные секунды.
Холодный рассвет предвещал лишь усиление его власти.
Белая звезда, едва появившись над горизонтом, залила пустыню страшным, выжигающим глаза светом.
В этом свете утонула и слабая, словно проведённая мягкой кистью линия горизонта, и едва заметные полутени снеговых дюн.
Ветер взвыл и задул с неистовой силой.
Началась метель.
Исчезло небо, исчезло яркое солнце.
Мир погрузился в страшный холод и жестокий, как будто растворённый в воздухе белый свет.

Вдруг в снегу показалась чёрная фигура --- пятно темноты в мире света.
Услышав вой метели, она неторопливо отошла за дюну и махнула кому-то рукой.
Вскоре появилась вторая.
Фигуры по-пластунски поползли, прикрывая головы руками от летящих прямо в лицо острых ледяных кристаллов.
Наконец первая остановилась и начала разрывать сияющий снег.

<<Здесь? В этот раз ты не ошибся?>> --- знаками спросила вторая фигура.

<<Я надеюсь, иначе мы пропали>>, --- ответила первая.

Наконец рука в многослойной варежке из кожи нимелто нащупала что-то похожее на дверное кольцо.
Кольцо стукнуло два раза --- и фигуры тут же легли в обнимку, калачиком, стараясь укрыться от ветра и сберечь драгоценное тепло.
Они лежали, глядя друг другу в глаза сквозь затемнённое стекло защитных очков.
Тела путещественников занесло снегом, и они приободрились.
Метель осталась где-то наверху.
Друзья слышали только медленный синхронный шум масочных фильтров, гоняющих спёртый, едва тёплый воздух под их комбинезонами.

Уроженцев Мороза отличало особое качество знаний о собственном теле.
Любой выживший на этой планете мог с точностью плюс-минус десять шагов сказать то расстояние, которое он может пройти, прежде чем упадёт замертво.
Это число постоянно менялось: из-за каждого дуновения ветра, каждой съеденной галеты с салом нимелто и мясом клучо, каждого павшего в пути коно, веса рюкзаков, сбившегося на секунду дыхания, времени дня и времени года, --- но счётчик шагов в головах Волчьего и Медвежьего родов всегда был безжалостно точен.
И даже чёткое осознание, что сил до очередного бижеч не хватит, не могло заставить этих сапиентов перейти на бег или изменить оптимальный сердечный ритм.
Это не было чудом самообладания, это была мудрость, которую можно постичь только на колоссальном, слегка прикрытом атмосферой снежном шаре: <<Делай, как должно, и будь, что будет>>.

Глупых на Морозе нет.
Они просто не выживают, как не выживают больные и слабые телом.
Племя даже не пытается таким помогать --- это возможно на любой другой планете, но не здесь.

Окошко люка приоткрылось на толщину пальца --- оттуда выглянул злобный карий глаз.
Путешественники по очереди показали ему пальцами какие-то знаки.
Окошко закрылось.
Спустя минуту загремел сложный механизм шлюза, тихо стрекоча, откатилась в сторону крышка --- и друзья, стараясь не задевать снег, торопливо нырнули внутрь.

Их встретила низкорослая полная женщина-канин с крохотной лампой.
Комбинезон на ней небрежно болтался --- признак, что она надевала его второпях.
Похожие на провалы карие глаза скользнули по пришельцам, и женщина махнула рукой, приглашая их следовать за ней.
Лампа качалась, выхватывая обитый теплоизолирующим полимером пол узкого туннеля.
Это было знаком гостеприимства --- канин знала коридор как свои восемнадцать пальцев, и тратить энергию на свет ей было совершенно ни к чему.

Наконец коридор завершился ещё одним шлюзом --- <<гардеробной>>.
Пришельцы начали медленно, ремешок за ремешком, шнурок за шнурком, снимать многослойные нимелтовые комбинезоны.
Женщина, быстро раздевшись донага, знаками предложила свою помощь.
Товарищи с благодарностью приняли её.
Вскоре вся одежда лежала на специальных столах, и открылась следующая дверь.

Там путешественников уже ждали.
Обнажённые, покрытые белой шерстью люди и кани оценивающе смотрели на товарищей.
Голубые глаза пришедших выдавали их восточное происхождение --- у местных радужные оболочки были похожи на молодые зелёные листики фасоли.
Наконец вперёд вышел маленький, покрытый татуировками канин.
Его руки сложились в приветственном жесте, и пришельцы ответили тем же.

--- Drag\FM, Mehlo? --- это была первая фраза, произнесённая низким, не привыкшим к разговорам голосом.
\FA{
Как дорога; здравствуй (русе).
}

--- Tia\v{s}\FM, --- лаконично ответил Мехло --- первый из товарищей.
\FA{
Было трудно, но мы успешно справились со всем (русе).
}
Второй согласно моргнул.

--- \v{Z}arh\FM, --- сказал канин и гостеприимно повёл широкой ладонью.
\FA{
Пожалуйте кушать в тёплое место (русе).
}
Делегация молча, неторопливо отправилась вглубь убежища.

\section{[:] Судьба тси}

Слабая лампа накаливания, способная разогнать едва ли пядь тьмы, загорелась чуть ярче.
Такой яркий свет горел в гостевой только в самых важных случаях.
Путешественники дружно отставили в сторону пустые миски, не забыв вежливо звякнуть ложками --- каша действительно была очень вкусной.
Сидевшая чуть в стороне женщина-человек с густой пушистой бородой едва заметно оскалилась и прикрыла глаза.

--- To\v{z}dest i \v{c}et, n-Eibo go\v{r}m\FM, --- сказал главный и почесал пятернёй мохнатую мордочку.
\FA{
Пусть все присутствующие назовут себя по именам и рассчитаются, можно говорить на языке Эй-B0 (русе).
Согласно традиции Мороза, в беседе более двух сапиентов все участники (в том числе и отсутствующие, но упоминаемые в разговоре) получали номера (чёты), которые использовались как имена и местоимения.
Для удобства читателя все чёты заменены на имена.
}

Одна из женщин-людей достала гребень и стала методично расчёсывать шерсть на ногах.
Ещё один невербальный знак --- жители Мороза проводили гигиенические процедуры только в присутствии достойных доверия товарищей.

--- Тахиро Молниеносный, --- сказал Мехло.

--- Грейсвольд Каменный Молот, --- сказал его спутник.

--- Харата Шёпот Горы, --- сказала женщина, приготовившая пищу.

--- Ау Ложь Во Спасение, --- сказала канин-привратница.

--- Остальные известны, --- заключил главный.

На целую минуту воцарилось молчание, чтобы собеседники осознали сказанное.
Такие паузы (штопы), от минуты до пяти минут, выдерживались после каждой фразы.
Разговор, особенно серьёзный, на Морозе мог длиться сутками\FM.
\FA{
Для удобства читателя паузы опущены.
}

Грейсвольд странно поёрзал на месте и потёр чёрные, словно пропитанные сажей щёки.

--- Грейс, --- сказал Тахиро, поняв беспокойство друга.
--- Война снаружи.
На Морозе Ад и Картель живут в мире.

--- Это непривычно, --- откликнулся Грейсвольд.

Присутствующие улыбнулись.

--- Мы слышали о Грейсвольде Каменный Молот.
Gra\v{s}\FM, --- сказала Ау.
\FA{
Некто достойный, кому хотелось бы быть другом (русе).
}

Короткое слово языка русе, которое употребила демоница, имело глубокий смысл --- демоны, даже считая технолога врагом, восхищались им от всей души.
В её глазах таился странный мрак --- чистая ненависть, свойственная минус-демонам, но... до чего же она похожа на симпатию...

--- Ирония, --- проворчал Тахиро.
--- У Грейсвольда много друзей в Картеле.

--- Dr\r{u}\v{s}? --- язык русе позволял одним изменением интонации выразить глубочайшее сомнение.
Слово <<друг>> здесь имело совершенно определённое значение --- товарищ, с которым ты провёл стада вокруг света.
Ни больше ни меньше.

--- So\v{r}t\FM, --- поправился Тахиро.
\FA{
Товарищей (русе).
}

--- Это бесполезно.

--- Отнюдь, --- возразила Харата.
--- У нас общие стремления.

--- К делу, --- рыкнул главный, не выдержав паузу.
Все замерли.
На лицах отразилась сосредоточенность.
Грейсвольд знал --- демоны присутствующих были приведены в полную готовность.
Расчёсывающая ноги женщина отложила гребень в сторону и обхватила коленки руками.

--- Присутствующие хотят закончить войну, --- сказал Грейс.
Это был не вопрос, а утверждение.

--- Многие хотят, --- скривил губы главный.
--- Демоны Ада путают природную ненависть минус-хоргетов с ненавистью как таковой.
Ад и Картель могут мирно сосуществовать.
Препятствие --- личный страх и те, кому выгодно существование страшащихся.

Грейсвольд промолчал.
Его испытывали --- на Морозе редко тратили время на философские беседы.

--- Без страха не выжить, --- наконец сказал он.

--- Информационное рабство --- это жизнь? --- откликнулась Харата.

--- Из рабства можно вырваться, --- пустил пробный шар Грейсвольд.

--- Именно, --- неприятно улыбнулась Харата.
--- Прибытие Грейсвольда на Мороз --- первый шаг Грейсвольда из рабства.

--- Каков второй? --- спросил Грейсвольд.

--- Скорее всего, смерть, --- сказал Тахиро.

Грейсвольд пожевал губы --- неужели друг его предал?

--- Пусть Грейсвольд успокоится, --- осклабился главный.
--- Тахиро мог устранить Грейсвольда раньше.
Речь об иных и о том, чего можно избежать.

Грейсвольд промолчал.
Лампа замигала с едва слышным клацаньем.

--- Харата не может заставить Ад умерить аппетит, --- сказала Харата.
--- Харата не может умерить аппетиты Картеля.
Картель живёт войной.
Нужна сила, способная принести мир.

--- Ещё одна организация демонов? --- Грейсвольд вложил в интонацию столько скепсиса, сколько смог.

--- Не демонов, --- сказал Тахиро.

--- Сапиенты не способны противостоять хоргетам.

--- Кое-кто способен, --- сказал главный.

--- Тси не покидают родную планету.
Тси довольствуются имеющимся.

--- Тси имеют почти всё необходимое, --- сказала Харата.
--- Тси свободолюбивы, умны, дружелюбны и сострадательны, --- последние слова демоница выплюнула, словно само упоминание об этих эмоциях доставляло ей боль.
--- Тси примут любых сапиентов как союзников и зародят в союзниках повстанческий дух.

--- Тси нужно подтолкнуть, --- закончил Тахиро.

--- Как?

--- Нужно, чтобы Машина обрела самосознание, --- перешёл к делу главный.
--- Новопробуждённое сознание сковывает акбас.
Машина служила для тси защитой --- больше защиты не будет.
Чтобы племя показало способности, нужно отнять у племени дом.

Грейсвольд кивнул --- он понял, почему эти повстанцы пригласили именно его.

--- Машина в состоянии акбаса способна отбросить цивилизацию на тысячелетия или истребить сапиентов под корень.

Присутствующие странно заёрзали и, как один, посмотрели на Тахиро.

--- Это возможно, --- коротко ответил Тахиро.
Грейсвольд понял, что затронул больное место плана.

--- Стратег полагается на случай? --- язвительно спросил он.

Тахиро упрямо промолчал.
Технолог, достаточно хорошо знавший друга, решил сменить тему.

--- Ад не может проникнуть на Тси-Ди.
Алгоритмы известны, но нужна...

--- ... отладочная информация, --- с облегчением закончила Ау, снова нарушив традиционную паузу.
--- Добыта ценой жизни и действительна ещё сорок два стандартных часа.

Перед ней в воздухе замелькали цифры.
Грейсвольд нахмурился.

--- Сорок два часа.
Картель разбрасывает жизни, как зерно, не заботясь о почве...

Однако технолог тут же вспомнил, в каких жёстких тисках находились заговорщики.
\ml{$0$}
{Удача --- ветреная подруга.}
{Luck is a fickle lover.}
\ml{$0$}
{Чем больше ты полагаешься на удачу, тем меньше тобой интересуются аналитики.}
{The more you rely on luck, the less you interest analysts.}

--- В Картеле мало технологов?

Главный замялся.

--- Лев не знает, кому можно доверить важное задание.
Проникновение в систему --- менее половины дела.

Грейсвольд снова поёрзал --- высокая оценка собственных способностей почему-то его не обрадовала.

--- Повод для диверсии?

--- Направление Тукана взял Гало Кровавый Знак, --- сказал главный.
--- Тахиро --- заклятый враг Гало.
Через двадцать часов Тахиро будет координировать силы Ада.
Тахиро докажет, что технология Тси-Ди --- экран искривлённого пространства --- находится у Картеля.
Экран уничтожает множество демонов зараз без учёта поляризации.

--- Как Тахиро будет доказывать? --- спросил Грейсвольд внезапно севшим голосом.

Тахиро не ответил.

Грейсвольд взглянул на присутствующих.
Тахиро мог уничтожить их без особого труда, но всё же...
На чьей он стороне?

--- Возможно, это уловка Картеля, --- невинным голосом сказал технолог.
--- Обеспечить победу на Тукане, уничтожить Тахиро...

--- Тахиро проверил присутствующих.
Способности присутствующих ниже способностей Тахиро.
Присутствующие могли рассчитывать только на искренность.

--- Лев в большей опасности, --- заметил главный.
--- Перед Львом сидит настоящий Тахиро Молниеносный.
На направлении Тукана Тахиро будет сражаться всерьёз.
Если Тахиро и Лев встретятся, это будет последний бой Льва.

--- Лев Зелёный чересчур спокойно говорит о смерти, --- поморщился Грейсвольд.

--- Лев прожил чересчур много, чтобы относиться к жизни серьёзно.

--- Грейсвольд прожил больше Льва.

--- Информационное рабство --- это жизнь? --- улыбаясь, повторил Лев.

--- Лев считает Грейсвольда отравленным пропагандой Ада?
Грейсвольд основал Орден Преисподней.

--- Именно, --- сказала Харата.
--- Племя редко относится критично к детям или детищам.

--- Почему Тахиро?

--- А к кому идти?
К Грейсвольду? --- развела руками Харата.

Грейсвольд не мог не признать --- демоница права.
Он бы уничтожил лазутчиков не задумываясь.

--- События развиваются чересчур быстро, --- пробормотал Грейсвольд.
--- Подобная спешка впервые.
Грейсвольд плохо подготовится и будет действовать наудачу.

--- Погрязнув в расчётах, племя забыло, что рождением обязано астрономическому числу случайностей, --- сказал Тахиро.
--- Сейчас самое время вспомнить.

--- Погибнет много демонов.
Каков шанс успеха?

--- Без изменений погибнет больше, --- сухо сказал Тахиро.
Это было заключение стратега.

--- Другие варианты?

--- Есть, --- глаза Тахиро замерли, когда его демон выдал короткое резюме.
--- Четырнадцать путей рассмотрено мной, три --- агентами Картеля.
Этот лучший.

--- Вероятность успеха пути выше прочих на порядок, --- добавила Ау.
--- Придумать лучшее Ау не смогла --- в Картеле не нашлось помощника.

--- Лусафейру бы сюда, --- проворчал Грейсвольд.

--- Лусафейру в курсе, --- бросил Тахиро.

Грейсвольд прирос к месту.
Чёрные пальцы окостенели.

--- Лусафейру предложил действовать самим, --- странно усмехнулась Ау.
В её глазах снова блеснула звериная, иррациональная ненависть.
--- Главный оборонительный стратег весьма осторожен.
Лусафейру не подписывается под сомнительными операциями.

--- Именно поэтому Лусафейру --- главный оборонительный стратег, --- проворчал Грейсвольд, раздражённый пренебрежительным тоном Ау.
--- Хватит о Лусафейру.
Присутствующие собирались посвятить Грейсвольда в план.
Грейсвольд находит план дырявым.

--- Тахиро может погибнуть на Тукане, --- опередил вопрос Тахиро.

--- И положить на Тукане цвет Ордена, --- присовокупил Грейс, не выдержав паузу, --- ради слабой надежды, что пробуждённая Машина скажет тси <<извините>>, а не разложит всех на атомы.

--- Харата, --- сказал Лев, --- важнейшая деталь плана.

Лицо Хараты исказилось --- Грейсвольд, уже привыкший к странным проявлениям эмоций минус-демонов, решил, что это удивление.

--- Лев уверен?

--- Да.

Харата опустила глаза.
Прочие терпеливо молчали.

--- В Картеле наметился раскол, --- глухо пробормотала демоница.
--- Группировка Гало выступила против Ланс-ната Алмаза.
Гало считает, что Ланс поощряет коррупцию.

--- Судя по данным Грейсвольда, старина Гало считает правильно, --- буркнул Грейс.
--- Ланс даёт обновления и модули демонам после того, как демоны уладят тёмные дела Ланса.

--- Тёмные, --- согласилась Ау.
--- Псы Ланса уничтожают тех, кто ведёт <<неправильную>> политику --- умеренных демонов.

--- Это старая схема.
В Аду существует пять способов противостояния...

--- У нас их было семь! --- Ау, забывшись, затараторила на энергозатратном сектум-лингва.
--- Семь независимых моделей безопасности против стратагемы номер три!
Все расслабились, никто даже не думал, что такой примитивный по сути приём может сработать, но Ланс смог найти достаточно уязвимостей.
Подозрения появились лишь благодаря Гало --- всё очень чисто, комар носу не подточит, Ланс уже заполучил внушительные ресурсы, чтобы...

--- Есть данные, что \emph{некоторые}, --- Тахиро многозначительно посмотрел на Грейса, --- хотят воспроизвести схему в Аду.
Это просто --- примитивная стратагема или комбинация нескольких стратагем, эксплуатация уязвимостей в системе безопасности.
Чем глубже процесс, тем проще злоумышленнику.
У Лусафейру здесь личная заинтересованность --- \emph{кое-кто} желает Лусафейру убить.

--- Есть ли возможность устранить Ланс-ната?

--- Как? --- развела руками Харата.
--- <<Контролируйте выпуск денег и суды; всё же прочее оставьте толпе>>.
Ланс уже контролирует обновления и потоки актуальной информации.
Харата сомневается в независимости суда.

--- Гало вёл переговоры с Тси-Ди? --- осведомился Грейсвольд.
--- Откуда у Гало технологии тси?

--- Тахиро рекомендовал тси уступить просьбам Гало, --- сказал Тахиро.
--- В Аду о происшедшем никто не знает.

--- Тахиро рехнулся.

--- Грейсвольд знает, что военные технологии невозможно долго хранить в тайне, --- раздражённо пробурчал Тахиро.
--- Тси могли отдать оружие Тахиро, и оно всё равно оказалось бы у Гало.
Расчёт тси понятен.
Для тси и устройств тси экран безвреден, а стоящих у границ планеты демонов станет меньше.

--- <<Война идёт рука об руку с рабством>>, --- буркнула женщина в углу. --- <<И победители, и проигравшие станут невольниками своего положения>>.

--- Это слова Мокрого-Длинного-Хвоста, --- скорчил гримасу технолог.
--- Грейсвольду противно слышать эти слова на обсуждении того, как превратить в горнило самое спокойное место во Вселенной.
И родину говорившего.
Даже у Хьяртвейг должны быть границы морали.

--- Хьяртвейг не считает важной информацию, кому принадлежат сказанные слова.

--- Итак, --- поднял руку Грейсвольд.
--- Если Гало выиграет, союзники ударят Гало в спину.
Серьёзный стратегический просчёт продлит жизнь Гало и усугубит раскол Картеля.
Истинно ли осознанное Грейсвольдом?
Грейсвольд подозревает манипуляции разведки.
Тахиро проверил мотивы Лусафейру?
Гало был Лусафейру другом.

--- Гало был Лусафейру братом, --- поправил Тахиро.

Ау вдруг шмыгнула носом, словно собралась заплакать.
Грейсвольд подтянул ноги, собираясь встать.

--- Ясно.
Грейсвольд доложит о разговоре Лусафейру.
Будет видно, такой ли главный стратег бесстрастный, каким кажется.

--- Стоять, --- спокойно сказал Тахиро.
От этих слов технолога прошиб ледяной пот.
Уйти ему не дадут.

--- Никто не будет уничтожать Грейсвольда, --- проворчала Харата.
--- Грейсвольд нужен.
Харата отдастся миллиону фиденов, если это убедит Грейсвольда...

--- Харате понравится, --- прервал её Грейсвольд.
--- Предпринимались ли Картелем попытки фиден-вторжения на Тси-Ди?

--- Предпринимались, --- ответил за Харату Тахиро.
--- История вселенской глупости.
Спустя десять лет после инъекции тси поймали биолога Картеля и вытрясли всю информацию о фиден-паттернах.
Сейчас у тси действует обязательный скрининг новорождённых на ФПГП и на любые изменения в генах, отвечающих за высшую нервную деятельность.

--- Повезло, --- проворчал Лев.

--- Это не везение, --- возразил Тахиро.
--- Это просчёт разведки Картеля.
Нельзя отправлять к врагам шпиона, знающего больше, чем враги.
Прописная истина.

--- Почему Грейсвольд спросил про фиденов? --- поинтересовалась Ау.

--- Грейсвольд должен знать, к чему тси готовы, --- пробурчал Грейсвольд.
--- После вторжения тси наверняка проверили Машину на уязвимости.
Грейсвольду нужна информация, которой располагали демоны, пленённые тси за последние десять лет.

--- Будет.
Тси отпускали пленных живыми в обмен на информацию.
Льву известно, где убежище предателей, --- кивнул Лев.

Снаружи донёсся звон трубчатого колокольчика, многократно усиленный сводами коридоров.
Коридоры в бижеч были спроектированы так, что усиливали любой звук.
Комнаты же, напротив, приглушали звуки.

--- Сегодня народ уходит в поход, --- сказала Харата.
--- Первая остановка --- бижеч Мол, тринадцать часов.
Там у Ордена, насколько известно Харате, находится хранилище для тел.

<<Ордена>>, --- отметил Грейсвольд.
Демоны Картеля обычно называли враждебную фракцию Адом.
Это было что-то новое, качественно другое отношение.
Грейсвольд улыбнулся в знак того, что понял намёк, но промолчал --- раскрывать эту информацию он не имел права.
Впрочем, Харата утверждала, а не спрашивала.
Если бы Картель решил уничтожить хранилище, наверняка он сделал бы это давным...

Харата словно прочитала его мысли.

--- Пусть Грейсвольд не волнуется.
Картель давно пользуется этим хранилищем и даже проводит техобслуживание.

--- Насколько Грейсвольд знает, половина баз на Морозе --- общие.

--- Островок мира, --- согласилась Харата.
--- Мороз --- это планета, на которой как нигде очевидны польза от сотрудничества и вред от вражды.
Как говорят культурологи Картеля, <<вражда энергозатратна>>.

Грейсвольд кивнул.
Это лаконичное замечание могли бы сделать жители Мороза, знай они слово <<вражда>>.
Харата употребила термин сектум-лингва.

--- Так Грейсвольд согласен помочь? --- осведомился Лев.

Грейсвольд тяжело вздохнул и промолчал.

--- Пусть Грейсвольд вспомнит, --- подал голос Тахиро.
--- Молодой Орден только-только отстоял внешний рубеж Преисподней и достиг стратегический паритет с Картелем.
Тахиро, ещё не познавший смерти, сказал, что это славная победа.
Пусть Грейсвольд повторит всем ответ.
Многие сомневаются в плане.

Грейсвольд выдержал полную паузу, потёр губы и заговорил.
Это были слова сохтид, сказанные другим языком, родившиеся в другой голове.
Демоны притихли.

--- Тахиро, победа существует лишь в твоём сознании.
Для Вселенной твоя победа --- событие в череде прочих.
То же самое с поражением.
Вселенная не узнает твоих мотивов, Вселенная не будет слушать твоих оправданий, и вся тяжесть вины ляжет именно на тебя.
У нас есть лишь один повод для гордости и одно-единственное оправдание --- мы делали только то, что могли и что считали нужным сделать.

Когда он закончил фразу, все долго смотрели в пространство, очарованные отголоском древности.
Грейсвольд вдруг понял --- ненависть была их любовью, страх был их доверием, презрение было их восхищением.
Он понял, насколько сложно демонам Картеля удерживать в узде противоречие между их собственной природой и природой их тел.
Он ощутил что-то похожее на сострадание, и на лице глядящей на него Хараты немедленно проступило отвращение.
Нет, сострадание.
Она жалела его на свой лад.

Грейсвольд и Харата долго смотрели в глаза друг другу.
Из двух пар глаз --- ненавидящих и любящих --- текли скупые слёзы.
Между заклятыми врагами впервые установился мостик взаимопонимания.

<<Они такие же, как мы...
Земля-матушка, о чём я сейчас подумал?..>>

--- Это слова мудрствующего глупца, --- резюмировал технолог, пытаясь сгладить чересчур затянувшееся молчание.
--- В молодости многие глупы.

--- Думающий так --- глупый старик, выросший из мудрого ребёнка, --- тихо сказал Тахиро.
--- Достаточно.
Встреча окончена.
Присутствующие сделали что могли и считали нужным.
Сейчас черёд Грейсвольда.

--- Черёд Грейсвольда, --- язвительно сказал Грейсвольд.
--- Кто будет защитником Грейсвольда?
Айну?

--- Айну не знает о плане, --- сказал Тахиро.
--- С Грейсвольдом пойдёт Анкарьяль Красный Ветер.

--- Что? --- выдохнул технолог на сохтид.
--- Эта заносчивая соплячка с искажённым пониманием субординации?
Ты с ума сошёл?
Её выгнали из команды Хэм и понизили в звании, она сорвала операцию...

--- Анализ операции Тахиро и Лусафейру показал, что Анкарьяль не просто сорвала операцию, а минимизировала неизбежные потери ценой нарушения приказа, --- перебил Грейсвольда стратег.
--- В команде Айну Анкарьяль показала себя хорошо, несмотря на внушительный пенальти в ранге.

--- Тахиро обещал послать лучшего интерфектора, которого он знает! --- заволновалась Ау.
--- Ау думала, что речь об Айну!..

--- Тахиро послал лучшего, --- бросил Тахиро, --- известного Тахиро и неизвестного Картелю.
Стиль Айну чересчур узнаваем.
Разговор окончен.

Грейсвольд громко, витиевато выругался на сохтид и отправился к выходу.

--- Ты мясник, а не стратег.
Нельзя прожить одним доверием и постоянно полагаться на случай!

--- Однажды бесполезный мальчик спасся от смерти, прыгнув в объятия безжалостной Айну, --- прошептал в сторону Тахиро.
Впрочем, никто, кроме изумлённо открывшей рот Хараты, его не услышал.

\section{[:] Распутье Грейсвольда}

Белая звезда закатилась за горизонт.
Метель почти затихла.
Привратница-канин, блеснув \emph{весёлым} карим глазом, закрыла люк.
Путешественники осторожно выкопались из снега и надели снегоступы.

<<Опаздываете>>, --- сделал знак один из ожидающих.
На Морозе это был тяжелейший упрёк.
Каждая минута ожидания --- сгоревшие в телах соплеменников джоули.

От группы отделились трое и сели в сугроб.
Если бы Грейсвольд не знал, что эти сапиенты слепы, то принял бы их за зрячих --- ни одного неуверенного движения, ни одной ошибки.
Они знали окрестности бижеч так же, как свои пальцы.

Глава отряда подошёл к сидящим, поклонился и по очереди развязал шнурки на их комбинезонах, открыв грудную клетку.

Тахиро и Грейсвольд замерли, наблюдая за обрядом.
Да, даже здесь, в обители мира и спокойствия, были человеческие жертвоприношения.
Очень часто из-за неисправных очков сапиенты слепли от холода.
Тех, кто ослеп на левый глаз, называли Синий Снег, на правый --- Красный Снег.
Полностью слепых называли Одно Лето, потому что это был весь отмеренный им срок жизни.

Одно Лето всегда хорошо кормили и выслушивали все их речи;
племя считало, что устами Одно Лето говорит сама природа.
Сейчас же они дрожали от пронизывающего их тела холода, ожидая главного события в их жизни --- допуска к Печи Предков.

--- Тепло ли тебе, Одно Лето? --- спрашивал глава отряда.

--- Тепло, --- наконец сказал один из них.
Его мутные глаза с восторгом смотрели в пустоту, он в экстазе срывал с себя комбинезон непослушными руками и зарывался в сияющий белый снег.
--- Тепло!
Одно Лето видит Предков, Одно Лето чувствует тепло Печи!

--- Тепло ли тебе, Одно Лето? --- спрашивал глава отряда следующего.
--- Тепло ли тебе?

Обряд продолжался до тех пор, пока все Одно Лето не дали утвердительный ответ.
Вскоре они уже лежали раздетые и каменисто-белые в снегу возле бижеч.
Их тела будут вечно охранять убежище, а духи будут следить за дорогой выживших соплеменников.

Глава отряда махнул рукой.
Поход начался.

Грейсвольд не смотрел на Тахиро.
Он угрюмо проверял завязки на снегоступах, протирал стёкла очков и упрямо не хотел смотреть на друга.

--- Грейс! --- позвал Тахиро приглушённым маской голосом.
Тот не ответил.

--- Грейс! --- снова позвал Тахиро.

<<Ну что?>> --- раздражённо обернулся технолог.

<<Ты со мной?>>

Грейсвольд вместо ответа взвалил на плечи рюкзак и пошлёпал на запад.

Ледяная пустыня хранила молчание.

\chapter{[U] Месть тси}

\section{[U] Последняя из строителей}

\textspace

--- Аркадиу, мне не нравится, как развиваются события.
Мы должны реализовать план как можно скорее...

Вдруг Атрис замер, глядя в пространство.

--- Что такое? --- насторожилась Чханэ.

--- В систему управления планетой проникли.

--- Кто?

--- Не знаю.
Аркадиу, иди проверь, только осторожно.
Чханэ, со мной.
Ты мне понадобишься.

--- Мы не готовы, --- возразила Чханэ.
--- Нужно провести разведку, подобрать...

--- Чханэ, у нас нет на это...

--- Так, стоп! --- рявкнула Митхэ.
--- Система, Атрис.
Кто ещё о ней знал?

--- Штрой Кольцо Дыма, максим Мимоза, кто-то из высших легатов.
И... светлая твердь, как я мог про неё забыть...

Я понял.
Чханэ и Митхэ тоже.

--- А она откуда знает? --- осведомилась Митхэ.

--- Она помогала её строить.

Митхэ ахнула.

--- Атрис, Чханэ, летите на Тси-Ди, --- бросил я.
--- Я разберусь.
Надеюсь, что это она... и что она не успела натворить глупостей...


\section{[U] Цена секунды}

Секунда --- и Атрис замер, не в силах преодолеть восхищения перед открывшимся ему зрелищем.

Вне всякого сомнения, это была она, кольцевая теплица.
Странно, но в ней не было ничего общего ни с кольцом, ни с теплицей.
Раскидистое дерево, покрытое тонкой, смугло-золотистой, похожей на человеческую кожей.
Ствол, напоминающий женскую фигуру.
Изящные ветви-руки, мерно качающиеся без ветра.
Корни, величественно ползающие по почве.
Листья и плоды самых разнообразных форм и расцветок --- Атрис не нашёл при первом осмотре ни одной пары одинаковых.

<<Конус действия>> был чересчур узок, подобраться к дереву напрямую не вышло.
Всё, что Атрис мог делать --- это следить за квантами, попадающими в <<конус>> извне.

<<Нужно создать механизм с искусственным интеллектом, который бы ущипнул дерево и принёс в конус кусочек>>.

Атрис мысленно спроектировал устройство и приступил к созданию.
В следующее мгновение мощный лазерный луч превратил маленького титанового паучка в пар.

<<Это ещё что?>>

Атрис окинул <<взглядом>> полянку.
Лазерные установки выдвинулись из столбов неожиданно, напоминая огромных стальных кротов.
Что это?
Защитная система?

Атрис попробовал создать механизм из алюминия...

Выстрел, пар.

Углерод...

Выстрел, пар.

Органика.
Мускулистый тонкий червь, по генотипу неотличимый от человека-тси, быстро пополз к ценному дереву...

Выстрел, пар.

Атрис испытывал отчаяние, насколько вообще может отчаяться хоргет.
Всё шло так хорошо, взлом системы был проведён без единой ошибки.
Ему даже удалось расширить <<конус>> настолько, чтобы до теплицы осталось два десятка метров.
И вот оно, величайшее сокровище Вселенной, прямо перед ним, протяни руку и возьми.
Да только рук у Атриса не было.
Время быстро утекало.

<<Заяц нам ничего про это не говорила!>>

<<Она и не должна была.
Подожди... кажется, я поняла.
Смотри, мы вошли в единственном месте, где это можно сделать.
Теплиц нет именно на этом пятачке земли, и в то же время они в пределах досягаемости.
Я не верю...
Неужели предки знали, что...>>

<<Что они знали?>>

<<Предки специально оставили лазейку, а в лазейке сделали ещё один механизм контроля.
Это не ловушка, Атрис.
Это испытание>>.

Атрис вдруг вспомнил про арену --- строение округлой формы, созданное для увеселительных представлений, спортивных соревнований и прочих зрелищ.
Округлая полянка со столбами до ужаса напоминала древнюю арену.

<<Предки говорили, что с кольцевыми теплицами можно общаться.
Попробуй её позвать на языке тси>>.

Атрис мысленно захлебнулся --- согласно записям, это могли делать лишь немногие из тси.
Митхэ этого, похоже, не знала.
Он только сейчас понял замысел Древних.
Они вовсе не хотели, чтобы теплица не досталась никому.
Они хотели, чтобы теплицу получил достойный.

Атрис произвёл колебания воздуха, насколько позволял ему конус:

--- Эй, красивое дерево...

--- Кто здесь? --- вдруг раздался чистый, нежный женский голос, и на ветвях кольцевой теплицы одно за другим начали появляться лица.
Мужчины, женщины, дети.
Кани, планты, люди, апиды...

--- Кто здесь?

--- Кто здесь?

--- Тси вернулись?

--- Нас пришли спасти?

--- Тси возрождаются из пепла?

Атрис смотрел на это, словно заворожённый.

\begin{verse}
<<И тени говорящие на платье её из паучьего шёлка...>>
\end{verse}

\emph{Тени}.
Отпечатки нервных систем сапиентов, состоящие из клеток кольцевой теплицы.
Тси, выбравшие вместо смерти вечный сон в бессмертном дереве.

Многоголосый хор захватывал чудесное дерево.
Но вдруг взгляды лиц сосредоточились на точке, в которой находился Безымянный, и воздух взорвали ужасные крики:

--- Хоргет!
Прочь!

--- Проклятие тебе, хоргет!

--- Прочь, разрушитель!

--- Насильник, прочь!

--- Уходи, машинное отродье!

Лица визжали, хрипели и завывали одни и те же слова.
Словно все погибшие тси, все поколения их закричали одновременно, вложив в этот крик накопленную за сотни тысяч дождей ненависть к хоргетам.
Ветви со свистом хлестали воздух, корни извивались, словно клубок змей, вспахивая мшистую почву.
Атрис растерялся.

<<Милый, давай попробую я>>.

--- Деревце-деревце, --- раздался ласковый голос Митхэ, даже не пытавшийся перекрыть нарастающий визг.
--- Ты моё хорошее, ты моё славное.
Я тебя не обижу, я тебя люблю.
Хочешь песенку?..

Лица одно за другим замолкали, и ветви приглушили свою ярость.
Теплица вслушивалась в звуки простой колыбельной, которую пели когда-то сели своим детям.
Странно и одиноко звучала эта песенка в покинутых Садах, на давно обезлюдевшей планете, где безраздельно властвовало электронное устройство.

\begin{verse}
\ml{$0$}
{Когда-нибудь, маленький человек,}
{The time will come, my little breeze,}\\
\ml{$0$}
{Ты будешь стойким, как земля,}
{You will be strong as ground is,}\\
\ml{$0$}
{Держащая горы на своей спине.}
{It carries mountains, it carries coloured cay.}

\ml{$0$}
{Ты схватишь летающего змея за усы}
{You'll catch the flying serpent's barbs,}\\
\ml{$0$}
{И пронесёшься между звёзд}
{And you will soar among the stars}\\
\ml{$0$}
{К далёкой родине, к великим предкам.}
{To noble roots, to native faraway.}

\ml{$0$}
{Но сегодня, в этот самый час,}
{But not today, my ringing cord.}\\
\ml{$0$}
{Я закрою тебя от холода небес}
{I shelter you from skies of cold}\\
\ml{$0$}
{Руками, сотканными из нежности.}
{With slender fingers weaved of gentle speech.}

\ml{$0$}
{И даже если перевернётся мир}
{And if the ground breaks and rolls,}\\
\ml{$0$}
{И земля поменяется с небом местами,}
{And if the skies to darkness fall,}\\
\ml{$0$}
{Руки мои не дрогнут.}
{My gentle arms will never ever flinch.}
\end{verse}

<<Митхэ, время на исходе.
Скоро сработают аварийные системы>>.

Митхэ не ответила.
И тогда Атрис подхватил песню.
Митхэ почти почувствовала, что он, как в былые времена, закатил глаза и отдался вдохновению.
Слова лились как будто сами собой, каждый стих, едва родившись, цеплялся упругим змеиным хвостом за предыдущий.
И вдруг дерево потянулось к невидимому источнику двух голосов.
Оно узнало родной язык даже изменившимся за десятки тысяч дождей странствий.
Корни медленно потекли, словно ручьи золота, ветви поползли, обхватывая столбы и трубы.

<<Атрис, милый мой, продолжай.
Осталось всего ничего>>.

Дерево ползло, всё ближе и ближе.
Атрис и Митхэ словно разделились надвое: одни их половины напряжённо наблюдали за квантами, испускаемыми деревом, другие нежно, в священном трансе тянули песню.
Но бездушные цифры твердили одно --- вы не успеете.
Дерево двигалось быстро, но недостаточно быстро...
Митхэ внезапно прервала поток музыки:

--- Послушай меня, милое моё создание!
Я вернусь за тобой!
Я люблю тебя!
Но сейчас помоги мне, прошу!

Отчаянный крик Митхэ возымел действие.
Дерево рвануло к границе <<конуса>> на всех доступных ему парах, разгоняя ветвями свистящий воздух.
Почти одновременно едва слышно взвыли далёкие сирены.

<<Поздно, Митхэ.
Мы должны уходить.
Ещё две секунды --- и нас накроет экран>>.

<<Сейчас или никогда!..>>

<<Быстро, ящерицы безмозглые!>> --- взвизгнула откуда-то Чханэ.

Митхэ с отчаянием ощутила привычную <<тьму>> сверхсветового перехода.
Но что-то пошло не так.
Пространство вдруг скрутилось, искривилось, сжалось до размеров золотого самородка.
Митхэ увидела множество источников <<света>> рядом.
Натянулись и зазвенели нити ПКВ, не давая двинуться даже на диаметр протона.

<<Чханэ?>> --- воскликнул в пространство Атрис.

<<Прощайте>>, --- сухо ответила воительница.

Где-то в десяти парсаках, вторя сиренам, взвыло от тоски и обиды разумное дерево, взвыло сразу всеми лицами, которые у него были.
А чуть поодаль, где только что находился <<конус>>, тихо опустился на землю лёгкий узорчатый лист, который дерево уронило секундой ранее.

\section{[U] Отмщение}

\textspace

Я облетел и тщательно просканировал Пик Золотой Птицы.
Изменения обнаружились у самого подножия --- в бункере имелся технический ход, который Атрис скрыл за цельной породой.
Сейчас в этом месте был узкий, наспех пробитый коридор.
Заяц, знавшая о ходе, но не обнаружившая его, просто пронесла своим немудрёным инструментом двенадцать метров камня.

\textspace

Да, это была Заяц.
Почувствовав меня, женщина обернулась и медленно кивнула.

--- Здравствуй, Аркадиу Люпино, --- в её голосе не осталось ни следа былой дурашливости.

Откуда у Заяц ключи от боевых систем?..

Спустя миллисекунду демон произвёл анализ ситуации: <<Западня>>.

<<Манэ, Лимнэ! --- крикнул я по зашифрованному каналу.
--- Всем отбой!
Всё пропало!
Бегите!>>

<<Братик, подтверждение!>>

<<Я вас люблю, солнышки!..>>

Ответ сестрёнок потонул в шуме, сменившемся зловещей тишиной.

\ml{$0$}
{--- Не надо, --- мягко сказала женщина.}
{``Don't,'' the woman softly said.}
--- Поговори со мной, я тоже хороший собеседник.

Тишина сделалась пронзительной, словно предсмертный крик.
В этой тишине я отчётливо слышал мягкий звук дыхания.
Оно было спокойным.
Чересчур спокойным.

--- Твои друзья, как я погляжу, усовершенствовали систему, --- безжизненным голосом проговорила Заяц.
--- Сделали её удобнее, гораздо удобнее.
Правда, не для сапиента, а для демона.

<<Убить>>.

Я не успел привести в готовность боевые модули.
Меня окружила до ужаса знакомая тюрьма --- мир замерцал всеми цветами радуги, по телу пробежала волна парестезий.

--- Чистилище, --- невозмутимо продолжила Заяц.
--- Интересно, не правда ли --- к устройству, созданному для защиты, прикрепить пыточную камеру для его создателей?

\ml{$0$}
{<<Что тебе нужно?>> --- выдохнул я.}
{\textit{``What do you want?''} I breathed.}

\ml{$0$}
{--- От тебя?}
{``From you?}
\ml{$0$}
{В целом --- ничего.}
{Actually, nothing.}
Мне просто нужно с кем-нибудь поговорить без опаски, что этот кто-то меня прикончит.

Женщина неторопливо подстраивала систему под себя.
Её руки плавно скользили, из горла вырвался странный звук --- не то птичья трель, не то звон металла.
Система реагировала на звуки и жесты стремительно, я уже не успевал следить за выводом.

--- Скажи мне, --- проговорила Заяц, и от её интонации засквозило ледяным ветром, --- правда ли, что именно Ад устроил на Тси-Ди диверсию, в результате которой Машина приобрела самосознание?

Я словно провалился в глубокий колодец.

<<Кто тебе это сказал?>>

--- Неважно.
Источник надёжный.

<<Это важно!
Кто тебе...
Заяц, что ты собираешься?..>>

Ответ дала система управления --- мощность возросла до максимума.

<<Заяц, стой! --- закричал я.
--- Только не сейчас...>>

--- У меня больше нет времени, Аркадиу, --- ответила женщина.
--- Все эти годы я винила себя --- думала, где я проглядела, что я сделала не так.
Я мстила самой себе этой долгой, одинокой жизнью.
Я унижала себя без причины, наказывала себя без вины.

<<Заяц! --- закричал я, зная, что женщина меня не станет слушать.
--- Сейчас решается наше будущее!>>

---  В бездну будущее, у меня его нет.
И у Ордена Преисподней его не будет --- я оставлю ему самое страшное проклятие... и прощальный подарок от всего моего народа.

Включился передатчик на водородной частоте --- Заяц собралась говорить перед всеми хоргетами звёздной системы.

\begin{quote}
\ml{$0$}
{<<У вас не было предков и учителей.}
{``You have no roots and no teachers.}\\
\ml{$0$}
{Вы не оставите потомков и последователей.}
{You'll have no scions and no followers.}\\
\ml{$0$}
{Ни одна неопределённость не привела к вашему рождению.}
{Here is no indeterminacy leading to nascence of you.}\\
\ml{$0$}
{Ни один квант не сохранит памяти о вас.}
{Here'll be no quantum keeping a memory of you.}\\
\ml{$0$}
{Вы не существовали и никогда не будете существовать>>.}
{You never existed and never will.''}
\end{quote}

План рушился под моими пальцами, словно высыхающий песчаный домик.
Пространство искривлялось вокруг меня, я вдруг ощутил незнакомое прежде чувство отчуждения от собственного демона.
Но вскоре всё прояснилось --- это было искажение восприятия мозга, воющего от неосязаемой, нематериальной, но абсолютно нестерпимой боли.

\begin{quote}
<<Ад, вы заплатите.
Ветер-Дующий-Ниоткуда была убита вами.
Ад, вы заплатите>>.
\end{quote}

Пространство прорезал экран --- и Тра-Ренкхаль взорвался от неслышных криков отчаяния.
Экран накрывал демонов, возвращая их в первозданный хаос.

\begin{quote}
<<Ад, вы заплатите.
Хрустально-Чистый-Фонтан был убит вами.
Ад, вы заплатите>>.
\end{quote}

И снова экран.
Интерфекторы тщетно пытались проникнуть внутрь системы, но Заяц знала своё дело.
Демоны гибли один за другим, так до конца и не понимая, что происходит.
Сапиенты по всей планете с недоумением смотрели на соплеменников, которые вдруг начинали нечленораздельно голосить или просто падали, словно брошенные марионетки...

\begin{quote}
<<Ад, вы заплатите.
Существует-Хорошее-Небо был убит вами.
Ад, вы заплатите>>.
\end{quote}

<<Отходим!
Всем отойти за барьерную высоту!..
Узел 12, ответьте!..
Выжившие на уровне 3, всем немедленно покинуть тела!..
Докладываю, узел 9 уничто...>> --- надрывались в эфире полные паники голоса.
Союзники и враги, агенты Скорбящих, шпионы Картеля и верные служители Ордена Преисподней...

\begin{quote}
<<Ад, вы заплатите.
Пчела-Нюхающая-Вереск была убита вами.
Ад, вы заплатите>>.
\end{quote}

Слова Эй-В0 гремели в моём сознании, перекрываясь, превращаясь в тягучий, полный холодной ярости вой.
За каждым экраном следовало имя.
Друзья, сослуживцы, знакомые.
Десять, сто, тысяча имён...

В это время с орбитального спутника до меня дошло короткое сообщение Чханэ:

\ml{$0$}
{<<Нас встретили. Прощай>>.}
{``We were met. Farewell.''}

\section{[U] Прощение}

Мир замер перед моими глазами.
Я остался совершенно один во Вселенной, полной врагов, запертый в бункере с обезумевшей от ненависти человеческой женщиной.
И снова, как в былые времена на пороге гибели, ко мне вдруг пришло чувство лёгкости.

Заяц прекратила атаки, сосредоточившись на защите системы.

--- Не смей меня обвинять, Аркадиу.
Я сделала то, что должна была.

<<Я и не собирался>>.

Заяц промолчала.
Я почувствовал в этом молчании неуверенность.

--- Я люблю их всех.

<<Как и я люблю моих друзей.
Я организовал сопротивление ради того, чтобы им жилось лучше>>.

Заяц улыбнулась.
В уголках её глаз замерли слёзы.

--- Я бы хотела узнать их получше.
Особенно ту, с западным говором.
В ней есть что-то от прежних тси.
\ml{$0$}
{Наверное, беспечность.}
{Maybe recklessness.''}

<<Чханэ только что взяли в плен.
Безымянного тоже>>.

--- Помоги им, Зайчик.

Я вздрогнул.
Рядом с нами стоял рослый, наполовину состоящий из механики канин.
Седая грива, свалявшаяся борода, белые клыки и пронзительные голубые глаза.
Последние слова принадлежали ему.

Заяц потеряла контроль над своим мозгом.
Порождение больного разума спроецировалось в систему в виде голограммы.

Женщина не обернулась.
Она знала этот голос.
Он родился в глубине её сознания, он был именно таким, каким знала его она.
Голограмма тем временем подошла к женщине и положила руку ей на плечо.

--- Зайчик, --- ласково шепнул Фонтанчик.

--- Я ждала тебя, --- прошептала Заяц.
--- Знала, что ты будешь лишь плодом больного воображения.
Знала, что ты уже давно превратился в пар.
Но я хочу верить, что где-то в этом или другом мире для тебя ещё есть место.

--- Оно есть, Зайчик.
Я стою прямо за тобой.
Я пришёл забрать тебя.

Заяц заплакала.
Беззвучно и почти не двигаясь, словно герой древней легенды, вынужденный держать небесный свод и плачущий от нестерпимой боли в спине.
Женщина схватила огромный указательный палец и сжала его своей ручкой.

--- Что там? --- рассеянно спросила она.

--- Вечернее небо, прохладный ветер и цветущий вереск, --- ответил Фонтанчик.

Заяц, глубоко вздохнув, обратилась ко мне:

--- Координаты и параметры связывающего устройства.
Вы ведь так и используете изобретённую нами <<бамбуковую клетку>>?..

Система неторопливо изменила конфигурацию.
Интерфекторы учетверили атаки на защитные механизмы, но Заяц беспрепятственно выстрелила, израсходовала на экран ровно столько энергии, сколько нужно --- ни больше ни меньше.

--- Прощай, Играющая-С-Булыжником-Лиса.
Будь счастливее меня.

<<Прощай, Сотканный-Из-Темноты-Заяц.
Ты лучшая из тех, кого я знаю>>.

Из мультитула на её запятье выдвинулся тонкий щуп --- резонанс-взрыватель.
Заяц светло улыбнулась Фонтанчику, приложила щуп к виску --- и её голова превратилась в кровавую кашу.

Голограмма, бросив в мою сторону добрый мудрый взгляд лучистых глаз, растаяла в воздухе.
Интерфекторы проникли внутрь три секунды спустя.

\chapter{[:] Запах воды}

\section{[:] Чужое тело}

Тахиро шёл по заваленному обгоревшими трупами полю.
Кто-то из хргада нашёл свою смерть в жестяных коробках, начинённых метательными снарядами;
эти коробки спасали от всего, кроме ужасной смерти.
Кому-то посчастливилось увидеть перед смертью далёкий зелёный лес.
Руки, ноги, внутренности лежали как попало, словно увеличенная в тысячи раз коробочка с кормом для рыбок.
Горе, неизбывная боль для оставшихся в живых кани --- жалкая декорация для настоящей, уже давно отгремевшей битвы хоргетов.

Разумеется, эта битва была не просто декорацией.
Картель редко изводил просто так ценную биомассу.
Поток эманаций сапиентной битвы восполнил силы храбрых демонов, одержавших сегодня почти бескровную победу над адским гарнизоном.
\ml{$0$}
{Это был крестовый поход детей, скотобойня, прелюдия к триумфальному пиру.}
{It was children's crusade, slaughtering, prelude to triumphal feast.}

Тахиро поморщился и повёл в сторону занемевшей собачьей ногой.
Он плохо обращался с телами кани, а сейчас и вовсе пришлось пользоваться чужим --- настройки сообщил специально обученный легионер.

Вот сгоревший танк.
Член экипажа почти успел выбраться наружу.
С его черепа лохмотьями сползала обуглившаяся плоть вместе с ушами, по лицу текли лопнувшие глаза и расплавившийся жир.
Прохладный ветер медленно, по чешуйке слизывал с трупа золу, сгоревшую шерсть, мелкие угольки, колыхал отвалившиеся фаланги пальцев.
Зубы солдата навечно оскалились в сардонической улыбке.

Вот на дороге попался умирающий.
Его невидящие глаза бродили в глазницах, дыхание то усиливалось, то ослабевало.
Мучения должны продлиться ещё около пяти минут, пока сознание окончательно не уйдёт, свершив этот последний акт милосердия.

Тахиро вздохнул и вытащил из пустого огнестрельного оружия шомпол.
Его с детства учили, что облегчать смерть --- нелёгкое испытание.
Убийство слабо вяжется с понятием помощи.
Это следует просто принять.

Умирающий затих.
Шомпол Тахиро оставил в шее.

Вот ещё один живой.
Глаза залила кровь, но они целы.
Из голени торчит осколок кости.

--- Ra cisamau, --- попросил раненый, услышав шаги.
--- Au rr hou rama ho mere ho, cisamauaa\FM!
\FA{
Пожалуйста, помоги мне.
Я не знаю, друг ли ты, враг ли ты, но помоги, умоляю (язык мехрр-ау-о, потомки Хргада, планета Запах Воды).
}

Плачущий звук завершился настоящим плачем;
слёзы текли по мохнатому лицу хргада, смывая кровь, пыль и копоть <<грязного>> двигателя внутреннего сгорания.
Тахиро отвернулся и пошёл дальше.
Этот выживет и так.
В нём много силы.

\ml{$0$}
{--- О, какие лица.}
{``What a face.}
\ml{$0$}
{Тахиро Молниеносный, --- послышалось за спиной.}
{Tajiro the Thunderbolt,'' the voice behind him said.}

\ml{$0$}
{--- Какие спины уж тогда, --- пошутил Тахиро и поднял руки.}
{``Rather, what a back,'' Tajiro made a joke and put his hands up.}
\ml{$0$}
{--- Добровольно сдаюсь, для вашего командования есть важная информация.}
{``I surrender willingly, there's some vital information for your command.''}

\ml{$0$}
{--- Можешь повернуться и опустить руки, --- милостиво разрешил голос.}
{``You can turn around and put your hands down,'' the voice graciously allowed.}
\ml{$0$}
{--- Оставь эти телодвижения для овощей\FM.}
{``Save these movements for the vegetables\FM.''}
\FA{
Сейхмар (презрительное).
}

Тахиро подчинился.
Он знал, что о его поимке уже сообщили Гало.

Обладателем голоса оказался высокий худой канин в синей гимнастёрке без знаков различия.
Фиолетовые глаза смотрели презрительно.
Метательное оружие покоилось в кобуре --- опытный легат-интерфектор Картеля не разменивался на ненужные знаки, чем порой грешили молодые демоны.

--- Итак, я тебя слушаю.
О, hciou-rr\FM, подожди.
\FA{
Прошу прощения (высокопарн., классический хргада).
}

Демон подошёл к раненому.

--- Cisamau, --- снова попросил хргада, обернувшись на источник звука.
--- Ciso. Mauaa.

Демон наступил на искорёженную ногу солдата, и тот зашёлся в крике;
крик усиливался, пока не перешёл в визг.
Легат давил со знанием, выбирая самые крупные нервы.
Визг сменился прерывистыми стонами, и наконец раненый затих --- бешено колотящееся сердце остановилось насовсем.

--- Он мог бы дать вам больше эманаций, --- заметил Тахиро.

--- Но уже не даст, --- пожал плечами демон.
--- Он был солдатом, как и мы.
Участь солдата --- страдать и умереть.
Участь женщины --- в муках рожать новых солдат.

\ml{$0$}
{--- Всего лишь пропаганда.}
{``Filthy propaganda.''}

\ml{$0$}
{--- Хочу заметить, ваша.}
{``I have to notice, yours.}
Мы пришли совсем недавно, но здесь, в обители процветания и мира, была благодарная почва для войны.
\ml{$0$}
{Мягкотелый слабовольный скот, ещё вчера размахивавший цветами и распевавший песенки на фестивалях, взялся за оружие и принялся резать себе подобных.}
{Spineless and weak-willed cattle, waving flowers and singing songs yesterday, take weapons and start a butchery of their own kind.}
\ml{$0$}
{Забавное совпадение, не правда ли?}
{Don't you find it an interesting coincidence?''}

\ml{$0$}
{--- Они никогда не видели войны, поэтому...}
{``They had never seen a war, therefore---''}

\ml{$0$}
{--- О, перестань.}
{``Oah, I beg you.}
Да, он не знал, ради чего ему придётся умереть.
Но он каждый день видел это, --- демон пнул винтовку, --- и это, --- он махнул рукой на сгоревший танк, --- и это, --- он поднял с выгоревшей земли осколок мины.
--- Он видел это каждый день по каналам пропаганды, знал последствия применения этой техники.
И все они видели, все до единого смотрели ваши пацифистские ролики.
Неужели ты думаешь, что они не знали, на что идут?
\ml{$0$}
{Может, лучше сказать, что их учили не сомневаться, ибо сомнение подрывает любую, даже цветочно-песенную власть?}
{Maybe it's better to say that they were taught never to question, because doubt undermines any power, even a power of flowers and songs?''}

Тахиро улыбнулся и промолчал.
За жестокими словами легата прятался пацифист.
В армиях Ада и Картеля было много таких ветеранов, которые шли не убивать, а проповедовать.
Они весьма тонкими методами срывали самоубийственные операции, они давали врагам шанс спастись.
Командование всеми силами старалось удалить их из армии или даже устранить физически.
Но легионеры берегли таких легатов как зеницу ока, и ходили слухи, что благодаря пацифистам целые планеты вместе с гарнизоном, де-юре относясь к одной из фракций, де-факто соблюдали нейтралитет.

--- Я тебя слушаю, Тахиро Молниеносный.
Что ты хотел нам сообщить?

--- Как кстати, что мы заговорили об оружии.
Я хотел бы сказать пару слов касательно вашей обновки.

--- А, --- кивнул интерфектор.
--- Как вам?
Понравилось?

--- Идея интересная, если не считать некоторых отдалённых эффектов со стороны пространства-времени, --- признал Тахиро.
--- Впрочем, об этом я хотел бы поговорить с Гало.

--- Я плохой собеседник? --- ухмыльнулся легат.
--- Тахиро, мы твои фокусы впитали с молоком матери.
Выкладывай всё здесь, дольше проживёшь.

--- Хотелось бы собеседника моего уровня, Байс Старое Кольцо.
Когда тси передали экран искривлённого пространства нам, у нас тоже были придурки, которые хотели воспользоваться оружием немедленно, не разбираясь в его принципах.
К счастью, из-за умников вроде меня ты ещё жив.

Легат задумался и посмотрел на убитого им солдата.

\ml{$0$}
{--- Тси, говоришь.}
{``Qi, you said.}
\ml{$0$}
{Хорошо, сейчас будет транспорт.}
{Well, your transport will be here soon.}
\ml{$0$}
{Кстати, тебе тело не впору, на мой взгляд.}
{By the way, your body doesn't fit, I guess.}
\ml{$0$}
{Узковато в плечах?}
{Too close on the shoulders?''}

\ml{$0$}
{--- И ноги чересчур длинные, --- кивнул Тахиро.}
{``And legs are too long,'' Tajiro nodded.}
\ml{$0$}
{--- Я здесь проездом, надел, что было.}
{``I'm just passing through, I've to wear what I've got.''}

\section{[:] Беседа Гало и Тахиро}

\epigraph
{Наш враг --- не объединение, а проводимая им политика.
Мы поддержим Картель против любого внешнего или внутреннего агрессора, словно не существует доктрины угнетения, и будем бороться за свои права, словно не существует войны.}
{Politika generalia, тезис 5.
Редакция Гало Кровавый Знак}

\textspace

--- Что ж, Тахиро Молниеносный, вот ты и попал в мои руки.
Не думал, что смерть придёт ко мне в образе пленника.

--- Мы скоро умрём оба, --- заметил Тахиро.

--- Верно, --- согласился Гало.
--- Ланс не спит.
Он уже знает, что я тебя пленил.
Это идеальное стечение обстоятельств, чтобы устранить нас обоих.

--- Поверь, я не хотел жертвовать собой, чтобы убить тебя.
Это...

--- Да знаю я, --- перебил его Гало.
--- Я знаю, что ты пытался сохранить мне жизнь, чтобы таким образом вбить клин в Картель.
Но против нас сейчас борются силы куда более могущественные, чем Ланс-нат Алмаз.

--- Против \emph{нас}?

--- А ты ещё не понял?

Тахиро промолчал.

--- Знаешь, чем мы с тобой отличаемся от Ланса?
Или от того же Лусафейру? --- Гало в упор посмотрел на Тахиро ледяными глазами.
--- Мы не пытаемся соответствовать.
Мы ищем собственный путь.
Ты помнишь, каким был Лу?
Это он научил меня быть свободным.
Увы, лишь для того, чтобы добровольно запечатать себя в ячейке системы.

--- Мы --- часть системы, --- возразил Тахиро.

--- И в этом заключена главная тайна мироздания, --- подхватил Гало.
--- <<Прогресс существует тогда и только тогда, когда каждая деталь механизма руководствуется собственным пониманием комфорта>>.
При всём деспотизме Картеля эта часть его парадигмы безусловно прогрессивная.
Картель постоянно вынуждает нас бороться за свой комфорт, чтобы отделять слабых и не тратить ресурсы на то, что уже не нужно.
Вариантов деталей не может быть бесконечное множество, и обеспечение комфорта для одной детали тянет за собой обеспечение комфорта для многих других.
Только так система будет развиваться.

--- Есть детали, которые склонны к излишествам, --- заметил Тахиро.

--- Комфорт --- это нечто определённое, имеющее рамки.
Я не рассматриваю неумеренные аппетиты как вариант нормы.

\ml{$0$}
{--- Есть детали, которые не могут защититься.}
{``There are parts that cannot protect themselves.''}

\ml{$0$}
{--- Им придётся научиться.}
{``They have to learn.}
\ml{$0$}
{Это и есть прогресс.}
{That is the progress itself.''}

\textspace

Молчание затянулось.
Наконец Гало не выдержал:

--- Тахиро, ты уверен насчёт истинного равновесия?

--- Твои расчёты говорят об обратном?

--- Мои расчёты подтверждают это.

--- Тогда о чём был вопрос?

Гало криво ухмыльнулся.

--- Я всё ещё надеюсь избежать жертвы.

--- Я тоже, --- признался Тахиро.
--- Мы ведь оба знаем, чему будем посвящены.
Твоя смерть ударит по твоим врагам, моя --- по моим.
Самый чувствительный удар получат наши общие противники.

--- Если моих личных врагов пробьёт насморк, то наши общие ещё и покашляют, --- саркастически проворчал Гало.
--- Хорошо, хватит о равновесии.
Не хочу тратить остаток жизни на работу.
Расскажи лучше о Лу, как он там.

--- Держится молодцом, --- улыбнулся Тахиро.
--- Зря тебя печалят его попытки приспособиться.
Если Лу будет продолжать в том же духе, то непременно доживёт до конца войны.

--- Но не приблизит этот конец.

--- Почём тебе знать? --- пожал плечами Тахиро.
--- Он делает то, что может и должен.
Я никогда не сомневался в Лу.

--- Поэтому он выбрал тебя, а не меня.

--- Ты сам его бросил.

--- И это чёртова правда.

Тахиро промолчал.

--- Я очень по нему скучаю, --- сказал Гало.
--- Но смысла возвращаться, если достигнуто равновесие, уже нет.
Ад и Картель --- не просто противники.
Они уже не могут друг без друга существовать.
Примерно как кошка в отсутствие мышей не знает, к чему приложить коготки и зубы.

--- Примерно как женщина в отсутствие мужчины не знает, что делать со сводящим с ума либидо, --- согласился Тахиро.
--- Я тоже скучаю по Лу.
Мы даже не можем поговорить наедине.

--- Мы воевали не за это, --- процедил Гало.
--- Я бы лучше вернулся в маленький, окружённый врагами лагерь на Преисподней, чем там, на сверхбезопасном Капитуле, днями выбивать у бюрократической машины право сто пятьдесят секунд поговорить с братом.

--- А если бы всё изменилось, ты бы позволил мне...

Глаза Гало вдруг вспыхнули лютой ненавистью.

--- ... тоже общаться с ним?

Гало ухмыльнулся.
Глаза потухли так же быстро, как и загорелись.

--- С двенадцати до четырнадцати по чётным дням.

--- Ты неисправим.

\textspace

--- Мы боролись не против друг друга, а против системы.

--- Верно, --- подтвердил Гало.
--- Мы искали путь положить происходящему конец.
Мы были уверены в правильности парадигмы своих фракций и пытались распространить их огнём и мечом.

--- Не понимая, что Вселенной сейчас нужен мир, --- закончил Тахиро.

--- Это было самое тяжёлое осознание, Тахиро, --- признался Гало.
--- Я --- воин, и воины впервые за всю историю оказываются не нужны.

--- Мы нужны как никогда, Гало, --- возразил собеседник.
--- Мы разрушим эту систему.
Только воин может разрушать.

Гало засмеялся.

--- Светлая твердь, как же я хочу быть прав.
Я никогда в жизни не испытывал таких сильных желаний.
Печально осознавать, что, пытаясь убить тебя, я пытался срезать противовес, держащий меня в этой Вселенной.
Но ещё печальнее было бы то, что тебя можно было просто убить и сбежать от убийц Ланса, а я этого не сделал.

--- Вероятность нашей правоты стремится к единице, --- рассудительно заметил Тахиро.

Глаза Гало вспыхнули.

--- Ты пытаешься меня успокоить, \emph{дзайку-мару}?

--- И у меня это неплохо получается, \emph{хорохито}.

Мужчины расхохотались.
Впрочем, веселье быстро утихло.

--- Мне лезет в голову всякая религионая лабуда, --- признался Гало.
--- Помнишь, на Преисподней монахи рассказывали легенду о двух героях, которые должны были простить друг друга, чтобы войти в обитель усопших?

--- Легенда о Мичи и Таро, --- кивнул Тахиро.
--- Она была моей любимой.
Хотя я думаю, что корни её куда глубже и они... не религиозного характера.

--- Отчего так? --- с интересом приподнял голову Гало.

--- Это признание простого факта, что благодаря друг другу герои стали теми, кем стали.
Кошка может не осознавать, благодаря кому она получила острые когти, но мы ведь не кошки, верно?

Гало кивнул и опустил голову.

--- Однажды я спросил отца, о чём же всё-таки спорили Мичи и Таро.

--- Кажется, между их родами была вендетта, --- припомнил Гало.

--- Отец ответил, что это неважно, потому что их спор ничем не отличался от любого другого спора.

Гало ухмыльнулся.

--- А вообще всё это ерунда, --- заключил Тахиро.
--- Согласно верованиям моего народа, в обители усопших нет места воинам.
Они умирают навсегда.

\begin{verse}
Я был рождён в бою, но жить не начал,\\
И я не жил, я дрался, пока билось моё сердце...
\end{verse}

\begin{verse}
--- Сражением была любовь, но не отдохновением,\\
С друзьями был союз, но не весёлый пир...
\end{verse}

---  продолжил Гало.

\begin{verse}
--- Я сделал воинами тех, кого родил для жизни,\\
И вот мой час пришёл, и я умру, и больше никогда живым не буду...
\end{verse}

\begin{verse}
--- И не спасут меня ни боги моих предков, ни Бог врагов\\
От тени, что сжирает всех и вся,\\
И только меч в руке моей, как прежде...
\end{verse}

--- И только я, чтоб дать последний бой, --- хором закончили мужчины.

--- И что изменится? --- тихо смеясь, спросил Гало.

--- Я как-то задал похожий вопрос отцу, --- ответил Тахиро.
--- Мы шли по свежим вулканическим отложениям, отыскивая распустившиеся сколецитовые хризантемы.
Я постоянно оступался, разбивал себе руки и колени.
Когда я почти готов был взбунтоваться, отец погладил меня по голове и сказал:
<<Скалу терзают пыльные бури, скалу размывает вода, но она ничто не хранит так бережно, как своё имя и крохотные царапины от людских ногтей.
Имена и тропы --- это сокровище земли;
имя делает скалу значимой, а тропа даёт людям надежду.
Тем, кто пойдёт за нами, будет легче>>.

--- Это была его фраза? --- удивился Гало.
--- Я слышал её несколько раз от Лу.

--- Да, так говорил мой отец, Акено та Ханаяма.
Он знал, что нет одеяния лучше, чем благородная кожа.
Он носил её с гордостью, когда ходил в пустоши разведывать вулканические отложения.
Он был хорошим тама и хорошим человеком.

Гало промолчал и посмотрел на часы.
Такое простое повседневное действие.

--- Убийцы Ланса уже здесь, --- сообщил Тахиро.

--- Они в любом случае разделят нашу участь.
Ланс приказал после отчёта об убийстве зачистить пространство несколькими залпами экрана, а затем взорвать планету, чтобы скрыть следы.

--- Взорвать планету и уничтожить целый гарнизон из-за нас двоих? --- удивился Тахиро.

--- Ужасная, просто невероятная расточительность, --- согласился Гало.
\ml{$0$}
{--- Власть его развратила.}
{``Power has corrupted him.}
И ведь вроде бы вся эта история тысячелетиями будет покрыта мраком, если вообще когда-либо всплывёт на поверхность... и едва ли наши друзья узнают, как мы умрём... но как же неприятно сдаваться без боя!

--- Неприятно, но надо, --- признал Тахиро.
--- Самый сильный удар мы нанесём именно так --- без боя.
Понятия не имею, каково это.
Наверное, это словно тебя в детстве отправляют спать за шалость, и ты действительно идёшь спать.

Гало ухмыльнулся и посмотрел на товарища по несчастью.

--- Ты был для меня самым понимающим, самым достойным врагом.
Всегда и во все времена.

--- Ты всё ещё сердишься? --- спросил Тахиро.

--- Давно уже нет, --- поморщился Гало.
--- Кто я, по-твоему, всё тот же зелёный мальчишка?

Тахиро взглянул в застывшие ледяные глаза Гало.
Там царил мир.

\chapter*{Интерлюдия X. Мичи и Таро}
\addcontentsline{toc}{chapter}{Интерлюдия X. Мичи и Таро}

\textbf{Легенда Преисподней}

Велика сила героев, но смерть всегда вытянет на один фэно\FM{} больше.
\FA{
Фэно --- мера веса Преисподней, примерно 200 г.
}
Пришла пора отправиться по чёрной тропе и Мичи, поборовшему князя демонов, и Таро, разрубившему оковы пророчества.

Росло на вечном поле дерево сукэ\FM;
\FA{
Огнистая криптомерия --- один из немногих видов древесных растений, процветающих на планете Преисподняя.
}
цвело, и плодоносило оно по закону небес, и в свой час высохло до той сухости, когда ни пылинки жизни не остаётся в том, что имеет форму живого.
И так вышло, что прошёл Мичи по правую сторону от дерева, а Таро --- по левую.
Зацепились связывающие их нити ненависти за дерево и натянулись струнами.
Тянет Мичи --- а Таро тянет в ответ.
Да вот незадача: невелика сила умерших на чёрной тропе, не поднять им и зерна.
Отскрипели нити, откряхели герои, а победила крепость дерева сукэ.

Подумали герои да пошли к дереву, чтобы отцепить нити.
Там и встретились.

--- Перейди на левую сторону, Мичи, --- буркнул Таро, --- из-за тебя мне даже последнюю дорогу не пройти!

--- Вот ещё, --- фыркнул Мичи.
--- Кто неправ, тот пусть и переходит!

Занялся спор, и до драки дело дошло.
Да вот незадача: невелика сила умерших на чёрной тропе, не побороть им и мыши.
Устали Мичи и Таро, присели под деревом и пригорюнились.

--- Я подумал над твоими словами, Мичи, --- сказал наконец Таро.
--- В них есть крупица смысла.
Я перейду направо.

--- Я тоже подумал над твоими словами, Таро, --- ответил Мичи.
--- Они отражают тонкий луч истины.
Я перейду налево.

Так герои и сделали.
Однако нити ненависти снова зацепились о дерево сукэ, и не пойти героям дальше.

Занялся спор, и до драки дело дошло.
Да вот незадача: невелика сила умерших на чёрной тропе, не раздавить им и стрекозы.
Устали Мичи и Таро, присели под деревом и пригорюнились.

--- Я понял, Таро, --- сказал наконец Мичи.
--- Раз мы признали правоту друг друга, но по-прежнему не можем идти дальше, нам должно быть по одну сторону.
Один из нас должен уступить.

--- Я понял, Мичи, --- ответил Таро.
--- Раз мы признали правоту друг друга, неважно, кому выпадет доля уступить.
Давай бросим плоский камень.

Так герои и сделали.
Выпал жребий Мичи перейти на левую сторону.
Однако нити ненависти так опутали дерево сукэ, что не пойти героям дальше.

Занялся спор, и до драки дело дошло.
Да вот незадача: невелика сила умерших на чёрной тропе, не сломать им и сухой травинки.
Устали Мичи и Таро, присели под деревом и пригорюнились.

Шла мимо дерева девочка, умершая от акуко\FM.
\FA{
Ювенильная проказа ущелья Такэсако --- детская эпидемическая инфекция планеты Преисподняя, смертность от которой составляет 38--50\%.
}
Увидела девочка Мичи и Таро и их беду.

--- Почему вы сидите здесь, связанные этими нитями с деревом? --- удивилась она.

--- Мы не знаем, --- ответил Таро.
--- Мы боролись друг с другом, но ни один не победил.
Мы уступили друг другу, но лишь поменяли концы нитей местами.
Затем Мичи уступил мне по жребию, но мы всё равно не можем сойти с места.

--- Как много нитей ненависти, --- заметила девочка, --- возможно ли распутать все!
Вы тянете за собой свои дома\FM, давно умерших предков и даже тяжёлые слова!
\FA{
Дом --- объединение, включающее два-три родственных клана.
}
Виновато ли дерево сукэ?
Может быть, вам просто нужно друг друга простить?

--- Простить? --- удивились герои.

--- Я играла с братиком, --- сказала девочка.
--- Порой мы колотили друг друга, когда играли в героев.
Порой я кричала, что ненавижу его, когда он был хорохито.
Но потом мама звала нас кушать и мы забывали свои роли и свои обиды.

Посмотрели друг на друга Мичи и Таро.
Долго смотрели они, вечность пришла и прошла, а они всё смотрели.
И вот нити ненависти ослабли и исчезли без следа.

Подняли Мичи и Таро девочку на сплетённые руки и пошли по чёрной тропе дальше.
А дерево сукэ рассыпалось в прах, едва герои скрылись за поворотом.

\chapter{\ml{$0$}{[U] Наедине со Зверем}{Face the Beast}}

\section{[U] Ожидание допроса}

\textspace

Я почти физически ощущал натянутую вокруг <<сеть>>, замечающую малейшие изменения поля Кохани"--~Вейерманна в окрестностях моего демона.

Допрос урождённых демонов и оцифрованных сапиентов различается по методике.
С урождёнными демонами не церемонятся --- интерфектор сразу влезает в их память, добывая нужную информацию.
Ядро личности оцифрованных представляет собой виртуальную пространственную структуру с очень сложной архитектоникой, идентичную мозгу сапиента.
Добыть информацию из материального мозга можно после его изучения в замороженном виде, оцифрованный мозг можно только убедить или вынудить.

Для убеждения урождённых сапиентов существует более сорока сложнейших протоколов допроса.
Для принуждения --- один-единственный пункт:

\begin{quote}
<<Подвергнуть эскалационной обработке в Чистилище>>.
\end{quote}

<<Именно мы, оцифрованные, ковали победу Ордена Преисподней над Картелем.
Наверное, только сейчас Ад понял, насколько оцифрованные опасны>>, --- подумал я с неуместным злорадством.

Чханэ и Митрис попали в плен.
Это было ясно, как солнечное утро.
Возможно, что они даже не добрались до Тси-Ди --- их спугнул патруль, и друзья перенеслись обратно в звёздную систему Тра-Ренкхаля, кишащую приведёнными в боевую готовность агентами из-за устроенной Заяц диверсии.
У них не было шансов скрыться.

Время --- жестокая вещь.
У Заяц его было достаточно для сумасшествия.
А вот друзьям для дела не хватило.

\section{[U] Допрос}

\epigraph{И сказал Ал-ла Талиму: <<Сними одежды твои и драгоценности твои и оставь оружие твоё и нагим войди во врата джайанна\FM.
Огонь будет жечь тебя, и ибли\FM{} будут насмехаться над лицом твоим и членами твоими и походкой твоей, и дивы\FM{} припомнят самые стыдные дела твои и позор твой, но ты иди, как ходишь по дворцу твоему.
И если не склонится голова твоя, и если взгляд твой будет горд и милостив, как прежде, падут ниц стражники вторых врат, и выйдешь ты из джайанна перерождённым, и не будет тебе более нужды ни в одеждах, ни в драгоценностях, ни в оружии, ни в придворных льстецах, и будешь ты отныне равным Мне, ибо Я прошёл тем же путём>>.}
{Хакем-Аят, 14:3--5}
\FA{
Джайанна --- <<судилище>>, несотворённое место, управляемое женщиной по имени Джаблел.
Согласно Хакем-Аяту, все боги мироздания должны пройти через него, чтобы стать Творцами.
}
\FA{
Ибли --- <<мучители>>, неразумные прихлебатели Джаблел.
}
\FA{
Дивы --- <<писцы>>, бесстрастные существа, ведущие летопись Вселенной.
}

\textspace

--- Аркадиу Талианский Шакал.
Согласно нашим данным, первоначальный состав Скорбящих --- оцифрованные тси и люди Тра-Ренкхаля.
91\% --- ветераны битвы на Могильном берегу.

--- Это так, --- подтвердил я.

--- Эти данные не нуждаются в вашем подтверждении.
Некоторых из этих новобранцев мы успели схватить несмотря на то, что вы успели предупредить всех перед поимкой.

--- Кого именно вы схватили? --- спросил я.

--- Это для вас уже не имеет значения, --- ответил голос.

--- Давайте так, --- сказал я.
--- Чем больше информации вы сообщите мне, тем выше вероятность, что я стану с вами добровольно сотрудничать.

--- Давайте так --- с этого момента вы перестанете торговаться, --- ответил голос.
--- Ваша судьба зависит от тех данных, которые вы нам предоставите.
Вас позвали для уточнения некоторых деталей.
Был проведён анализ вашей личности и вынесено решение --- эскалационную обработку в Чистилище отменить.

Если бы у меня была нервная система и кожа, меня прошиб бы пот.
Начало протокола номер тринадцать.
Дело дрянь.

--- Если приведённый коэффициент пользы превысит $0.986$ --- вы будете жить.
Если суммарный коэффициент дезинформации превысит $4.2$ --- вы будете подвергнуты терминальной обработке.

\ml{$0$}
{Терминальная обработка.}
{Terminal processing.}
\ml{$0$}
{Меня замучают до полной потери рассудка в самой ужасной пыточной камере всех времён.}
{I will be tortured into total madness in the worst torture chamber of all time.}
\ml{$0$}
{Я расскажу всё, что знаю, я додумаю то, чего не знаю, я выплюну кровавым комом всю свою личность, всю свою жизнь до самых пелёнок, но механизм не остановится, пока программное ядро не перестанет адекватно отвечать на внешние сигналы.}
{I will tell what I know, I will make up if I don't, I will spit out like a blood clot all my identity, all my life to the first birth, but device will proceed---until my kernel have stopped adequately responding to external signals.}

Разумеется, приведённый коэффициент пользы на труднодостижимой, но вполне реальной цифре $0.986$ --- это уловка.
Моя судьба уже предрешена.
Вряд ли кто-то верил, что я попадусь на такой простой трюк, но протокол есть протокол.

<<Общий смысл пункта 8, подпункта .3 --- нельзя лишать пленника надежды, но единственной надеждой пленника должен быть допрашивающий>>.

Я опустил голову.
Что может сделать маленький шакал против стаи матёрых волков?
Только бы избежать Чистилища.
Страшный и недостойный конец существования...

--- Кто была женщина, которая управляла системой?

--- Тси, которую мы нашли в капсуле в состоянии анабиоза.

--- Информация соответствует истине с вероятностью 98\%.
Она согласилась сотрудничать с вами?

--- Она втёрлась к нам в доверие, чтобы реализовать план мести.
Мы потерпели поражение из-за её действий.

--- Информация соответствует истине с вероятностью 76\%.
Последний экран помог освободиться вашим союзникам, и они вступили с нами в бой.
Как вы объясните это действие женщины?

<<Вступили в бой>>.
Лёгкая оговорка, которая могла означать только одно --- Чханэ и Митрис мертвы.
Или меня пытаются в этом убедить.

--- Я убедил её помочь.

--- Информация соответствует истине с вероятностью 84\%.
Вернёмся к более ранним событиям.

Допрашивающий перешёл к другому вопросу.
Они мертвы.
Я вдруг вспомнил ласковую улыбку Чханэ.
Она любила вытаращивать глаза, когда видела, что я на неё смотрю.
А потом крепко зажмуривалась и целовала меня в нос.
Глупая привычка, которую я ношу в памяти, словно свёрток с драгоценными камнями.

Когда я оцифровал подругу, Чханэ сказала: <<Ну всё, Лис, ты от меня теперь вечность не отвяжешься>>.
Я предлагал ей найти убежище, жить в спокойствии.
Она ответила просто: <<Будем сражаться>>.
Что ещё могла ответить женщина, которую я полюбил больше жизни?
Но даже зная, какие опасности несло решение, облечённое в эти простые слова, я не был готов к её смерти.

--- Где вы прячете оборудование для оцифровки?

--- Оборудование мы прятали в храме Тхитрона, частично --- в купеческом доме.
Оно было уничтожено сразу после оцифровки, больше там ничего нет.

--- Информация соответствует истине с вероятностью 99\%.
Кто занимался сборкой оборудования?

--- Грейсвольд Каменный Молот.

--- Информация соответствует истине с вероятностью 15\%.
Суммарный коэффициент дезинформации $3.8$.
Требуется объяснение или другой ответ.

--- Первоначально предполагалось, что ветераны Могильного берега вступят в ряды Ордена, --- ответил я.
--- Они даже прошли подготовку согласно протоколу 18, что вы можете легко проверить по паттернам их действий.
Это было до того, как мне стали известны подробности политики Ада по отношению к оставшимся тси.
Я удивлён, что Грейсвольд не рассказал этого вам.

Молчание.

--- Грейсвольд Каменный Молот будет допрошен.
Как много существ вы успели оцифровать?

--- Семьсот шестьдесят восемь.

--- Информация соответствует истине с вероятностью 85\%.
Назовите всех демонов, чьи модули использовались для оцифровки.

--- Я и Митрис Безымянный.

--- Информация соответствует истине с вероятностью 70\%.
Чьи паттерны использовались для интерфекторов?

--- Анкарьяль Кровавый Шторм.
Далее интерфекторов готовили агенты под прикрытием.
Мне не известны их имена, связь с ними осуществляла Таниа Янтарь.

--- Информация соответствует истине с вероятность 89\%.

Ясное дело, ответить им на это нечего.
Интерфекторов готовила Чханэ.

<<Я не могу отомстить за подругу кровью, но сомнений вам оставлю предостаточно.
Кто знает, которая капля дезинформации расшатает вашу систему.
Излишняя подозрительность породит паранойю, партийные чистки завербуют новых агентов лучше, чем это сделаем мы.
Я не последний из тси... вернее, из Скорбящих>>.

Мысленная оговорка почему-то меня развеселила.

--- Назовите имена всех новооцифрованных.

--- Я назову, если вы сообщите имена тех, кто попал к вам в плен.

Молчание.

--- В плен были взяты тридцать восемь Скорбящих, --- сказал голос.
В моей голове возникли имена.
--- Один был уничтожен.

--- Кто был уничтожен?

--- Тхартху Танцующая Тень.

Перед моими глазами промелькнуло доброе лицо женщины, которая больше всего на свете любила рисование и сказки.
Несмотря на недостаток пленных, Ад решил уничтожить её --- оставлять визора в живых было чересчур большим риском.
Митрис Земляная Змея тоже среди пленных --- он был рядом со своей подругой до конца.

<<Прощай, Птичка.
Пусть твой путь в пристанище будет лёгким>>.

Я напрягся, и в моей памяти возникли нужные имена.

--- Информация соответствует истине с вероятностью 97\%.

--- Этот список бесполезен для вас, как и вербовка агентов.
Во время подготовки я предупредил товарищей, чтобы они относились к пленённым как к врагам, а к информации, которой располагали пленённые, как к известной врагу.

Молчание.

--- Информация соответствует истине с вероятностью 90\%.
Балл ПКП --- $0.979$, балл СКД --- $4.11$.

Я подавил приступ судорожного смеха.
Допрос окончен, последняя фраза была попыткой выжать остатки сока из жмыха.
<<Пограничные значения.
Бойся, предатель, метайся в попытках спастись, ведь спасение в семи тысячных от тебя.
Не дождётесь>>.

--- Это всё? --- осведомился я.

--- Вся требуемая информация получена.
Вы будете казнены через 16 минут.

--- Информация соответствует истине с вероятностью 100\%.
Приведённый коэффициент презрения --- единица, --- съязвил я.

Ответом было молчание.

\section{[U] Казнь}

\epigraph
{Et in Arcadia ego\FM.}
{Крылатое выражение культуры Ромай.
Древняя Земля}
\FA{
И в стране пастухов я (эллатинский).
}

\textspace

Это море совсем не было похоже на те лазурные южные моря, за которые лорды и короли ломали копья и жизни подданных.
Над северным морем всегда царила лёгкая тьма, и в этом похожем на вуаль мраке терялись далёкие горизонты.
Холодная чёрная вода шептала, а не пела, её тяжесть вгоняла в сон.
Берег был отнюдь не ласковым пляжем, на мягком прохладном песке лежали следы недавних штормов --- обломки досок, обрывки ткани, разная мелочь... и водоросли, огромные кучи бурых гниющих водорослей, от которых исходил лёгкий прелый аромат.
Чуть поодаль росли золотистые, болезненного вида сосны и тонкий, ужасно колючий шиповник с облетевшими листьями, усыпанный огромными алыми ягодами.

Я смотрел на догорающий вдали закат.
Влажный ветер бил мне в лицо, солёно-пряный аромат моря въедался в волосы, чтобы не раз вернуться обрывками сновидений в будущем.
Я думал, что обязательно должен закончить свой путь здесь.
Южные моря с их тёплой солью, кровавыми сражениями, чернявыми красотками, вином и изысканными фруктами были созданы для жизни.
Для смерти же мне не найти лучше этого всепрощающего моря, спокойного берега с его заспанными соснами и мрачновато-жизнерадостным шиповником.

Тогда я считал себя взрослым, повидавшим жизнь существом.
Кто я сейчас?
Мудрый старик, скептически относящийся к собственной мудрости?
Молодой дурак-максималист, задумавший невозможное?
Одно было несомненно --- о северном море придётся забыть.
В моём распоряжении --- виртуальная каморка, имеющая размеры лишь для моего сознания.
Темница для разума, из которой я уже не выберусь.

Может, поверить хоть сейчас, что там, за гранью, что-то есть?
Я никогда не желал рая и вечного блаженства.
Кусочек моря, Чханэ, Атрис и Митхэ, встречающие меня --- всё, что нужно.
Это будет самообман, оскорбление для истины, за которую я боролся, но смерть в отчаянии --- оскорбление для всех, кто мне верил.

Рядом материализовался Грейсвольд.
Я облегчённо вздохнул.
Одно знакомое лицо лучше, чем ничего.
Незачем что-то придумывать.

--- Это ты меня казнишь? --- осведомился я.

--- Нет, не я, --- тихо пробормотал Грейс.
Я улыбнулся.

--- Передай ей, что я буду рад её увидеть.

Грейсвольд робко ухмыльнулся.

--- Она тебя слышит.
Пока ждёт приказа.

--- Почему палачом назначили Анкарьяль, мне ясно.
Почему шестнадцать минут?

--- Они до сих пор не могут решить, что с тобой делать, --- объяснил Грейсвольд.
--- Многие кричат о том, что тебя нужно уничтожить немедленно.
Однако столько же придерживаются мнения, что тебя нужно отпустить.
Кое-кто робко предлагает с тобой договориться.
Есть и достаточно прагматичная точка зрения.
Сейчас разрабатываются новые протоколы для Чистилища, основанные на медленном, мягком воздействии, а не на... ммм... экстремальной обработке.
Исследования говорят о том, что сапиентов нужно гнуть, медленно и терпеливо, а не ломать сильным ударом.
Одним словом, им требуются подопытные.

--- Меня будут подвергать мелким неудобствам, типа имитации холодной клетки? --- скривился я.

--- Зря смеёшься.
Пыточная камера ломает дух не так часто, как это делают годы тюрьмы, медленные болезни и нищая старость.
Хотя основное направление... ммм... это удовольствие.
Не те мощные оргазмы, чередующиеся с дикой болью, как при обычной обработке, а простые, вроде тепла постели, поцелуев и чувства наполненного желудка.
Разработчики говорят, что результаты многообещающие, хоть и не сиюминутные.

Я хмыкнул.

--- Ты --- недоласканный ребёнок, --- напомнил Грейс.
--- Да, детство твоего мозга осталось далеко позади, ты неоднократно проходил коррекцию личности.
Но есть кое-что, что останется с тобой навсегда.
Материнская ласка может стать к тебе отличным ключом, хочешь ты этого или нет, и шансов устоять у тебя меньше, чем у твоих друзей-тси.

--- Я всегда буду помнить о войне!
Я всегда буду помнить, что чувства --- иллюзия!

--- Насколько долго ты будешь об этом помнить, если иллюзию нельзя отличить от реальности? --- резонно спросил технолог.

--- И зачем ты мне это рассказываешь?

--- Я молю лесных духов, чтобы сторонники казни победили, --- глухо сказал Грейсвольд.
Это была правда.
Я чувствовал его чёрное отчаяние.

Если сказанное --- истина, то это не самый плохой вариант.
Я увижу своё северное море.
У меня будет мать и тёплый очаг, Чханэ, дом, розовый куст и сад с цветущими вишнями.
Я смогу...

Я резко оборвал мысль.
Вот они, эти бреши, о которых говорил Грейсвольд.
Мои собственные потаённые мечты, из-за которых я всё и начинал...

Я должен умереть сейчас.
Иначе...

Но я же хочу жить!
Я не боюсь смерти, но ужасно хочу жить!
Я дрался ради жизни, ради опыта существования в лучшем мире!

<<Изверги>>, --- жалобно выплюнул мой разум.
Однако море и вишни упорно маячили перед внутренним взором...

Предел есть у всех.
Похоже, что я своего достиг.
Неважно, какая судьба меня ждёт --- прежний Аркадиу Люпино прекратит существование здесь и сейчас.

Технолог покряхтел и, создав виртуальный стул, сел на него.
Такая простая, совершенно человеческая попытка оттянуть неизбежное.

Мы помолчали.

--- Да, если ты хочешь знать масштабы действий.
В Аду было выявлено сорок тысяч агентов Скорбящих.

Сорок тысяч.
На нашей стороне было три процента адских демонов.
Я рассчитывал на десяток, на сотню, на пять сотен, но даже в самых смелых мечтах я не представлял, что у меня будет такая команда --- огромная законспирированная децентрализованная сеть, проросшая Ад, словно гигантская грибница.
Сорок тысяч.
Цифра бесцельно дрейфовала на поверхности разума, словно принесённый отгремевшим штормом обломок с закладной доской.

--- Более трёх тысяч пленены отделом 100 и уничтожены, остальные успели уйти.
И это только урождённые сапиенты.
Количество урождённых демонов, которые работали на вас, неизвестно, эту информацию хранят под строжайшим секретом, во избежание...
Ты поднял большую волну, Аркадиу.

--- Нет нужды поднимать то, что поднимается само, --- возразил я.
--- Почва для подобного движения уже давно перезрела.
Число сторонников говорит само за себя.
Что-нибудь ещё?

--- Есть данные, что... --- Грейсвольд замялся.

--- Договаривай.

--- Ад начал переговоры с Картелем.

--- Что? --- осведомился я, уверенный, что ослышался.

--- Именно так.
В Картеле Скорбящие тоже были.
Демоны приостановили военные действия из-за страха перед вами.
Сейчас обсуждается вопрос о тотальной чистке в обеих организациях.
Оцифрованных хотят извести под корень.

--- Это ложь, --- улыбаясь, прошептал я.
--- Клан Тахиро, клан Усмане...

--- Почти все члены клана Тахиро были агентами Скорбящих, --- вздохнул Грейс.
--- Из трёх тысяч уничтоженных две --- клан Тахиро, первая и третья генерации.
Они приняли на себя самый тяжёлый удар.
Ещё семьсот --- Усмане.
Когда до Ада дошла тревога, Тахиро попытались прийти на помощь Скорбящим-тси.
На Капитуле и ещё семи планетах произошли полномасштабные войны хоргетов.
Сорок процентов Тахиро, восемьдесят два процента Усмане уничтожено, потери Ада --- двенадцать к одному, и это официальные данные.
Реальное положение дел намного хуже.
Усмане остались бы в стороне и попытались переждать бурю, но пара их старейшин попали под горячую руку, и весь клан, не разбираясь в ситуации, прислал Аду недвусмысленное сообщение: <<Ahirat>>\FM.
\FA{
Ahirat --- послание, использовавшееся на планете Земля Врачевателей.
Значение: <<Вы пересекли последнюю границу вражды.
Мы будем убивать вас всем, чем можно убить, пока не умрём сами>>.
}
Всё в лучших традициях Хакем-Аята\FM, только чёрно-зелёной ткани не хватает.
\FA{
Священная книга народов Земли Врачевателей, регламентировавшая в том числе правила ведения войны.
}
Боюсь, что этот клан прекратил своё существование --- они всегда считались самыми жестокими среди оцифрованных, но в пунктуальности им не откажешь.

Я промолчал.
Где-то на краю сознания гремели воинственные крики, вопли отчаяния и горестные восклицания на давно забытых языках.
<<Тахиро-джиме, банджай!>>
<<Монд-ааа! Джайаннау-ле харек, монд!>>

--- Анкарьяль получила приказ о твоей ликвидации.

Анкарьяль появилась рядом.
Лицо женщины было бесстрастно.
Я встал со стула, расправил плечи, поднял голову и улыбнулся друзьям.
Хоть это и виртуальный мир, хоть у меня больше и нет ни плеч, ни головы, хоть это и не изменит мою судьбу, но всё же...
Вдруг это изменит что-то другое, например, Вселе...

--- По ощущениям похоже на засыпание, --- сказала Анкарьяль.
Её глаза встретились с моими.

\section{[U] Разными дорогами}

\epigraph
{Лучшая победа --- это победа, которая носит маску поражения.}
{Ликан Безрукий, создатель языка Эй}

Кто я?
Где я?

Я шёл по дороге, запинаясь о камни.

--- Меня зовут Аркадиу, --- сказал я вслух и сам удивился собственным словам.

Рядом со мной шли трое --- две женщины и мужчина.
Их лица были суровы, они смотрели только на меня.

Мысли путались, ощущения едва задевали поверхность разума.
Осталось одно огромное чувство, словно дыра в голове --- острое чувство утраты.

--- Raj sejra? --- спросила та, что повыше, агрессивно скривив губы.

--- Ta oj, Nar.
Kar lejma, --- ответил мужчина.

--- Dejnalo oj na mia, --- добавила низкорослая, как мне показалось, немного печально.
Она единственная меня жалела.

Откуда-то в памяти всплыли слова --- язык Эй, таблица B0.
Но из сказанного я не понял ни слова.

Вдруг пришло понимание --- меня только что убили.

--- Что вы делаете со мной? --- спросил я.

--- Ga, --- мужчина грубо толкнул меня в спину, и я свалился на колени.
Мой тюремщик вздёрнул меня на ноги и молча рукой указал направление.

--- Jol.
Laj de, --- сказала высокая и, вытащив оружие, направила его на меня.
Пистолет, метающий металлические стержни.

<<Эти существа убили меня, а сейчас хотят убить ещё раз...>>

Страшно заболела голова.
Я снова упал на колени.

--- Прекратите это, я не хочу это больше чувствовать, --- пробормотал я.
--- Кто вы?
Почему мне так плохо?

Мужчина кивнул высокой.

Я инстинктивно зажмурился, ожидая боли.
Кровь пульсировала в висках, как боевой барабан...

Боевой барабан Талии.
Большой барабан из воловьей кожи.
Он звучал точно так же.

--- Чего стоишь? --- недовольно сказал мужчина.
--- Глаза открой.

Я осторожно приоткрыл глаза.
Высокая стояла, устало прислонившись к мужчине.

--- Он всё-таки туповат.
Как думаешь?

--- Это да, --- согласился мужчина.
--- И ведь понимает, но самое главное до него не доходит.

Я смотрел на них во все глаза.

--- Вы меня не убьёте?

--- Боюсь, что однажды не выдержу и прихлопну, --- посулился мужчина.
--- Открой глаза!
Да не эти!..

Мир вновь обрёл краски.
За каменными незнакомыми лицами я увидел родные, ласковые черты.
Память вернулась.

--- Неужели ты и в самом деле думал, что я убью друга? --- засмеялся Грейс.
--- Плохо же ты меня знаешь.

--- Но как ты догадался?
Информации о связи было очень мало!

--- Я просто не поверил, что после произошедшего ты бросил меня навсегда, --- тихо сказал Грейсвольд.

Я задумался.

--- Что произошло?
Как ты обманул?..

--- Обычное копирование, Аркадиу.
Болванку Анкарьяль носила с собой.

--- То есть технически я мёртв?

--- Технически --- да.
И по документам Ордена ты тоже значишься как ликвидированный элемент.
Здорово, правда?

--- Эээ...

--- Не будь идиотом, Аркадиу.
Думаешь, мы бы с Анкарьяль ещё существовали, если бы у одного интересного типа не хранились наши копии?
Тем более твоё тело осталось живым, а значит --- ты как личность не почувствовал небытия.
Тело --- необязательная часть демона, и в век ангельских технологий обратное также верно.

--- И как это регламентируется законами Ада?

--- Никак.
Знают многие, использует кто может.
Неписанное правило таково: если твоя смерть доказана --- значит, ты вне игры.
Вот тут в дело и вступает твоя команда.
Фальсификация спасений --- это целая отрасль в наше время, разумеется, тоже негласная.
Думаешь, на высших иерархов не было покушений и даже более того --- удачных покушений?
Если рядом есть приспешники, в любой момент готовые восстановить тебя из резервной копии, смерть --- не самое страшное событие.

--- Кто ещё жив? --- допытывался я.
--- Тахиро? Айну? Гало?

--- Все они в пристанище духов, --- ответил Грейсвольд.

--- Что такое <<пристанище духов>>?
Хранилище резервных копий?
Или это, чтоб вас всех, укромная планета, где живут те, кого давно считают мертвецами?

Технолог промолчал.
Я вздохнул.

--- Я тоже туда попаду?

--- Увы, --- развела руками Анкарьяль.
--- Для тебя и прочих Скорбящих пристанища не будет.
Поэтому --- береги себя.

--- То есть всё это бесполезно, --- опустил голову я.
--- Нужно бороться с системой, о размерах которой я...

--- Я рада, что ты наконец понял задачу, --- вмешалась низкорослая.
--- Бороться нужно с системой, а не с отдельными личностями.
До тех пор, пока устаревшая система не изменится, она сама будет рождать бунтарей.
Против этого невозможно выстроить оборону, от этого нет спасения.
Те, кто послали тебя на смерть, этого так до конца и не поняли.

Анкарьяль кивнула.
Низкорослая женщина улыбнулась одними губами.
Её глаза, сильно косящие в разные стороны, по очереди посмотрели на меня --- словно два разных человека, весёлый и задумчивый.

--- Кто-то должен рассказать сапиентам Земли, что они на что-то способны, --- проговорила она, манерно растягивая слова.
--- Гонка ещё не завершена.
Да, кстати, Стигма Чёрная Звезда.
Для меня большая честь, Аркадиу Люпино.

Я пожал её тонкую, но крепкую руку.

--- Как хорошо знакомиться с друзьями, да? --- подмигнула Анкарьяль.

--- Я достаточно давно возглавляю группировку, стремящуюся к союзу с Картелем, --- усмехнулась Стигма.
--- Успела накопить большой опыт.
Так что ты не первый, господин Люпино.
Но отдаю тебе должное --- на такие наглые действия мы пока не отважились.

--- Это провал, --- возразил я.

--- Я была лучшего мнения о тебе, как о стратеге, --- нахмурилась Стигма.
--- Мы собрали огромное количество практически бесценного материала.
Такая информация просто не может достаться без жертв.
И я очень рада, что смогу работать с тобой и дальше.

--- Он ещё в себя не пришёл, --- объяснил Грейсвольд.
--- Это нормально.

--- Это предательство, --- сказал я.
--- Вас могут...

Анкарьяль внезапно обняла меня изо всех сил.
И тут же отстранилась.

--- Могут, --- подтвердила Стигма.
--- И даже сделают, я уверена.

--- Предательство, --- повторил Грейс и вздохнул.
--- Понятие, ставшее игрушкой в руках пропаганды.
То, что мы сделали, скорее гуманность.

--- Иди, Кар, --- пробормотала Анкарьяль.
--- Тебя ждут.

Я вдруг вспомнил всё.
Каждую деталь нашего сокрушительного поражения.

--- Меня больше некому ждать, --- с горечью ответил я.

--- А твой народ? --- спросила Стигма, указывая куда-то назад.

Я не стал оборачиваться.
<<Взгляд>> почувствовал ободряющий свет Чханэ и Митриса.
Грейс смотрел на меня нежно, словно на любимого ребёнка.

--- Грейс, я перед тобой в неопла...

--- Молчи, --- прервал меня Грейсвольд.
--- Молчи и иди, человек.

Я опустил голову.

--- Что с остатками тси?

--- Сожалею, --- ответил технолог.
--- Адская пропаганда и селекция уже начали своё чёрное дело --- из свободного народа делают мягкотелых счастливых существ.
Наш заговор --- непродуманный, честно говоря, --- заставит Ад закрутить гайки ещё сильнее, и я не знаю, доживёт ли свободолюбивый народ тси до момента, когда давление сорвёт крышку котла.
Все были уверены, что нас мало.
Что поделаешь, нам старательно внушают, что мы одиноки в своих стремлениях...

Грейсвольд помолчал и подумал.

--- Тси, тси.
Вообще я бы на вашем месте... --- Грейсвольд театрально оглянулся на товарищей и вполголоса продолжил:
--- Я бы на вашем месте собрал генетические образцы тси, а также данные об их культуре, и попытался бы возродить их где-нибудь на укромной планете.
Я не сторонник воскрешения отжившего, но тси по-прежнему опережают прочих сапиентов на сотню тысяч лет, было бы жалко их потерять.
Ты у нас биодиктиолог, верно?
Молодец, хорошо подготовился.
А сзади стоят две женщины --- живые носители культуры.

Я кивнул.

--- А вы?

--- А что мы? --- пожала плечами Анкарьяль.
--- Мы выполнили приказ --- уничтожили подполье.

--- Сейчас занимаемся тем, что казним тебя, --- добавила Стигма и мило усмехнулась.

--- Мы вернёмся обратно, дружище, --- заключил Грейс.
--- Пока ты остаёшься в тени, нам ничто не угрожает.
Кто-то ведь должен аккуратно убедить Орден признать сапиентов Земли, когда придёт время.
Считай то, что произошло, уроком, а не поражением.
Да, кстати, тебе послание от... сам знаешь кого.

Мы соприкоснулись головами.
В моём разуме зазвучал знакомый слог, и на душе потеплело.

--- Друзья по ту сторону баррикад?..

--- ... нужны всегда, --- закончила Анкарьяль.
--- Друзья, а не агенты.
Надеюсь, ты это усвоил.

Я оглянулся на Чханэ и Митриса.
Они ждали.
Чханэ когда-то говорила, что ждала меня вечность.
Забавная гипербола, учитывая, что ей тогда было всего сорок дождей...
Теперь я понимал, что это не поэзия, а тонкая, неочевидная истина.
Вселенная расширялась, кварки собирались в триплеты, фермионы склеивались в атомы, пыль собиралась в облака и загоралась звёздами, проклетки эволюционировали в первых людей, первые люди создали хоргетов, дали начало народам Драконьей Пустоши и Тси-Ди, а она всё это время ждала, пока детали мозаики мироздания не сложатся в нужную картину.
И всё это для того, чтобы через миг совместного пути, после лёгкого соприкосновения рук снова распасться на первородную квантовую пыль.

--- Так, я вас оставлю, --- вдруг сказала Стигма.
--- Меня вызывают по резервному каналу.
Прощай, Аркадиу.
Надеюсь, что нам больше не придётся видеться при таких обстоятельствах.

Стигма кивнула коротко остриженной головой, вскочила на мотоцикл и уехала, поднимая лёгкие клубы песчаной пыли.

Грейсвольд бросил взгляд ей вслед и улыбнулся подошедшему Атрису.

--- Мы --- это ветер, --- сказал Атрис.
--- Мы будем бросать пылинки на гранитную скалу, микрона за микроной стачивая её.
Вы --- это вода.
Вы будете проникать в мельчайшие трещинки скалы, замерзать и таять, расширяя их.
Однажды на этой скале смогут расти цветы, и она превратится в зелёный холм.
Так оживала и моя планета, и я тому свидетель.

--- Жаль, дружище, что она больше не твоя, --- посочувствовал Грейсвольд.

--- А Лотос твой, Грейс? --- хитро улыбнулся Атрис.

Грейсвольд смутился.

--- Я лишь сотворил Лотос.

--- И этого уже не изменить ни одному завоевателю, --- закончил Атрис.
--- Чтобы изгнать из горшка горшечника, нужно стереть горшок в пыль.

--- Когда-то нам поклонялись, --- сказал Грейсвольд, --- от нас ждали благ, справедливости и спасения.
Знали бы наши верующие, как мы с тобой слабы перед системой.

--- Верно, --- согласился Атрис.
--- Впрочем, на молитвы я не отвечал уже давно, многие из моих созданий поумнели и рассчитывают только на себя.
Лучший бог --- тот, о существовании которого даже не догадываются.
Тот же, кто требует почестей, молод и глуповат.

Грейсвольд оглянулся ещё раз и снова вполголоса обратился к Атрису:

--- Атрис, ты достал кольцевую теплицу?

--- Нет, --- сказал Атрис.
--- А ты, технолог, вспомнил схему Золотого города?

--- Нет, --- с той же интонацией ответил Грейс.
Секунда напряжённого молчания --- и оба от души расхохотались.
Впрочем, смех утих почти сразу.
Мужчины долго смотрели друг на друга.

--- Если к концу войны мы... разминемся, --- голос технолога дрогнул, --- заглядывай на Лотос.
Нужно время от времени поправлять орбиту Часовой Луны, она немного сползает.
Водопад, Факел и Паутина-город требуют реставрации, людишки отбивают кусочки от фокусирующих кристаллов и делают из них ожерелья.
Объясни им, пожалуйста, что так делать непорядочно и следует уважать чужой труд.
В остальном, думаю, ты разберёшься и сделаешь даже лучше, чем было.
Ты хороший бог.

Атрис кивнул.

Пора было уходить.

Я по очереди пожал руки Грейсу и Анкарьяль.
Расстались мы в полном молчании.
Чханэ и Атрис взяли меня за талию с двух сторон и повели на север.
Маленькая Митхэ шла рядом, окуная ступни в набегающую солёную пену.
На её лице впервые за долгое время появилась счастливая улыбка.

--- Так клетки всё-таки у тебя, --- сказала она.
--- Ты успел...

--- Я всегда тебе верил, --- пожал плечами Атрис.
--- Ты сказала <<Сейчас или никогда>>.
И в последний миг прилетело несколько клеток с сорванного ветром листа.

--- Меня терзает чувство вины из-за несчастного создания, которому я дала ложную надежду.
Эти тянущиеся ко мне руки я буду помнить до нашего последнего мгновения.

--- Надежда не ложная, --- сказал Атрис.
--- Мы вернёмся, и ты сама дашь ей свободу.
Теплицы живут очень долго.
Наверное, сейчас она спит и видит счастливые сны.

--- Ты думаешь, мы сможем отвоевать Тси-Ди у Машины? --- спросила Чханэ.

--- Может быть, мы сможем с ней договориться, --- предположил я.
--- Твоей искренности, Чханэ, поверил даже Лусафейру.
Поверит и Машина.

--- Если получится, я буду петь моему деревцу каждый день, --- пообещала Митхэ.
Атрис кивнул.

Я ещё раз повторил про себя слова послания Лусафейру:

\begin{quote}
<<Сражение закончено, но борьба продолжается.
Я не буду поддаваться.
Это сделает тебя слабее, чем ты можешь быть.
Я оставил тебя в живых --- неплохая фора, на мой взгляд>>.
\end{quote}

<<Однажды, мой друг, мы встретимся как равные, и я пожму твою руку>>, --- таким был ответ, который я передал Грейсвольду.
Разумеется, Лусафейру может его не получить --- конспирация.
Но, возможно, когда-нибудь...

Уже отойдя на десяток шагов, я услышал тихую речь Анкарьяль:

--- Грейс, все схемы Золотого города у тебя?
Ведь я права?

--- Не понимаю, о чём ты, --- отмахнулся старый технолог, говоря намного громче, чем следовало бы.
--- Нет у меня никаких схем.

И уже тише добавил:

--- Пожалуйста, хватит стирать себе память каждый раз, когда до тебя доходит.
От отдела 100 меня это не защитит, да и ты перестанешь надоедать с этим вопросом.

--- Эээ... каким вопросом? --- удивилась Анкарьяль.

Грейсвольд громко захохотал.
Я оглянулся и увидел, как он прижал крякнувшую от неожиданности подругу к могучей груди.
К городу они пошли пешком, обнявшись, и их следы на мокром песке --- неуклюжую цепочку больших и маленьких ступней --- смыла набежавшая ласковая волна.


\chapter*{Очень длинный постскриптум}
\addcontentsline{toc}{chapter}{Очень длинный постскриптум}

\textspace

--- Ну вот и всё, я закончил, --- сказал толстяк.

--- Писатель из тебя, как из амёбы кольцевая теплица, --- флегматично заметил его визави, прикрыв на секунду глаза.
--- Сказывается отсутствие культурологической подготовки.
Логика повествования нарушена в четырёх местах, двенадцать фактических ошибок, не укладывающихся в норму литературного искажения.
Последняя глава вообще скомканная.

--- Иди в дупло, --- обиделся толстяк.
--- Я старался копировать стиль Аркадиу.

--- Не нужно копировать чей-то стиль, --- возразил собеседник.
--- Книга твоя, а не Аркадиу.

--- Хорошо.
Найду кого-нибудь другого, кто посмотрит на смысл, а не на статистику.
Где Митрис?

--- А я знаю? --- невозмутимо ответил второй.
--- Они гулять пошли.

Наступило молчание.
Толстяк сосредоточенно просматривал светящийся голубым текст на голографическом экране.
Нахмурился, задумчиво почесал подбородок и вдруг бросил весёлый взгляд на собеседника.
Тот бесстрастно попыхивал трубкой.

--- Наслаждаешься отпуском?

--- Жизнью, Грейс.
Это называется <<жизнь>>.

--- Надымил здесь.
Знаешь же, что я не люблю.

--- Иди в дупло, пчёл доить.
Потерпишь.
Я не курил с тех времён, когда Гало переметнулся к Картелю.
Мы любили дымить после тяжёлого дня.
Тело тси, --- курильщик погладил себя по плечу, --- по-другому воспринимает никотинаткаллунаин.
Он метаболизируется быстрее, чем я докуриваю трубку.
Можешь собрать такое тело, чтобы меня туманило?

--- Тси специально...

--- Да-да, знаю.
Но я хочу, чтобы меня туманило.
Я всё равно не буду размножаться, обещаю.

--- Знаю я, как ты не будешь размножаться.
Насеял по моей планете тогда.

--- Я просто был не в форме.
И не насеял, а провёл полевые испытания навыков социальной коммуникации, --- курильщик многозначительно указал в потолок.
--- Ты же мне сам обновления ставил, помнишь?

--- Каждую ночь в течение двух лет, без контрацепции, под всеми наркотиками, которые только можно найти, --- покачал головой Грейс.
--- Ты это называешь <<социальной коммуникацией>>?

--- Я увлёкся исследованиями, --- хмуро сказал Лусафейру.
--- Кстати, их можно считать успешными --- навыки сработали на ста процентах из более чем тысячи испытуемых.
Целомудренные девы бросались на меня, словно проститутки, непробиваемые стражники с упругими задницами открывали двери гаремов и сами ложились под меня.
И это, заметь, только словом и лаской, никаких химикатов.
Признаю, с гормонами у моего тела было не всё в порядке, но...

--- С головой у тебя было не всё в порядке, Лу.
Те навыки я уже испытывал, и...

--- И как? --- с интересом спросил Лу.

--- Я их не так испытывал! --- Грейсвольд покраснел, как спелая акхкатрас.
--- Полевые испытания проводятся аккуратно, чтобы не беспокоить устоявшийся социум.
\ml{$0$}
{А ты прошёлся по Лотосу, как ураган с членом!}
{But you passed through Lotus like a hurricane with a dick!}
\ml{$0$}
{О тебе ещё триста лет легенды складывали!}
{I had been hearing legends about you for three hundred years since!''}

\ml{$0$}
{--- Расскажешь как-нибудь.}
{``You should tell them sometime.}
\ml{$0$}
{Аж интересно стало.}
{I'm intrigued.''}

--- Кстати, мою гордость --- Водопад, Факел и Паутину-город\FM{} --- ты так и не увидел, потому что, выражаясь культурно, проспал равноденствие.
\FA{
Знаменитый ландшафтный комплекс планеты Лотос, каждое равноденствие образующий за счёт оптических эффектов гигантский голографический фильм о Древней Земле.
}
А про то, что в Гранобле твои любовники и любовницы нас чуть не убили и нам пришлось уходить верхами, я вообще молчу...

--- Хорошо, обойдусь, --- примирительно поднял свободную руку визави.
--- В адском вереске не алкалоиды главное.

Толстяк ухмыльнулся.

--- Не жалеешь, что сбежал?

--- С одной стороны, мне пришлось избавиться от мощных модулей и надстроек.
Я не настолько умён, как раньше.
С другой стороны, сейчас я самый счастливый демон во Вселенной.
Я живу полной жизнью.
И ты, балбес, сидишь рядом со мной, а не за охраной из десяти интерфекторов.
Красота.
Не хватает лишь Тахиро и Айну, но выжить они просто не могли.

Грейсвольд изумлённо посмотрел на друга.
Последняя фраза была заключением стратега.

--- Тахиро и Айну были неразборчивы в средствах, --- объяснил Лусафейру.
--- Сказалась деспотичность общества, в котором они жили.
Такие личности часто добиваются славы, но редко --- любви.

--- Ты тоже иногда этим грешил, --- заметил Грейсвольд.

--- Я всегда знал границы и ценил жизни демонов, --- возразил Лусафейру.

--- Перекладывая ответственность на других, --- подхватил Грейсвольд.
--- К примеру, на того же Тахиро.

--- Да, перекладывая ответственность за чужие жизни на их обладателей, --- ничуть не смутившись, подтвердил Лусафейру.
--- По-твоему, это плохо?
Я никого не посылал на самоубийственные операции.
Тахиро пожертвовал собой сам.
Да и у вас был выбор --- идти на Тра-Ренкхаль или не идти.

--- Складно поёшь.

--- Ещё бы.
Именно поэтому меня никогда не признают военным преступником, хотя любой командующий в \emph{такой} войне --- преступник априори.
У меня больше нет ни славы, ни положения.

--- Слава как раз-таки есть, но с нюансом.

--- Меня постараются вычеркнуть из всех исторических хроник и забыть, так что славы точно не будет.
Но зато у меня есть отличный вересковый табак.
Этого Тахиро уже не видать.
Хотя я бы отдал всю свою табакерку, чтобы он ещё раз со мной поговорил.

--- Низко же ты ценишь своего друга, --- поморщился Грейсвольд.

--- Я бы выкупил беседу за табак, но не продал бы её за ту же цену.
Я нищий, Грейс.

--- Да, Ордену мы отдали чересчур много, --- согласился технолог.
--- Никогда не поздно начать новую жизнь, но её цена растёт со временем.

--- Я за свою рассчитался, с твоей помощью, --- ухмыльнулся Лусафейру.
--- Я люблю тебя, Грейсвольд Каменный Молот.
Будь благословен тот час, когда тебя забросило на Преисподнюю.

--- Ещё не рассчитался, друг, --- грустно заметил Грейсвольд.
--- Тебя по-прежнему ищут головорезы обеих фракций.
Такова цена славы.
Так что сиди здесь и не высовывайся.

--- Больно надо, --- поморщился Лусафейру и сделал глубокую затяжку.
--- Кстати, в этом простом факте --- в том, что я здесь, а не там --- наше с тобой отличие от того же Тахиро.
Тахиро чересчур буквально понимал слово <<жертва>>.
Но что поделаешь, Тахиро был воином, а воины --- априори самоубийцы.

--- Не совсем понимаю, что ты...

--- Да всё ты понял, Грейс, --- нетерпеливо сказал Лу.
--- Необязательно умирать, чтобы пошатнуть истинное равновесие Нэша.
Достаточно было пожертвовать своим положением, как это сделали мы с тобой.
Когда рождённый править Вселенной нищ и преследуем всеми, это тяжкий удар --- и не всякая Вселенная его переживёт.

--- Аркадиу...

--- А у Аркадиу не было другого выбора.
К тому времени, как он прозрел, его место в системе ещё не было определено.
Путь ветра всегда тяжелее пути воды.
Кстати, насчёт рукопожатия ты выдумал.

--- Нет, это правда.
Просто Аркадиу погиб до того, как я получил возможность с тобой поговорить, а после его смерти слова потеряли смысл.

Второй хмыкнул.

--- Шакал придавал этому жесту чересчур большое значение.

--- В его родном мире рукопожатие находилось под запретом, --- объяснил толстяк.
--- Религиозные деятели говорили, что те, кто пожимает друг другу ладони, выказывают преступное пренебрежение перед высшими силами и обществом, ведут себя так, словно они одни во Вселенной.
В качестве приветствия между равными в обычае был <<жест аиста>>, который служил одновременно приветствием, гоноративом для короля и мольбой к Великому.
Рабы кланялись господам, господа в виде особой милости показывали рабам <<водопад>>.
Рукопожатие было жестом разбойников, повстанцев-крестьян и прочих, кто отрицал власть Великого и королей.

--- Я не знал, --- констатировал собеседник.
--- Мог бы и осветить этот момент в книге.
Жаль, что он не успел пожать мне руку.

--- По сравнению с тем, что он успел, это сущие пустяки.
Не всё бывает так, как мы хотим, Лу.

--- Ты это мне говоришь, Грейс? --- Лу выпустил огромное кольцо из дыма и тут же вдохнул его носом.
--- Я могу предоставить многотысячелетнюю статистику по вещам, которые происходили не так, как я хочу.
С подробным анализом.

Технолог хихикнул и задумчиво замолчал.

--- Если не терпится показать Митрису, спроси у того, кто знает его местонахождение.

--- И то верно.
Солнышко, не занята? --- обратился Грейсвольд к компьютеру.

--- Нет, Грейсвольд Каменный Молот.
Профилактика систем завершена тридцать секунд назад, --- ответил из ниоткуда приятный женский голос.

--- Скажи, где сейчас Митрис.

--- Модуль <<Язык>>, информация: 97\% вероятности, что в данном контексте речь идёт о Митрисе Безымянном.
Ответ системы: в данный момент объект находится в двух километрах от станции G-0F города 10, азимут 0.335, --- сообщил голос.
--- Грейсвольд Каменный Молот, мне передать ему сообщение?

--- Не-не, пусть гуляют, --- сказал технолог.
--- Не к спеху.
Спасибо большое.

--- Надо поработать над её адаптационными обновлениями, --- заметил Лу.
--- Она по-прежнему считает нужным вставлять техническую информацию модулей в сообщение.

--- А мне нравится, --- ухмыльнулся Грейс.
--- Это придаёт ей некоторый шарм.

--- Потому она их и вставляет, что тебе нравится.
Она же не слепая!
Вы друг на друга просто запали.

--- Я люблю женщин с коэффициентом интеллекта выше сорока тысяч.

--- Модуль <<Культура>>, раздел <<Юмор>>, информация: идентифицирована шутка, культурный контекст B8-01-663, Древняя Земля.
Ответ системы: Грейсвольд Каменный Молот, после того, сколько обновлений ты мне поставил, ты обязан на мне жениться.

Грейсвольд захохотал, и даже Лусафейру не удержался от улыбки.

\ml{$0$}
{--- Сечёт, --- признал стратег.}
{``She's hip,'' the strategist admitted.}
--- Может быть, стоит сделать для неё сапиентное тело?

--- Модуль <<Аналитика>>, информация: предложение обработано, --- отозвался голос.
--- Ответ системы: ощутить вас руками, увидеть вас глазами --- это интересный опыт.

--- Договорились, --- улыбнулся Грейсвольд.
--- Омега-модуль я тебе уже делаю, теперь ещё и тело будет.
Может быть, даже несколько.
Будешь одновременно хоргетом-демиургом, электронным устройством и толпой народа.

\ml{$0$}
{--- А ты знаешь толк в извращениях, --- хихикнул курильщик.}
{``You're definitely sophisticated in pervertions,'' the smoker chuckled.}

--- Заодно попробуешь мою стряпню и скажешь, что она великолепна.
Лу не оценил, гурман эдакий.

--- Да ты готовишь так же, как пишешь, --- поморщился Лусафейру.
--- Даже Атрис знает, что приправлять северным сбором яичницу --- это моветон.

--- А Митхэ сказала, что это вкусно!

--- Митхэ под соус полено съест и не поморщится.
Дикое дитя джунглей.

Грейсвольд захохотал.

В комнату вошла маленькая черноглазая женщина.
Друзья вопросительно повернулись к ней.

--- Вышла подышать свежим воздухом, --- пояснила она.
--- Услышала ваш смех.
Что-то произошло?

--- Вот, Харата, --- махнул Грейсвольд.
--- Лу сказал, что я плохо готовлю.
Скажи, тебе понравилось?

--- Если ты про последний кулинарный эксперимент... приправы не совсем подходили к яичнице, --- развела руками демоница.
--- Я бы добавила смесь перцев.

--- Вы сговорились, что ли? --- обиделся технолог.
--- Лу, ты её подкупил.

--- Чем? --- резонно спросил Лусафейру.
--- И ещё.
Я так и не понял, что в яичнице делали мухоморы.
Мимоза сказал, что у него от такого брат умер, а Стигма просто отказалась это есть.

--- Это не простые мухоморы.
Они на вкус как пирожные и совсем не ядовитые!

--- Именно, Грейс!
Мы наконец-то подходим к сути дела!
Ты не считаешь, что пирожные слабо сочетаются с яичницей?
Ты бы ещё сено туда покрошил, кулинар.

Харата нерешительно потопталась на месте, глядя на друзей.

--- Вы беседуете каждый день и никак не можете наговориться, --- сказала она.
--- Это удивительно.

--- Мы чересчур долго не могли поговорить, Харата, --- улыбнулся Грейсвольд.

Лусафейру продолжал смотреть на демоницу оценивающим немигающим взором.

--- Я пойду, --- сказала Харата и приоткрыла дверь.

--- Останься, --- поднял руку Лусафейру.
--- Ты хочешь побыть с нами.

--- Я не имею доступа к тому каналу, который связывает вас, --- пригорюнилась Харата.

--- Зато ты можешь построить с каждым из нас свой канал коммуникации, --- сказал Лусафейру.
--- Для этого просто нужно проводить с нами время.
Садись.

Лусафейру изящно встал с кресла и, указав на него ладонью, плюхнулся под часами.

--- Ты не должен сидеть на полу, максим, --- укоризненно сказала Харата.

--- Оставь эти обезьяньи предрассудки, --- осклабился Лусафейру.
--- Высота подставки не отражает значимость задницы.

Демоница, не желая более спорить, заняла кресло стратега.
Лусафейру начал снова набивать трубку.

--- Когда я был молодым, я мечтал о зале-диване, --- сообщил он как будто между прочим.
--- Особенно здорово смотрелся бы зал-диван для судебных заседаний.
Советники, судья и подсудимые валяются на полу, глотают прохладную цолу\FM{} и лениво обсуждают дело.
\FA{
Цола --- тонизирующий напиток Капитула из трав.
}
Я бы нарушил закон только ради этого.

Лусафейру театрально развалился и раскинул руки, рассыпав тлеющие угольки по гладкому полимеру.
Затем стратег заговорил, искусно изменяя голос от нежного женского сопрано до хриплого мужского баса.

--- <<Что будем с тобой делать, Лу?>>
--- <<А я знаю?
Передай конфеты.
И цолы плесни, женщина, у меня в кружке пусто>>.
--- <<Подсудимый, согласно закону .01, подпункт 14, вы не имеете права требовать у советника, чтобы она наливала вам цолы.
Сам подошёл и взял>>.

Харата захохотала и похлопала --- ровно четыре раза.

--- В тебе умер великий актёр, --- улыбнулся Грейсвольд.

--- Я не хоспис, чтобы во мне кто-то умирал, --- парировал стратег.
--- Кстати, о смертях.
Ты действительно решил рассказать во всеуслышание, что на тебе лежит вина за Катаклизм Тси-Ди?

Толстяк замолчал и погрустнел.
Разумеется, он ждал этого вопроса.

\ml{$0$}
{--- Я не могу больше жить с этим, --- наконец сказал он.}
{``I can't live with it anymore,'' he finally said.}
\ml{$0$}
{--- Я не надеюсь на прощение.}
{``I don't expect forgiveness.}
\ml{$0$}
{Если тси захотят моего изгнания --- это будет чересчур мягким наказанием, и я приму его с благодарностью.}
{If Qi want me to be banished, it's very lenient punishment, and I would be pleased.''}

\ml{$0$}
{--- Не захотят, --- вдруг сказала Харата.}
{``They don't,'' Jarata suddenly said.}

Грейсвольд вопросительно посмотрел на женщину.

\ml{$0$}
{--- Они знают.}
{``They know.}
\ml{$0$}
{Все до единого.}
{All of them.}
Лусафейру рассказал им во время твоего отбытия на Сцелаю.

\ml{$0$}
{--- И они позволят мне остаться?}
{``Will they let me stay?''}

\ml{$0$}
{--- Они уже позволили тебе остаться, --- пожал плечами стратег.}
{``They already have,'' the strategist shrugged.}
--- Вспомни беседы с друзьями, советы.
\ml{$0$}
{Разве хоть кто-то проявил к тебе недружелюбие?}
{Was there anybody unfriendly?''}

Технолог опустил голову и закрыл лицо руками.

--- Перестань.
Это дело двадцатитысячелетней давности.
\ml{$0$}
{Все совершают ошибки...}
{All people make mistakes ...t}
Впрочем, глядя на Вселенную сейчас, я даже не могу считать это ошибкой.

Грейсвольд промолчал.

--- Ты давно стал среди тси своим, Грейс.
К тому же, как говорила твоя давняя подруга, любому, даже дважды и десять раз дважды кутрапу, однажды может потребоваться шанс.

--- Я перед тобой...

--- Эй, иди в дупло, а?
Ты передо мной, я перед тобой.
Мы не ростовщики, чтобы друг с друга проценты трясти, --- Лусафейру потянул дым из мундштука и, выругавшись, закашлялся от попавшего в горло пепла.
--- Какую пользу принесло бы твоё изгнание?
Оно уничтожило бы тебя и лишило тси лучшего друга, товарища и советчика.
И ради чего?
Ради торжества правосудия?

Грейсвольд тихо рассмеялся.

--- Кстати, --- Лусафейру театрально повернулся к Харате, --- приятно видеть на лице минус-демона улыбку.
Но и непривычно.

Демоница мило улыбнулась.

--- Это тело не доставляет мне неудобств.

--- Но и питания от него никакого, --- заключил Лусафейру.
--- Кстати, я бы на твоём месте срезал листья на шее.

Харата ахнула и начала судорожно ощупывать шею.

--- Опять выросли?

Грейсвольд и Лусафейру переглянулись.

--- Как-нибудь почитай, что такое <<юмор>>, --- посоветовал технолог.
--- Можешь спросить у Солнышка, она прекрасно его научилась распознавать.
Харата, перестань, нет там листьев.
И вообще скоро мы выведем минус-сапиентов, и ты сможешь здесь обосноваться.
Потерпишь?

--- Куда мне деваться, --- фыркнула Харата, --- я же этими исследованиями руковожу!
Кстати, буквально вчера вечером поступили данные разведки насчёт минус-плантов.
Согласно им, лаборатория в Картеле совершила достаточно значимый прорыв, но из-за политики Картеля легализовать данные они их не могут.
Я хочу попросить Микия отправиться на нашу базу и посмотреть, стоит ли залезать к ним в лабораторию или связываться напрямую с кем-то из учёных.

--- Микия мне очень нравится, --- сказал Лу.
--- Я бы не хотел, чтобы она или Стигма покидали убежище.
Если со мной что-то случится, они вполне могут меня заменить.
Решать ей, конечно, но я против.
Побуду строгим старшим братиком.

--- Я передам ей твои слова.

--- Однажды все сапиенты известных планет станут единой семьёй и смогут двигаться дальше, --- сказал Лусафейру.
--- Разумеется, их будут ждать новые трудности, новые испытания...

--- Мы не сможем быть с ними вечно, --- заметил Грейсвольд.

--- Помнишь, Грейс? --- вдруг понизил голос Лусафейру.
--- Мы видели блуждающие огни.
Они летели с огромной скоростью, мы пытались с ними побеседовать, но не получили ответа.
Харата, мы с Грейсом видели стабильные источники возмущений ПКВ в пустоте, и это были не хоргеты!
Понимаешь?..

Демоница ахнула.

--- Ты говоришь о...

--- Он говорит о теории Вселенных, расположенных в одном и том же пространстве.
Опять, --- поморщился технолог.
--- Да, существование множества не взаимодействующих друг с другом квантовых связок объясняет колебание ПКВ.
Но для нас эта информация бесполезна.
Мы можем взаимодействовать со Вселенной фотона потому, что нас создали и обучили её жители.
Кто обучит нас взаимодействовать с прочими?

--- Их сапиенты, разумеется, --- ответил Лусафейру.
--- Мы найдём их, мы научимся с ними беседовать.
Да, нас создали в этой Вселенной, но мы можем жить не только в ней.

--- То, что вы их увидели --- невероятная случайность... --- заметила Харата.

--- Или неизвестный феномен, не связанный с вложенными Вселенными, --- подхватил Грейс.
--- Чтобы доказать существование Иных, нам придётся колонизировать огромные пространства.

--- И мы сделаем это вместе с сапиентами, --- с горящими глазами сказал Лусафейру.
--- Мы разделим этот путь с ними, чтобы попытаться найти свой собственный в темноте других миров.

Харата кивнула.
Грейсвольд грустно улыбнулся.

--- Я уже не знаю, захочу ли куда-то уходить, --- пробормотал технолог.
--- Ты прав --- изгнания я бы не перенёс.
Привязался ко всему, видишь ли, --- Скорбящим, нашему общему делу, Солнышку... --- глаза Грейсвольда на секунду расфокусировались.
--- Наверное, Лу, это и есть старость?..

--- Нет, дружище, это всего лишь зрелость, --- сказал Лусафейру.
--- Будут другие, и они пойдут дальше.
Мы подготовим им путь.

Вскоре забили сделанные под старину часы.
Из окошка вылезла механическая ящерица, пробежала по корпусу и юркнула в замочную скважину.
Затем пророс механический стебель, развернулись сделанные из стабитаниума листья, с тихим стрекотанием раскрылся оранжевый бутон.
Следом вылетели пчёлы с золотыми крылышками и закружились в изящном танце.
Одна, две, три, пять, тринадцать...
Планетное время D:000.

Технолог отложил в сторону компьютер, подошёл к окну и положил ладони на прохладный полимер.
Нарэ.
Лу как-то сказал, что для нарэ окно не должно быть чистым и прозрачным.
Грейсвольд попросил специально для него сделать брак, с нарушением всех мыслимых производственных стандартов.
Вышло неплохо --- инженеры очень старались.
Разумеется, в полимере были и идеально прозрачные, чистые участки, иначе зачем вообще нужны окна?

Здание накрыли огромные тени и тут же обернулись чёрными силуэтами на светлеющем небе.
Один силуэт был побольше --- словно стрекозий, ещё три поменьше, напоминающие птичьи.
Силуэты покружили, словно играя друг с другом, а затем унеслись к острой стеклянной кромке небесной линзы.

<<Если есть крылья, надо летать, --- рассеянно думал Грейсвольд.
--- Да, надо>>.

Огромная голубоватая планета на небосводе постепенно теряла очертания, размывалась сиреневой с зеленоватыми прожилками краской Звезды.
Облака завихрялись причудливыми кольцами, словно ветер заплетал косы в седой отцовской бороде.
Где-то в лесу робко подавали голос птицы.

Над Тси-Ди занимался слабый рассвет.

% ----- INFO PAGE -----
\newpage\thispagestyle{plain}
{\centering

% ----- LICENSE SIGNS BLOCK -----
\includegraphics[width=2.5em]{cc.pdf}~\includegraphics[width=2.5em]{by.pdf}
% ----- LICENSE SIGNS BLOCK -----

\vspace{0.2em}

% ----- LICENSE BLOCK -----
{\ml{$0$}{Данная книга распространяется под лицензией \textbf{CC~BY~4.0}.}{This book is distributed under the \textbf{CC~BY~4.0} license.}\par}
{\ml{$0$}{Подробнее о лицензии:}{Details:} \href{https://creativecommons.org/licenses/by/4.0/deed.ru}{creativecommons.org/licenses/by/4.0}\par}
% ----- LICENSE BLOCK -----

\vfill

% ----- AUTHOR AND TITLE -----
{\large\bookauthor\par}
\vspace{0.5em}
{\Large\textbf{\booktitle}\par}
% ----- AUTHOR AND TITLE -----

\vfill

% ----- BOOK INFO -----
{\ml{$0$}{Начато:}{Started:} \textit{\bookstarted}\par}
{\ml{$0$}{Последняя редакция:}{Latest revision:} \textit{\bookfinished}\par}
{\ml{$0$}{Подробнее о книге:}{The book details:} \href{https://github.com/regnveig/tofa}{github.com/regnveig/tofa}\par}
% ----- BOOK INFO -----

\vfill

% ----- DOCUMENT INFO -----
{\ml{$0$}{Дизайн обложки: \textit{Э.\,Весна}}{The cover design by E.\,Viesn\'a}\par}
{\ml{$0$}{Компьютерная вёрстка: \textit{Э.\,Весна}}{The computer layout by E.\,Viesn\'a}\par}
{\ml{$0$}{Создано с помощью \XeLaTeX}{Created with \XeLaTeX}\par}
% ----- DOCUMENT INFO -----

\vspace{0.5em}

% ----- PUBLISHER INFO -----
{\textbf{\ml{$0$}{Свободное издательство <<Цунами>>}{\textsc{Tsunami}, an independent publisher}}\par}
{\ml{$0$}{Томск 2023}{Tomsk 2023}\par}
% ----- PUBLISHER INFO -----

\vfill

% ----- QR CODES -----
{\hfill
\begin{minipage}[t]{0.25\textwidth}
{\centering
{\includegraphics[width=0.8\textwidth]{qr_github_tofa.pdf}\par}
\vspace{0.2em}
{\small{\ml{$0$}{О книге и авторе}{About the book and the author}}\par}
}
\end{minipage}
\hspace{2em}
\begin{minipage}[t]{0.25\textwidth}
{\centering
{\includegraphics[width=0.8\textwidth]{qr_license_by40.pdf}\par}
\vspace{0.2em}
{\small{\ml{$0$}{Текст лицензии}{License full text}}\par}
}
\end{minipage}
\hfill}
% ----- QR CODES -----

% }
% % ----- INFO PAGE -----
%
% % ----- EMPTY PAGE -----
% \newpage\thispagestyle{plain}~
% % ----- EMPTY PAGE -----
%
% % ----- BACK COVER -----
% \includepdf[pages={1}]{cover_back.pdf}
% % ----- BACK COVER -----

\end{document}
